\documentclass[11pt]{amsbook}

\usepackage{amsmath,amssymb}
\usepackage{booktabs}
\usepackage[colorlinks=true,linkcolor=black,citecolor=black,urlcolor=black]{hyperref}
\usepackage{cleveref}

% ---------- theorem environments ----------
\theoremstyle{plain}
\newtheorem{theorem}{Theorem}[chapter]
\newtheorem{proposition}[theorem]{Proposition}
\newtheorem{lemma}[theorem]{Lemma}
\newtheorem{corollary}[theorem]{Corollary}
\newtheorem{hypothesis}[theorem]{Hypothesis}
\newtheorem{eresult}[theorem]{Empirical input}
\theoremstyle{definition}
\newtheorem{definition}[theorem]{Definition}
\theoremstyle{remark}
\newtheorem{remark}[theorem]{Remark}
\crefname{hypothesis}{hypothesis}{hypotheses}
\Crefname{hypothesis}{Hypothesis}{Hypotheses}
\crefname{eresult}{empirical input}{empirical inputs}
\Crefname{eresult}{Empirical input}{Empirical inputs}

% ---------- status lines ----------
\newcommand{\uses}[1]{\par\nopagebreak\noindent{\footnotesize\emph{Uses: #1.}\par}}
\newcommand{\evidence}[1]{\par\nopagebreak\noindent{\footnotesize\emph{Evidence: #1.}\par}}
\newcommand{\reference}[1]{\par\nopagebreak\noindent{\footnotesize\emph{Reference: #1.}\par}}
\newcommand{\motivation}[1]{\par\nopagebreak\noindent{\footnotesize\emph{Motivation: #1.}\par}}
\newcommand{\scope}[1]{\par\nopagebreak\noindent{\footnotesize\emph{Scope: #1.}\par}}

\title[Quantum Mechanics and General Relativity under the Energy--Time Hypothesis]{Quantum Mechanics and General Relativity\\ under the Energy--Time Hypothesis}
\author{}
\date{}

\begin{document}

\begin{abstract}
This paper investigates one hypothesis: time advances in elementary steps of Planck duration, and the action transferred in each step is an integer multiple of $\hbar/2$. The only other inputs are empirical: special relativity, quantum mechanics, Newtonian gravity, and standard thermodynamics. Every inference is stated as a theorem or corollary of an explicitly listed premise set. The framework reproduces the oscillator zero-point energy, a minimal measurable time of order the Planck time, and the wave equation with exact light-cone propagation on a spacetime lattice. On larger scales it selects a unique expanding solution, $R = ct$, of a Newtonian gravitating ball with constant mass injection, and it yields a boundary thermodynamics with entropy $S = \pi k_{\mathrm B} n^{2}$ and temperature $T = \hbar/2\pi k_{\mathrm B} t$, in which $E = TS$ holds identically. The resulting constraint $H^{2} = 1/t^{2}$ coincides formally with the Friedmann equation and contains no term for a cosmological constant. Uniqueness theorems fix the injection rate $c^{3}/2G$, the conductance law of the three-component energy transfer, the Cayley form of the unitary step, and the Planck scale triad itself. Open problems are listed explicitly. Confrontation with observational data is deferred to a companion volume.
\end{abstract}

\keywords{energy--time hypothesis, discrete time, Planck units, boundary thermodynamics, holographic entropy, open systems}

\maketitle
\tableofcontents

\part{Foundations}

\chapter{The energy--time hypothesis}\label{ch:eth}

This chapter fixes the units before any hypothesis is made. Section~\ref{sec:units} shows that the constants $c$, $\hbar$ and $G$ determine a unique system of time, length and mass scales, with no free dimensionless parameter. Nothing in that section asserts that nature advances in steps of the Planck time; that statement is Hypothesis~A1 of Section~\ref{sec:hypotheses}, and an independent argument for its magnitude is given in Section~\ref{sec:measurement}. Throughout the paper, empirical inputs are stated first, hypotheses second, derived results last; no result is proved from a statement that depends on it.

\section{The dimensional basis and Planck units}\label{sec:units}

\subsection{The three experimental inputs}

\begin{eresult}[Invariance of the speed of light]\label{er:light}
There exists a universal finite constant $c$ with $[c] = L T^{-1}$, the same in all inertial frames, which bounds the speed of energy and signal transfer.
\evidence{Michelson and Morley (1887) and the null experiments that followed; Einstein (1905). By the SI definitions of the metre (1983) and the second (2019), the numerical value of $c$ is exact.\footnote{The exactness is a definition of the unit system, not a claim of unlimited precision for any particular velocity measurement; the empirical content is the frame-independence.}}
\scope{All relativistic phenomenology to date; in this paper $c$ enters every conversion between time and length.}
\end{eresult}

\begin{eresult}[Quantum of action]\label{er:action}
There exists a universal constant $\hbar$ with $[\hbar] = M L^{2} T^{-1}$ such that action is exchanged in quanta of size $\hbar$.
\evidence{Planck (1900, black-body radiation); Einstein (1905, photoelectric effect); de Broglie (1923, matter waves, $p = h/\lambda$).\footnote{de Broglie's relation appears here only as part of the historical chain leading to $\hbar$; wherever it is used as a premise in this paper, the use is stated explicitly at that point.} $\hbar$ is exact in the SI (fixed Josephson and von Klitzing constants, Kibble-balance watt realisation).}
\scope{All quantum phenomenology.}
\end{eresult}

\begin{eresult}[Newtonian constant of gravitation]\label{er:grav}
There exists a universal constant $G$ with $[G] = M^{-1} L^{3} T^{-2}$ such that the non-relativistic gravitational field of a point mass $M$ at distance $r$ has strength $GM/r^{2}$.
\evidence{Cavendish (1798) and the torsion-balance tradition; CODATA 2018 value $G = 6.674\,30(15) \times 10^{-11}\ \mathrm{m^{3}\,kg^{-1}\,s^{-2}}$, relative uncertainty $2.2 \times 10^{-5}$.}
\scope{Non-relativistic gravity; in this paper $G$ enters through exactly one combination, the Planck scales of Proposition~\ref{prop:triad}.}
\end{eresult}

We record the working values:
\begin{equation}\label{eq:constants}
c = 2.997\,924\,58 \times 10^{8}\ \mathrm{m\,s^{-1}},\qquad
\hbar = 1.054\,571\,817 \times 10^{-34}\ \mathrm{J\,s},
\end{equation}
with $c$ and $\hbar$ exact by definition, and $G$ carrying the uncertainty quoted in \cref{er:grav}.

\subsection{Dimensions as exponent vectors}

Every mechanical dimension is a monomial in the three base dimensions $M, L, T$; we encode it as an exponent triple. The dimensions of the three input constants are, row by row,
\begin{equation}\label{eq:dimtable}
\begin{gathered}
[\hbar] = M^{1} L^{2} T^{-1} \;\longleftrightarrow\; (1, 2, -1),\qquad
[c] = M^{0} L^{1} T^{-1} \;\longleftrightarrow\; (0, 1, -1),\\
[G] = M^{-1} L^{3} T^{-2} \;\longleftrightarrow\; (-1, 3, -2).
\end{gathered}
\end{equation}
A general combination $\hbar^{a} c^{b} G^{d}$ therefore has dimension
\begin{equation}\label{eq:general}
[\hbar^{a} c^{b} G^{d}] = M^{\,a - d}\; L^{\,2a + b + 3d}\; T^{\,-a - b - 2d}.
\end{equation}
Equation \cref{eq:general} is pure exponent arithmetic. No numerical constant enters at this stage.

\begin{lemma}[Dimensional independence]\label{lem:dim}
If $\hbar^{a} c^{b} G^{d}$ is dimensionless, then $a = b = d = 0$.
\uses{\cref{er:light,er:action,er:grav}}
\end{lemma}

\begin{proof}
Dimensionless means the exponent of each of $M, L, T$ in \cref{eq:general} vanishes:
\begin{equation}\label{eq:system}
\begin{aligned}
a - d &= 0 && (M)\\
2a + b + 3d &= 0 && (L)\\
-a - b - 2d &= 0 && (T)
\end{aligned}
\end{equation}
This homogeneous system has only the trivial solution if and only if its coefficient matrix is nonsingular. Reading the matrix off \cref{eq:system} column by column (one column per constant):
\begin{equation}\label{eq:matrix}
A = \begin{pmatrix}
1 & 0 & -1\\
2 & 1 & 3\\
-1 & -1 & -2
\end{pmatrix},
\end{equation}
with determinant expanded along the first row:
\begin{equation}\label{eq:detexpand}
\det A
= 1\cdot\bigl(1\cdot(-2) - 3\cdot(-1)\bigr)
- 0\cdot\bigl(2\cdot(-2) - 3\cdot(-1)\bigr)
+ (-1)\cdot\bigl(2\cdot(-1) - 1\cdot(-1)\bigr).
\end{equation}
The two surviving cofactors are
\begin{equation}\label{eq:cofactors}
1\cdot(-2) - 3\cdot(-1) = -2 + 3 = +1,
\qquad
2\cdot(-1) - 1\cdot(-1) = -2 + 1 = -1,
\end{equation}
and therefore
\begin{equation}\label{eq:det}
\det A = 1\cdot(+1) + (-1)\cdot(-1) = 1 + 1 = 2 \neq 0.
\end{equation}
Hence \cref{eq:system} has only the trivial solution.
\end{proof}

The value $\det A = 2$ has one consequence used repeatedly below: when we solve inhomogeneous systems of the same form, all solutions are unique and have rational exponents. No root of a measured number, and no $2\pi$, enters through this route.

\begin{proposition}[Unique scale triad]\label{prop:triad}
Up to powers, the only scales of time, length and mass constructible from $\{\hbar, c, G\}$ alone are
\begin{equation}\label{eq:triad}
t_{\mathrm P} = \sqrt{\frac{\hbar G}{c^{5}}},\qquad
\ell_{\mathrm P} = \sqrt{\frac{\hbar G}{c^{3}}},\qquad
M_{\mathrm P} = \sqrt{\frac{\hbar c}{G}}.
\end{equation}
\uses{\cref{lem:dim}}
\end{proposition}

\begin{proof}
A scale of time is a combination \cref{eq:general} with dimension $M^{0}L^{0}T^{1}$: the right-hand sides of \cref{eq:system} become $(0, 0, 1)$, and uniqueness follows from $\det A \neq 0$ (\cref{lem:dim}). We solve by elimination:
\begin{equation}\label{eq:solvet}
\begin{aligned}
&(M)\;:\quad a = d,\\
&(L)\;:\quad 2a + b + 3d = 0 \;\Longrightarrow\; b = -2a - 3d \xrightarrow{\ a=d\ } b = -5a,\\
&(T)\;:\quad -a - b - 2d = 1 \;\Longrightarrow\; -a + 5a - 2a = 1 \;\Longrightarrow\; 2a = 1,
\end{aligned}
\end{equation}
so $a = \tfrac12$, $d = \tfrac12$, $b = -\tfrac52$.
Hence $t = \hbar^{1/2} c^{-5/2} G^{1/2} = \sqrt{\hbar G / c^{5}}$.

For a length, the right-hand sides of \cref{eq:system} become $(0, 1, 0)$, and the same elimination gives
\begin{equation}\label{eq:solvel}
a = d \;\;\text{from } (M),\qquad
b = -a - 2d \;\;\text{from } (T),\qquad
2a + b + 3d = 1 \xrightarrow{\ b=-3a\ } 2a = 1,
\end{equation}
so $(a, b, d) = (\tfrac12, -\tfrac32, \tfrac12)$ and $\ell = \sqrt{\hbar G / c^{3}}$.

For a mass, the right-hand sides become $(1, 0, 0)$, and
\begin{equation}\label{eq:solvem}
\begin{gathered}
(T)\;:\; b = -a - 2d,\qquad
(M)\;:\; a - d = 2a = 1,\\
(L)\;:\; 2a + b + 3d = 0 \xrightarrow{\ (T)\ } a + d = 0 \;\Longrightarrow\; d = -a,
\end{gathered}
\end{equation}
so $(a, b, d) = (\tfrac12, \tfrac12, -\tfrac12)$ and $M = \sqrt{\hbar c / G}$. The three exponent vectors are distinct, and the uniqueness of each is \cref{lem:dim}; powers of a scale are the only remaining freedom.
\end{proof}

Working values, using \cref{eq:constants} and \cref{er:grav}:
\begin{equation}\label{eq:values}
\begin{gathered}
t_{\mathrm P} = 5.391\,247(60) \times 10^{-44}\ \mathrm{s},\\
\ell_{\mathrm P} = 1.616\,255(18) \times 10^{-35}\ \mathrm{m},\qquad
M_{\mathrm P} = 2.176\,434(24) \times 10^{-8}\ \mathrm{kg}.
\end{gathered}
\end{equation}
The parenthesised digits are the uncertainty inherited from $G$ alone (see \cref{rem:guncertainty}).

\subsection{The Planck relations}

\begin{definition}[Planck units and the Planck relations]\label{def:units}
We adopt
\begin{equation}\label{eq:planckrelations}
t_{\mathrm P} := \sqrt{\frac{\hbar G}{c^{5}}},\qquad
\ell_{\mathrm P} := c\, t_{\mathrm P},\qquad
M_{\mathrm P} c^{2} := \frac{\hbar}{t_{\mathrm P}}.
\end{equation}
The last two lines are the Planck relations: they fix the conversion rate between the time axis and the length axis, and between the time axis and the mass axis.
\end{definition}

The Planck relations are conventions of unit choice, not physical claims. What must be checked is their consistency with \cref{prop:triad} --- that the right-hand sides of \cref{eq:planckrelations} coincide with the constructed scales. The check is two lines of algebra; both are written out, because the rest of the paper depends on the absence of hidden factors here.

\begin{lemma}[Consistency of the Planck relations]\label{lem:relations}
The relations \cref{eq:planckrelations} hold as identities of the constructed scales \cref{eq:triad}:
\begin{equation}\label{eq:check}
c\, t_{\mathrm P} = \ell_{\mathrm P},
\qquad
M_{\mathrm P} c^{2} = \frac{\hbar}{t_{\mathrm P}}.
\end{equation}
\uses{\cref{def:units}, \cref{prop:triad}}
\end{lemma}

\begin{proof}
For the first relation, square both sides (an equivalence, both sides being positive):
\begin{equation}\label{eq:check1}
c^{2} t_{\mathrm P}^{2} = c^{2}\cdot\frac{\hbar G}{c^{5}} = \frac{\hbar G}{c^{3}} = \ell_{\mathrm P}^{2}
\quad\Longrightarrow\quad
c\, t_{\mathrm P} = \ell_{\mathrm P}.
\end{equation}
For the second relation, again square both sides:
\begin{equation}\label{eq:check2}
M_{\mathrm P}^{2} c^{4} = \frac{\hbar c}{G}\, c^{4} = \frac{\hbar c^{5}}{G},
\qquad
\left(\frac{\hbar}{t_{\mathrm P}}\right)^{2} = \hbar^{2}\cdot\frac{c^{5}}{\hbar G} = \frac{\hbar c^{5}}{G},
\end{equation}
where the second equality uses $1/t_{\mathrm P}^{2} = c^{5}/(\hbar G)$. The two right-hand sides coincide; taking positive roots gives \cref{eq:check}.
\end{proof}

\begin{definition}[Planck energy]\label{def:epsilon}
\begin{equation}\label{eq:epsilon}
\varepsilon_{\mathrm P} := \frac{\hbar}{t_{\mathrm P}}.
\end{equation}
\end{definition}

\begin{corollary}[Two names, one energy]\label{cor:epmp}
$\varepsilon_{\mathrm P} = M_{\mathrm P} c^{2}$ exactly.
\uses{\cref{def:epsilon,def:units}}
\end{corollary}

\begin{proof}
The second Planck relation \cref{eq:check} reads $M_{\mathrm P}c^{2} = \hbar/t_{\mathrm P}$; the right-hand side is $\varepsilon_{\mathrm P}$ by \cref{def:epsilon}.
\end{proof}

Working value:
\begin{equation}\label{eq:epsvalue}
\varepsilon_{\mathrm P} = 1.956\,08 \times 10^{9}\ \mathrm{J} \approx 1.220\,9 \times 10^{28}\ \mathrm{eV},
\end{equation}
with the same inherited uncertainty as $M_{\mathrm P}$.

\begin{remark}[Physical content of the Planck relations]\label{rem:reading}
The relation $\ell_{\mathrm P} = c\,t_{\mathrm P}$ says that the spatial spacing and the time step are not independent settings: one conversion rate, $c$, turns one axis into the other. The relation $M_{\mathrm P}c^{2} = \hbar/t_{\mathrm P}$ does the same between the mass axis and the time axis. Together they mean that the lattice of Hypothesis~A1 has a single calibration: there is no adjustable ratio between its temporal, spatial and energetic scales. What has not been shown is that any physical process realises a step of size $t_{\mathrm P}$. That is a separate statement, made only in Section~\ref{sec:hypotheses}, with its magnitude discussed independently in Section~\ref{sec:measurement}.
\end{remark}

\begin{remark}[No $2\pi$ in this section]\label{rem:twopi}
Every relation in this section is a rational function of $\{\hbar, c, G\}$ with rational exponents; no factor of $2\pi$ appears, and none can, because none of the three inputs contains an angle. This begins the paper's tracking of numerical factors: every factor of $2\pi$ encountered later is traced to a specific point of entry --- the conversion between angular frequency and cycles in Section~\ref{sec:quantum}, the phase lattice of the oscillator, or the Brillouin boundary of the propagation lattice --- and each such point is marked where it first occurs.
\end{remark}

\begin{remark}[Uncertainty carried by $G$]\label{rem:guncertainty}
$t_{\mathrm P}$ and $\ell_{\mathrm P}$ inherit exactness from $\hbar$ and $c$; $M_{\mathrm P}$ and $\varepsilon_{\mathrm P}$ inherit the relative uncertainty $2.2 \times 10^{-5}$ of $G$ (\cref{er:grav}). Every dimensionless statement, and every combination in which $G$ cancels, is unaffected. In particular, the boundary thermodynamics of Chapter~\ref{ch:boundary}, which is $G$-free, carries no uncertainty from this source, while any statement involving $\varepsilon_{\mathrm P}$ explicitly (for example the threshold of Chapter~\ref{ch:light}) carries it in full.
\end{remark}

% [document continues]

\section{The energy quantum and the oscillator floor}\label{sec:quantum}

This section collects the quantum-mechanical inputs used later in the chapter. Everything in it is standard; the presentation separates what is taken from experiment from what is derived, because the hypotheses of Section~\ref{sec:hypotheses} depend on this distinction.

\subsection{Two inputs of quantum kinematics}

\begin{eresult}[Canonical commutation relation]\label{er:ccr}
For every conjugate pair $(\hat x, \hat p)$ of a quantum system,
\begin{equation}\label{eq:ccr}
[\hat x, \hat p] = \hat x \hat p - \hat p \hat x = i \hbar\,\mathbb{1}.
\end{equation}
\evidence{Born and Jordan (1925), Dirac (1925), as the defining commutator of matrix mechanics; verified ubiquitously since, for example in the spectrum of the harmonic oscillator (below) and in Josephson tunnelling.}
\scope{All non-relativistic quantum systems; nothing in this paper uses relativistic refinements of \cref{eq:ccr}.}
\end{eresult}

\begin{eresult}[Energy--frequency relation]\label{er:efreq}
A quantum of energy $E$ and a frequency $\omega$ stand in the relation
\begin{equation}\label{eq:efreq}
E = \hbar \omega,
\end{equation}
with $\omega$ an angular frequency throughout this paper.
\evidence{Planck (1900, black-body radiation, $E = h \nu$ with $\hbar = h/2\pi$); Einstein (1905, photoelectric effect); Compton (1923, scattering).}
\scope{All quantum systems; in this paper \cref{eq:efreq} is used only as an input and is never derived --- deriving it from the quantization rule of Section~\ref{sec:hypotheses} would reverse the dependency order and is explicitly not done.}
\end{eresult}

\Cref{er:efreq} carries the first $2\pi$ of the paper. Planck's original form was $E = h\nu$ with $\nu$ a cycle frequency; writing $\omega = 2\pi\nu$ converts it to \cref{eq:efreq}. All formulas below use angular frequencies. The convention matters quantitatively at exactly one place --- the frame-rate oscillator of Section~\ref{sec:hypotheses}, where the two candidates $\omega_0 = 1/t_{\mathrm P}$ and $\omega_0 = 2\pi/t_{\mathrm P}$ differ by $2\pi$, and only the first is consistent with the injection theorem of Chapter~\ref{ch:unitarity} (\cref{rem:framerate}).

\subsection{The uncertainty relation, derived once}

\begin{lemma}[Robertson]\label{lem:robertson}
For Hermitian operators $\hat A$, $\hat B$ and any state,
\begin{equation}\label{eq:robertson}
\Delta A \, \Delta B \;\geq\; \tfrac{1}{2}\,\bigl|\langle [\hat A, \hat B] \rangle\bigr|.
\end{equation}
\uses{nothing beyond the algebra of the state space}
\end{lemma}

\begin{proof}
Remove the means, which do not affect spreads: $\hat A' := \hat A - \langle \hat A \rangle$, $\hat B' := \hat B - \langle \hat B \rangle$, so that $\Delta A = \| \hat A' \psi \|$, $\Delta B = \| \hat B' \psi \|$. Cauchy--Schwarz applied to $\hat A' \psi$ and $\hat B' \psi$ gives
\begin{equation}\label{eq:schwarz}
\bigl|\langle \hat A' \psi \,|\, \hat B' \psi \rangle\bigr| \;\leq\; \| \hat A' \psi \| \, \| \hat B' \psi \| \;=\; \Delta A \, \Delta B.
\end{equation}
Split the matrix element into real and imaginary parts:
\begin{equation}\label{eq:split}
\langle \hat A' \hat B' \rangle \;=\; \tfrac{1}{2}\langle \{ \hat A', \hat B' \} \rangle + \tfrac{1}{2}\langle [ \hat A', \hat B' ] \rangle .
\end{equation}
The anticommutator expectation is real; the commutator expectation is purely imaginary, because $[\hat A, \hat B]^{\dagger} = -[\hat A, \hat B]$. Hence the modulus of the sum is at least the modulus of its imaginary part:
\begin{equation}\label{eq:imag}
\bigl|\langle \hat A' \hat B' \rangle\bigr| \;\geq\; \tfrac{1}{2}\,\bigl|\langle [\hat A, \hat B] \rangle\bigr|,
\end{equation}
where $[\hat A', \hat B'] = [\hat A, \hat B]$ since constants commute. Chaining \cref{eq:imag} with \cref{eq:schwarz} gives \cref{eq:robertson}.
\end{proof}

\begin{lemma}[Kennard]\label{lem:kennard}
\begin{equation}\label{eq:kennard}
\Delta x \, \Delta p \;\geq\; \frac{\hbar}{2}.
\end{equation}
\uses{\cref{er:ccr} via \cref{lem:robertson}}
\end{lemma}

\begin{proof}
Substitute $\hat A = \hat x$, $\hat B = \hat p$ in \cref{lem:robertson}: $\langle [\hat x, \hat p] \rangle = i\hbar$ by \cref{eq:ccr}, whose modulus is $\hbar$.
\end{proof}

\begin{remark}
In this presentation \cref{lem:kennard} is a consequence of \cref{er:ccr}, one level below an independent postulate; every use of \cref{eq:kennard} traces back to the commutation relation.
\end{remark}

\subsection{The oscillator floor}

\begin{proposition}[Ground-state floor]\label{prop:floor}
For the harmonic oscillator
\begin{equation}\label{eq:hamiltonian}
\hat H = \frac{\hat p^{2}}{2m} + \frac{1}{2} m \omega^{2} \hat x^{2},
\end{equation}
every state satisfies
\begin{equation}\label{eq:floor}
\langle \hat H \rangle \;\geq\; \frac{\hbar\omega}{2},
\end{equation}
with both terms of \cref{eq:hamiltonian} contributing $\hbar\omega/4$ at the minimum.
\uses{\cref{lem:kennard}}
\end{proposition}

\begin{proof}
Split each squared operator into fluctuation and mean parts, $\langle \hat x^{2} \rangle = (\Delta x)^{2} + \langle \hat x \rangle^{2} \geq (\Delta x)^{2}$, and identically for $\hat p$; dropping the non-negative mean squares can only lower the right-hand side:
\begin{equation}\label{eq:meanssplit}
\langle \hat H \rangle \;=\; \frac{\langle \hat p \rangle^{2} + (\Delta p)^{2}}{2m} + \frac{1}{2} m \omega^{2}\bigl( \langle \hat x \rangle^{2} + (\Delta x)^{2} \bigr) \;\geq\; \frac{(\Delta p)^{2}}{2m} + \frac{1}{2} m \omega^{2} (\Delta x)^{2}.
\end{equation}
Now use \cref{lem:kennard} in the form $\Delta p \geq \hbar / (2\,\Delta x)$, valid because $\Delta x > 0$ for any normalisable state:
\begin{equation}\label{eq:substituted}
\langle \hat H \rangle \;\geq\; \frac{\hbar^{2}}{8 m (\Delta x)^{2}} + \frac{1}{2} m \omega^{2} (\Delta x)^{2} \;=:\; g(\Delta x).
\end{equation}
Minimise $g$ over $\Delta x > 0$. The function tends to $+\infty$ both as $\Delta x \to 0$ and as $\Delta x \to \infty$, so any stationary point is the global minimum. Differentiate:
\begin{equation}\label{eq:derivative}
g'(\Delta x) = -\frac{\hbar^{2}}{4 m (\Delta x)^{3}} + m \omega^{2}\, \Delta x \;=\;0
\;\;\Longrightarrow\;\;
(\Delta x)^{2} = \frac{\hbar}{2 m \omega},
\end{equation}
where the root is unambiguous because $\Delta x > 0$. Evaluate the two terms of \cref{eq:substituted} at this point, separately:
\begin{equation}\label{eq:twoquarters}
\frac{\hbar^{2}}{8 m \cdot \frac{\hbar}{2 m \omega}} = \frac{\hbar \omega}{4},
\qquad
\frac{1}{2} m \omega^{2} \cdot \frac{\hbar}{2 m \omega} = \frac{\hbar \omega}{4}.
\end{equation}
Both equal $\hbar\omega/4$: position and momentum fluctuations contribute equally at the floor, a balance forced by the minimisation. Adding the two gives \cref{eq:floor}.
\end{proof}

\begin{remark}[The two quarters]\label{rem:quarters}
The equal split of \cref{eq:twoquarters} is the classical virial statement $\langle \hat T \rangle = \langle \hat V \rangle$ for the oscillator; the virial theorem is not invoked, since \cref{eq:twoquarters} is already the direct evaluation.
\end{remark}

\begin{proposition}[Saturation by Gaussian states]\label{prop:gauss}
The bound \cref{eq:floor} is attained, with $\langle \hat H \rangle = \hbar\omega/2$, by every minimum-uncertainty Gaussian
\begin{equation}\label{eq:gaussian}
\psi(x) = \bigl(2\pi (\Delta x)^{2}\bigr)^{-1/4} \exp\Bigl( -\frac{(x - \langle \hat x \rangle)^{2}}{4 (\Delta x)^{2}} + \frac{i \langle \hat p \rangle x}{\hbar} \Bigr),
\qquad (\Delta x)^{2} = \frac{\hbar}{2 m \omega},\quad \langle \hat p \rangle = 0,
\end{equation}
and only by such states.
\uses{\cref{lem:kennard} (equality condition), \cref{eq:hamiltonian}}
\end{proposition}

\begin{proof}
Equality in \cref{eq:floor} requires two conditions simultaneously: equality in \cref{eq:meanssplit}, i.e.\ $\langle \hat x \rangle = \langle \hat p \rangle = 0$; and equality in \cref{eq:kennard}, which by the equality condition of Cauchy--Schwarz requires $\hat A' \psi \propto \hat B' \psi$, i.e.\ $\hat p \psi \propto \hat x \psi$ --- a first-order differential equation whose normalisable solutions are exactly \cref{eq:gaussian}. It remains to verify the saturation. For a real centred Gaussian, integration by parts gives $\langle \hat p^{2} \rangle = \hbar^{2} \int |\psi'|^{2}\,dx$, and $\psi' = -\frac{x}{2(\Delta x)^{2}}\psi$, so
\begin{equation}\label{eq:gaussmoments}
\langle \hat p^{2} \rangle = \frac{\hbar^{2}}{4 (\Delta x)^{4}} \langle \hat x^{2} \rangle = \frac{\hbar^{2}}{4 (\Delta x)^{2}} = \frac{\hbar^{2}}{4} \cdot \frac{2 m \omega}{\hbar} = \frac{1}{2} m \hbar \omega,
\end{equation}
using $\langle \hat x^{2} \rangle = (\Delta x)^{2}$. Then
\begin{equation}\label{eq:gaussvalue}
\langle \hat H \rangle = \frac{\langle \hat p^{2} \rangle}{2m} + \frac{1}{2} m \omega^{2} \langle \hat x^{2} \rangle = \frac{\hbar\omega}{4} + \frac{\hbar\omega}{4} = \frac{\hbar\omega}{2},
\end{equation}
by the same two evaluations as \cref{eq:twoquarters}.
\end{proof}

\subsection{The ladder spectrum}

\begin{definition}[Ladder operators]\label{def:ladder}
\begin{equation}\label{eq:ladder}
\hat a := \sqrt{\frac{m\omega}{2\hbar}}\Bigl( \hat x + \frac{i \hat p}{m\omega} \Bigr),
\qquad
\hat a^{\dagger} := \sqrt{\frac{m\omega}{2\hbar}}\Bigl( \hat x - \frac{i \hat p}{m\omega} \Bigr).
\end{equation}
\end{definition}

\begin{lemma}[{$[\hat a, \hat a^{\dagger}] = 1$}]\label{lem:aadag}
The ladder operators satisfy
\begin{equation}\label{eq:aadag}
[\hat a, \hat a^{\dagger}] = \mathbb{1}.
\end{equation}
\uses{\cref{er:ccr}, \cref{def:ladder}}
\end{lemma}

\begin{proof}
Expand the commutator term by term:
\begin{equation}\label{eq:aadagexpand}
[\hat a, \hat a^{\dagger}] = \frac{m\omega}{2\hbar}\, \Bigl[ \hat x + \frac{i\hat p}{m\omega},\; \hat x - \frac{i\hat p}{m\omega} \Bigr]
= \frac{m\omega}{2\hbar}\Bigl( - \frac{i}{m\omega} [\hat x, \hat p] + \frac{i}{m\omega} [\hat p, \hat x] \Bigr),
\end{equation}
where $[\hat x, \hat x] = [\hat p, \hat p] = 0$ killed the first and last terms. Substituting $[\hat x, \hat p] = i\hbar$ and $[\hat p, \hat x] = -i\hbar$ (\cref{eq:ccr}):
\begin{equation}\label{eq:aadagcancel}
-\frac{i}{m\omega}(i\hbar) = +\frac{\hbar}{m\omega},
\qquad
\frac{i}{m\omega}(-i\hbar) = +\frac{\hbar}{m\omega},
\end{equation}
both signs resolving to $+$, and
\begin{equation}\label{eq:aadagvalue}
[\hat a, \hat a^{\dagger}] = \frac{m\omega}{2\hbar}\cdot\frac{2\hbar}{m\omega} = \mathbb{1}.
\end{equation}
The factor $2$ in the numerator is the sum of the two equal cross terms; dropping either is the most common way to lose a factor of $2$ here, so both are written out in \cref{eq:aadagexpand}.
\end{proof}

\begin{proposition}[Hamiltonian in number form]\label{prop:numberform}
\begin{equation}\label{eq:numberform}
\hat H = \hbar\omega\,\Bigl( \hat a^{\dagger} \hat a + \tfrac{1}{2} \Bigr).
\end{equation}
\uses{\cref{def:ladder}, \cref{lem:aadag}}
\end{proposition}

\begin{proof}
Invert \cref{eq:ladder} by addition and subtraction:
\begin{equation}\label{eq:invert}
\hat x = \sqrt{\frac{\hbar}{2 m \omega}}\, (\hat a + \hat a^{\dagger}),
\qquad
\hat p = -i \sqrt{\frac{m \hbar \omega}{2}}\, (\hat a - \hat a^{\dagger}).
\end{equation}
Square each line, keeping operator order:
\begin{equation}\label{eq:squares}
\hat x^{2} = \frac{\hbar}{2 m \omega}\bigl( \hat a + \hat a^{\dagger} \bigr)^{2},
\qquad
\hat p^{2} = -\frac{m \hbar \omega}{2}\bigl( \hat a - \hat a^{\dagger} \bigr)^{2}.
\end{equation}
Substitute into \cref{eq:hamiltonian}:
\begin{equation}\label{eq:hsub}
\hat H = \frac{\hbar\omega}{4}\Bigl( \bigl(\hat a + \hat a^{\dagger}\bigr)^{2} - \bigl(\hat a - \hat a^{\dagger}\bigr)^{2} \Bigr),
\end{equation}
the minus sign coming from the $-$ in $\hat p^{2}$. Expand both squares and cancel the diagonal terms $\hat a^{2}$, $\hat a^{\dagger 2}$:
\begin{equation}\label{eq:expand}
\bigl(\hat a + \hat a^{\dagger}\bigr)^{2} - \bigl(\hat a - \hat a^{\dagger}\bigr)^{2}
= 2\bigl( \hat a \hat a^{\dagger} + \hat a^{\dagger} \hat a \bigr).
\end{equation}
Finally use \cref{eq:aadag} in the form $\hat a \hat a^{\dagger} = \hat a^{\dagger} \hat a + \mathbb{1}$:
\begin{equation}\label{eq:hnumber}
\hat H = \frac{\hbar\omega}{2}\bigl( 2\,\hat a^{\dagger}\hat a + \mathbb{1} \bigr) = \hbar\omega\Bigl( \hat N + \tfrac{1}{2} \Bigr),
\qquad \hat N := \hat a^{\dagger} \hat a,
\end{equation}
which is \cref{eq:numberform}.
\end{proof}

\begin{theorem}[Oscillator spectrum]\label{thm:spectrum}
The spectrum of $\hat H$ is
\begin{equation}\label{eq:spectrum}
E_{n} = \hbar\omega\Bigl( n + \tfrac{1}{2} \Bigr),\qquad n = 0, 1, 2, \ldots
\end{equation}
\uses{\cref{prop:numberform}, \cref{lem:aadag}, positivity of the norm}
\end{theorem}

\begin{proof}
From the Jacobi identity and \cref{eq:aadag}:
\begin{equation}\label{eq:nshift}
[\hat N, \hat a] = [\hat a^{\dagger}, \hat a]\, \hat a = -\hat a,
\qquad
[\hat N, \hat a^{\dagger}] = +\hat a^{\dagger}.
\end{equation}
Hence if $\hat N |n\rangle = n |n\rangle$, then $\hat N \hat a |n\rangle = (n-1)\hat a |n\rangle$ and $\hat N \hat a^{\dagger}|n\rangle = (n+1)\hat a^{\dagger}|n\rangle$: $\hat a$ lowers and $\hat a^{\dagger}$ raises the eigenvalue by one. Positivity bounds the descent: for any state,
\begin{equation}\label{eq:positivity}
\langle \psi | \hat N | \psi \rangle = \| \hat a \psi \|^{2} \geq 0,
\end{equation}
so $n \geq 0$ for every eigenvalue. A non-terminating descent would produce $n = -1, -2, \ldots$ in contradiction with \cref{eq:positivity}; therefore the chain terminates at some $|0\rangle$ with
\begin{equation}\label{eq:ground}
\hat a |0\rangle = 0 \quad\Longrightarrow\quad \hat N |0\rangle = 0 \quad\Longrightarrow\quad E_{0} = \tfrac{1}{2}\hbar\omega,
\end{equation}
using \cref{eq:numberform}. The terminal state is the ground state: in position space $\hat a \psi_0 = 0$ is the first-order equation $\sqrt{m\omega/2\hbar}\,(x + \frac{\hbar}{m\omega}\partial_x)\psi_0 = 0$, whose unique normalisable solution is the Gaussian \cref{eq:gaussian} with $(\Delta x)^2 = \hbar/2m\omega$ --- the same state that saturated \cref{prop:floor}. Repeated application of $\hat a^{\dagger}$ reaches every non-negative integer, so the spectrum is exactly \cref{eq:spectrum}.
\end{proof}

\begin{corollary}[Two derivations, one floor]\label{cor:consistency}
The ground-state energy obtained from the spectrum, $E_{0} = \hbar\omega/2$ (\cref{eq:ground}), equals the variational floor of \cref{prop:floor}.
\uses{\cref{thm:spectrum}, \cref{prop:floor}}
\end{corollary}

\begin{proof}
Both equal $\hbar\omega/2$; \cref{eq:ground} and \cref{eq:floor} are numerically identical statements about the same operator. The agreement is a consistency of two derivations from the same input, not an independent confirmation.
\end{proof}

\subsection{The half-unit of action}

\begin{definition}[Half of the quantum of action]\label{def:alpha0}
\begin{equation}\label{eq:alpha0}
\alpha_{0} := \frac{\hbar}{2}.
\end{equation}
\end{definition}

\begin{remark}[Origin of the half-unit]\label{rem:half}
The factor $1/2$ in \cref{eq:alpha0} is not conventional. It is the floor of \cref{prop:floor} written as $E_{0} = \alpha_{0}\,\omega$: the smallest action a bound system can carry under the commutation relation \cref{eq:ccr}. The quantization rule of Section~\ref{sec:hypotheses} will count action in integer multiples of this $\alpha_{0}$.
\end{remark}

\begin{remark}[$2\pi$ convention at the frame rate]\label{rem:framerate}
All frequencies below are angular. The first place where the angular/cycle convention changes a number is the frame-rate oscillator of Section~\ref{sec:hypotheses}: the convention $\omega_0 = 1/t_{\mathrm P}$ gives the zero-point cell $E_{0} = \alpha_0\omega_0 = \hbar/2t_{\mathrm P} = \varepsilon_{\mathrm P}/2$, while the cycle reading, $\omega_0 = 2\pi/t_{\mathrm P}$, would give $E_0 = \pi\varepsilon_{\mathrm P}$. The choice is not free: only the first is consistent with the injection theorem (Chapter~\ref{ch:unitarity}), as checked there.
\end{remark}

% [document continues]

\section{The hypothesis layer}\label{sec:hypotheses}

The inputs so far are empirical. This section states what the framework supposes on top of them: three hypotheses, A1, A2, and P. They are assumptions in the plain sense of the word --- provisional, falsifiable in principle, and derived from nothing in this paper. Two rules govern their use. First, nothing below may serve as a premise for them; every statement they support is, for that reason, conditional. Second, wherever possible each hypothesis is split into a structure part and a magnitude part, justified separately: magnitudes may receive independent experimental support (as Hypothesis~A1 does in Section~\ref{sec:measurement}), while structures stand or fall with the framework as a whole.

\begin{hypothesis}[Discrete time, A1]\label{as:a1}
Physical processes are registered on a discrete sequence of steps (frames) indexed by
\begin{equation}\label{eq:frames}
n \in \mathbb{Z}^{+} = \{1, 2, 3, \ldots\},
\end{equation}
and time enters every formula of this paper only through
\begin{equation}\label{eq:tick}
t(n) = n\, t_{\mathrm P}.
\end{equation}
\uses{\cref{def:units} (the step size $t_{\mathrm P}$)}
\motivation{The energy--time hypothesis requires a time axis whose steps can carry a quantum of action; on a continuous axis, $E\,t_{\mathrm P}$ would range over a continuum and the rule of \cref{as:a2} would have no discrete object to count. The step is a sequence marker, not the reading of any clock: two observers in arbitrary relative motion can disagree about every clock rate, yet index their processes against one and the same sequence. Whether nature in fact steps at the scale $t_{\mathrm P}$ is a question of magnitude, discussed independently in Section~\ref{sec:measurement}; the structure of discreteness is not derived anywhere in this paper.}
\scope{All statements of this paper. The exclusion of $n = 0$ is taken up in Chapter~\ref{ch:boundary}, where the first frame carries a special condition.}
\end{hypothesis}

\begin{hypothesis}[Action quantization, A2]\label{as:a2}
Every exchange of energy $E$ registered on the time axis of \cref{as:a1} is quantized in integer multiples of the half-unit of action:
\begin{equation}\label{eq:quantization}
E\, t_{\mathrm P} = k\, \alpha_{0},\qquad k \in \mathbb{Z}^{+},
\end{equation}
equivalently $E = \frac{k}{2}\, \varepsilon_{\mathrm P}$.
\uses{\cref{as:a1}, \cref{def:alpha0}}
\motivation{Section~\ref{sec:quantum} showed that bound systems already live in half-units of action: their floor is $E_{0} = \alpha_{0}\omega$ and their spectrum is spaced by $\hbar\omega = 2\alpha_{0}\omega$ (\cref{thm:spectrum}). Equation \cref{eq:quantization} extends the half-unit from the oscillator to the time axis itself: action occurs in quanta of size $\alpha_{0}$, and energy is the quantum number $k$ read in units of $1/t_{\mathrm P}$. Historically, rules of this form descend from the Sommerfeld--Wilson condition $\oint p\,dq = n h$ (1915); the present rule quantizes the energy--time product instead of a phase-space loop.}
\scope{Energy \emph{exchanges}: interactions, injections, emissions and absorptions. Whether a freely propagating quantum must itself sit on the grid of \cref{eq:quantization} during its flight is not decided here; it is listed as an open question and revisited in Chapter~\ref{ch:lattice}.}
\end{hypothesis}

\begin{hypothesis}[Frame-rate oscillator, P]\label{as:p}
There exists a reference oscillator whose angular frequency equals the step rate:
\begin{equation}\label{eq:framerate}
\omega_{0} = \frac{1}{t_{\mathrm P}}.
\end{equation}
\motivation{The hypothesis layer needs one quantum system whose zero-point energy calibrates the frame scale. With $\omega_{0}$ as in \cref{eq:framerate}, the oscillator floor of \cref{thm:spectrum} becomes $E_{0} = \alpha_{0}\omega_{0} = \varepsilon_{\mathrm P}/2$ (\cref{rem:zeropoint-cell}): a single cell of \cref{eq:quantization} at the smallest $k$. The choice between $\omega_{0} = 1/t_{\mathrm P}$ and $\omega_{0} = 2\pi/t_{\mathrm P}$ is not free --- they differ by $2\pi$ and assign different zero-point cells; only the first is consistent with the injection theorem (Chapter~\ref{ch:unitarity}), as flagged in \cref{rem:framerate}.}
\scope{The zero-point and propagation chapters; every numerical factor of $\pi$ traced in \cref{rem:twopi,rem:framerate}.}
\end{hypothesis}

\begin{remark}[The zero-point cell]\label{rem:zeropoint-cell}
Combining \cref{as:p} with the oscillator floor and the Planck relations, in steps:
\begin{equation}\label{eq:zeropointcell}
E_{0} = \alpha_{0}\,\omega_{0} = \frac{\hbar}{2}\cdot\frac{1}{t_{\mathrm P}} = \frac{\hbar}{2 t_{\mathrm P}} = \frac{\varepsilon_{\mathrm P}}{2},
\end{equation}
where the last equality is \cref{cor:epmp}. The zero-point energy of the reference oscillator is exactly the $k = 1$ quantum of \cref{eq:quantization}. The same number reappears by a different route --- the per-frame injection of Chapter~\ref{ch:unitarity} --- and the agreement is recorded there as a consistency check, with the shared premise stated openly: both routes are expressed in the units of Section~\ref{sec:units}, and independence is claimed only for the derivation chains, not for the unit system.
\end{remark}

\begin{remark}[Status of the input layer]\label{rem:layer-scope}
The input layer of this paper is now closed, and it is small:
\begin{center}
\begin{tabular}{@{}p{3.6cm}p{4.6cm}p{3.2cm}@{}}
\toprule
Layer & Content & Status \\
\midrule
Empirical inputs (\cref{sec:units,sec:quantum}) & $c$; $\hbar$; $G$; $[\hat x,\hat p]=i\hbar$; $E=\hbar\omega$ & measured or standard quantum kinematics \\
Hypotheses (\cref{sec:hypotheses}) & A1 discrete time; A2 action quantization; P frame rate & working hypotheses \\
\bottomrule
\end{tabular}
\end{center}
Everything later is derivation from these rows, plus explicitly cited external results. Two disciplines are stated once and enforced throughout. (i) \emph{Status discipline}: no statement derived under A1--P may be quoted as evidence for A1--P. (ii) \emph{Scope discipline}: \cref{as:a2} governs exchanges, not free propagation; mechanism questions --- what the source of the per-frame injection is, why the reference oscillator ticks at the frame rate --- are listed as open problems and are discussed, if at all, in remarks headed \emph{Reading}, never as claims about mechanism.
\end{remark}

\section{Odd quantum numbers for bound states}\label{sec:oddquantum}

One consequence of the hypothesis layer can be recorded immediately. It is small, but it is used repeatedly later: bound spectra quantize only on the odd sites of the action axis.

\begin{theorem}[Odd quantum numbers]\label{thm:oddquantum}
Let $E_{n} = \hbar\omega_{0}\bigl( n + \tfrac{1}{2} \bigr)$, $n = 0, 1, 2, \ldots$, be the levels of the reference oscillator of \cref{as:p}. Then
\begin{equation}\label{eq:oddquantum}
E_{n}\, t_{\mathrm P} = (2n+1)\,\alpha_{0},
\end{equation}
so every bound level carries an odd quantum number $k = 2n+1 \in 2\mathbb{Z}^{+}+1$, and no bound level realizes an even multiple of $\alpha_{0}$.
\uses{\cref{as:p}, \cref{thm:spectrum}, \cref{def:alpha0}, \cref{as:a2} (the reading of $k$)}
\end{theorem}

\begin{proof}
Substitute the spectrum and the frame rate, one step at a time:
\begin{equation}\label{eq:oddsubstitute}
E_{n}\, t_{\mathrm P} = \hbar\,\omega_{0}\Bigl( n + \tfrac{1}{2} \Bigr) t_{\mathrm P} = \hbar\Bigl( n + \tfrac{1}{2} \Bigr) = 2\alpha_{0}\Bigl( n + \tfrac{1}{2} \Bigr) = (2n+1)\,\alpha_{0},
\end{equation}
where the second equality uses $\omega_{0} t_{\mathrm P} = 1$ (\cref{eq:framerate}) and the third uses $\hbar = 2\alpha_{0}$ (\cref{eq:alpha0}). The integer $2n+1$ is odd for every $n \geq 0$, and an even multiple $2m\,\alpha_{0}$ would require $2n + 1 = 2m$, i.e.\ $n = m - \tfrac12 \notin \mathbb{Z}$, which no level satisfies.
\end{proof}

\begin{remark}[Reading of the odd structure]\label{rem:oddreading}
The structure has a single source. The zero-point floor $E_{0} = \alpha_{0}\omega_{0} = \varepsilon_{\mathrm P}/2$ (\cref{rem:zeropoint-cell}) is the half-quantum contribution that \cref{rem:half} traced to the commutation relation itself; every excited level adds full quanta of $2\alpha_{0}$ on top of that floor. On the action axis, floor plus quanta lands exclusively on odd sites: the oddness is the zero-point contribution made visible. Two consistency checks are immediate. Each $k = 2n+1$ lies in $\mathbb{Z}^{+}$, so the oscillator levels are admissible under the quantization rule \cref{eq:quantization}; and the smallest of them, $n = 0$, $k = 1$, is the zero-point cell of \cref{rem:zeropoint-cell}.
\end{remark}

\begin{remark}[Open question: the even sites]\label{rem:evensites}
\Cref{thm:oddquantum} says that even multiples of $\alpha_{0}$ are not populated by bound-state quantization. It does not say what, if anything, populates them. Whether processes without a binding potential --- scattering processes, the boundary injection of Chapter~\ref{ch:unitarity} --- realize even quantum numbers is an open problem, to be decided where those processes are defined. No mechanism is asserted in either direction.
\end{remark}

\section{The measurement argument}\label{sec:measurement}

\cref{as:a1} posits that nature registers processes on a discrete sequence whose step size is $t_{\mathrm P}$. The \emph{structure} of that claim is a hypothesis; its \emph{magnitude} can be supported independently, by an argument of Wigner's type that uses only the empirical inputs of \cref{sec:units} and the results cited below --- in particular, not A1 itself, since support for A1 must not presuppose A1. The argument is a heuristic in the tradition of Wigner (1957)\footnote{E.~P. Wigner, \emph{Relativistic invariance and quantum phenomena}, Rev.\ Mod.\ Phys.\ \textbf{29}, 255 (1957).} and Salecker--Wigner (1958);\footnote{H.~Salecker and E.~P. Wigner, \emph{Quantum limitations of the measurement of space-time distances}, Phys.\ Rev.\ \textbf{109}, 571 (1958).} it yields the \emph{scale} and the \emph{form of the inequality}, and it does not yield a proof that nature ticks. The coefficient of the inequality moves with measurement conventions; the combination $\hbar G / c^{5}$ does not.

\begin{eresult}[Energy--time uncertainty, operational form]\label{er:etuncertainty}
To resolve a process whose state changes observably within a time $\tau$, the participating states must carry an energy width
\begin{equation}\label{eq:etuncertainty}
\Delta E \;\geq\; \frac{\hbar}{2\tau}.
\end{equation}
\evidence{Mandelstam and Tamm (1945),\footnote{L.~Mandelstam and I.~Tamm, \emph{The uncertainty relation between energy and time in non-relativistic quantum mechanics}, J.\ Phys.\ (USSR) \textbf{9}, 249 (1945); the relation follows from \cref{er:ccr} for states evolving under a time-independent Hamiltonian, with $\tau$ the characteristic evolution time.} adopted here as an operational criterion for measurement resolution rather than as a statement about self-adjoint time operators.}
\scope{The measurement-resolution statements of this section.}
\end{eresult}

\begin{eresult}[Mass--energy equivalence and Schwarzschild radius]\label{er:gr}
\begin{equation}\label{eq:massenergy}
E = M c^{2},
\qquad
r_{\mathrm s}(M) = \frac{2 G M}{c^{2}}.
\end{equation}
\reference{Einstein (1905); Schwarzschild (1916). Conditions: the second formula is the static, spherically symmetric exterior solution of general relativity, used here at the level at which it has been tested (weak field, slow motion).}
\end{eresult}

\begin{lemma}[Spatial extent of a resolving apparatus]\label{lem:apparatussize}
An apparatus that resolves a time $\tau$ has spatial extent
\begin{equation}\label{eq:apparatussize}
L \;\leq\; c\,\tau.
\end{equation}
\uses{\cref{er:light}}
\end{lemma}

\begin{proof}
A measurement is a causal process that must complete within its own resolving time: at least one signal has to cross the apparatus --- prepare a subsystem here, read the effect there --- and the fastest possible crossing of an extent $L$ takes $L/c$ (\cref{er:light}). Requiring $L/c \leq \tau$ gives \cref{eq:apparatussize}.
\end{proof}

\begin{lemma}[Gravitational readability condition]\label{lem:readability}
An apparatus of mass $M$ and extent $L$ can deliver a readout to its exterior only if
\begin{equation}\label{eq:readability}
L \;\geq\; r_{\mathrm s}(M).
\end{equation}
\uses{\cref{er:gr} (Schwarzschild radius), \cref{er:light} (causal structure)}
\end{lemma}

\begin{proof}
Suppose $L < r_{\mathrm s}(M)$. By the standard collapse behaviour of general relativity, the worldlines making up the apparatus then terminate at the singularity within a time of order $L/c$,\footnote{This is the one step of the section that relies on the interior Schwarzschild geometry rather than on measured weak-field physics; the argument is of Wigner type and the support it gives is correspondingly soft. The dependence is flagged rather than hidden.} while a readout, by \cref{lem:apparatussize}, needs at least one completed signal crossing. The collapse timescale is shorter than the readout timescale, so no readout reaches the exterior; contraposition gives \cref{eq:readability}.
\end{proof}

\begin{lemma}[Width of the energy support]\label{lem:spectralwidth}
If $\Delta E$ is the standard deviation of the energy of a state whose spectral measure is supported on an interval of width $W$, then
\begin{equation}\label{eq:spectralwidth}
W \;\geq\; 2\,\Delta E.
\end{equation}
\end{lemma}

\begin{proof}
Let the support be $[\lambda_{\min}, \lambda_{\max}]$, $W = \lambda_{\max} - \lambda_{\min}$, and let $\mu$ be the mean. Deviations from the mean cannot exceed the distances to the endpoints in their respective directions, so
\begin{equation}\label{eq:bhattacharya}
\mathrm{Var}(E) = \bigl\langle (E - \mu)^{2} \bigr\rangle \;\leq\; (\lambda_{\max} - \mu)(\mu - \lambda_{\min}),
\end{equation}
the product of the two one-sided spans. Their sum is fixed, $W = (\lambda_{\max} - \mu) + (\mu - \lambda_{\min})$, and two non-negative numbers with fixed sum have product at most $(\mathrm{sum}/2)^{2}$ (equivalently, $(u - v)^{2} \geq 0$ rearranges to $4uv \leq (u+v)^{2}$). Hence
\begin{equation}\label{eq:amgm}
(\Delta E)^{2} \;\leq\; \frac{W^{2}}{4}
\quad\Longrightarrow\quad
W \geq 2\,\Delta E.
\end{equation}
\end{proof}

\begin{proposition}[Scale of minimal time resolution]\label{prop:tau}
No measurement scheme built from the inputs of \cref{sec:units} resolves times below
\begin{equation}\label{eq:tau}
\tau \;\geq\; t_{\mathrm P}.
\end{equation}
\uses{\cref{er:etuncertainty}, \cref{er:gr}, \cref{lem:apparatussize}, \cref{lem:readability}, \cref{lem:spectralwidth}, \cref{def:units}. Explicitly not used: \cref{as:a1,as:a2,as:p}.}
\end{proposition}

\begin{proof}
Resolve a time $\tau$. \emph{Step 1 (energy width).} By \cref{er:etuncertainty},
\begin{equation}\label{eq:mstep1}
\Delta E \;\geq\; \frac{\hbar}{2\tau}.
\end{equation}
\emph{Step 2 (mass of the apparatus).} By \cref{lem:spectralwidth}, realizing the width \cref{eq:mstep1} requires an energy scale of at least $\Delta E$, hence a mass
\begin{equation}\label{eq:mstep2}
M \;=\; \frac{E_{\mathrm{app}}}{c^{2}} \;\geq\; \frac{\Delta E}{c^{2}} \;\geq\; \frac{\hbar}{2 \tau c^{2}},
\end{equation}
using \cref{eq:massenergy} for the first equality.
\emph{Step 3 (gravitational radius of the apparatus).} Substitute \cref{eq:mstep2} into the Schwarzschild radius and let the factor $2$ cancel the factor $\tfrac12$ explicitly:
\begin{equation}\label{eq:mstep3}
r_{\mathrm s}(M) = \frac{2 G M}{c^{2}} \;\geq\; \frac{2 G}{c^{2}}\cdot\frac{\hbar}{2 \tau c^{2}} \;=\; \frac{G \hbar}{\tau c^{4}}.
\end{equation}
\emph{Step 4 (window for the apparatus).} Chaining \cref{lem:readability} with \cref{lem:apparatussize},
\begin{equation}\label{eq:mstep4}
\frac{G \hbar}{\tau c^{4}} \;\leq\; r_{\mathrm s}(M) \;\leq\; L \;\leq\; c\,\tau,
\end{equation}
so the lower bound of \cref{eq:mstep3} must not exceed the upper bound of \cref{eq:apparatussize}.
\emph{Step 5 (solve for $\tau$).} Multiply \cref{eq:mstep4} through by the positive quantity $\tau c^{4}$:
\begin{equation}\label{eq:mstep5}
G \hbar \;\leq\; \tau^{2} c^{5}
\quad\Longleftrightarrow\quad
\tau^{2} \;\geq\; \frac{\hbar G}{c^{5}} \;=\; t_{\mathrm P}^{2},
\end{equation}
the last equality being the definition of $t_{\mathrm P}$ (\cref{def:units}). Both sides are positive, so the root preserves the inequality: $\tau \geq t_{\mathrm P}$.
\end{proof}

\begin{remark}[What the argument does and does not establish]\label{rem:license}
Three statements close the argument. (i) \emph{What is supported:} the magnitude. A discrete time axis with step size $t_{\mathrm P}$ is consistent with measurement, because any attempt to resolve a smaller time founders on the apparatus's own gravitational radius. (ii) \emph{What is not established:} the structure. \cref{prop:tau} neither derives discreteness nor selects $t_{\mathrm P}$ over any larger step; the choice of the sequence remains \cref{as:a1}. (iii) \emph{Coefficient mobility:} the factor $\tfrac12$ in \cref{eq:etuncertainty} and the support-width factor $2$ in \cref{lem:spectralwidth} both move with convention (standard deviation versus half-width, characteristic time versus resolving time); what is convention-independent is the combination $\hbar G / c^{5}$ itself. Every convention-dependent factor is displayed above rather than absorbed, so the reader can see exactly where each one enters.
\end{remark}

\begin{remark}[Genealogy]\label{rem:genealogy}
Within the heuristic family to which this argument belongs, the earliest node is de Broglie's matter-wave relation $p = h/\lambda$ (1923):\footnote{L.~de Broglie, \emph{Ondes et quanta}, C.\ R.\ Acad.\ Sci.\ \textbf{177}, 507 (1923).} the energy--frequency relation \cref{er:efreq} and the wave kinematics $\omega = ck$ combine as its descendants. The genealogy is recorded for orientation; it plays no logical role below.
\end{remark}

\begin{corollary}[Frames are counted from $n = 1$]\label{cor:counting}
If time steps exist at the supported scale, the frame index of \cref{eq:frames} runs from $n = 1$: the attempt to register $n = 0$ is the attempt to resolve $\tau < t_{\mathrm P}$ itself.
\uses{\cref{prop:tau}, \cref{as:a1}}
\end{corollary}

\section{Three kinds of statements}\label{sec:interpretation}

The foundation is now closed, and it is worth stating once, at the point where it closes, what kind of statement the physical prose of this paper is allowed to make.

\begin{remark}[Formal statements, operational prose, readings]\label{rem:register}
This paper uses three kinds of statements, and they do not mix. \emph{Formal statements} (hypotheses, empirical inputs, definitions, lemmas, propositions, theorems, corollaries) carry the logical content and are checked by their proofs or citations. \emph{Operational prose} carries the physics: what a formula says about measurement, what an inequality excludes, what a scale means. \emph{Readings} --- the remarks headed \emph{Reading} --- state the most useful physical picture of a result \emph{under the hypotheses}; they are interpretations, not mechanism claims. The standing analogy is the equivalence principle: ``inertia is gravity'' is two descriptions of one acceleration, recorded because it is the most useful form of the principle; it is not a mechanism, and was never read as one. Every reading in this paper is bounded in the same way: it says what a result is most usefully \emph{like}, never what it is \emph{made of}.
\end{remark}

\begin{remark}[Enforcement; first list of open problems]\label{rem:mechanismregistry}
The rule of \cref{rem:register} is enforced by placement: interpretive content appears only in remarks, never inside hypotheses and never as a step of a proof. At the end of this chapter, the following mechanism questions are open, in the sense that this paper makes no claim about them in either direction:
\begin{enumerate}
\item population of the even action sites (\cref{rem:evensites});
\item whether a freely propagating quantum sits on the quantization grid \cref{eq:quantization} during its flight (scope of \cref{as:a2});
\item the internal composition of whatever drives the per-frame injection of Chapter~\ref{ch:unitarity} (no mechanism claim about a source is made anywhere in this paper);
\item the reason the reference oscillator of \cref{as:p} ticks at the frame rate (the hypothesis records the convention; the check that fixes it comes later, and nothing deeper is claimed).
\end{enumerate}
The list is cumulative: Chapter~\ref{ch:conclusion} collects all entries, and every entry is a question of mechanism, not of derivation.
\end{remark}

\begin{remark}[What Chapter 1 has established]\label{rem:ch1summary}
The input layer of the paper is exactly the table of \cref{rem:layer-scope}: five empirical inputs, three hypotheses, and nothing else. Upon it stand: the Planck scale and its uniqueness (\cref{sec:units}); the quantum of action and the oscillator floor, from which the half-unit $\alpha_{0}$ arises (\cref{sec:quantum}); the hypothesis layer A1--P, whose magnitudes and structures are supported separately (\cref{sec:hypotheses}); the odd quantum numbers of bound spectra (\cref{sec:oddquantum}); and the independent measurement argument for the step size (\cref{sec:measurement}). Every premise used in arriving at these results appears earlier in the chapter than the result itself, and the two uses that would have violated this order --- deriving \cref{er:efreq} from below, or supporting \cref{as:a1} from itself --- are named and refused where they would have occurred. The rest of the paper builds on this table alone.
\end{remark}

\chapter{Light at the Planck scale}\label{sec:light}\label{ch:light}

This chapter locates the ceiling above which a pair of notions --- a \emph{localized entity of energy} $E$, and a \emph{wave of wavelength} $\lambda$ --- can no longer both apply. The derivation is short algebra standing on two empirical inputs and one external result; its output, a threshold wavelength and a threshold energy, is a boundary of the description, not a new dynamics: the chapter says where a way of speaking fails, not what replaces it. Mechanism questions are handled under the rule of \cref{rem:register}: nothing below asserts what a wave becomes past the threshold.

\section{Inputs and conventions}\label{sec:light-inputs}

Three inputs enter. The first is quantum kinematics already recorded as an empirical input:

\begin{center}
\begin{tabular}{@{}lll@{}}
\toprule
Input & Statement & Status \\
\midrule
energy--frequency relation & $E = \hbar\omega$ & \cref{er:efreq} \\
dispersion of massless waves & $\omega = c k$, \quad $k = 2\pi/\lambda$ & \cref{er:dispersion} \\
gravitational radius & $r_{\mathrm s} = 2 G E / c^{4}$ & \cref{er:gr} \\
\bottomrule
\end{tabular}
\end{center}

The dispersion relation is not assumed ad hoc:

\begin{eresult}[Dispersion of massless waves]\label{er:dispersion}
Massless excitations in vacuum satisfy
\begin{equation}\label{eq:dispersion}
\omega = c\, k, \qquad k = \frac{2\pi}{\lambda}.
\end{equation}
\reference{Standard wave kinematics. The relation is rederived independently in Chapter~\ref{ch:lattice}, where it holds exactly throughout the Brillouin zone; the two derivations meet as a consistency check, and nothing in the present chapter uses the lattice result.}
\end{eresult}

The third input, the Schwarzschild radius in the form $r_{\mathrm s} = 2GE/c^{4}$ (the mass--energy form of \cref{eq:massenergy}), is quoted from \cref{er:gr} and is used exactly at the level at which that result was recorded: the field of a localized energy, treated non-relativistically.

\begin{remark}[The $2\pi$ convention, first working instance]\label{rem:twopi-light}
This chapter keeps the standard convention $k = 2\pi/\lambda$ of \cref{eq:dispersion}. An equivalent convention absorbs $2\pi$ into the constant, i.e.\ writes $\omega \equiv c/\lambda$; under it, the numerical coefficients of \cref{sec:light-threshold} move, and the next section computes both variants side by side so that the movement is visible term by term. What no convention can move is the \emph{order of magnitude}: both variants place the threshold within one power of $\ell_{\mathrm P}$ and of $\varepsilon_{\mathrm P}$, and it is the order of magnitude, not the coefficient, that the physics depends on. This remark is the first working instance of the tracking promised in \cref{rem:twopi}: the factor $2\pi$ enters \emph{here}, at the convention, and every factor of $\pi$ below is its descendant.
\end{remark}

With the inputs recorded, the next section performs the derivation in three displayed steps, tracking each cancellation.

\section{The collapse threshold}\label{sec:light-threshold}

The physical question is: how far can the pair of notions ``localized entity'' and ``wave'' be pushed before they exclude each other? Resolution finer than $\lambda$ requires energy of order $\hbar c/\lambda$; energy gravitates; and once the energy needed to resolve $\lambda$ carries a gravitational radius larger than $\lambda$ itself, nothing can both be localized to $\lambda$ and carry that energy. The scale at which this happens is computable in three steps from the recorded inputs.

\begin{proposition}[Collapse threshold]\label{prop:threshold}
The condition $r_{\mathrm s} = \lambda$, with $E$ and $\lambda$ related by the inputs of \cref{sec:light-inputs}, determines a unique threshold
\begin{equation}\label{eq:threshold}
\lambda_{\mathrm c} = 2\sqrt{\pi}\; \ell_{\mathrm P},
\qquad
E_{\mathrm c} = \sqrt{\pi}\; \varepsilon_{\mathrm P},
\end{equation}
separating $\lambda > \lambda_{\mathrm c}$ ($r_{\mathrm s} < \lambda$) from $\lambda < \lambda_{\mathrm c}$ ($r_{\mathrm s} > \lambda$).
\uses{\cref{er:efreq}, \cref{er:dispersion}, \cref{er:gr}, \cref{def:units}, \cref{cor:epmp}}
\end{proposition}

\begin{proof}
\emph{Step 1 (energy at wavelength).} Chain the inputs, one substitution at a time:
\begin{equation}\label{eq:stepenergy}
E = \hbar \omega \;\xrightarrow{\ \cref{eq:dispersion}\ }\; \hbar c k \;\xrightarrow{\ k = 2\pi/\lambda\ }\; E(\lambda) = \frac{2\pi \hbar c}{\lambda}.
\end{equation}
The factor $2\pi$ is the convention factor of \cref{rem:twopi-light}, entering through $k$.

\emph{Step 2 (gravitational radius, and the threshold equation).} Substitute \cref{eq:stepenergy} into the Schwarzschild radius of \cref{er:gr} and track the powers of $c$ explicitly:
\begin{equation}\label{eq:stepradius}
r_{\mathrm s} = \frac{2 G}{c^{4}}\, E(\lambda) = \frac{2 G}{c^{4}} \cdot \frac{2\pi \hbar c}{\lambda} = \frac{4 \pi \hbar G}{c^{3} \lambda},
\end{equation}
where $c / c^{4} = c^{-3}$ is the only cancellation. Setting $r_{\mathrm s} = \lambda$:
\begin{equation}\label{eq:stepthreshold}
\frac{4 \pi \hbar G}{c^{3} \lambda} = \lambda
\;\;\Longrightarrow\;\;
\lambda^{2} = \frac{4 \pi \hbar G}{c^{3}} = 4\pi\, \ell_{\mathrm P}^{2},
\end{equation}
using $\ell_{\mathrm P}^{2} = \hbar G / c^{3}$ (\cref{prop:triad}). Both sides are positive, so the root is unambiguous:
\begin{equation}\label{eq:lamc}
\lambda_{\mathrm c} = \sqrt{4\pi}\; \ell_{\mathrm P} = 2\sqrt{\pi}\; \ell_{\mathrm P}.
\end{equation}

\emph{Step 3 (threshold energy).} Evaluate \cref{eq:stepenergy} at $\lambda_{\mathrm c}$, and simplify the numerical factor separately from the units:
\begin{equation}\label{eq:estep}
E_{\mathrm c} = \frac{2\pi \hbar c}{2\sqrt{\pi}\, \ell_{\mathrm P}} = \underbrace{\frac{2\pi}{2\sqrt{\pi}}}_{=\;\sqrt{\pi}}\; \frac{\hbar c}{\ell_{\mathrm P}} = \sqrt{\pi}\; \frac{\hbar}{t_{\mathrm P}} = \sqrt{\pi}\; \varepsilon_{\mathrm P},
\end{equation}
where the unit conversion $\hbar c / \ell_{\mathrm P} = \hbar / t_{\mathrm P}$ is the first Planck relation $\ell_{\mathrm P} = c\, t_{\mathrm P}$ of \cref{def:units}, and the last equality is \cref{cor:epmp}. The arithmetic factor alone: $2\pi / 2\sqrt{\pi} = \pi/\sqrt{\pi} = \sqrt{\pi}$.

\emph{Monotonicity.} $E(\lambda)$ is strictly decreasing in $\lambda$ (\cref{eq:stepenergy}), $r_{\mathrm s}$ is strictly increasing in $E$ (\cref{eq:massenergy}), so $r_{\mathrm s}(E(\lambda)) = 4\pi\hbar G / c^{3}\lambda$ is strictly decreasing in $\lambda$; equality holds exactly at $\lambda_{\mathrm c}$, which is therefore the unique crossing.
\end{proof}

\begin{remark}[Where the factors of $\pi$ come from]\label{rem:sqrtpi}
The coefficients of \cref{eq:threshold} are descendants of the single convention factor $2\pi$ of \cref{rem:twopi-light}: the threshold equation squares the $2\pi$ of \cref{eq:stepenergy}, giving $\sqrt{4\pi} = 2\sqrt{\pi}$ in $\lambda_{\mathrm c}$, and the surviving $\sqrt{\pi}$ in $E_{\mathrm c}$ is that same factor after halving. No $\pi$ enters from \cref{sec:units} --- the Planck relations are $\pi$-free (\cref{rem:twopi}) --- and none enters from the gravitational input, which is linear in $E$. Every $\pi$ of this chapter is the convention's descendant, and if the convention is changed they all move together, as the next remark verifies.
\end{remark}

\begin{remark}[Convention dependence of the coefficients]\label{rem:convention-move}
Under the alternative convention $\omega \equiv c/\lambda$ of \cref{rem:twopi-light}, the same three steps give
\begin{equation}\label{eq:altconvention}
E(\lambda) = \frac{\hbar c}{\lambda}
\;\;\Longrightarrow\;\;
r_{\mathrm s} = \frac{2\hbar G}{c^{3}\lambda}
\;\;\Longrightarrow\;\;
\lambda_{\mathrm c}' = \sqrt{2}\; \ell_{\mathrm P},\qquad
E_{\mathrm c}' = \frac{\varepsilon_{\mathrm P}}{\sqrt{2}}.
\end{equation}
The coefficients moved ($2\sqrt{\pi} \leftrightarrow \sqrt{2}$ in length, $\sqrt{\pi} \leftrightarrow 1/\sqrt{2}$ in energy); what did not move is the statement that the threshold lies within one power of $\ell_{\mathrm P}$ and of $\varepsilon_{\mathrm P}$. The threshold \emph{coefficient} is not an invariant of the description; its \emph{order of magnitude} is. All quantitative claims below are made at the level of the invariant.
\end{remark}

\begin{remark}[Physical content]\label{rem:threshold-content}
Below $\lambda_{\mathrm c}$, the energy required for resolution carries $r_{\mathrm s} > \lambda$: the energy that was meant to probe the region surrounds it. The two descriptions of one and the same physical situation --- ``a localized entity of energy $E$'' and ``a wave of wavelength $\lambda$'' --- exclude each other there. The threshold is therefore a boundary of this pair of descriptions, not a dynamical transition: the derivation itself used the continuous inputs \cref{er:efreq,er:dispersion,er:gr}, and its conclusion marks the scale at which that combined way of speaking can no longer be applied to itself. The same boundary appears independently in the measurement argument: \cref{lem:readability} demands $r_{\mathrm s} \leq L$ for any readable apparatus, and \cref{lem:apparatussize} demands $L \leq c\tau$; at the resolution of a wavelength $\lambda$ the two bounds meet at exactly the condition $r_{\mathrm s} = \lambda$ of \cref{prop:threshold}. One boundary, two derivations; the proof of \cref{prop:tau} did not use this chapter, and this chapter did not use \cref{prop:tau}, so the two are independent as far as this section is concerned.
\end{remark}

\begin{remark}[Commensurability with the step size]\label{rem:commensurate}
The ratio $\lambda_{\mathrm c} / \ell_{\mathrm P} = 2\sqrt{\pi}$ is irrational: no integer number of step sizes equals the threshold, and the one-step wavelength $\lambda = \ell_{\mathrm P}$ lies strictly inside the collapse region ($\ell_{\mathrm P} < \lambda_{\mathrm c}$, since $1 < 2\sqrt{\pi}$). The first statement matters later, where the lattice chapter derives a second wall (the Brillouin boundary) and the ordering of the two walls is compared; the second records, in coefficient-free form, that a single-step wave sits deep inside the collapse region. Both statements are recorded here in the invariant, coefficient-free form.
\end{remark}

\section{The threshold against the quantization rule}\label{sec:light-tension}

The threshold energy was derived from continuous concepts, while \cref{as:a2} quantizes energy on the grid $E = \frac{k}{2}\,\varepsilon_{\mathrm P}$, $k \in \mathbb{Z}^{+}$. The two must be compared explicitly.

\begin{remark}[The tension, stated exactly]\label{rem:tension}
If the threshold energy is forced onto the grid of \cref{eq:quantization}, it would require
\begin{equation}\label{eq:tension}
E_{\mathrm c} = \sqrt{\pi}\; \varepsilon_{\mathrm P} \;\overset{?}{=}\; \frac{k}{2}\,\varepsilon_{\mathrm P}
\quad\Longleftrightarrow\quad
k = 2\sqrt{\pi},
\end{equation}
and $2\sqrt{\pi}$ is irrational --- no integer $k$ realizes the threshold. The two statements are not in contradiction, because they do not speak about the same kind of object: \cref{eq:quantization} quantizes \emph{exchanges} (scope of \cref{as:a2}), while \cref{prop:threshold} states where a \emph{pair of descriptions} fails. The threshold inherits the validity domain of its own inputs (\cref{er:efreq,er:dispersion,er:gr}), all of them continuous; it marks the boundary of that continuous way of speaking, and it is not claimed to be itself a state of the grid. In the taxonomy of \cref{sec:oddquantum}, the threshold belongs to the same open problem as the even action sites (\cref{rem:evensites}): a non-grid point whose population, if any, this paper does not describe. The question is stated, not answered, in both places.
\end{remark}

\begin{remark}[Jurisdiction]\label{rem:jurisdiction}
``What does a wave become past the threshold?'' is a mechanism question, and the rule of \cref{rem:register} applies verbatim: this paper asserts no mechanism in either direction. What is asserted is weaker and exact: below $\lambda_{\mathrm c}$ the descriptions ``localized entity of energy $E$'' and ``wave of wavelength $\lambda$'' exclude each other (\cref{rem:threshold-content}); any replacement description is outside the competence of this chapter, and the chapter's only output is the boundary \cref{eq:threshold} itself.
\end{remark}

\begin{remark}[Convergence of independent routes]\label{rem:convergence}
The scale $\ell_{\mathrm P}$ has now been reached by three routes that do not share premises: (i) entity confinement, \cref{prop:threshold}; (ii) readability of measurement, \cref{prop:tau} together with \cref{rem:threshold-content}; (iii) commensurability with the step size, \cref{rem:commensurate}, whose second wall is derived in \cref{ch:lattice}. The convergence is recorded as a \emph{robustness statement} about the scale: three independent derivations of one scale are stronger than one. No mechanism is inferred from the convergence --- the three routes may rest on different physics behind the same scale, and this paper records the convergence without drawing that inference.
\end{remark}

\begin{remark}[What Chapter 2 has established]\label{rem:ch2summary}
A unique threshold pair $(\lambda_{\mathrm c}, E_{\mathrm c}) = (2\sqrt{\pi}\,\ell_{\mathrm P},\; \sqrt{\pi}\,\varepsilon_{\mathrm P})$, derived in three steps from two empirical inputs and one external result; its coefficients are convention-dependent, its order of magnitude is not (\cref{rem:convention-move}); it is a boundary of description, not a transition (\cref{rem:threshold-content}); its non-integrality on the quantization grid belongs to the same open problem as the even action sites (\cref{rem:tension}); and its scale is reached independently by two further routes (\cref{rem:convergence}). No dynamics was added; none was needed.
\end{remark}

\part{The open chain}

\chapter{The open three-component chain}\label{ch:cascade}

This chapter studies an open three-component system: energy enters from an exterior source, flows one way through three components, and is tracked frame by frame. Two kinds of content must be kept apart from the start, and they are marked as they occur. \emph{Structural content} --- hypotheses, exact solutions, uniqueness of the steady state --- is derived and proved like anything else in this paper. \emph{Interpretive content} --- the spatial picture of where the components live, the physical origin of the radiation component, the mechanism of the one-step delay --- was guided by the author's physical intuition at the time of construction and remains \textbf{speculative}; it is recorded in remarks labelled as such, and no derivation below depends on it.

\section{The hypotheses of the chain}\label{sec:cascade-axioms}

\begin{definition}[Content vector]\label{def:contentvector}
The energy content at frame $n$ is the vector
\begin{equation}\label{eq:contentvector}
\boldsymbol{\mathcal{E}}^{(n)} = \bigl( E_{\mathrm B}^{(n)},\, E_{\mathrm r}^{(n)},\, E_{\mathrm m}^{(n)} \bigr)^{\!\mathrm T} \in \mathbb{R}^{3},
\end{equation}
whose components are named the \emph{bulk} ($\mathrm B$), \emph{radiation} ($\mathrm r$), and \emph{matter} ($\mathrm m$) contents. The order of the names records the direction of energy flow (\cref{as:b1}); the spatial reading of the order --- a bulk region adjoining the exterior, a radiation layer, a matter core --- is speculative and carries no logical weight (\cref{rem:speculation}).
\end{definition}

\begin{hypothesis}[One-way conduction, B1]\label{as:b1}
Energy flows only along $\mathrm B \to \mathrm r \to \mathrm m$; there is no reverse channel at any frame.
\motivation{Thermodynamic in origin: the entropy quadratic in the content (\cref{ch:boundary}) admits no regular partition function, hence no equilibrium state, hence no channel along which a reverse flow could equilibrate. The root itself is a theorem of \cref{ch:boundary}, and B1 is stated here as a hypothesis because the chain must exist before the boundary theorem that explains it can be proved; the two meet as a consistency check, not as a shared premise.}
\end{hypothesis}

\begin{hypothesis}[Conductance share, B2]\label{as:b2}
At frame $n$, each component conducts downstream a fraction $1/n$ of its own content:
\begin{equation}\label{eq:share}
\text{conducted amount} = \frac{1}{n}\, E_{i}^{(n)},\qquad i \in \{\mathrm B, \mathrm r\}.
\end{equation}
\motivation{As a reading, $1/n$ is the share of a component that has resided for the expected time of $n$ frames; as a necessity, $1/n$ is the unique conductance law for which a non-degenerate linear steady state exists at all --- \cref{thm:forced-cascade} of \cref{sec:forced-cascade} proves this, and until then the law stands as a hypothesis.}
\end{hypothesis}

\begin{hypothesis}[Constant injection, B3]\label{as:b3}
At every frame, a constant amount $J$ is injected into the bulk component:
\begin{equation}\label{eq:injection}
\boldsymbol{s} = (J, 0, 0)^{\mathrm T}, \qquad J = \text{const} > 0,
\end{equation}
independent of the content $\boldsymbol{\mathcal{E}}^{(n)}$ and of $n$.
\motivation{Structure and magnitude are separated, as throughout this paper. The \emph{structure} --- constant, state-independent injection --- is the content of B3 and is used immediately. The \emph{magnitude} $J$ is not chosen here; it is fixed by the injection theorem of \cref{ch:unitarity} as $J = \varepsilon_{\mathrm P}/2$, and the fixing is one-way (the injection theorem uses B3's structure, never its value, so no circle forms).}
\end{hypothesis}

\begin{definition}[Chain equations and boundary condition]\label{def:cascade}
The evolution of the content vector is
\begin{equation}\label{eq:cascade}
\begin{aligned}
E_{\mathrm B}^{(n+1)} &= \Bigl(1 - \tfrac{1}{n}\Bigr) E_{\mathrm B}^{(n)} + J,\\
E_{\mathrm r}^{(n+1)} &= \Bigl(1 - \tfrac{1}{n}\Bigr) E_{\mathrm r}^{(n)} + \frac{1}{n} E_{\mathrm B}^{(n)},\\
E_{\mathrm m}^{(n+1)} &= E_{\mathrm m}^{(n)} + \frac{1}{n} E_{\mathrm r}^{(n)},
\end{aligned}
\end{equation}
with the boundary condition
\begin{equation}\label{eq:firstframe}
E_{\mathrm r}^{(1)} = E_{\mathrm m}^{(1)} = 0, \qquad E_{\mathrm B}^{(1)} = \frac{J}{2}.
\end{equation}
\end{definition}

The boundary condition is \emph{not} an independent hypothesis: it is the first application of the steady-state split rule of \cref{sec:forced-cascade} --- each passing component retains half of its inflow and transmits half --- evaluated at the first frame, where the retention factor $1 - \tfrac1n$ vanishes and the recursion cannot generate the initial content by itself.

\begin{remark}[The one-step delay, structural reading]\label{rem:delay}
In \cref{eq:cascade}, every inflow term draws on the \emph{previous} frame's content: $E_{\mathrm r}^{(n+1)}$ is fed by $E_{\mathrm B}^{(n)}$, never by $E_{\mathrm B}^{(n+1)}$. No response is instantaneous. The formal ground --- that no signal crosses a frame boundary faster than one step, i.e.\ that the principle $v \leq c$ is a property of the step structure --- is established in \cref{sec:steps} and is \emph{not} used here; the chain stands on \cref{as:b1,as:b2,as:b3} alone. A second reading of the same delay: radiation produced at frame $n$ appears as content only from frame $n+1$ on; at any single frame, only inherited contents are present, never the production in progress.
\end{remark}

\begin{remark}[Speculative readings of this chapter]\label{rem:speculation}
Recorded once, so that no remark below needs to repeat it. (i) \emph{Spatial picture}: bulk as the region adjoining the exterior boundary, radiation as a layer, matter as the core --- intuition about \emph{where} the components live; speculative. (ii) \emph{Origin of the radiation component}: whether anything radiates in the generalised sense intended here, or whether the radiation channel is closer in character to an Unruh-type excitation, is \textbf{unknown}; the Unruh resemblance is an unverified analogy, and this paper asserts neither it nor any other mechanism. (iii) \emph{The halving pattern} of the shares below was reached by physical intuition; what is proved below is only that the hypotheses force it and that it is the unique non-degenerate steady state. Wherever an original intuition is found to have stated something the derivation does not support, the derivation governs and the wording here supersedes it.
\end{remark}

\begin{remark}[Truncation caution]\label{rem:truncation}
The conductance law of \cref{as:b2} continues mathematically: nothing stops one from writing components $16, 32, \ldots$ steps down the halving sequence. Whether such a continuation corresponds to a physical object, is the shadow of a deeper structure, or is an artifact of the mathematical form, is \textbf{not known}, and this paper makes no claim in either direction. The chain \emph{of this paper} has three components; the matter component is terminal and retains everything it receives. The working rule is conservative rather than ontological: no formula below extends the chain beyond its stated content, and this discipline asserts nothing about what does or does not lie beyond it. The reader is cautioned once, at the hypothesis layer, so that no later formula is read as extending the chain by itself.
\end{remark}

\section{Exact solutions, verified line by line}\label{sec:exact-solutions}

The chain equations \cref{eq:cascade} are linear recurrences with the variable coefficient $1 - \tfrac1n$. The boundary condition \cref{eq:firstframe} fixes the first frame, and for $n \geq 1$ the map $\boldsymbol{\mathcal{E}}^{(n)} \mapsto \boldsymbol{\mathcal{E}}^{(n+1)}$ is affine with linear part of determinant $\bigl(1-\tfrac1n\bigr)^{2} > 0$ for $n \geq 2$ and prescribed value at $n = 1$; induction therefore gives a \emph{unique} solution, and the task reduces to identifying it. All formulas in this section are written in $J$; the magnitude of $J$ is fixed later (\cref{rem:jform}).

\begin{proposition}[Exact solution of the chain]\label{prop:exact-solutions}
The unique solution of \cref{eq:cascade} with boundary condition \cref{eq:firstframe} is
\begin{equation}\label{eq:solutions}
\begin{gathered}
E_{\mathrm B}^{(n)} = \frac{nJ}{2},\qquad
E_{\mathrm r}^{(1)} = 0,\;\; E_{\mathrm r}^{(n)} = \frac{nJ}{4}\ (n \geq 2),\\
E_{\mathrm m}^{(1)} = 0,\;\; E_{\mathrm m}^{(n)} = \frac{(n-2)J}{4}\ (n \geq 2).
\end{gathered}
\end{equation}
\uses{\cref{eq:cascade}, \cref{eq:firstframe}}
\end{proposition}

\begin{proof}
Direct substitution, with every arithmetic step displayed and the frame transitions treated separately.

\emph{Bulk component.} For general $n \geq 1$:
\begin{equation}\label{eq:checkB}
\textsc{RHS}_{\mathrm B}(n) = \Bigl(1 - \frac{1}{n}\Bigr)\frac{nJ}{2} + J = \frac{(n-1)J}{2} + \frac{2J}{2} = \frac{(n+1)J}{2},
\end{equation}
matching the formula at frame $n+1$. At $n = 1$ the retention factor vanishes and the boundary condition supplies $E_{\mathrm B}^{(1)} = J/2$, which the formula also gives ($1 \cdot J/2$).

\emph{Radiation component, transition $n = 1$.} The recursion gives
\begin{equation}\label{eq:checkR1}
E_{\mathrm r}^{(2)} = \Bigl(1 - \frac{1}{1}\Bigr) E_{\mathrm r}^{(1)} + \frac{E_{\mathrm B}^{(1)}}{1} = 0 \cdot 0 + \frac{J/2}{1} = \frac{J}{2} = \frac{2J}{4},
\end{equation}
matching the formula at $n = 2$. For $n \geq 2$:
\begin{equation}\label{eq:checkR2}
\textsc{RHS}_{\mathrm r}(n) = \Bigl(1 - \frac{1}{n}\Bigr)\frac{nJ}{4} + \frac{1}{n}\cdot\frac{nJ}{2} = \frac{(n-1)J}{4} + \frac{J}{2} = \frac{(n-1)J}{4} + \frac{2J}{4} = \frac{(n+1)J}{4},
\end{equation}
matching the formula at frame $n+1$.

\emph{Matter component, transitions $n = 1, 2$.} The recursion gives
\begin{equation}\label{eq:checkM1}
\begin{gathered}
E_{\mathrm m}^{(2)} = E_{\mathrm m}^{(1)} + \frac{E_{\mathrm r}^{(1)}}{1} = 0 + 0 = 0 = \frac{(2-2)J}{4},\\
E_{\mathrm m}^{(3)} = E_{\mathrm m}^{(2)} + \frac{E_{\mathrm r}^{(2)}}{2} = 0 + \frac{J/2}{2} = \frac{J}{4} = \frac{(3-2)J}{4},
\end{gathered}
\end{equation}
both matching the formula. For $n \geq 2$:
\begin{equation}\label{eq:checkM2}
\textsc{RHS}_{\mathrm m}(n) = \frac{(n-2)J}{4} + \frac{1}{n}\cdot\frac{nJ}{4} = \frac{(n-2)J}{4} + \frac{J}{4} = \frac{(n-1)J}{4},
\end{equation}
matching the formula at frame $n+1$. Every case checks; by uniqueness, \cref{eq:solutions} is the solution.
\end{proof}

\begin{remark}[$J$-form as discipline about the order of derivation]\label{rem:jform}
These solutions could equally be written with $\varepsilon_{\mathrm P}/4$ in place of $J/4$ --- the value $J = \varepsilon_{\mathrm P}/2$ being the content of the injection theorem (\cref{ch:unitarity}). Quoting that value here would use, as a premise, a statement whose own derivation uses \cref{as:b3}: a circle in the dependency order, of exactly the kind this paper refuses (\cref{rem:register}). The $J$-form of \cref{eq:solutions} is therefore the exact statement at this point of the text; the $\varepsilon_{\mathrm P}$-form is recovered once, in \cref{ch:unitarity}, when the value has been derived.
\end{remark}

\begin{proposition}[Total content and its increment]\label{prop:total}
The total content of the chain
\begin{equation}\label{eq:totaldef}
N(n) := E_{\mathrm B}^{(n)} + E_{\mathrm r}^{(n)} + E_{\mathrm m}^{(n)}
\end{equation}
equals
\begin{equation}\label{eq:total}
N(n) = \frac{(2n-1)\,J}{2},
\end{equation}
and grows by exactly the injected amount per frame:
\begin{equation}\label{eq:increment}
N(n+1) - N(n) = J \qquad \text{for all } n \geq 1.
\end{equation}
\uses{\cref{eq:cascade}, \cref{eq:solutions}}
\end{proposition}

\begin{proof}
\emph{Summation form of the increment.} Add the three equations of \cref{eq:cascade} frame by frame. The conducted terms cancel in pairs: $E_{\mathrm B}^{(n)}/n$ appears with a minus sign in the bulk equation and a plus sign in the radiation equation, and $E_{\mathrm r}^{(n)}/n$ with a minus in the radiation equation and a plus in the matter equation. Explicitly,
\begin{equation}\label{eq:telescope}
N(n+1) - N(n) = -\frac{E_{\mathrm B}^{(n)}}{n} + J + \frac{E_{\mathrm B}^{(n)}}{n} - \frac{E_{\mathrm r}^{(n)}}{n} + \frac{E_{\mathrm r}^{(n)}}{n} = J,
\end{equation}
each pair cancelling as written; only the injection survives. This proves \cref{eq:increment} directly from the hypotheses, without the solution.

\emph{Closed form.} Sum \cref{eq:increment} from frame $1$ to $n-1$ and add the boundary value $N(1) = E_{\mathrm B}^{(1)} = J/2$ (the other two components vanish there by \cref{eq:firstframe}):
\begin{equation}\label{eq:sumincrement}
N(n) = \frac{J}{2} + (n-1)J = \frac{J}{2} + \frac{(2n-2)J}{2} = \frac{(2n-1)J}{2}.
\end{equation}
Independently, substitution of \cref{eq:solutions} gives the same value:
\begin{equation}\label{eq:sumcheck}
\frac{nJ}{2} + \frac{nJ}{4} + \frac{(n-2)J}{4} = \frac{2nJ + nJ + (n-2)J}{4} = \frac{(4n-2)J}{4} = \frac{(2n-1)J}{2},
\end{equation}
the two routes agreeing as they must.
\end{proof}

\begin{remark}[Shares]\label{rem:shares}
Within the chain, the shares are
\begin{equation}\label{eq:shares}
E_{\mathrm B} : E_{\mathrm r} : E_{\mathrm m} = \frac{1}{2} : \frac{1}{4} : \frac{1}{4} \qquad (\text{exactly, for } n \geq 2),
\end{equation}
since $E_{\mathrm r}^{(n)} = nJ/4$ and $E_{\mathrm m}^{(n)} = (n-2)J/4 = nJ/4 - J/2$ differ from each other by the constant $J/2$, which is $O(1)$ against contents growing linearly in $n$; the statement is therefore exact at every frame from $n = 2$ up to a bounded, decaying relative correction. Against the injected total $nJ$ (the running sum of \cref{eq:injection}), the shares read $1 : \tfrac12 : \tfrac14 : \tfrac14$ asymptotically. The halving pattern $\tfrac12, \tfrac14, \tfrac18, \tfrac18$ is the same table normalised to the injected total with the residual included; both readings are recorded, and neither adds content.
\end{remark}

\begin{remark}[The residual, deferred]\label{rem:residual}
The total injected by frame $n$ is $nJ$; the chain holds $(2n-1)J/2$ of it. The difference,
\begin{equation}\label{eq:residual}
D(n) := nJ - N(n) = \frac{J}{2} \qquad \text{for all } n,
\end{equation}
is frame-independent. One number receives three readings later in this paper (\cref{sec:cascade-scope}): the per-frame transmitted amount, the first-frame retained amount, and this residual. The identification of $D$ with the energy content of the exterior boundary belongs to \cref{ch:boundary} and is not anticipated here; \cref{eq:residual} is kept as pure arithmetic.
\end{remark}

\begin{remark}[Symbol discipline and a mathematical aside]\label{rem:kernel-symbol}
Two notes. (i) The response kernel of \cref{sec:cascade-kernel} carries two frame arguments, source and target; its source argument is denoted $m$ throughout, not $k$, to keep the symbol $k$ reserved for the quantum number of \cref{eq:quantization} --- one symbol, one meaning, and renamings of this kind are recorded rather than silent. (ii) The product structure behind the kernel, $\prod (1 - \tfrac1j)$, is a telescoping product, and its moments generate series of the harmonic/$\zeta$ type; such series are pure mathematics and are kept out of the physical layer of this paper, per the truncation discipline of \cref{rem:truncation}.
\end{remark}

\section{The response kernel and stability}\label{sec:cascade-kernel}

The chain equations \cref{eq:cascade} are affine: their right-hand sides are linear in $\boldsymbol{\mathcal{E}}^{(n)}$ plus the constant injection. Deviations from the exact solution \cref{eq:solutions} therefore obey a closed linear system, and every structural fact of this section is a property of that linear system alone.

\begin{definition}[Deviation system and response kernel]\label{def:kernel}
Let $\boldsymbol{\delta}^{(n)}$ be a deviation from the exact solution. The \emph{deviation system} is the homogeneous part of \cref{eq:cascade}:
\begin{equation}\label{eq:deviation}
\begin{aligned}
\delta_{\mathrm B}^{(n+1)} &= \Bigl(1 - \tfrac{1}{n}\Bigr) \delta_{\mathrm B}^{(n)},\\
\delta_{\mathrm r}^{(n+1)} &= \Bigl(1 - \tfrac{1}{n}\Bigr) \delta_{\mathrm r}^{(n)} + \tfrac{1}{n}\,\delta_{\mathrm B}^{(n)},\\
\delta_{\mathrm m}^{(n+1)} &= \delta_{\mathrm m}^{(n)} + \tfrac{1}{n}\,\delta_{\mathrm r}^{(n)}.
\end{aligned}
\end{equation}
The \emph{response kernel} $\mathcal{W}(n; m)$ is the scalar multiplier carrying a bulk deviation from frame $m$ to frame $n \geq m$:
\begin{equation}\label{eq:kerneldef}
\mathcal{W}(n; m) := \prod_{j = m}^{n-1} \Bigl(1 - \frac{1}{j}\Bigr), \qquad \mathcal{W}(n; n) = 1 \ \text{(empty product)}.
\end{equation}
The source argument is denoted $m$; the symbol $k$ remains reserved for the quantum number of \cref{eq:quantization} (\cref{rem:kernel-symbol}).
\end{definition}

\begin{proposition}[Closed form of the kernel]\label{prop:kernel}
\begin{equation}\label{eq:kernel}
\mathcal{W}(n; m) = \frac{m-1}{n-1} \qquad (n \geq m \geq 1).
\end{equation}
\uses{\cref{eq:kerneldef}}
\end{proposition}

\begin{proof}
Write the factors in un-cancelled form and cancel step by step:
\begin{equation}\label{eq:kernel-telescope}
\mathcal{W}(n; m) = \frac{m-1}{m}\cdot\frac{m}{m+1}\cdot\frac{m+1}{m+2}\cdots\frac{n-3}{n-2}\cdot\frac{n-2}{n-1}.
\end{equation}
Each numerator cancels the preceding denominator, and only the first numerator and the last denominator survive: $\mathcal{W}(n; m) = (m-1)/(n-1)$. At $n = m$ the product is empty and the formula gives $(m-1)/(m-1) = 1$, as required by \cref{eq:kerneldef}.
\end{proof}

\begin{corollary}[Annihilation of first-frame deviations]\label{cor:annihilation}
$\mathcal{W}(n; 1) = 0$ for all $n \geq 2$: a deviation introduced at the first frame is erased from the bulk channel entirely.
\uses{\cref{prop:kernel} at $m = 1$}
\end{corollary}

The vanishing factor is the retention factor $1 - \tfrac1n$ at $n = 1$: the same zero that forced the boundary condition \cref{eq:firstframe} also protects everything downstream from whatever the first frame was seeded with. This is the chain's internal boundary layer (\cref{rem:delay}).

\begin{proposition}[Cocycle law]\label{prop:cocycle}
For $n \geq \ell \geq m$,
\begin{equation}\label{eq:cocycle}
\mathcal{W}(n; \ell)\,\mathcal{W}(\ell; m) = \mathcal{W}(n; m).
\end{equation}
\uses{\cref{prop:kernel}}
\end{proposition}

\begin{proof}
Substitute the closed form and cancel:
\begin{equation}\label{eq:cocycleproof}
\mathcal{W}(n; \ell)\,\mathcal{W}(\ell; m) = \frac{\ell - 1}{n - 1}\cdot\frac{m - 1}{\ell - 1} = \frac{m-1}{n-1} = \mathcal{W}(n; m).
\end{equation}
Equivalently, the product of \cref{eq:kerneldef} splits as $\prod_{j=\ell}^{n-1}\cdot\prod_{j=m}^{\ell-1}$ with the two halves sharing no factor. The law is the semi-group property of the deviation flow: composition of transports equals the direct transport.
\end{proof}

\begin{definition}[Frame-over-frame ratio]\label{def:normalization}
\begin{equation}\label{eq:gamma}
\gamma_{n} := \frac{N(n)}{N(n+1)} = \frac{2n-1}{2n+1},\qquad 1 - \gamma_{n} = \frac{2}{2n+1} \sim \frac{1}{n},
\end{equation}
using \cref{eq:total}. The complementary fraction $1 - \gamma_{n} = J / N(n+1)$ is the injected cell against the next total.
\end{definition}

\begin{proposition}[Stability of the chain]\label{prop:stability}
Every solution of the deviation system \cref{eq:deviation} decays relative to the growing contents:
\begin{equation}\label{eq:decay}
\delta_{\mathrm B}^{(n)} \propto \frac{1}{n},\qquad
\delta_{\mathrm r}^{(n)} \propto \frac{\ln n}{n},\qquad
\delta_{\mathrm m}^{(n)} \ \text{bounded},
\end{equation}
and the zero-sum $\delta_{\mathrm B}^{(n)} + \delta_{\mathrm r}^{(n)} + \delta_{\mathrm m}^{(n)} = \text{const}$ holds. In the continuum reading $t = n t_{\mathrm P}$ ($n \gg 1$), the function
\begin{equation}\label{eq:lyapunov}
V = \tfrac{1}{2}\,\delta_{\mathrm B}^{2} + \delta_{\mathrm r}^{2}
\quad\text{satisfies}\quad
\dot V = -\frac{(\delta_{\mathrm B} - \delta_{\mathrm r})^{2} + \delta_{\mathrm r}^{2}}{t} \leq 0,
\end{equation}
a Lyapunov function for the deviation flow.\footnote{A.~M. Lyapunov, \emph{The general problem of the stability of motion} (1892): a scalar function decreasing along trajectories certifies stability. The function \cref{eq:lyapunov} is constructed for the present system; its decay rate is not optimised.}
\uses{\cref{prop:kernel}, \cref{eq:deviation}}
\end{proposition}

\begin{proof}
\emph{Bulk channel.} Directly from \cref{prop:kernel}: $\delta_{\mathrm B}^{(n)} = \mathcal{W}(n; m)\,\delta_{\mathrm B}^{(m)} = \frac{m-1}{n-1}\delta_{\mathrm B}^{(m)} \propto \frac{1}{n}$.

\emph{Radiation channel.} Multiply the radiation equation by $n$ and define $\rho^{(n)} := (n-1)\,\delta_{\mathrm r}^{(n)}$; then
\begin{equation}\label{eq:rho}
n\,\delta_{\mathrm r}^{(n+1)} = (n-1)\,\delta_{\mathrm r}^{(n)} + \delta_{\mathrm B}^{(n)}
\quad\Longrightarrow\quad
\rho^{(n+1)} = \rho^{(n)} + \delta_{\mathrm B}^{(n)},
\end{equation}
so $\rho$ telescopes with summand $\delta_{\mathrm B}^{(j)} = \frac{m-1}{j-1}\delta_{\mathrm B}^{(m)}$:
\begin{equation}\label{eq:rhotelescope}
\rho^{(n)} = \rho^{(m)} + (m-1)\,\delta_{\mathrm B}^{(m)} \sum_{j = m}^{n-1} \frac{1}{j-1} = \rho^{(m)} + (m-1)\,\delta_{\mathrm B}^{(m)}\,\bigl( H_{n-2} - H_{m-2} \bigr),
\end{equation}
where $H_{r} = \sum_{i=1}^{r} 1/i$ is the $r$-th harmonic number. Dividing by $n - 1$:
\begin{equation}\label{eq:exactdr}
\delta_{\mathrm r}^{(n)} = \frac{m-1}{n-1}\Bigl[ \delta_{\mathrm r}^{(m)} + (m-1)\,\delta_{\mathrm B}^{(m)}\,\bigl(H_{n-2} - H_{m-2}\bigr) \Bigr] \;\propto\; \frac{\ln n}{n},
\end{equation}
using the classical asymptotic $H_{r} = \ln r + \gamma_{\mathrm{E}} + o(1)$, with $\gamma_{\mathrm{E}}$ Euler's constant.\footnote{L.~Euler (1734) showed $H_r - \ln r$ converges; the limit $\gamma_{\mathrm E} \approx 0.5772$ is named after him. Only the logarithmic growth matters here.}

\emph{Matter channel.} From \cref{eq:deviation}, $\delta_{\mathrm m}^{(n)} = \delta_{\mathrm m}^{(m)} + \sum_{j=m}^{n-1} \delta_{\mathrm r}^{(j)}/j$; with $\delta_{\mathrm r}^{(j)} \leq C\,\frac{\ln j}{j}$, comparison with the integral $\int^{\infty} \frac{\ln x}{x^{2}}\,dx = \Bigl[ -\frac{\ln x + 1}{x} \Bigr]^{\infty}$, which converges, bounds $\delta_{\mathrm m}$ independently of $n$.

\emph{Zero-sum.} Add the three equations of \cref{eq:deviation}: every conducted term cancels against its inflow partner, $-\delta_{\mathrm B}/n + \delta_{\mathrm B}/n - \delta_{\mathrm r}/n + \delta_{\mathrm r}/n = 0$, so the total deviation is constant: perturbations do not change the summed content, they migrate downstream through the channels.

\emph{Lyapunov function (continuum reading).} For $n \gg 1$ the deviation system reads $\dot\delta_{\mathrm B} = -\delta_{\mathrm B}/t$, $\dot\delta_{\mathrm r} = (\delta_{\mathrm B} - \delta_{\mathrm r})/t$, $\dot\delta_{\mathrm m} = \delta_{\mathrm r}/t$. Differentiate \cref{eq:lyapunov} and complete the square in the numerator:
\begin{equation}\label{eq:vdot}
\dot V = \delta_{\mathrm B}\Bigl(-\frac{\delta_{\mathrm B}}{t}\Bigr) + 2\delta_{\mathrm r}\Bigl(\frac{\delta_{\mathrm B} - \delta_{\mathrm r}}{t}\Bigr)
= \frac{-\delta_{\mathrm B}^{2} + 2\delta_{\mathrm B}\delta_{\mathrm r} - 2\delta_{\mathrm r}^{2}}{t}
= -\frac{(\delta_{\mathrm B} - \delta_{\mathrm r})^{2} + \delta_{\mathrm r}^{2}}{t} \leq 0,
\end{equation}
where the last equality is the completing-the-square step $-\delta_{\mathrm B}^{2} + 2\delta_{\mathrm B}\delta_{\mathrm r} = -\bigl[(\delta_{\mathrm B}-\delta_{\mathrm r})^{2} - \delta_{\mathrm r}^{2}\bigr]$, written out rather than quoted.
\end{proof}

\begin{remark}[Where the logarithm lives]\label{rem:harmonic-anchor}
The logarithm of \cref{eq:decay} is not an approximation: it is the harmonic number of \cref{eq:exactdr}, exact at every frame. The harmonic/$\zeta$-type series recorded as pure mathematics in \cref{rem:kernel-symbol} thus finds its physical anchor here --- the $H_{n-2}$ of the radiation channel --- and the remainder of the chain carries no further transcendental entry.
\end{remark}

\begin{remark}[Reading of the stability]\label{rem:stability-reading}
The deviation system is homogeneous, so \cref{eq:decay} is decay of \emph{relative} deviations; absolute deviations in the matter channel need not shrink. What the chain forgets is shape: injected perturbations are swept downstream and diluted as $1/n$ (bulk), $\ln n/n$ (radiation), while their sum rides along untouched. The system is stable not because disturbances vanish but because they are organised out of the visible channels --- the physical content of \cref{cor:annihilation} at every frame, not only the first.
\end{remark}

\section{The forced chain}\label{sec:forced-cascade}

The conductance law \cref{eq:share} was introduced as a hypothesis, with a promise attached: the law $1/n$ is not one choice among many, but the \emph{only} law compatible with a non-degenerate steady state. This section redeems the promise. The argument runs in two parts --- a split rule forced by the frame algebra, and an existence theorem that fixes the exponent of the conductance law --- and both are stated with $J$ symbolic, per \cref{rem:jform}.

\subsection{The forced split}

Assume the chain has reached linear growth,
\begin{equation}\label{eq:steadystate}
E_{i}^{(n)} = \kappa_{i}\, n \qquad (i \in \{\mathrm B, \mathrm r, \mathrm m\}),
\end{equation}
which is the natural steady form here: the total content increases by exactly $J$ per frame (\cref{prop:total}), and linear growth of the total is the unique steady shape consistent with a constant increment. At each node, the inflow $J_{i}$ splits into growth and conduction:
\begin{equation}\label{eq:nodebalance}
J_{i} = \underbrace{E_{i}^{(n+1)} - E_{i}^{(n)}}_{\text{growth}} + \underbrace{\frac{1}{n} E_{i}^{(n)}}_{\text{conducted, by \cref{as:b2}}},
\qquad\text{so}\qquad
J_{i} = \kappa_{i} + \kappa_{i} = 2\kappa_{i},
\end{equation}
because for linear growth the growth term is $\kappa_{i}\bigl((n+1) - n\bigr) = \kappa_{i}$ and the conducted term is $\tfrac1n\cdot\kappa_{i}n = \kappa_{i}$ --- the share \cref{eq:share} makes the two contributions \emph{equal at every node}, whatever $\kappa_{i}$ is. The split is forced, not chosen: each passing node retains exactly half of its inflow. The coefficients are now determined recursively, step by step, from the injection $J_{0} = J$:
\begin{equation}\label{eq:splitchain}
\begin{aligned}
&\kappa_{\mathrm B} = \frac{J_{0}}{2} = \frac{J}{2}, \qquad \mathcal T_{\mathrm B} = \frac{J}{2} =: J_{1},\\
&\kappa_{\mathrm r} = \frac{J_{1}}{2} = \frac{J}{4}, \qquad \mathcal T_{\mathrm r} = \frac{J}{4} =: J_{2},\\
&\kappa_{\mathrm m} = J_{2} = \frac{J}{4} \qquad (\text{terminal node: retains all}).
\end{aligned}
\end{equation}
The flow sequence $J_{i} = J / 2^{i}$ halves at every hop, and the coefficients are unique. Comparison with \cref{eq:solutions}: the exact solution has $E_{\mathrm B}^{(n)} = nJ/2$ and $E_{\mathrm r}^{(n)} = nJ/4$ exactly, while $E_{\mathrm m}^{(n)} = (n-2)J/4 = \kappa_{\mathrm m}n - J/2$ differs from the steady form by the constant $J/2$ --- a bounded offset inherited from the boundary condition \cref{eq:firstframe}. The forced state is the attractor, reached up to $O(1)$ offsets; this is the stability of \cref{prop:stability} seen from the solution side.

\subsection{The conductance exponent is forced}

\begin{theorem}[Existence of a non-degenerate steady state]\label{thm:forced-cascade}
Consider the conductance family
\begin{equation}\label{eq:family}
E^{(n+1)} = \Bigl(1 - \frac{1}{n^{\gamma}}\Bigr) E^{(n)} + J, \qquad \gamma > 0,\ J > 0,
\end{equation}
of which \cref{as:b2} is the member $\gamma = 1$. A solution with a non-degenerate linear asymptote --- bulk, radiation and matter contents all $\propto n$ with non-zero coefficients --- exists \textbf{if and only if} $\gamma = 1$; the coefficients are then those of \cref{eq:splitchain}, uniquely.
\uses{\cref{eq:cascade}, \cref{eq:share}, power comparison (\emph{dominant balance})}
\end{theorem}

\begin{proof}
Seek an asymptotic solution $E^{(n)} = g_{0}\, n^{\eta}\bigl(1 + o(1)\bigr)$ with $g_{0} > 0$ and insert it into \cref{eq:family}. The increment on the left and the loss on the right have the respective leading behaviours
\begin{equation}\label{eq:balance}
E^{(n+1)} - E^{(n)} = g_{0}\,\eta\, n^{\eta - 1}\bigl(1 + o(1)\bigr),
\qquad
-\frac{1}{n^{\gamma}}E^{(n)} = -g_{0}\, n^{\eta - \gamma}\bigl(1 + o(1)\bigr),
\end{equation}
so the equation demands
\begin{equation}\label{eq:compare}
g_{0}\eta\, n^{\eta-1} \;=\; -g_{0}\, n^{\eta-\gamma} + J \qquad (\text{to leading order}).
\end{equation}
Three exponent comparisons exhaust the possibilities.\footnote{The method is \emph{dominant balance}: an asymptotic equation can hold only if its largest terms cancel, which fixes the exponents first and the coefficients afterwards. Standard reference: C.~M. Bender and S.~A. Orszag, \emph{Advanced Mathematical Methods for Scientists and Engineers} (1978), Ch.~3.}

\emph{Case $\gamma < 1$.} Then $\eta - \gamma > \eta - 1$: the loss term grows strictly faster than the increment, whatever $\eta$ is. If $\eta \geq 1$, the loss eventually exceeds the constant injection $J$ and \cref{eq:compare} cannot balance at large $n$; if $\eta < \gamma$, the increment decays while the loss dominates; the borderline $\eta = \gamma$ forces the constant loss $-g_{0}$ to cancel $J$ exactly, leaving the decaying increment unbalanced. No linear asymptote exists; the true asymptote is sublinear, $E^{(n)} \sim J\,n^{\gamma}$, drained by a conductance faster than the injection can sustain linear growth.

\emph{Case $\gamma = 1$.} The two sides of \cref{eq:compare} carry the same power $n^{\eta - 1}$ iff $\eta = 1$; with $\eta > 1$ the increment grows while the right side stays bounded (impossible), and with $\eta < 1$ the increment decays while $J$ stands (impossible). At $\eta = 1$ the coefficients must satisfy
\begin{equation}\label{eq:coefficient}
g_{0} = -g_{0} + J \quad\Longrightarrow\quad g_{0} = \frac{J}{2} \quad (\text{uniquely}).
\end{equation}

\emph{Case $\gamma > 1$.} Then $\eta - \gamma < \eta - 1$: the loss decays relative to the increment, and $\eta = 1$ balances with $g_{0} = J$ --- the bulk channel alone grows linearly. But the conducted flow is $n^{-\gamma} E^{(n)} \sim J\, n^{1-\gamma} \to 0$: the downstream channels starve, and the three-component steady state is degenerate (non-zero coefficients fail).

Uniqueness of the coefficients for $\gamma = 1$ is the recursion \cref{eq:splitchain}: the split rule $J_{i} = 2\kappa_{i}$ determines each $\kappa_{i}$ from its inflow, and the terminal node retains all. Hence a non-degenerate linear steady state exists iff $\gamma = 1$, with coefficients \cref{eq:splitchain} uniquely.
\end{proof}

\begin{remark}[What fixes the law, and what does not]\label{rem:gamma-one}
\Cref{thm:forced-cascade} fixes the exponent $\gamma = 1$ by \emph{existence alone} --- no consistency condition, no fitting, no second assumption; the shares need no second fixing, the internal bridge of \cref{ch:boundary} carrying what a coincidence requirement would have carried. The law \cref{eq:share} thereby rises from convention to forced structure: \cref{as:b2} remains the official statement, and the theorem is its justification.
\end{remark}

\section{Scope of the chain}\label{sec:cascade-scope}

Three statements fix what the chain does and does not claim. They are collected here, at the exit of the chapter, because each of them is a guard on the way the earlier formulas may be read.

\begin{remark}[Truncation, restated without commitment]\label{rem:truncation-restated}
The working content of this paper is the three-component chain; the matter component is terminal. The mathematical continuation $16, 32, \ldots$ of the halving sequence exists as algebra, and its \emph{status} is open: it may be a pure artifact of the mathematical form, it may have a physical correspondent, or it may be the visible shadow of a structure this paper does not contain. Nothing in this chapter decides among these, and the chapter claims nothing in either direction. What is enforced is only a working discipline: no formula of this paper is to be extended beyond the stated content, and the discipline itself is silent about what lies beyond.
\end{remark}

\begin{remark}[The one-frame invisibility of production]\label{rem:invisibility}
By \cref{eq:solutions}, the radiation content at frame $n$ draws on bulk contents of frames $n-1$ and earlier only: the current-frame production
\begin{equation}\label{eq:production}
\frac{1}{n}E_{\mathrm B}^{(n)} = \frac{1}{n}\cdot\frac{nJ}{2} = \frac{J}{2}
\end{equation}
is a frame-intermediate quantity, not a content; it acquires content status only from frame $n+1$ onward, through the radiation equation of \cref{eq:cascade}. Two quantities are therefore offset by exactly one frame: the \emph{content at a frame} (what a single-frame observation sees) and the \emph{production during a frame transition} (what the transition generates). They cannot be merged into one simultaneous reading, and no formula of this paper merges them: each content equation uses only inherited contents, by the delay structure of \cref{rem:delay}. As a reading: at any single frame, only what was produced at least one frame ago is present.
\end{remark}

\begin{remark}[One constant, three roles]\label{rem:threenames}
The quantity
\begin{equation}\label{eq:Jhalf}
\frac{J}{2}
\end{equation}
appears in three roles: as the per-frame conducted amount $\frac1n E_{\mathrm B}^{(n)}$ of \cref{eq:production}; as the first-frame retained content $E_{\mathrm B}^{(1)}$ of \cref{eq:firstframe}; and as the frame-independent residual $D(n) = nJ - N(n)$ of \cref{eq:residual}. The equality of the three is arithmetic, not structural: they belong to different frame windows --- a production rate, an initial condition, and a running difference --- and reading them as one held object would ascribe to a number an identity the frame structure does not give it. What is structural is the steady-state balance behind them: each frame retains $J/2$ of its inflow and transmits $J/2$, so the residual is re-created at every step rather than carried by anything. The identification of $D$ with the exterior content is deferred to \cref{ch:boundary} (\cref{rem:residual}); this chapter keeps the arithmetic.
\end{remark}

\begin{remark}[What Chapter 3 has established]\label{rem:ch3summary}
On the three hypotheses B1--B3 and the boundary condition alone: the exact frame-by-frame solution \cref{eq:solutions}; the constant increment $J$ of the total content; the response kernel $\mathcal{W}(n;m) = (m-1)/(n-1)$ with its annihilation of first-frame deviations and its cocycle law; stability in the sense of \cref{eq:decay}; and the forced split, with the existence theorem that the conductance exponent $\gamma = 1$ is the unique non-degenerate choice. The chapter's own open items, each with a named successor: the thermodynamic root of B1 (\cref{ch:boundary}); the magnitude $J$ (\cref{ch:unitarity}); the even action sites (\cref{rem:evensites}); and the speculative readings of \cref{rem:speculation}, which remain speculation. Nothing in the chapter required dynamics, and nothing in it assumed any.
\end{remark}

\chapter{The Newtonian leg}\label{ch:newton}

This chapter introduces the only spatial degree of freedom the framework has: the radius $R$ of a uniform ball. Its treatment is in two parts. The first part, \cref{sec:steps}, is the \emph{stepping structure}: what the radius \emph{is} on the frame sequence, before any dynamics. The second part, the remaining sections, subjects the same radius to Newtonian energetics and shows that the energetics, rather than adding a dynamics, is forced into one compatible shape. A reader who accepts the stepping structure on first reading may take \cref{prop:uniform-rate} as the chapter's central fact and treat the rest as its Newtonian shadow.

\section{The stepping structure}\label{sec:steps}

\begin{hypothesis}[Stepping]\label{as:stepping}
The radius is read off the frame count by
\begin{equation}\label{eq:stepping}
R(n) = n\,\ell_{\mathrm P}, \qquad n \in \mathbb{Z}^{+},
\end{equation}
with the same frame index as \cref{as:a1}.
\motivation{Same source as \cref{as:a1}: the frame sequence is the primary quantity, and the spatial axis is not an independent degree of freedom but the same sequence read in units of $\ell_{\mathrm P}$, the conversion rate being the first Planck relation. The hypothesis joins the hypothesis layer as its fourth member; hypotheses introduced after \cref{sec:hypotheses} are recorded individually with their motivation, and the concluding chapter (\cref{ch:conclusion}) collects the full table.}
\end{hypothesis}

\begin{proposition}[Uniform rate, exact]\label{prop:uniform-rate}
For every frame,
\begin{equation}\label{eq:uniformrate}
\frac{\Delta R}{\Delta t} = \frac{R(n+1) - R(n)}{t_{\mathrm P}} = \frac{\ell_{\mathrm P}}{t_{\mathrm P}} = c.
\end{equation}
\uses{\cref{as:stepping}, \cref{def:units} (first Planck relation)}
\end{proposition}

\begin{proof}
Numerator first, by \cref{eq:stepping}: $R(n+1) - R(n) = (n+1)\ell_{\mathrm P} - n\ell_{\mathrm P} = \ell_{\mathrm P}$, with no remainder, because $R$ is linear in $n$: the difference quotient of a linear function is exact and carries no correction term. The denominator is the step, $\Delta t = t_{\mathrm P}$ (\cref{eq:tick}). The ratio is $\ell_{\mathrm P}/t_{\mathrm P}$, and the first Planck relation $\ell_{\mathrm P} = c\,t_{\mathrm P}$ (\cref{def:units}) converts it to $c$.
\end{proof}

\begin{corollary}[No instantaneous cross-frame response]\label{cor:delay}
No physical response can transmit information across a frame boundary in less than one step: any chain of responses crossing a distance $L$ requires at least $L / \ell_{\mathrm P}$ frames, and in particular a response between consecutive frames requires at least one step. The maximal displacement per frame is $\ell_{\mathrm P}$, attained by the radius itself.
\uses{\cref{prop:uniform-rate}}
\end{corollary}

\begin{proof}
By \cref{prop:uniform-rate}, the fastest available displacement per frame is the step $\ell_{\mathrm P}$ taken at rate $c$; a response chain transmitting information across $L$ must therefore accumulate at least $\lceil L/\ell_{\mathrm P}\rceil$ frames, one step each by \cref{eq:tick}. Nothing faster is available on the frame sequence of \cref{as:a1,as:stepping}.
\end{proof}

This corollary discharges the two forward references of \cref{ch:cascade}: the delay that entered that chapter as motivation (\cref{rem:delay}) and the one-frame offset between production and content (\cref{rem:invisibility}) are here formal consequences of the stepping structure, not additional assumptions. The chain stands, as it always did, on its own hypotheses; the delay is now known to be forced.

\begin{remark}[Two readings of $R = ct$]\label{rem:tworeadings}
\Cref{eq:uniformrate} equips the identity $R = ct$ with two readings, and the rule of \cref{rem:register} assigns them different status. As \emph{kinematics}: the radius advances one step per tick, and $R = ct$ is the conversion of that statement into continuous units --- a reading that is exact, derivation-free, and already complete. As \emph{dynamics}: $R = ct$ could be imagined as a solution of an equation of motion, sustained against deceleration by some driving --- a reading that requires a theory this paper does not contain, and which \cref{sec:newton-ball} will show to be not even available: after the matching, the would-be equation of motion is an identity. The first reading is the content of the framework; the second is not asserted, anywhere.
\end{remark}

\section{The Newtonian ball, and other geometries}\label{sec:newton-ball}

\subsection{Setting and energetics}\label{sec:ball-setting}

The setting is a uniform, isotropic ball: radius $R(t)$, homogeneous density $\rho(t)$, total content
\begin{equation}\label{eq:ballmass}
M = \rho\, V(R), \qquad V(R) = \frac{4\pi}{3} R^{3},
\end{equation}
with a unit test particle at the edge. Two classical results enter as external inputs.

\begin{eresult}[Shell theorem and energy integral]\label{er:ball}
(i) A spherically symmetric shell exerts no force on a body inside it; the ball acts on its edge as if its mass were concentrated at the centre. (ii) The edge particle's specific energy is conserved:
\begin{equation}\label{eq:energy}
\frac{1}{2}\dot R^{2} - \frac{G M}{R} = E = \text{const}.
\end{equation}
\reference{Newton, \emph{Principia} (1687), Book I, Props.~70--71 (shell theorem); the conserved quantity \cref{eq:energy} is the vis viva integral of Leibniz (1686), for a central potential. For the Newtonian treatment of an expanding cosmological ball, see Milne and McCrea (1934),\footnote{E.~A. Milne and W.~H. McCrea, \emph{Newtonian universes and the curvature of space}, Quart.\ J.\ Math.\ \textbf{5}, 73 (1934): the Newtonian ball as a derivation of the Friedmann equation. The present chapter uses the same device at the frame level.} whose lineage this section follows.}
\end{eresult}

With $M$ constant for the moment, \cref{eq:energy} has three branches, and it is worth displaying them explicitly before anything is matched to the frame sequence, because each branch is a \emph{genuinely distinct fate} of the same equation:
\begin{equation}\label{eq:branches}
\dot R^{2} = \frac{2 G M}{R} + 2E,
\qquad
\begin{cases}
E < 0: & \text{halts at } R_{\max} = \tfrac{G M}{|E|},\ \text{then recollapse};\\[6pt]
E = 0: & \dot R^{2} = \tfrac{2 G M}{R};\\[6pt]
E > 0: & \dot R^{2} = \tfrac{2 G M}{R} + 2E.
\end{cases}
\end{equation}
On the third branch $\dot R \to \sqrt{2E}$ at late times. The halting radius in the first line follows from $\dot R = 0$ in \cref{eq:branches}: $0 = 2GM/R_{\max} + 2E$ with $E = -|E|$, solved for $R_{\max} = GM/|E|$. The three branches are parabolic, elliptic and hyperbolic in the orbital sense, and within Newtonian mechanics alone no principle selects among them: $E$ is a free constant of integration, fixed by initial conditions.

The stepping structure of \cref{sec:steps} now acts on all three branches at once, because it fixes $\dot R = c$ exactly (\cref{prop:uniform-rate}). Substituting into \cref{eq:energy},
\begin{equation}\label{eq:branches-stepped}
E = \frac{c^{2}}{2} - \frac{G M(R)}{R}
\qquad\Longleftrightarrow\qquad
M(R) = \frac{c^{2} R}{2 G} - \frac{E R}{G},
\end{equation}
so the three branches become three \emph{mass profiles}: $E = 0$ is the unique branch whose mass content is linear in the radius,
\begin{equation}\label{eq:linear-profile}
E = 0 \;\Longleftrightarrow\; M(R) = \frac{c^{2} R}{2 G};
\end{equation}
$E < 0$ requires $M < c^{2}R/2G$ at every radius, and $E > 0$ requires $M > c^{2}R/2G$. The linear profile \cref{eq:linear-profile} is the deposited-shell profile of \cref{ch:boundary}; the agreement is recorded here and verified there, one-way (\cref{ch:boundary} does not use this section). Which branch the matching ultimately selects is settled in \cref{sec:frame-readout} --- the formal statement that the definition of the reading forces $E = 0$ is \cref{prop:compatibility} there --- and until then all three branches remain on the table as what they are: a real freedom of the Newtonian layer, not yet a decided question.

\begin{remark}[No numerical values, by design]\label{rem:no-numbers}
This chapter computes shapes and exponents only. Numerical values enter this paper in exactly two places: the working values of the Planck units (\cref{eq:values}, \cref{eq:epsvalue}), and the temperatures of the boundary thermodynamics where they are unavoidable (\cref{ch:boundary}). Everything between those two places is symbolic, because every relation below is a relation of forms, and forms are what the frame structure determines.
\end{remark}

\subsection{Other geometries, and what would change}\label{sec:other-geometries}

This subsection is optional and exploratory: the ball of \cref{sec:ball-setting} is not the only geometry on which the same formalism can be run, and the exercise is worth recording because it shows which conclusions are the ball's and which are not. Its status must be stated first and plainly: the energy equation \cref{eq:energy} is a theorem of flat Newtonian mechanics (\cref{er:ball}); on a curved or higher-dimensional space it is an \emph{ansatz}, a formal extension whose credentials this subsection does not claim. Everything below is conditional on that ansatz, and the conclusions are correspondingly soft.

The general-$d$ ansatz has two ingredients: a homogeneous content $M = \rho\, V_{d}(R)$ and a shell-theorem potential at the edge, which in $d$ spatial dimensions follows from the flux of gravity through the $(d-1)$-sphere of area $\Omega_{d-1} R^{d-1}$:
\begin{equation}\label{eq:d-dim}
\begin{gathered}
g(R)\,\Omega_{d-1} R^{d-1} = \Omega_{d-1} G M
\;\Longrightarrow\;
g(R) = \frac{G M}{R^{d-2}}\\
\;\Longrightarrow\;
U(R) = -\frac{G M}{(d-2)\,R^{d-2}}\ (d \geq 3),
\end{gathered}
\end{equation}
the potential obtained by integrating $-g$ from infinity, with the dimensional constant absorbed into $G$. For $d = 3$ this returns \cref{eq:energy} exactly: $\Omega_{2} = 4\pi$ gives $g = GM/R^{2}$ and $U = -GM/R$. The energy equation reads
\begin{equation}\label{eq:energy-d}
\frac{1}{2}\dot R^{2} - \frac{G M}{(d-2) R^{d-2}} = E,
\qquad\text{so at } E = 0:\qquad
H^{2} = \frac{\dot R^{2}}{R^{2}} = \frac{2 G \rho\, V_{d}(R)}{(d-2)\,R^{d}}.
\end{equation}
The volumes follow from the ball-volume recursion $V_{d}(R) = \int_{0}^{R} \Omega_{d-1} r^{d-1}\,dr$ with $\Omega_{2} = 4\pi$, $\Omega_{3} = 2\pi^{2}$:
\begin{equation}\label{eq:volumes}
V_{3}(R) = \int_{0}^{R} 4\pi r^{2}\,dr = \frac{4\pi}{3}R^{3},
\qquad
V_{4}(R) = \int_{0}^{R} 2\pi^{2} r^{3}\,dr = \frac{\pi^{2}}{2}R^{4}.
\end{equation}
For the $3$-sphere as an entire space (the cosmological reading), the homogeneous content is $M = \rho \cdot 2\pi^{2} R^{3}$, the volume of $S^{3}$ of radius $R$.

Three geometries, then, under the same ansatz. The constraint coefficient at $E = 0$, from \cref{eq:energy-d}:
\begin{equation}\label{eq:coefficients}
H^{2} = 2 G \rho\, \frac{V_{d}(R)}{(d-2)R^{d}}:
\qquad
\begin{array}{l}
\text{flat 3-ball:} \quad H^{2} = 2G\rho\cdot\dfrac{4\pi/3}{1} = \dfrac{8\pi G}{3}\rho,\\[10pt]
\text{$S^{3}$:} \quad H^{2} = 2G\rho\cdot\dfrac{2\pi^{2}}{1} = 4\pi^{2} G\rho,\\[10pt]
\text{flat 4-ball:} \quad H^{2} = 2G\rho\cdot\dfrac{\pi^{2}/2}{2} = \dfrac{\pi^{2} G}{2}\rho.
\end{array}
\end{equation}
And the injected solution $M = \mu t$ at $E = 0$:
\begin{equation}\label{eq:pumped-d}
\begin{aligned}
d = 3:\quad & \dot R^{2} = \frac{2G\mu t}{R} \;\Longrightarrow\; R^{\frac12} dR = \sqrt{2G\mu}\,t^{\frac12} dt
\;\Longrightarrow\; \tfrac{2}{3}R^{\frac32} = \tfrac{2}{3}\sqrt{2G\mu}\,t^{\frac32}\\
& \;\Longrightarrow\; R = (2G\mu)^{\frac13} t;\\
d = 4:\quad & \dot R^{2} = \frac{2G\mu t}{2R^{2}} = \frac{G\mu t}{R^{2}} \;\Longrightarrow\; R\,dR = \sqrt{G\mu}\,t^{\frac12} dt
\;\Longrightarrow\; \tfrac{R^{2}}{2} = \tfrac{2}{3}\sqrt{G\mu}\,t^{\frac32}\\
& \;\Longrightarrow\; R \propto t^{\frac34},
\end{aligned}
\end{equation}
with $S^{3}$ sharing the $d = 3$ line because its ansatz potential is the same $1/R$. The comparison:
\begin{center}
\begin{tabular}{@{}llll@{}}
\toprule
Geometry & Pumped dynamics & Constraint coefficient & Outcome \\
\midrule
flat 3-ball & $R = (2G\mu)^{1/3}t$, linear & $8\pi G/3$ & matches \cref{ch:gr} \\
$S^{3}$ & linear (same) & $4\pi^{2} G$ & coefficient fails \\
flat 4-ball & $R \propto t^{3/4}$ & $\pi^{2}G/2$ & matching itself fails \\
\bottomrule
\end{tabular}
\end{center}

\begin{remark}[What the exercise establishes]\label{rem:geometry-verdict}
The Newtonian ball is not an arbitrary choice, and the table shows what selects it: the stepping matching requires the pumped dynamics to be \emph{linear} in $t$, which the $d = 3$ potential $1/R$ delivers and the $d = 4$ potential $1/R^{2}$ does not; and the reading of \cref{ch:gr} requires the constraint coefficient $8\pi G/3$, which the flat 3-ball delivers and $S^{3}$ does not. The ball is selected \emph{retroactively} by two requirements it is the only one of the three to meet. This subsection claims nothing for the ansatz beyond the conditional: with the credentials of \cref{eq:energy-d} unexamined, the table records what changes, and the main text's standing on the flat ball is strengthened, not weakened, by knowing the alternatives.
\end{remark}

\section{The pumped family}\label{sec:pumped-family}

The ball of \cref{sec:newton-ball} was treated with constant content $M$. This section opens the system: content is added from the outside at a constant rate, the same structural choice as \cref{as:b3} at the level of the whole ball. Which energy branch is used must be declared at the start: the section works on the $E = 0$ branch of \cref{eq:branches}, because \cref{sec:ball-setting} identified that branch as the one whose mass profile can be linear in $R$ --- the shape the stepping matching will require --- and \cref{sec:frame-readout} formalizes the selection. The restriction is stated, not hidden; the $E \neq 0$ pumped solutions are not developed here.

\begin{hypothesis}[Constant injection]\label{as:pump}
The ball's content grows at a constant rate:
\begin{equation}\label{eq:pump}
\frac{dM}{dt} = \mu, \qquad \mu = \text{const} \geq 0,
\end{equation}
with $R(0) = 0$ (the first frame is empty, in line with the counting of \cref{cor:counting}).
\motivation{The open system needs a source, and a constant rate is the minimal such structure --- one new number, no further shape. The mechanism behind the injection is an open question (\cref{rem:mechanismregistry}); \cref{eq:pump} names the structure, and nothing deeper is implied by the name. The rate $\mu$ itself is a structure parameter here; its value is fixed by the stepping matching in \cref{sec:frame-readout}, the same separation of magnitude and structure as in \cref{as:b3}.}
\end{hypothesis}

Integrating \cref{eq:pump} once, with initial content $M_{0} \geq 0$:
\begin{equation}\label{eq:mt}
M(t) = M_{0} + \mu t.
\end{equation}
The $E = 0$ energy equation \cref{eq:branches} with \cref{eq:mt} reads
\begin{equation}\label{eq:pumpeq}
\dot R^{2} = \frac{2G (M_{0} + \mu t)}{R}.
\end{equation}

\begin{proposition}[General pumped solution]\label{prop:generalsolution}
The unique solution of \cref{eq:pumpeq} with $R(0) = 0$ is
\begin{equation}\label{eq:generalsolution}
R(t) = (2G)^{1/3}\,\mu^{-2/3}\,\Bigl[ \bigl(\mu t + M_{0}\bigr)^{3/2} - M_{0}^{3/2} \Bigr]^{2/3}.
\end{equation}
\uses{\cref{eq:pumpeq}, \cref{eq:pump}}
\end{proposition}

\begin{proof}
Separate variables in \cref{eq:pumpeq}, taking the positive root since $R$ grows:
\begin{equation}\label{eq:separate}
R^{1/2}\,dR = \sqrt{2G\,(M_{0} + \mu t)}\;dt.
\end{equation}
Substitute $u = M_{0} + \mu t$, so that $du = \mu\,dt$ and $dt = du/\mu$:
\begin{equation}\label{eq:usub}
\int R^{1/2}\,dR = \int \sqrt{2G}\;u^{1/2}\,\frac{du}{\mu}
\quad\Longrightarrow\quad
\frac{2}{3} R^{3/2} = \frac{\sqrt{2G}}{\mu}\cdot\frac{2}{3}\,u^{3/2} + C,
\end{equation}
the two elementary integrations written separately. Determine $C$ from $R(0) = 0$, where $u = M_{0}$:
\begin{equation}\label{eq:constant}
0 = \frac{\sqrt{2G}}{\mu}\cdot\frac{2}{3}\,M_{0}^{3/2} + C
\quad\Longrightarrow\quad
C = -\frac{2\sqrt{2G}}{3\mu}\,M_{0}^{3/2}.
\end{equation}
Substitute back and cancel the common factor $\tfrac23$:
\begin{equation}\label{eq:presolved}
R^{3/2} = \frac{\sqrt{2G}}{\mu}\Bigl[ \bigl(\mu t + M_{0}\bigr)^{3/2} - M_{0}^{3/2} \Bigr].
\end{equation}
Raise both sides to the power $2/3$ --- unambiguous, both sides positive --- tracking the exponents of each factor:
\begin{equation}\label{eq:raise}
R = \underbrace{\bigl(\sqrt{2G}\bigr)^{2/3}}_{(2G)^{1/3}}\ \underbrace{\bigl(\mu^{-1}\bigr)^{2/3}}_{\mu^{-2/3}}\ \Bigl[ \bigl(\mu t + M_{0}\bigr)^{3/2} - M_{0}^{3/2} \Bigr]^{2/3},
\end{equation}
which is \cref{eq:generalsolution}. Uniqueness follows from \cref{eq:pumpeq} by standard Picard iteration: the right-hand side $\sqrt{2G(M_{0}+\mu t)/R}$ is Lipschitz in $R$ away from $R = 0$, and the solution \cref{eq:generalsolution} satisfies $R(t) \sim t^{2/3}$ as $t \to 0$, hence approaches the singularity with the integrable exponent; the verification is the formula itself.
\end{proof}

Two limits organise the whole family, and both are computed from \cref{eq:presolved} rather than quoted.

\begin{proposition}[Dust limit and linear limit]\label{prop:twolimits}
\begin{equation}\label{eq:twolimits}
\mu t \ll M_{0}:\quad R(t) \to \Bigl(\frac{9 G M_{0}}{2}\Bigr)^{1/3} t^{2/3},
\qquad
\mu t \gg M_{0}:\quad R(t) \to (2G\mu)^{1/3}\,t.
\end{equation}
\uses{\cref{eq:generalsolution}, binomial expansion to first order}
\end{proposition}

\begin{proof}
\emph{Dust limit.} Factor $M_{0}$ out of the bracket in \cref{eq:presolved} and expand with $\varepsilon = \mu t / M_{0} \ll 1$, writing the binomial coefficient explicitly:
\begin{equation}\label{eq:dustexpand}
\bigl(\mu t + M_{0}\bigr)^{3/2} - M_{0}^{3/2} = M_{0}^{3/2}\Bigl[ \bigl(1 + \varepsilon\bigr)^{3/2} - 1 \Bigr]
\approx M_{0}^{3/2}\Bigl[ 1 + \frac{3}{2}\varepsilon - 1 \Bigr]
= \frac{3}{2}\,\mu t\, M_{0}^{1/2},
\end{equation}
to first order in $\varepsilon$; the next term is $O(\varepsilon^{2})$ and is dropped with the expansion recorded. Substitute into \cref{eq:presolved}:
\begin{equation}\label{eq:dustsolve}
\begin{gathered}
R^{3/2} = \frac{\sqrt{2G}}{\mu}\cdot\frac{3}{2}\,\mu t\, M_{0}^{1/2} = \frac{3}{2}\sqrt{2 G M_{0}}\; t\\
\quad\Longrightarrow\quad
R = \Bigl( \frac{3}{2}\sqrt{2 G M_{0}} \Bigr)^{2/3} t^{2/3} = \Bigl( \frac{9 G M_{0}}{2} \Bigr)^{1/3} t^{2/3},
\end{gathered}
\end{equation}
where the last simplification is $\bigl(\tfrac32\bigr)^{2/3} (2GM_{0})^{1/3} = \bigl( \tfrac94 \cdot 2 G M_{0}\bigr)^{1/3} = (9GM_{0}/2)^{1/3}$. This is the constant-content solution of \cref{sec:newton-ball} with mass $M_{0}$: the dust limit recovers the closed ball, as it must.

\emph{Linear limit.} Now $M_{0} \ll \mu t$; factor $(\mu t)^{3/2}$ out instead, and note the second term is negligible by ratio $M_{0}/\mu t$:
\begin{equation}\label{eq:linearexpand}
\bigl(\mu t + M_{0}\bigr)^{3/2} - M_{0}^{3/2} \approx (\mu t)^{3/2}\Bigl( 1 + O\bigl(\tfrac{M_{0}}{\mu t}\bigr) \Bigr).
\end{equation}
Then \cref{eq:presolved} gives
\begin{equation}\label{eq:linearsolve}
R^{3/2} = \frac{\sqrt{2G}}{\mu}\,(\mu t)^{3/2} = \sqrt{2G\mu}\;t^{3/2}
\quad\Longrightarrow\quad
R = (2G\mu)^{1/3}\,t,
\end{equation}
using $(\mu^{3/2}/\mu)^{2/3} = \mu^{1/3}$. The expansion is linear in $t$ with constant velocity
\begin{equation}\label{eq:velocity}
v = \dot R = (2G\mu)^{1/3}.
\end{equation}
\end{proof}

The two limits meet where the two terms of $M(t)$ are comparable, at the crossing time
\begin{equation}\label{eq:crossing}
t_{*} = \frac{M_{0}}{\mu}.
\end{equation}
A ball that begins with content $M_{0}$ behaves as closed dust up to $t \sim t_{*}$ and as an open, linearly expanding ball afterwards; one equation with one switch --- $\mu = 0$ or $\mu > 0$ --- and two asymptotic shapes exhaust the family.

\begin{remark}[One equation, one switch]\label{rem:oneswitch}
One equation with one switch gives two asymptotic regimes --- \cref{prop:twolimits} is the formal content, and what survives the formalisation is the asymptotics and nothing else: the switch is a statement about the \emph{source} ($\mu$ vanishes or not), not about the geometry, and the transient between the two limits is the single family \cref{eq:generalsolution}, continuous with continuous $\dot R$ across $t_{*}$.
\end{remark}

\begin{remark}[Finiteness at the origin selects the pure case]\label{rem:pure}
The pure pumped solution $M_{0} = 0$ is distinguished within the family by behaviour at $t \to 0$: for $M_{0} > 0$, $\dot R^{2} = 2G M_{0}/R \to \infty$ as $R \to 0$, while for $M_{0} = 0$, $\dot R^{2} = 2G\mu t/R = (2G\mu)^{2/3}$ stays finite --- the velocity \cref{eq:velocity} holds from the first frame on. Requiring the rate to be well-defined at the origin --- the natural regularity condition on the frame sequence --- sets $M_{0} = 0$. It is a regularity selection, not a covert extra assumption.
\end{remark}

\section{The frame-algebra reading}\label{sec:frame-readout}

The sections above treated $M$ as a quantity with its own dynamics --- constant, or growing at a rate $\mu$. This section reverses the direction: on the frame sequence, $M$ is \emph{read off} the frame count, and the would-be equation of motion becomes an identity. The reversal is the chapter's central move, and it is made in one definition.

\begin{definition}[Mass reading]\label{def:readout}
At frame $n$, the content of the ball is \emph{defined} by
\begin{equation}\label{eq:readout}
M(n) := \frac{R\,(\Delta R / \Delta t)^{2}}{2G}.
\end{equation}
\end{definition}

The definition requires commentary of exactly one kind: \cref{eq:readout} has the \emph{form} of the energy equation \cref{eq:energy} with $E = 0$ solved for $M$, and that form is the whole point. In \cref{sec:newton-ball} the same expression sat on the left of a dynamical equation, with $M$ an independent variable and $E$ a free constant; here the expression is a \emph{definition}, and nothing is solved. The summary of this section --- that the equation is a conversion rate, not a dynamical law --- is now a provable statement rather than a description (\cref{prop:identity} below).

\begin{proposition}[Content at frame $n$]\label{prop:readoutvalue}
The reading \cref{eq:readout} with the stepping structure gives
\begin{equation}\label{eq:readoutvalue}
M(n) = \frac{n\,c^{3} t_{\mathrm P}}{2G}, \qquad M(n)\,c^{2} = \frac{n\,\varepsilon_{\mathrm P}}{2}.
\end{equation}
\uses{\cref{def:readout}, \cref{as:stepping}, \cref{prop:uniform-rate}, \cref{def:units}, \cref{cor:epmp}}
\end{proposition}

\begin{proof}
Substitute the stepping values $R = n\ell_{\mathrm P}$ (\cref{eq:stepping}) and $\Delta R/\Delta t = c$ (\cref{eq:uniformrate}) into \cref{eq:readout}:
\begin{equation}\label{eq:readoutstep1}
M(n) = \frac{n\,\ell_{\mathrm P}\,c^{2}}{2G} \;\xrightarrow{\ \ell_{\mathrm P} = c\,t_{\mathrm P}\ }\; \frac{n\,c^{3} t_{\mathrm P}}{2G},
\end{equation}
the conversion using the first Planck relation (\cref{def:units}). For the energy form, multiply by $c^{2}$ and reduce $c^{5}t_{\mathrm P}/G$ to $\hbar/t_{\mathrm P}$. The needed identity is verified by squaring, as in \cref{lem:relations}:
\begin{equation}\label{eq:reduce}
\Bigl(\frac{c^{5} t_{\mathrm P}}{G}\Bigr)^{2} = \frac{c^{10}}{G^{2}}\cdot\frac{\hbar G}{c^{5}} = \frac{\hbar c^{5}}{G},
\qquad
\Bigl(\frac{\hbar}{t_{\mathrm P}}\Bigr)^{2} = \hbar^{2}\cdot\frac{c^{5}}{\hbar G} = \frac{\hbar c^{5}}{G},
\end{equation}
using $t_{\mathrm P}^{2} = \hbar G / c^{5}$ (\cref{def:units}) for the second equality; both squared quantities are positive, so the unsquared identity holds:
\begin{equation}\label{eq:reducefinal}
\frac{c^{5} t_{\mathrm P}}{G} = \frac{\hbar}{t_{\mathrm P}} = \varepsilon_{\mathrm P},
\end{equation}
the last equality being \cref{cor:epmp} read together with \cref{def:epsilon}. Hence $M(n)\,c^{2} = n\,c^{5}t_{\mathrm P}/2G = n\,\varepsilon_{\mathrm P}/2$.
\end{proof}

\begin{proposition}[The energy equation becomes an identity]\label{prop:identity}
With $M$ given by the reading \cref{eq:readout}, the left and right sides of the energy equation \cref{eq:energy} coincide at every frame with $E = 0$:
\begin{equation}\label{eq:identity}
\frac{1}{2}\Bigl(\frac{\Delta R}{\Delta t}\Bigr)^{2} - \frac{G\,M(n)}{R(n)} = \frac{c^{2}}{2} - \frac{c^{2}}{2} = 0.
\end{equation}
\uses{\cref{prop:readoutvalue}, \cref{as:stepping}}
\end{proposition}

\begin{proof}
The first term is $\tfrac12 c^{2}$ by \cref{prop:uniform-rate}. For the second, substitute \cref{eq:readoutstep1} and cancel frame by frame:
\begin{equation}\label{eq:identityproof}
\frac{G\,M(n)}{R(n)} = \frac{G}{n\,\ell_{\mathrm P}}\cdot\frac{n\,c^{3} t_{\mathrm P}}{2G} = \frac{c^{3} t_{\mathrm P}}{2\,\ell_{\mathrm P}} \;\xrightarrow{\ \ell_{\mathrm P} = c\,t_{\mathrm P}\ }\; \frac{c^{3} t_{\mathrm P}}{2\,c\,t_{\mathrm P}} = \frac{c^{2}}{2},
\end{equation}
where the numerator $G$ cancels the denominator $G$, the frame index $n$ cancels, and one power of $c$ survives against $\ell_{\mathrm P} = c\,t_{\mathrm P}$. The two terms are equal; the difference is $0$ identically.
\end{proof}

\begin{proposition}[Compatibility: the matching forces $E = 0$]\label{prop:compatibility}
In pure Newtonian mechanics, $E$ in \cref{eq:energy} is a free constant of integration with three genuinely distinct branches (\cref{eq:branches}). Once the content is defined by the reading \cref{eq:readout}, no choice remains:
\begin{equation}\label{eq:compat}
E = \frac{1}{2}\Bigl(\frac{\Delta R}{\Delta t}\Bigr)^{2} - \frac{G\,M(n)}{R(n)} \;\equiv\; 0,
\end{equation}
identically in $n$, by \cref{prop:identity}. The branch selection is therefore forced by the matching, not chosen: $E = 0$ was a freedom of the Newtonian layer alone, and it disappears the moment the layer is matched to the frame sequence.
\uses{\cref{prop:identity}}
\end{proposition}

This is the precise content of the $E = 0$ statement, and its precise boundary: as a statement of Newtonian mechanics it is a real freedom with real alternatives; as a statement of the matched framework it is an identity. Presenting it as a standing \emph{option} would have implied a discarded alternative that the framework does not in fact possess; presenting it as forced from the start would have hidden that the Newtonian layer on its own is genuinely richer. Both errors are avoided by keeping the two layers in the open.

\subsection{The fixing of the injection rate}

The pumped family of \cref{sec:pumped-family} reaches the stepping structure only at one rate. Matching the linear limit \cref{eq:velocity}, $v = (2G\mu)^{1/3}$, with the exact frame rate \cref{eq:uniformrate}, $v = c$:
\begin{equation}\label{eq:pinning}
(2G\mu)^{1/3} = c \quad\xrightarrow{\ \text{cube both sides}\ }\ 2G\mu = c^{3} \quad\Longrightarrow\quad \mu = \frac{c^{3}}{2G},
\end{equation}
uniquely. The injection rate is not chosen: it is the unique constant rate at which the pumped ball moves exactly one step per tick, and with it the whole family collapses to a single trajectory,
\begin{equation}\label{eq:ct}
R(t) = c\,t, \qquad M(t) = \frac{c^{3} t}{2G},
\end{equation}
with $M_{0} = 0$ by the regularity selection of \cref{rem:pure}.

\begin{definition}[Power and per-frame injection]\label{def:power}
The injection power is the rate of energy input,
\begin{equation}\label{eq:power}
\mathcal P := \frac{d(Mc^{2})}{dt} = \mu c^{2} = \frac{c^{5}}{2G},
\end{equation}
and the per-frame injection is its one-tick content,
\begin{equation}\label{eq:jtick}
J := \mathcal P\, t_{\mathrm P} = \frac{c^{5} t_{\mathrm P}}{2G} = \frac{\varepsilon_{\mathrm P}}{2},
\end{equation}
using \cref{eq:reducefinal}.
\end{definition}

\begin{remark}[Two quantities named $J$, one theorem apart]\label{rem:tworeadings-j}
\Cref{eq:jtick} defines the per-frame injection of the source as $\varepsilon_{\mathrm P}/2$; \cref{as:b3} posits a constant per-frame injection $J$ into the chain, its magnitude deliberately left open (\cref{rem:jform}). That the two are the same quantity is the content of the injection theorem of \cref{ch:unitarity}, and the identification is one-way: the injection theorem will use \cref{as:b3}'s structure and \cref{eq:power}, and conclude $J = \varepsilon_{\mathrm P}/2$ --- never assuming it.
\end{remark}

\begin{remark}[Kinematics, not dynamics, resolved]\label{rem:resolved}
\Cref{prop:compatibility} resolves the second reading of \cref{rem:tworeadings}. The would-be equation of motion, after the matching, has degenerated into an identity: there is no equation left of which $R = ct$ could be the solution. ``Expansion'' in this framework is the reading of frame counting in spatial units, and the question of what dynamics drives it --- the question the second reading presupposes --- has no slot in which to be posed (\cref{rem:ch2summary} for the scale, and \cref{ch:gr} for the formal coincidence of the resulting constraint with the Friedmann equation). Nothing in this section forbids a reader from embedding the frame sequence in a larger theory that has dynamics; it records that this paper is not that theory.
\end{remark}

\section{Density reading and equation of state}\label{sec:density-readout}

\begin{definition}[Hubble parameter of the framework]\label{def:hubble}
\begin{equation}\label{eq:hubble}
H(t) := \frac{\dot R}{R}.
\end{equation}
This is the standard cosmological usage of the symbol, specialised to the present framework: $H$ is a defined ratio, not a dynamical variable, and $H_{0} := H(t_{0})$ will denote its value at the present age $t_{0}$. The dimensionless $h = H_{0} / (100\ \mathrm{km\,s^{-1}\,Mpc^{-1}})$ of observational cosmology is not used in this paper.
\end{definition}

\begin{proposition}[Hubble law of the reading]\label{prop:hubblelaw}
\begin{equation}\label{eq:hubblelaw}
H(t) = \frac{1}{t} \qquad \text{for all } t > 0.
\end{equation}
\uses{\cref{eq:ct}, \cref{def:hubble}}
\end{proposition}

\begin{proof}
By \cref{eq:ct}, $\dot R = c$ and $R = ct$, so $H = c / (ct) = 1/t$; the constant $c$ cancels against itself. There is no exceptional instant: the law holds at every frame and in the continuum reading alike.
\end{proof}

\begin{proposition}[Density of the reading]\label{prop:density}
The homogeneous density implied by \cref{eq:ct} is
\begin{equation}\label{eq:density}
\rho(t) = \frac{M(t)}{V(R(t))} = \frac{c^{3}t / 2G}{\tfrac{4}{3}\pi c^{3} t^{3}} = \frac{3}{8\pi G\, t^{2}}.
\end{equation}
\uses{\cref{eq:ct}, \cref{eq:ballmass}}
\end{proposition}

\begin{proof}
Substitute the trajectory \cref{eq:ct} into the mass and the volume, and cancel the powers of $c$ and $t$ term by term: numerator $c^{3}t/2G$; denominator $\tfrac43\pi c^{3}t^{3}$; quotient $\frac{3 c^{3} t}{8\pi G c^{3} t^{3}} = \frac{3}{8\pi G t^{2}}$, the factors $c^{3}$ cancelling in the first step.
\end{proof}

Two consistency statements follow immediately, and both are worth recording because they are frequently confused with dynamical content.

\begin{proposition}[Constraint self-check]\label{prop:selfcheck}
The density \cref{eq:density} satisfies
\begin{equation}\label{eq:selfcheck}
H^{2} = \frac{8\pi G}{3}\,\rho,
\end{equation}
identically, with both sides equal to $1/t^{2}$.
\uses{\cref{prop:hubblelaw}, \cref{prop:density}}
\end{proposition}

\begin{proof}
Left side: $H^{2} = 1/t^{2}$ by \cref{eq:hubblelaw}. Right side: $\frac{8\pi G}{3}\cdot\frac{3}{8\pi G t^{2}} = \frac{1}{t^{2}}$, every constant cancelling term by term. The equality is an identity of the reading, not a dynamical constraint imposed on it.
\end{proof}

\Cref{eq:selfcheck} has the form of the Friedmann constraint of general relativity, and \cref{ch:gr} is devoted to the precise sense in which the two coincide (formal coincidence, one-way, with the general-relativistic side quoted as an external result). Nothing of that enters here.

\begin{eresult}[Conservation scaling of a perfect fluid]\label{er:conservation}
A homogeneous perfect fluid with constant equation-of-state parameter $w = p/(\rho c^{2})$ dilutes as
\begin{equation}\label{eq:scaling}
\rho \propto a^{-3(1+w)},
\end{equation}
with $a(t)$ the scale factor of the space it fills.
\reference{Standard cosmology; the scaling follows from stress-energy conservation $\dot\rho + 3H(\rho + p/c^{2}) = 0$ with constant $w$, by one quadrature.}
\end{eresult}

\begin{proposition}[Effective equation of state]\label{prop:eos}
With the identification $a \propto t$ (\cref{ch:gr} formalises the bridge; here it is the identity $R = ct$ itself), the dilution law \cref{eq:density} and the scaling \cref{eq:scaling} force
\begin{equation}\label{eq:eos}
-3(1+w) = -2 \quad\Longrightarrow\quad w = -\frac{1}{3},
\end{equation}
and correspondingly $\rho c^{2} + 3p = \rho c^{2}(1 + 3w) = 0$.
\uses{\cref{prop:density}, \cref{er:conservation}}
\end{proposition}

\begin{remark}[Effective, not instantaneous]\label{rem:instantaneous}
The exponent $w = -\tfrac13$ is an \emph{effective} index: it belongs to the dilution law \cref{eq:density} together with the external conservation scaling, not to any material of which the injected content is made. The injected content is deposited at rest, as the shell profile of \cref{ch:boundary} shows; a sourced continuity equation assigns the injected stream its own instantaneous index, and the two indices answer different questions --- what the fluid is, versus how its averaged density falls. The distinction is drawn here in one paragraph and enforced throughout \cref{ch:boundary,ch:gr}: the framework's content is pressureless as deposited material, and $w = -\tfrac13$ is the effective shape of its dilution. Neither statement is weakened by the other.
\end{remark}

\section{Equivalence of the forms, and the integrability discipline}\label{sec:equivalence}

Three propositions close the chapter's mathematical content; a discipline statement closes its method.

\begin{proposition}[The dust ball is the $\mu \to 0$ limit]\label{prop:dustlimit}
At fixed initial content $M_{0} > 0$, the pumped family \cref{eq:generalsolution} converges to the constant-content solution of \cref{sec:newton-ball}:
\begin{equation}\label{eq:dustlimit}
\mu \to 0:\qquad R(t) \;\longrightarrow\; \Bigl(\frac{9 G M_{0}}{2}\Bigr)^{1/3} t^{2/3},
\end{equation}
the $t^{2/3}$ law of the closed ball with the coefficient \cref{eq:dustsolve}.
\uses{\cref{eq:generalsolution}, the expansion \cref{eq:dustexpand}}
\end{proposition}

\begin{proof}
The limit $\mu \to 0$ at fixed $M_{0}$ sends $\varepsilon = \mu t / M_{0} \to 0$ at every fixed $t$, which is precisely the regime of the expansion \cref{eq:dustexpand}; the result was computed there and is quoted here only to record that the limit exists and lands on the closed ball.
\end{proof}

\begin{proposition}[Lattice consistency]\label{prop:latticeconsistency}
The continuous trajectory $R(t) = ct$ and the frame reading $R(n) = n\ell_{\mathrm P}$ agree exactly at the matching times:
\begin{equation}\label{eq:match}
t = n\,t_{\mathrm P}:\qquad R(t) = c\,n\,t_{\mathrm P} = n\,\ell_{\mathrm P} = R(n),
\end{equation}
using $\ell_{\mathrm P} = c\,t_{\mathrm P}$ (\cref{def:units}). Moreover, because $R$ is linear in $n$, every frame-sum of $R$ coincides with its integral exactly: the Euler--Maclaurin corrections, which measure the difference between sum and integral, involve only derivatives of order $\geq 2$ of the summand, and a linear function has none. The consistency is therefore at the level of identities, not of asymptotics.
\uses{\cref{eq:stepping}, \cref{eq:ct}, \cref{def:units}}
\end{proposition}

\begin{corollary}[The closed ball is incompatible with stepping]\label{cor:incompatible}
No constant-content branch $M = \text{const}$ survives the matching: the reading \cref{def:readout} with stepping gives $M(n) = n\,c^{3}t_{\mathrm P}/2G \propto n$ (\cref{prop:readoutvalue}), while a constant content is independent of $n$; the two conditions are contradictory at every frame from $n = 2$ on. The compatibility of the Newtonian energetics with the frame sequence is carried exclusively by the pumped, linearly growing content.
\uses{\cref{prop:readoutvalue}, \cref{as:stepping}}
\end{corollary}

\begin{remark}[An independent check, recorded]\label{rem:triangulation}
The value $\mu = c^{3}/2G$ was fixed in \cref{sec:frame-readout} from the stepping matching alone. An independent consistency check on the resulting content $E = c^{5}t/2G$ will be constructed in \cref{ch:boundary}, where the entropy and temperature of the boundary are built from separate inputs and the identity $E = TS$ is verified rather than assumed --- a consistency check, not a premise. The dependence is one-way: nothing in this chapter uses the boundary construction.
\end{remark}

\subsection{The integrability discipline}

\begin{remark}[What may be integrated]\label{rem:integrability}
Two kinds of quantity appear in this paper, and the operations they admit differ. A \emph{state quantity} is defined at frame level and additive over frames: radius, content, entropy below --- summing its frame values is an operation of the frame sequence itself. A \emph{ratio quantity} is a quotient of such quantities per frame: the Hubble parameter $H = 1/t$ is the example in hand. The discipline: \textbf{state quantities may be integrated; ratio quantities may not}, because integrating a ratio manufactures an object of the continuum reading that has no frame-level meaning --- an accumulation of rates rather than of contents. The test is exact: $X$ is integrable on the frame sequence iff there is a frame-defined $Y$ with $X = \Delta Y / \Delta n$ identically, for then $\sum X\,\Delta n$ telescopes into $Y$; otherwise the formal integral is continuum-level shorthand and is flagged as such wherever it appears. The example makes the test concrete: $H(n) = 1/(n t_{\mathrm P})$ fails it, since
\begin{equation}\label{eq:lnfails}
\frac{\Delta \ln n}{\Delta n} = \ln\Bigl(1 + \frac{1}{n}\Bigr) \;\neq\; \frac{1}{n} \qquad (n \geq 1),
\end{equation}
the inequality being strict because $\ln(1 + x) < x$ for $x > 0$; $\ln t$ is thus a continuum object, not a frame accumulation, and no formula below integrates $H$. The integrals actually performed in this chapter --- of $R$ in \cref{sec:pumped-family,sec:newton-ball} --- are integrals of the state quantity $R$ and are exact by linearity; the bridge between frame-sums and integrals is $\Delta n / \Delta t = 1/t_{\mathrm P}$ (\cref{eq:tick}). The discipline binds every later chapter.
\end{remark}

\begin{remark}[What Chapter 4 has established]\label{rem:ch4summary}
The stepping structure (\cref{as:stepping}) and its exact rate $c$ (\cref{prop:uniform-rate}); the delay it forces on every cross-frame response (\cref{cor:delay}), discharging the forward references of \cref{ch:cascade}; the Newtonian ball with its three honest branches (\cref{eq:branches}) and the selection of $E = 0$ by the matching (\cref{prop:compatibility}); the pumped family, its dust and linear limits, and the fixing of the injection rate to $\mu = c^{3}/2G$; the power $\mathcal P = c^{5}/2G$ and the per-frame injection $J = \varepsilon_{\mathrm P}/2$, defined and awaiting identification with the chain's $J$ in \cref{ch:unitarity}; the density reading and its effective equation of state $w = -\tfrac13$; and the integrability discipline that governs what may be done with all of the above. Open items, each with a named successor: the mechanism of the source (open, \cref{rem:mechanismregistry}); the boundary identifications of the content and the residual (\cref{ch:boundary}); the formal coincidence with the Friedmann equation (\cref{ch:gr}); and the value of $J$ as a theorem about the chain (\cref{ch:unitarity}). The chapter used five external results, all cited; it added two members to the hypothesis layer (\cref{as:stepping,as:pump}); and it introduced no dynamics of its own.
\end{remark}

\part{Unitarity and openness}

\chapter{Unitarity and the open system}\label{ch:unitarity}

The framework carries two evolution laws that superficially contradict each other: a \emph{unitary} quantum layer, in which norms are conserved, and an \emph{open} content layer, in which the total grows by a constant injection every frame. This chapter shows the two are compatible --- because they act on different objects --- and then proves that the open layer's constant increment is not free: it is exactly one minimal action quantum per frame. This is the value fixed by the injection theorem, and it identifies the two quantities named $J$ of \cref{rem:tworeadings-j}.

\section{The two layers and their symbols}\label{sec:two-layers}

\begin{definition}[Quantum layer]\label{def:quantum-layer}
The quantum layer consists of a state vector $\Psi_{n} \in \mathcal{H}$ at each frame, with
\begin{equation}\label{eq:quantum-layer}
\| \Psi_{n} \| = \text{const},
\end{equation}
and evolution by unitary steps $\Psi_{n+1} = U_{n}\Psi_{n}$ whose form (Cayley steps of a Hermitian generator) and verification are the subject of \cref{sec:cayley-audit}. The symbol $\Psi$ is reserved: whatever a unitary step acts on is a quantum state.
\end{definition}

\begin{definition}[Content layer]\label{def:content-layer}
The content layer is the classical affine flow of \cref{ch:cascade}: the vector $\boldsymbol{\mathcal E}^{(n)} = \bigl(E_{\mathrm B}^{(n)}, E_{\mathrm r}^{(n)}, E_{\mathrm m}^{(n)}\bigr)^{\mathrm T}$ with the chain equations \cref{eq:cascade} and constant injection \cref{eq:injection}. Its evolution is not norm-preserving --- the total $N(n)$ grows by $J$ per frame (\cref{prop:total}) --- and it has no amplitude interpretation: its components are energy contents, not probability amplitudes.
\end{definition}

\begin{proposition}[Compatibility of the layers]\label{prop:layers}
The two layers do not conflict, because they act on different objects: $\Psi_{n}$ is a vector in a Hilbert space with a norm, $\boldsymbol{\mathcal E}^{(n)}$ is a vector in $\mathbb{R}^{3}$ with a sum. A conflict between ``unitary'' and ``growing'' appears only when one symbol is made to do both jobs, and the split \cref{def:quantum-layer,def:content-layer} removes precisely that ambiguity.
\uses{\cref{def:quantum-layer}, \cref{def:content-layer}}
\end{proposition}

The relation between the layers is stated once and flagged as borrowed:

\begin{eresult}[Measurement bridge]\label{er:bridge}
Energy readings of the quantum layer reproduce the contents: for the observable $\hat E_{i}$ associated with component $i$,
\begin{equation}\label{eq:measurement-bridge}
E_{i} = \langle \Psi_{n} | \hat E_{i} | \Psi_{n} \rangle .
\end{equation}
\reference{Born's measurement postulate; the framework does not derive it and does not extend it. The bridge is used exactly where a content must be read from the quantum layer, and nowhere else.}
\end{eresult}

\begin{eresult}[Stinespring dilation theorem]\label{er:stinespring}
Every completely positive trace-preserving map on a system arises as the reduction of a unitary evolution of the system tensor an environment: non-unitarity of the reduced dynamics is the signature of openness, not a defect of the dynamics.
\reference{W.~F. Stinespring, \emph{Positive functions on $C^{*}$-algebras}, Proc.\ Amer.\ Math.\ Soc.\ \textbf{6}, 211 (1955). The content layer's explicit environment is constructed in \cref{ch:dilation}, and the existence statement here is quoted, not used, until then.}
\end{eresult}

\begin{remark}[Symbol discipline]\label{rem:symbol-split}
The split is enforced as a symbol discipline: $\Psi$ for whatever unitary steps act on, $\boldsymbol{\mathcal E}$ for whatever the chain acts on, and no shared symbols between them. The contents are never square-rooted; amplitudes are never summed as energies: the two objects are of different types, and the notation keeps the types apart.
\end{remark}

\begin{remark}[What unitarity does and does not guard]\label{rem:unitarity-scope}
The quantum layer's constant norm (\cref{eq:quantum-layer}) is a statement about $\Psi$, and it survives everything this paper does to $\boldsymbol{\mathcal E}$: the growth of $N(n)$ is not a violation of unitarity but a property of a different object. The framework-wide statement, proved in \cref{sec:injection-theorem}, is sharper: the growth is not arbitrary --- it is quantised in units of $\alpha_{0}$ per frame. Unitarity guards the amplitudes; the injection theorem guards the energy balance; and the two never guard the same thing twice.
\end{remark}

\section{The injection theorem}\label{sec:injection-theorem}

\begin{theorem}[Injection theorem]\label{thm:injection}
Per frame, the growth of the total content carries exactly one minimal action quantum:
\begin{equation}\label{eq:inj-theorem}
\Delta E \cdot t_{\mathrm P} = \alpha_{0} = \frac{\hbar}{2},
\qquad
\Delta E := N(n+1) - N(n).
\end{equation}
\uses{the linking premise \cref{eq:link} below; \cref{prop:total} (step 1); \cref{def:power} (step 2); \cref{def:units} and \cref{def:alpha0} (step 3)}
\end{theorem}

The linking premise is stated separately, because it is an identification joining two definitions and deserves to be visible as such:
\begin{equation}\label{eq:link}
\begin{gathered}
\text{the injection constant $J$ of \cref{as:b3}}\\
\text{is the one-tick content of the source power $\mathcal P$ of \cref{def:power}.}
\end{gathered}
\end{equation}
\cref{eq:link} is one-way: it says what the chain is fed by, and it will never be used in the reverse direction to characterise the source through the chain.

\begin{proof}
\emph{Step 1 (increment equals the injection).} By \cref{prop:total}, $\Delta E = N(n+1) - N(n) = J$ for every frame, independently of $n$ and of the contents themselves: the conducted terms cancel pairwise in the frame sum, and only the injection survives (\cref{eq:telescope}).

\emph{Step 2 (injection in terms of power).} By \cref{eq:link} and the definition of the per-frame injection (\cref{def:power}),
\begin{equation}\label{eq:step2}
\Delta E = J = \mathcal P\, t_{\mathrm P} = \frac{c^{5}}{2G}\, t_{\mathrm P}.
\end{equation}

\emph{Step 3 (multiply by one tick, and reduce).}
\begin{equation}\label{eq:step3}
\Delta E \cdot t_{\mathrm P} = \mathcal P\, t_{\mathrm P}^{2} = \frac{c^{5}}{2G}\cdot\frac{\hbar G}{c^{5}} = \frac{\hbar}{2} = \alpha_{0},
\end{equation}
where $t_{\mathrm P}^{2} = \hbar G / c^{5}$ is the definition of $t_{\mathrm P}$ squared (\cref{def:units}), and the reduction $c^{5}\cdot\hbar G/(2G c^{5})$ displays the cancellation of $c^{5}$ and $G$ explicitly.
\end{proof}

\begin{corollary}[The value of $J$]\label{cor:twoj}
The injection constant is fixed at
\begin{equation}\label{eq:jpinned}
J = \frac{\varepsilon_{\mathrm P}}{2},
\end{equation}
and the exact solution \cref{eq:solutions} takes its energy-unit form
\begin{equation}\label{eq:epsform}
E_{\mathrm B}^{(n)} = \frac{n\,\varepsilon_{\mathrm P}}{4},\qquad
E_{\mathrm r}^{(n)} = \frac{n\,\varepsilon_{\mathrm P}}{8},\qquad
E_{\mathrm m}^{(n)} = \frac{(n-2)\,\varepsilon_{\mathrm P}}{8},
\end{equation}
obtained from \cref{eq:solutions} by the single substitution $J = \varepsilon_{\mathrm P}/2$, applied once and recorded here once.
\uses{\cref{thm:injection}, \cref{eq:reducefinal}}
\end{corollary}

\begin{proof}
From \cref{eq:step3}, $J = \Delta E = \alpha_{0}/t_{\mathrm P} = (\hbar/2)/t_{\mathrm P} = \hbar/2t_{\mathrm P} = \varepsilon_{\mathrm P}/2$, the last equality by \cref{cor:epmp}. The forms \cref{eq:epsform} follow term by term: $nJ/2 = n\varepsilon_{\mathrm P}/4$, and so on for each row.
\end{proof}

\Cref{eq:jpinned} discharges the open magnitude of \cref{as:b3} (\cref{rem:jform}): the value was not chosen in \cref{ch:cascade}, and it is not chosen here --- it is derived, once, from premises each of which was established earlier.

\begin{corollary}[Realisation of the smallest quantum number]\label{cor:quantum-realized}
Per frame, the content layer realises exactly the $k = 1$ quantum of \cref{eq:quantization}: $\Delta E = \tfrac{1}{2}\varepsilon_{\mathrm P}$ is the smallest admissible value of \cref{eq:quantization}, and it occurs at every frame.
\uses{\cref{thm:injection}, \cref{as:a2}}
\end{corollary}

\begin{remark}[The meeting with the zero-point cell]\label{rem:meeting}
The number $\varepsilon_{\mathrm P}/2$ now stands at the end of two derivation chains: the oscillator chain (\cref{thm:spectrum} with \cref{as:p}: $E_{0} = \alpha_{0}\omega_{0} = \varepsilon_{\mathrm P}/2$) and the injection chain (\cref{eq:step3}). The meeting is recorded as a consistency of the framework, with the shared premise declared as in \cref{rem:zeropoint-cell}: both chains are expressed in the units of \cref{sec:units}, and independence is claimed only for the derivations, not for the unit system. What the meeting adds beyond either chain alone: the zero-point floor of the reference oscillator and the per-frame growth of the content layer are the same quantity read in two different languages --- the $k = 1$ cell is realised by both, simultaneously and by construction.
\end{remark}

\begin{remark}[Reading of the theorem]\label{rem:injection-reading}
The natural unit of per-frame growth is action: the content layer admits no smaller nonzero step than $\alpha_{0}$, and the energy reading $\varepsilon_{\mathrm P}/2$ is that step expressed in energy units through $\alpha_{0}/t_{\mathrm P}$. The reading is exact and conservative: no statement is made about \emph{why} the source injects at this rate (\cref{rem:register}), only that whatever the source is, its one-tick content is one minimal action quantum --- the strongest non-mechanistic statement the framework permits itself about the injection.
\end{remark}

\section{The Cayley step, verified}\label{sec:cayley-audit}

The quantum layer's evolution law (\cref{def:quantum-layer}) is now specified and checked. The check matters because the framework leans on this step being \emph{exactly} unitary, and because the degree to which the step is unique is a real question with a two-sided answer (\cref{prop:unique}, \cref{eq:pade22}).

\begin{definition}[Cayley step]\label{def:cayley}
Let $\hat H$ be Hermitian (the generator of the frame) and define
\begin{equation}\label{eq:cayleyparam}
\vartheta := \frac{\hat H\, t_{\mathrm P}}{2\hbar} \qquad (\text{Hermitian}),
\qquad
U := (1 - i\vartheta)(1 + i\vartheta)^{-1}.
\end{equation}
\end{definition}

\begin{proposition}[Unitarity of the step]\label{prop:unitarity}
$U^{\dagger} U = \mathbb{1}$.
\uses{\cref{eq:cayleyparam} (Hermiticity of $\vartheta$ only)}
\end{proposition}

\begin{proof}
Both factors of $U$ are functions of the single Hermitian operator $\vartheta$ --- one a polynomial, one a resolvent --- and functions of one operator commute. Hence the adjoint inverts the order of the commuting factors, and the factors cancel pairwise:
\begin{equation}\label{eq:unitarity}
\begin{gathered}
U^{\dagger}U = \bigl[(1 + i\vartheta)^{-1}\bigr]^{\dagger}(1 - i\vartheta)^{\dagger}(1 - i\vartheta)(1 + i\vartheta)^{-1}\\
= (1 + i\vartheta)^{-1}(1 + i\vartheta)\,(1 - i\vartheta)(1 - i\vartheta)^{-1} = \mathbb{1},
\end{gathered}
\end{equation}
where the Hermiticity of $\vartheta$ was used in the second equality to turn adjoints into sign flips of $i$. The resolvent exists: $1 + i\vartheta$ is invertible because the spectrum of $i\vartheta$ is purely imaginary and never equals $-1$.
\end{proof}

\begin{remark}[Direction guard]\label{rem:direction-guard}
The proof of \cref{prop:unitarity} uses nothing but the Hermiticity of $\vartheta$: no evolution equation, no Schr\"odinger dynamics, no external benchmark. The unitarity of the Cayley step is therefore not borrowed from quantum mechanics --- it is algebra --- and conversely the step is \emph{not} derived from any physical evolution law. The comparison with the exponential below is one-way, as the rule of \cref{rem:register} requires: a circle of exactly the kind refused in \cref{er:efreq} (an earlier result quoted as premise for its own ancestor) is not merely avoided here, it is structurally absent.
\end{remark}

\begin{proposition}[Uniqueness among degree-one rational steps]\label{prop:unique}
Among rational maps $R(\vartheta) = P(\vartheta)/Q(\vartheta)$ with $\deg P, \deg Q \leq 1$, the conditions
\begin{equation}\label{eq:threeconds}
\text{(i) exact unitarity } |R(i\theta)| = 1 \ \forall \theta \in \mathbb{R},
\qquad
\text{(ii) matching } e^{-2i\vartheta} \text{ through } O(\vartheta^{2}),
\end{equation}
select $R(\vartheta) = (1 - i\vartheta)/(1 + i\vartheta)$ uniquely.
\uses{\cref{eq:cayleyparam}; pure algebra}
\end{proposition}

\begin{proof}
Normalise the constant term: $P = a(1 + b\vartheta)$, $Q = a(1 + d\vartheta)$ with $a \neq 0$; the overall scale cancels from \cref{eq:threeconds}. Matching to $e^{-2i\vartheta} = 1 - 2i\vartheta + O(\vartheta^{2})$ gives one complex condition,
\begin{equation}\label{eq:match-cayley}
(1 + b\vartheta)(1 + d\vartheta)^{-1} = 1 + (b - d)\,\vartheta + O(\vartheta^{2})
\;\Longrightarrow\; b - d = -2i .
\end{equation}
Exact unitarity on the imaginary axis, $|1 + ib\theta| = |1 + id\theta|$ for all real $\theta$, requires the moduli-squared to agree at every order:
\begin{equation}\label{eq:moduli}
1 + (ib + id^{*})\theta - b d^{*}\theta^{2} = \cdots \ \text{real and even in } \theta
\;\Longrightarrow\; b = \pm d ,
\end{equation}
the sign coming from the $\theta^{2}$-coefficient. The branch $b = +d$ gives $b - d = 0 \neq -2i$, contradicting \cref{eq:match-cayley}; the branch $b = -d$ gives $-2d = -2i$, i.e.\ $d = i$, $b = -i$, and $R(\vartheta) = (1 - i\vartheta)/(1 + i\vartheta)$ --- the Cayley step. Both branches were written out, so the loss of either is visible.
\end{proof}

\Cref{prop:unique} carries its degree restriction in its assumptions, where it belongs: without it, the conditions admit higher-order examples such as \cref{eq:pade22}, and uniqueness fails.

\begin{remark}[Approximation status, convention fixed]\label{rem:approx}
The physical evolution of standard quantum mechanics over one tick is $e^{-i\hat H t_{\mathrm P}/\hbar} = e^{-2i\vartheta}$ \emph{at the level of the external kinematics this paper quotes but does not derive}. The Cayley step is the $[1,1]$ Pad\'e approximant of $e^{-2i\vartheta}$:
\begin{equation}\label{eq:pade11}
\frac{1 - i\vartheta}{1 + i\vartheta} = 1 - 2i\vartheta - 2\vartheta^{2} + 2i\vartheta^{3} + O(\vartheta^{4}),
\qquad
e^{-2i\vartheta} = 1 - 2i\vartheta - 2\vartheta^{2} + \tfrac{4}{3}i\vartheta^{3} + O(\vartheta^{4}),
\end{equation}
matching through $\vartheta^{2}$; the per-step discrepancy first appears at $O(\vartheta^{3})$ in the map, i.e.\ $O(\vartheta^{2})$ in the exponent, and over a fixed elapsed time (a fixed number of steps per unit time) the accumulated discrepancy is $O(\vartheta^{2})$. Higher-order exactly-unitary steps exist: the $[2,2]$ Pad\'e $P/Q$ with
\begin{equation}\label{eq:pade22}
P = 1 - i\vartheta - \frac{\vartheta^{2}}{3}, \qquad Q = 1 + i\vartheta - \frac{\vartheta^{2}}{3},
\qquad
|P(i\theta)|^{2} = 1 + \frac{\theta^{2}}{3} + \frac{\theta^{4}}{9} = |Q(i\theta)|^{2} \ \forall\theta,
\end{equation}
is exactly unitary and matches through $O(\vartheta^{4})$; the equality in \cref{eq:pade22} is verified by expanding both moduli-squared and observing they coincide term by term. Two notes close the check. First, the convention: throughout this section $\vartheta = \hat H t_{\mathrm P}/2\hbar$ and the expansion is that of $e^{-2i\vartheta}$, giving the coefficients $1/3$, $1/9$ of \cref{eq:pade22}; in the alternative convention $\theta = 2\vartheta$ (exponential $e^{-i\theta}$) the same series reads $1 + \theta^{2}/12 + \theta^{4}/144$ --- the two are one change of parameter, and both appear in the literature. Second, the existence of \cref{eq:pade22} confirms that \cref{prop:unique} cannot be strengthened beyond degree one: the Cayley step is the unique \emph{lowest-order} unitary rational step, not the unique unitary step. The framework's use of the Cayley step is thereby an explicit choice of the lowest order, with higher orders available and not taken.
\end{remark}

\begin{remark}[Time-dependent generators]\label{rem:timedependent}
If the generator is $\hat H(n)$, Hermitian at each frame, then each frame's Cayley step is legitimate by \cref{prop:unitarity}, and the evolution over $n$ frames is the ordered product $U_{n}\cdots U_{1}$, the discrete time-ordered product \emph{as in standard kinematics}. This is compatible with the content layer's growth: unitarity guards the amplitudes of the quantum layer at every frame, while the content layer grows --- the two statements concern different objects (\cref{prop:layers}), and the system-side non-conservation together with the expansion-side conservation is examined in \cref{ch:symmetries}, not here.
\end{remark}

\begin{remark}[Counterexample: explicit Euler]\label{rem:euler}
The naive step $U = 1 - i\vartheta$ (explicit Euler) fails the check immediately:
\begin{equation}\label{eq:euler}
|1 - i\vartheta|^{2} = (1 - i\vartheta)(1 + i\vartheta) = 1 + \vartheta^{2} > 1,
\end{equation}
the norm grows at every nonzero step, and the map is not unitary. The framework does not adopt it; the record of its failure is kept so that ``exactly unitary'' above is read as a discriminating property of the Cayley choice, not a property of any rational step.
\end{remark}

\section{The transfer kernel is a legitimate channel}\label{sec:cptp}

The content layer is affine: a linear transfer plus a constant source. This section checks the linear part with the tools of open quantum systems, and records what the check does and does not cover.

\begin{definition}[Transfer matrix]\label{def:transfer}
The homogeneous part of the chain \cref{eq:deviation} acts through
\begin{equation}\label{eq:transfermatrix}
\mathbb{M}_{n} =
\begin{pmatrix}
1 - \frac{1}{n} & 0 & 0\\[4pt]
\frac{1}{n} & 1 - \frac{1}{n} & 0\\[4pt]
0 & \frac{1}{n} & 1
\end{pmatrix},
\qquad
\boldsymbol{\delta}^{(n+1)} = \mathbb{M}_{n}\,\boldsymbol{\delta}^{(n)},
\end{equation}
and the full chain reads $\boldsymbol{\mathcal E}^{(n+1)} = \mathbb{M}_{n}\,\boldsymbol{\mathcal E}^{(n)} + \boldsymbol{s}$ with the constant source $\boldsymbol{s}$ of \cref{eq:injection}.
\end{definition}

\begin{lemma}[Column-stochasticity]\label{lem:stochastic}
Every column of $\mathbb{M}_{n}$ is a probability vector: entries non-negative and
\begin{equation}\label{eq:stochastic}
\begin{gathered}
\text{column 1: } \Bigl(1 - \frac{1}{n}\Bigr) + \frac{1}{n} + 0 = 1,
\qquad
\text{column 2: } 0 + \Bigl(1 - \frac{1}{n}\Bigr) + \frac{1}{n} = 1,\\
\text{column 3: } 0 + 0 + 1 = 1.
\end{gathered}
\end{equation}
\uses{\cref{eq:transfermatrix}}
\end{lemma}

\begin{proof}
Each sum is written out rather than inferred: the first two columns pair the retention factor $1 - \tfrac1n$ with the conducted share $\tfrac1n$, which sum to $1$ by \cref{eq:share}; the third column is the terminal node, which retains all. Non-negativity holds for $n \geq 1$ since $1 - \tfrac1n \geq 0$ and $\tfrac1n > 0$.
\end{proof}

\begin{lemma}[Kraus form of the transfer]\label{lem:kraus}
The map on operators defined by
\begin{equation}\label{eq:kraus}
K_{ji} = \sqrt{(\mathbb{M}_{n})_{ji}}\; |j\rangle\langle i|, \qquad i, j \in \{\mathrm B, \mathrm r, \mathrm m\},
\end{equation}
is a completely positive trace-preserving channel on the states diagonal in the content basis: trace preservation follows from \cref{lem:stochastic},
\begin{equation}\label{eq:complete}
\sum_{i,j} K_{ji}^{\dagger} K_{ji}
= \sum_{i,j} (\mathbb{M}_{n})_{ji}\; |i\rangle\langle i|
= \sum_{i} \Bigl( \sum_{j} (\mathbb{M}_{n})_{ji} \Bigr) |i\rangle\langle i|
= \sum_{i} |i\rangle\langle i| = \mathbb{1},
\end{equation}
where the reordering of the double sum --- Kraus adjoints reassemble by input index $i$ --- is the only step, and each bracketed column sum is $1$ by \cref{eq:stochastic}. Complete positivity is immediate from the Kraus form itself. The action on a diagonal operator reproduces the transfer: $K_{ji} |i\rangle\langle i| K_{ji}^{\dagger} = (\mathbb{M}_{n})_{ji} |j\rangle\langle j|$, so diagonal states evolve by the stochastic matrix and off-diagonal terms are erased --- the transfer, as a quantum process, is exactly the classical content flow.
\uses{\cref{lem:stochastic}}
\end{lemma}

\begin{remark}[Scope of the check]\label{rem:cptp-scope}
Three boundaries are stated once. (i) The channel \cref{eq:kraus} covers the \emph{linear transfer} only; the source $\boldsymbol{s}$ is an affine constant, not a channel --- it grows the trace by design and has no Kraus representation. (ii) The complete positivity of \cref{eq:kraus} is certified by its Kraus form, which is to say the check is self-contained: no external theorem is invoked. (iii) The environment whose reduction would generate this channel exists by \cref{er:stinespring}; its explicit, minimal construction is \cref{ch:dilation}, and nothing in this section presupposes it.
\end{remark}

\begin{remark}[What Chapter 5 has established]\label{rem:ch5summary}
The two layers of the framework are compatible, and in what precise sense: distinct objects (\cref{prop:layers}), distinct symbols (\cref{rem:symbol-split}), connected by one borrowed bridge (\cref{er:bridge}). The injection theorem (\cref{thm:injection}): the per-frame growth of the content layer is exactly one minimal action quantum, fixing $J = \varepsilon_{\mathrm P}/2$ (\cref{cor:twoj}) and realising the $k = 1$ quantum of \cref{eq:quantization} at every frame (\cref{cor:quantum-realized}). The quantum step is exactly unitary (\cref{prop:unitarity}), uniquely so among degree-one rational steps (\cref{prop:unique}), with the approximation status fixed on a single convention (\cref{rem:approx}), and time-dependent generators admitted frame by frame (\cref{rem:timedependent}). The internal transfer of the content layer is a certified channel (\cref{lem:kraus}). Two questions receive two-sided answers: the freedom structure of the environment, settled by explicit construction in \cref{ch:dilation}; and the forcedness of the Cayley step --- forced at the lowest order (\cref{prop:unique}), not forced in general (\cref{eq:pade22}).
\end{remark}

\chapter{The explicit dilation}\label{ch:dilation}

What environment does the content layer's non-unitarity imply? This chapter answers constructively: an explicit environment of minimal size is built, its minimality is proved, and the construction is shown to leave no unnamed degrees of freedom. The chapter's first task is to say precisely what kind of dilation the content layer requires, because the answer is classical, and the classicality is itself a structural fact.

\section{The content layer is a classical affine flow}\label{sec:classical-flow}

\begin{proposition}[Positivity of the flow]\label{prop:positivity}
If $\boldsymbol{\mathcal E}^{(n)}$ has non-negative components, so does $\boldsymbol{\mathcal E}^{(n+1)} = \mathbb{M}_{n}\,\boldsymbol{\mathcal E}^{(n)} + \boldsymbol{s}$. The positive octant is invariant: energy contents never become negative at any frame.
\uses{\cref{eq:transfermatrix}, \cref{eq:injection}, \cref{lem:stochastic}}
\end{proposition}

\begin{proof}
Each component of $\boldsymbol{\mathcal E}^{(n+1)}$ is a sum of products of non-negative quantities: the entries of $\mathbb{M}_{n}$ are non-negative for $n \geq 1$ ($1 - \tfrac1n \geq 0$ and $\tfrac1n > 0$, by \cref{eq:transfermatrix}), the source $\boldsymbol{s} = (J, 0, 0)^{\mathrm T}$ is non-negative, and the components of $\boldsymbol{\mathcal E}^{(n)}$ are non-negative by hypothesis. Every term of every sum is non-negative, so every component of $\boldsymbol{\mathcal E}^{(n+1)}$ is non-negative. The induction starts at the boundary condition \cref{eq:firstframe}, whose components are non-negative.
\end{proof}

\begin{remark}[What classicality means here]\label{rem:what-dilation}
No operation of this paper ever forms a superposition of content vectors: the chain acts on $\mathbb{R}^{3}$, not on a Hilbert space, and the components are measured contents (\cref{er:bridge}), not amplitudes. The absence of phase, interference and linearity-over-$\mathbb{C}$ is therefore not an approximation or a restriction to be relaxed --- it is the type of the object. Two consequences follow. First, the dilation question for the content layer as a whole is a question of \emph{classical} probability: how a Markov flow with a constant source is represented as the marginal of a larger stochastic system. Second, the quantum dilation machinery (Kraus operators, environments, tensor products) is needed only where a genuinely quantum object is present --- and there is exactly one such object in this layer, the transfer channel \cref{eq:kraus}. The chapter's work splits accordingly: a classical construction for the affine flow, and a quantum construction, already verified in \cref{lem:kraus}, for the transfer alone.
\end{remark}

\begin{remark}[Reading of the classicality]\label{rem:classicality-reading}
A conservative reading, per \cref{rem:register}: the content layer is classical because it \emph{records} --- its components are the values a measurement would return, and records do not superpose. Nothing here asserts that the world behind the records is classical; the framework is silent about that, in both directions. What is asserted is only that no formula of this paper requires the content layer to carry quantum structure, and that the dilation below respects this by being a classical construction with one precisely located quantum corner (\cref{sec:transfer-dilation}).
\end{remark}

\section{The explicit quantum dilation of the transfer}\label{sec:transfer-dilation}

The one quantum object of the content layer is the transfer channel \cref{eq:kraus}. This section builds its dilation explicitly --- one fresh register per frame --- and verifies the two properties the dilation must have: the reduced dynamics must reproduce the transfer, and the extended map must be an isometry. Only one earlier fact is used, the column-stochasticity of \cref{lem:stochastic}; nothing below is borrowed from the later chapters, and nothing is assumed about the register beyond its definition.

\begin{definition}[Frame register and dilation isometry]\label{def:isometry}
Let $\mathcal H_{\mathrm s}$ carry the content basis $\{|i\rangle\}$, $i \in \{\mathrm B, \mathrm r, \mathrm m\}$, and let frame $n$ carry a register $\mathcal H_{e_{n}}$ with an orthonormal family of record states $\{|ji\rangle\}$, one for each ordered pair $(j, i)$. The \emph{dilation isometry} is
\begin{equation}\label{eq:isometry}
V_{n}^{\mathrm{tr}} |i\rangle = \sum_{j} \sqrt{(\mathbb{M}_{n})_{ji}}\; |j\rangle_{\mathrm s} \otimes |ji\rangle_{e_{n}}.
\end{equation}
\end{definition}

\begin{proposition}[$V_{n}^{\mathrm{tr}}$ is an isometry]\label{prop:isometry}
$(V_{n}^{\mathrm{tr}})^{\dagger} V_{n}^{\mathrm{tr}} = \mathbb{1}_{\mathrm s}$.
\uses{\cref{eq:isometry}, \cref{lem:stochastic}}
\end{proposition}

\begin{proof}
Compute the matrix element directly, using the orthonormality of the record states in both labels:
\begin{equation}\label{eq:vdagv}
\begin{aligned}
\langle i' | (V_{n}^{\mathrm{tr}})^{\dagger} V_{n}^{\mathrm{tr}} | i \rangle
&= \sum_{j, j'} \sqrt{(\mathbb{M}_{n})_{ji}\,(\mathbb{M}_{n})_{j'i'}}\;\; \langle j' | j \rangle_{\mathrm s}\; \langle j' i' | j i \rangle_{e_{n}}\\
&= \sum_{j, j'} \sqrt{(\mathbb{M}_{n})_{ji}\,(\mathbb{M}_{n})_{j'i'}}\;\; \delta_{j' j}\; \delta_{j' j}\,\delta_{i' i}\\
&= \delta_{i' i} \sum_{j} (\mathbb{M}_{n})_{ji}
= \delta_{i' i},
\end{aligned}
\end{equation}
where the second equality uses $\langle j'|j\rangle_{\mathrm s} = \delta_{j'j}$ and $\langle j'i'|ji\rangle = \delta_{j'j}\delta_{i'i}$ (orthonormal record family), the third collapses the double sum over $j'$, and the last equality is the column-stochasticity of \cref{lem:stochastic} --- each column sums to $1$. The result is the identity matrix, which is the claim.
\end{proof}

\begin{proposition}[Reduction to the transfer]\label{prop:reduction}
For every input operator, tracing out the register reproduces the transfer channel \cref{eq:kraus}:
\begin{equation}\label{eq:reduction}
\operatorname{Tr}_{e_{n}}\Bigl[ V_{n}^{\mathrm{tr}} |i\rangle\langle i'| (V_{n}^{\mathrm{tr}})^{\dagger} \Bigr]
= \delta_{i i'} \sum_{j} (\mathbb{M}_{n})_{ji}\; |j\rangle\langle j|.
\end{equation}
\uses{\cref{eq:isometry}, \cref{lem:stochastic}}
\end{proposition}

\begin{proof}
Expand the product from \cref{eq:isometry}:
\begin{equation}\label{eq:dilation-expand}
V_{n}^{\mathrm{tr}} |i\rangle\langle i'| (V_{n}^{\mathrm{tr}})^{\dagger}
= \sum_{j, j'} \sqrt{(\mathbb{M}_{n})_{ji}\,(\mathbb{M}_{n})_{j'i'}}\; |j\rangle\langle j'|_{\mathrm s} \otimes |ji\rangle\langle j'i'|_{e_{n}}.
\end{equation}
The partial trace over the register kills every tensor factor whose record states are orthogonal:
\begin{equation}\label{eq:ptrace}
\operatorname{Tr}_{e_{n}}\bigl( |ji\rangle\langle j'i'| \bigr) = \langle j'i' | j i \rangle = \delta_{j' j}\,\delta_{i' i},
\end{equation}
so only terms with $j' = j$ and $i' = i$ survive, and their square roots become the entries themselves:
\begin{equation}\label{eq:collapse}
\operatorname{Tr}_{e_{n}}[\cdots] = \delta_{i i'} \sum_{j} (\mathbb{M}_{n})_{ji}\; |j\rangle\langle j|.
\end{equation}
For $i = i'$ this is the stochastic transition of \cref{lem:stochastic} on the diagonal state $|i\rangle\langle i|$; for $i \neq i'$ the result is zero, so off-diagonal input terms are erased. Comparing with \cref{eq:kraus}: $K_{ji}|i\rangle\langle i|K_{ji}^{\dagger} = (\mathbb{M}_{n})_{ji}|j\rangle\langle j|$, and the off-diagonal erasure is the same. The reduced dynamics is the transfer channel, exactly.
\end{proof}

\begin{remark}[What the register records]\label{rem:record}
The record state $|ji\rangle$ says: at frame $n$, content moved from $i$ to $j$. Because the family is orthonormal, different records are perfectly distinguishable: the environment ends each frame holding the complete classical record of the transition that occurred, and no interference between alternatives survives the partial trace --- which is why the quantum dilation reproduces a \emph{classical} transfer. Retaining all registers, the trajectory of the system together with the string of records is a classical joint Markov chain, and its consistency under composition is the cocycle law \cref{eq:cocycle} already proved at the level of the kernel: composing two frames of records composes their transitions, with no new content.
\end{remark}

\begin{remark}[Freshness of the registers]\label{rem:fresh}
Each frame receives a \emph{fresh} register: the definition \cref{eq:isometry} couples the system only to $\mathcal H_{e_{n}}$ at frame $n$, and no register is reused. This is the Markov property in the environment language --- the transfer has no memory beyond the current content, so its dilation needs none --- and the resulting picture, a system colliding with one fresh register per frame, is developed as the input--output reading in \cref{sec:collision}. The source $\boldsymbol{s}$ is not part of this dilation: it is a classical constant of the affine flow, treated separately in \cref{sec:source-branch}.
\end{remark}

\section{Minimality: the environment has dimension five}\label{sec:minimality}

The dilation \cref{eq:isometry} uses one register state per transition that can actually occur. This section proves that nothing smaller can work, and that the bound is achieved exactly.

\begin{lemma}[Nonzero pattern of the transfer]\label{lem:pattern}
For $n \geq 2$, exactly five entries of $\mathbb{M}_{n}$ are non-zero:
\begin{equation}\label{eq:pattern}
\begin{gathered}
(\mathbb{M}_{n})_{\mathrm{BB}} = 1 - \frac{1}{n},\quad
(\mathbb{M}_{n})_{\mathrm{rB}} = \frac{1}{n},\quad
(\mathbb{M}_{n})_{\mathrm{rr}} = 1 - \frac{1}{n},\\
(\mathbb{M}_{n})_{\mathrm{mr}} = \frac{1}{n},\quad
(\mathbb{M}_{n})_{\mathrm{mm}} = 1,
\end{gathered}
\end{equation}
corresponding to retention of the bulk, bulk-to-radiation conduction, retention of the radiation, radiation-to-matter conduction, and retention of the terminal component. At $n = 1$ the two retention factors $1 - \tfrac1n$ vanish, leaving four non-zero entries --- the boundary case already recorded in \cref{cor:annihilation}.
\uses{\cref{eq:transfermatrix}}
\end{lemma}

\begin{proof}
Read off \cref{eq:transfermatrix} row by row: row $\mathrm B$ has one entry; row $\mathrm r$ has two; row $\mathrm m$ has two (one conduction, one retention). All five listed values are strictly positive for $n \geq 2$ ($1 - \tfrac1n \geq \tfrac12 > 0$); at $n = 1$, $1 - \tfrac1n = 0$ removes the two diagonal retention entries of $\mathrm B$ and $\mathrm r$, leaving four.
\end{proof}

\begin{lemma}[The Kraus operators are Hilbert--Schmidt orthogonal]\label{lem:hs}
For the Kraus operators \cref{eq:kraus} with $(\mathbb{M}_{n})_{ji} > 0$ and $(\mathbb{M}_{n})_{j'i'} > 0$,
\begin{equation}\label{eq:hs}
\langle K_{ji}, K_{j'i'} \rangle_{\mathrm{HS}}
= \operatorname{Tr}\bigl( K_{ji}^{\dagger} K_{j'i'} \bigr)
= \sqrt{(\mathbb{M}_{n})_{ji}\,(\mathbb{M}_{n})_{j'i'}}\;\operatorname{Tr}\bigl( |i\rangle\langle j | j' \rangle\langle i'| \bigr)
= \sqrt{\cdots}\; \delta_{j j'}\,\delta_{i i'},
\end{equation}
so distinct Kraus operators are orthogonal in Hilbert--Schmidt inner product.
\uses{\cref{eq:kraus}}
\end{lemma}

\begin{proof}
One computation, written out: the adjoint is $K_{ji}^{\dagger} = \sqrt{(\mathbb{M}_{n})_{ji}}\;|i\rangle\langle j|$, the product inside the trace is $|i\rangle\langle j|j'\rangle\langle i'| = \delta_{j j'} |i\rangle\langle i'|$, and the trace of $|i\rangle\langle i'|$ is $\langle i'|i\rangle = \delta_{i i'}$. The two Kronecker deltas are the whole content of the lemma.
\end{proof}

\begin{theorem}[Minimal environment dimension is five]\label{thm:minimal}
For $n \geq 2$, the minimal environment dimension of any dilation of the transfer channel \cref{eq:kraus} is exactly $5$.
\uses{\cref{lem:pattern}, \cref{lem:hs}, elementary rank counting}
\end{theorem}

\begin{proof}
\emph{Lower bound.} The Choi matrix of the channel --- formed from any Kraus representation as a sum of rank-one terms, one per Kraus operator,\footnote{M.-D.~Choi, \emph{Completely positive linear maps on complex matrices}, Linear Algebra Appl.\ \textbf{10}, 285 (1975): the matrix now called the Choi matrix, and the rank of that matrix is the least number of Kraus operators any representation can have. Only the direction ``number of Kraus operators $\geq$ rank'' is used, and it is immediate from the sum of rank-one terms.} --- has rank at most the number of Kraus operators in that representation. Here the five operators $K_{ji}$ of \cref{lem:pattern} are Hilbert--Schmidt orthogonal (\cref{lem:hs}), hence linearly independent (orthogonal non-zero vectors are independent), so they form \emph{a} representation with exactly five operators; any other representation needs at least five. Equivalently: the Choi matrix built from them is diagonal in the five orthonormal vectors $|j\rangle \otimes |i\rangle$ with eigenvalues $(\mathbb{M}_{n})_{ji} > 0$ --- its rank is exactly $5$, counted by counting the positive diagonal entries.

\emph{Achievement.} The dilation \cref{eq:isometry} uses the record states $|ji\rangle$ for precisely the five pairs of \cref{lem:pattern} at $n \geq 2$: its register has dimension $5$. The lower bound is met.
\end{proof}

\begin{remark}[The number five, fixed twice]\label{rem:five}
The count is forced from two independent directions: the transfer matrix's sparsity pattern \cref{lem:pattern} and the minimality theorem \cref{thm:minimal}. Both are checked against the same five entries, so the environment's dimension is not a choice but a theorem of the chain's structure; at $n = 1$ the corresponding number is $4$, matching the boundary layer's reduced pattern.
\end{remark}

\section{The source branch}\label{sec:source-branch}

The dilation of the transfer is complete. What remains of the content layer is the source $\boldsymbol{s}$, and its treatment is shorter: it requires no dilation at all.

\begin{remark}[The source is not a reduced dynamics]\label{rem:source-no-dilation}
The source term $\boldsymbol{s} = (J, 0, 0)^{\mathrm T}$ is a constant classical vector added at every frame (\cref{as:b3}), with the fixed value $J = \varepsilon_{\mathrm P}/2$ (\cref{cor:twoj}). It is not a channel: its trace grows by design, it admits no Kraus representation (\cref{rem:cptp-scope}), and it has no off-diagonal content to preserve. Asking for its ``environment'' would therefore be a category error --- a dilation answers the question ``of what larger unitary system is this map the reduction?'', and the source, not being a reduction, does not pose that question. What the source has instead of an environment is a status: it is a primitive affine term of the flow, with its mechanism an open question (\cref{rem:mechanismregistry}). Its ``environment'', insofar as the word applies at all, is the injection itself.
\end{remark}

\begin{proposition}[State-independence, re-verified in the dilation]\label{prop:state-independence}
Both structural facts of the chain --- that the transfer preserves the trace (hence acts identically regardless of the state it acts on) and that the increment is a constant --- hold in the dilation automatically: the first by \cref{prop:isometry} (the isometry property involves no state), the second because the source is a constant vector.
\uses{\cref{prop:isometry}, \cref{as:b3}}
\end{proposition}

\begin{remark}[No unnamed degrees of freedom]\label{rem:no-unnamed}
The chapter's construction is now complete and is closed by an exhaustion statement. The dilation of the content layer consists of exactly: (i) the frame-register stream of \cref{def:isometry} --- explicit, one fresh register per frame, of dimension $5$ at $n \geq 2$ and $4$ at $n = 1$, minimal by \cref{thm:minimal}; and (ii) the source, a hypothesis-layer constant whose mechanism is open. Between these two, \textbf{no unnamed degree of freedom remains}: any additional register, field, or hidden variable would have to couple to the system through one of them, and both are accounted for --- the register stream is minimal, and the source is not a freedom at all but a primitive. The statement is structural, not ontological: it certifies that the framework's explicit content is complete, and it is silent about what lies beyond it in the sense of \cref{rem:truncation-restated}.
\end{remark}

\section{The input--output picture, and the closure}\label{sec:collision}

The construction of \cref{sec:transfer-dilation,sec:minimality,sec:source-branch} stands on its own, but it fits into a standard and physically intuitive picture, which this section records as a reading with its citation.

\begin{eresult}[Collision model]\label{er:collision}
An open system may be represented as colliding with one fresh environmental unit per time step: each unit arrives in a standard state, interacts (``collides'') with the system, and is then either discarded or retained. The reduced dynamics is Markovian if and only if the units are fresh and uncorrelated with the system at arrival.
\reference{J.~Rau, \emph{Relaxation phenomena in spin and harmonic oscillator systems}, Phys.\ Rev.\ \textbf{129}, 1880 (1963), origin of the collision-model device; the register-string formulation used below is the standard one.}
\end{eresult}

\begin{remark}[The chain as a collision model]\label{rem:collision-reading}
The dilation fits the picture exactly, with the registers of \cref{def:isometry} as the colliding units: fresh register at frame $n$, one collision $V_{n}^{\mathrm{tr}}$, then discard (the reduced dynamics \cref{eq:reduction}) or retain (the full record of \cref{rem:record}). The model is \emph{time-non-homogeneous} --- the register's dimension and the collision operator change with $n$ through $\mathbb{M}_{n}$ --- which is the collision-model reading of the chain's $1/n$ structure. With the records retained, the system trajectory plus the register string is a classical joint Markov chain whose transition composition is exactly the cocycle law \cref{eq:cocycle} already proved at kernel level: the two descriptions meet, and no new content appears in the meeting.
\end{remark}

\begin{remark}[Closure of the open item]\label{rem:closure}
The question of the environment's structure --- what degrees of freedom it has --- is closed on four independent grounds, each established in its own section: \emph{existence}, by the explicit construction \cref{def:isometry}; \emph{dimension}, fixed at exactly five by \cref{thm:minimal}; \emph{explicit form}, the frame-register stream with its record states; and \emph{exhaustion}, the absence of unnamed degrees of freedom by \cref{rem:no-unnamed}. What remains open is only what was always mechanism: the physical realisation of the registers and the internal structure of the source. Both are mechanism questions under the rule of \cref{rem:register}, and the closure is structural: complete as far as the framework's explicit content goes, silent beyond it.
\end{remark}

\begin{remark}[What Chapter 6 has established]\label{rem:ch6summary}
The content layer is classical in type (\cref{prop:positivity,rem:what-dilation}), and its one quantum corner --- the transfer channel --- admits an explicit dilation (\cref{def:isometry}) verified in both required directions (\cref{prop:isometry,prop:reduction}), minimal in environment dimension (\cref{thm:minimal}), with a source that is a primitive rather than a reduction (\cref{rem:source-no-dilation}) and an exhaustion statement leaving no unnamed freedom (\cref{rem:no-unnamed}). The open item is thereby closed in every sense the framework's content allows (\cref{rem:closure}); the input--output picture organises the construction as a time-non-homogeneous collision model without adding to it (\cref{rem:collision-reading}). Together with \cref{lem:kraus}, the framework's open dynamics is now fully certified: lawful as a channel, explicit as a dilation, minimal as a construction.
\end{remark}

\part{Boundary and readings}

\chapter{Boundary thermodynamics}\label{ch:boundary}

This chapter constructs the entropy and temperature of the boundary from two experimental anchors, and shows that the relation $E = TS$ is here an identity. Every temperature in this chapter is a value derived from the energy--time hypothesis and its companions; no temperature below depends on any cosmological model, and no cosmological model enters the chapter.

\section{Geometric preliminaries}\label{sec:geometry}

\begin{lemma}[The boundary is a sphere]\label{lem:boundary-sphere}
Under the stepping structure and the uniformity of \cref{ch:newton}, the boundary at frame $n$ is the sphere of radius
\begin{equation}\label{eq:boundaryradius}
L(n) = n\,\ell_{\mathrm P},
\end{equation}
with area
\begin{equation}\label{eq:area}
A(n) = 4\pi L(n)^{2}.
\end{equation}
\uses{\cref{as:stepping}, uniformity (\cref{er:ball})}
\end{lemma}

\begin{proof}
Uniformity (isotropy about the centre) makes the stepped front an equidistant set: all points reached in $n$ steps lie at the same distance from the centre, and equidistance from a point is the defining property of a sphere. The distance is $n$ step lengths by \cref{eq:stepping}, with the step $\ell_{\mathrm P}$; the area of the sphere of radius $L$ is $4\pi L^{2}$ by the solid geometry of three-dimensional Euclidean space.
\end{proof}

\begin{lemma}[Volume]\label{lem:volume}
\begin{equation}\label{eq:volume}
V(L) = \int_{0}^{L} 4\pi r^{2}\,dr = \frac{4\pi}{3} L^{3}.
\end{equation}
\end{lemma}

\begin{proof}
One line: the integrand is the area of the sphere of radius $r$ (\cref{eq:area}), and $\int_{0}^{L} r^{2}\,dr = L^{3}/3$, giving the factor $4\pi/3$.
\end{proof}

\begin{lemma}[Deposited profile]\label{lem:deposited}
Each frame deposits the content $J$ at the boundary, and the deposit becomes a shell of thickness $\ell_{\mathrm P}$ at rest. The mass profile is
\begin{equation}\label{eq:profile}
M(r) = \frac{c^{2} r}{2 G},\qquad \frac{dM}{dr} = \frac{c^{2}}{2G},
\end{equation}
with density $\rho(r) \propto r^{-2}$, finite total mass at every radius, and no singularity at the centre.
\uses{\cref{cor:twoj}, \cref{cor:epmp}, \cref{def:units}, \cref{as:stepping}, the rest default below}
\end{lemma}

\begin{proof}
\emph{Deposit per frame.} The injected energy per frame is $J = \varepsilon_{\mathrm P}/2$ (\cref{cor:twoj}); by \cref{cor:epmp}, $\varepsilon_{\mathrm P} = M_{\mathrm P}c^{2}$, so the deposited mass per frame is
\begin{equation}\label{eq:deposit}
m_{\mathrm{dep}} = \frac{J}{c^{2}} = \frac{M_{\mathrm P}}{2}.
\end{equation}
\emph{Shell thickness.} One step per frame (\cref{as:stepping}) makes each deposited shell one step thick, $\ell_{\mathrm P}$; the shells do not move after deposition --- rest is the framework's default, since internal motion of the content would be an extra structure this paper nowhere introduces: the shells are deposited, not moving.
\emph{Mass per unit radius.} Dividing \cref{eq:deposit} by $\ell_{\mathrm P}$ and reducing with \cref{def:units}:
\begin{equation}\label{eq:surface-density}
\frac{m_{\mathrm{dep}}}{\ell_{\mathrm P}} = \frac{M_{\mathrm P}}{2\,\ell_{\mathrm P}}
= \frac{1}{2}\sqrt{\frac{\hbar c}{G}\cdot\frac{c^{3}}{\hbar G}}
= \frac{1}{2}\sqrt{\frac{c^{4}}{G^{2}}}
= \frac{c^{2}}{2G},
\end{equation}
each square-root reduction written separately. Integrating from the centre: $M(r) = \frac{c^{2}}{2G}\,r$.
\emph{Density and regularity.} With $A(r) = 4\pi r^{2}$ (\cref{eq:area}),
\begin{equation}\label{eq:densityprofile}
\rho(r) = \frac{dM/dr}{A(r)} = \frac{c^{2}}{8\pi G\, r^{2}},
\end{equation}
and the total mass inside any radius is $\int_{0}^{R} \rho(r)\,4\pi r^{2}\,dr = \frac{c^{2}R}{2G}$: the $r^{2}$ of the volume element cancels the $r^{-2}$ of the density, so the centre contributes a finite amount --- the profile is integrable, and there is no mass singularity.
\end{proof}

\begin{remark}[First-frame check]\label{rem:firstframe-check}
The check is on the \emph{rate}: the per-frame deposited mass \cref{eq:deposit} equals $J/c^{2}$, so the deposition reproduces the injection exactly, frame by frame, from the first frame on. (The per-frame deposit $\varepsilon_{\mathrm P}/2$ and the chain's first-frame retained content $E_{\mathrm B}^{(1)} = \varepsilon_{\mathrm P}/4$ are different quantities, related by the split rule's factor $2$; only the rate is a check of the structure.)
\end{remark}

\begin{remark}[Profiles, and what the constraint sees]\label{rem:profiles}
\Cref{eq:profile} is not the only conceivable mass profile: the uniform-density reading $M(r) = r^{3}/2Gt^{2}$ of the Newtonian ball is a fluid-dilution profile, and others can be written. What the framework's constraint can discriminate is less than a profile: \cref{prop:selfcheck} involves only the \emph{total} content over the \emph{total} volume, so profiles with the same totals give the same constraint, and the deposited profile of \cref{lem:deposited} is the framework's native default --- the zero-assumption content distribution of \cref{rem:register}. Alternatives are recorded, not adopted.
\end{remark}

\section{The entropy path}\label{sec:entropy-path}

The boundary's entropy is constructed from two experimental anchors and nothing else. Both are stated as empirical inputs with their domains; the construction is then pure substitution.

\begin{eresult}[Holographic bit density]\label{er:bitdensity}
One bit of information is carried per $4\ell_{\mathrm P}^{2}$ of boundary area:
\begin{equation}\label{eq:bitdensity}
\frac{d\mathcal N}{dA} = \frac{1}{4\ell_{\mathrm P}^{2}}.
\end{equation}
\evidence{Bekenstein (1973) and Hawking (1975) established the thermodynamics of horizons, of which \cref{eq:bitdensity} is the local, area-density form; it is used here as an empirical counting rule.}
\scope{Strict for de Sitter-type horizons; approximate for the present boundary, with the approximation carried by the anchor itself. Every entropy derived below inherits this domain, and nothing else in the framework does.}
\end{eresult}

\begin{eresult}[Boltzmann entropy]\label{er:boltzmann}
\begin{equation}\label{eq:boltzmann}
S = k_{\mathrm B}\ln W,
\end{equation}
with each bit an independent two-state unit, so that $W_{\mathrm{st}} = e^{\mathcal N}$ in the natural-logarithm convention.
\evidence{Boltzmann (1877);\footnote{L.~Boltzmann, \emph{\"Uber die Beziehung zwischen dem zweiten Hauptsatze der mechanischen W\"armetheorie und der Wahrscheinlichkeitsrechnung} (1877); the counting formula $S = k_{\mathrm B}\ln W$ and its two-state special case are the content of the input, not its derivation here. The information-theoretic reading ($\mathcal N$ bits, $W = 2^{\mathcal N}$ on the bit scale) is the Shannon form, one $\ln 2$ conversion away.} the formula is used as an empirical counting rule.}
\scope{Systems of weakly interacting units; the independence assumption is part of the anchor.}
\end{eresult}

\begin{proposition}[Information content of the boundary]\label{prop:information}
\begin{equation}\label{eq:information}
\mathcal N(n) = \frac{A(n)}{4\ell_{\mathrm P}^{2}} = \pi n^{2},
\qquad
\frac{d\mathcal N}{dA} = \frac{1}{4\ell_{\mathrm P}^{2}}.
\end{equation}
\uses{\cref{er:bitdensity}, \cref{lem:boundary-sphere}}
\end{proposition}

\begin{proof}
Substitute the area \cref{eq:area} with the boundary radius \cref{eq:boundaryradius}, $A = 4\pi n^{2}\ell_{\mathrm P}^{2}$, and divide by the cell $4\ell_{\mathrm P}^{2}$:
\begin{equation}\label{eq:ncancel}
\mathcal N = \frac{4\pi n^{2}\ell_{\mathrm P}^{2}}{4\ell_{\mathrm P}^{2}} = \pi n^{2},
\end{equation}
the factor $4\ell_{\mathrm P}^{2}$ cancelling term by term --- numerator cell against denominator cell --- with no independent constant entering. The area-law form $d\mathcal N/dA = 1/4\ell_{\mathrm P}^{2}$ is the input itself.
\end{proof}

\begin{corollary}[Boundary entropy]\label{cor:entropy}
\begin{equation}\label{eq:entropy}
S_{\partial} = k_{\mathrm B}\,\pi n^{2}.
\end{equation}
\uses{\cref{prop:information}, \cref{er:boltzmann}}
\end{corollary}

\begin{proof}
$\ln W_{\mathrm{st}} = \ln e^{\mathcal N} = \mathcal N$ by \cref{er:boltzmann}; multiplying by $k_{\mathrm B}$ and substituting \cref{eq:information} gives \cref{eq:entropy}.
\end{proof}

\begin{eresult}[Bekenstein--Hawking entropy]\label{er:bh}
\begin{equation}\label{eq:bh}
S_{\mathrm{BH}} = \frac{k_{\mathrm B}\,c^{3}\,A}{4\hbar G}.
\end{equation}
\reference{J.~D. Bekenstein, Phys.\ Rev.\ D \textbf{7}, 2333 (1973); S.~W. Hawking, Comm.\ Math.\ Phys.\ \textbf{43}, 199 (1975). Quoted as the external benchmark for the check below.}
\end{eresult}

\begin{proposition}[Verification against the benchmark]\label{prop:bh-check}
The boundary entropy \cref{eq:entropy} equals $S_{\mathrm{BH}}$ evaluated at the boundary area \cref{eq:area}, term by term.
\uses{\cref{eq:entropy}, \cref{er:bh}, \cref{lem:boundary-sphere}, \cref{def:units}}
\end{proposition}

\begin{proof}
Evaluate \cref{eq:bh} at $A = 4\pi n^{2}\ell_{\mathrm P}^{2}$ and reduce:
\begin{equation}\label{eq:bh-check}
S_{\mathrm{BH}} = \frac{k_{\mathrm B}c^{3}\cdot 4\pi n^{2}\ell_{\mathrm P}^{2}}{4\hbar G}
= \frac{k_{\mathrm B}\,\pi n^{2}\,c^{3}\ell_{\mathrm P}^{2}}{\hbar G}
= k_{\mathrm B}\,\pi n^{2},
\end{equation}
where the last equality uses $c^{3}\ell_{\mathrm P}^{2} = \hbar G$, itself the definition $\ell_{\mathrm P}^{2} = \hbar G/c^{3}$ of \cref{def:units} rearranged. Every factor cancels against its counterpart; nothing is fitted.
\end{proof}

\begin{remark}[Anchor and benchmark, and the honesty of the check]\label{rem:anchor-vs-verification}
The bit density \cref{er:bitdensity} and the Bekenstein--Hawking formula \cref{eq:bh} have the same content --- \cref{eq:bitdensity} \emph{is} \cref{eq:bh} divided by $k_{\mathrm B}\ln 2$ per bit. The distinction between them is the distinction between \emph{use} and \emph{verification}, not between two independent contents, and this paper states that openly rather than packaging the check as independent confirmation: \cref{prop:bh-check} verifies that the framework's construction reproduces the benchmark it was built to reproduce, which is a consistency check of the arithmetic and of the unit conversions --- valuable, and not more than that.
\end{remark}

\begin{remark}[Where the $\pi$ comes from]\label{rem:pi-origin}
The $\pi$ in $\pi n^{2}$ is geometric, not conventional: it is the $4\pi$ of the sphere's area (\cref{eq:area}) after the cell's $4$ cancelled (\cref{eq:ncancel}). It is not a descendant of the $2\pi$ convention tracked since \cref{rem:twopi} --- no angle entered \cref{sec:geometry} --- and it stands in the tracking of numerical factors as a separate entry, of geometric origin.
\end{remark}

\section{The temperature path}\label{sec:temperature-path}

\begin{definition}[Boundary energy]\label{def:eboundary}
\begin{equation}\label{eq:eboundary}
E_{\partial}(n) := M(n)\,c^{2} = \frac{n\,\varepsilon_{\mathrm P}}{2},
\end{equation}
the energy content of \cref{prop:readoutvalue}. As with every value in this chapter, it is derived under the hypotheses of Chapters \ref{ch:eth}--\ref{ch:newton} and depends on no cosmological model.
\end{definition}

\begin{hypothesis}[Temperature as a ratio]\label{as:ratio-temp}
The boundary temperature is defined by
\begin{equation}\label{eq:ratio-temp}
T_{\partial} := \frac{E_{\partial}}{S_{\partial}}.
\end{equation}
\motivation{The microcanonical derivative does not exist here, and the failure is worth one precise look before it is bypassed. Differentiating \cref{eq:entropy,eq:eboundary} with respect to the frame index,
\begin{equation}\label{eq:derivatives}
\frac{dS_{\partial}}{dn} = 2\pi k_{\mathrm B}\,n,
\qquad
\frac{dE_{\partial}}{dn} = \frac{\varepsilon_{\mathrm P}}{2},
\end{equation}
the first from $\frac{d}{dn}\pi n^{2}$ and the second from the linear content, so
\begin{equation}\label{eq:microcanonical}
\Bigl(\frac{\partial S_{\partial}}{\partial E_{\partial}}\Bigr)^{-1} = \frac{dE_{\partial}/dn}{dS_{\partial}/dn} = \frac{\varepsilon_{\mathrm P}/2}{2\pi k_{\mathrm B}n} = \frac{\varepsilon_{\mathrm P}}{4\pi k_{\mathrm B}n} = \frac{T_{\partial}}{2},
\end{equation}
using \cref{eq:ratio-temp} for the last equality: the microcanonical object is off from the ratio by a factor of $2$, and no refinement of the derivative repairs the mismatch, because the mismatch is structural --- the entropy is quadratic in $n$ while the content is linear, and a linear content against a quadratic entropy has no consistent conjugate temperature in the derivative sense. The ratio definition is therefore not a convenience but the only operational definition available; it is the sixth member of the hypothesis layer, and it is flagged where it is used.}
\end{hypothesis}

\begin{proposition}[Boundary temperature]\label{prop:temperature}
\begin{equation}\label{eq:temperature}
T_{\partial}(n) = \frac{\varepsilon_{\mathrm P}}{2\pi k_{\mathrm B}\,n} = \frac{\hbar}{2\pi k_{\mathrm B}\,t},
\qquad t = n\,t_{\mathrm P}.
\end{equation}
\uses{\cref{eq:ratio-temp}, \cref{eq:eboundary}, \cref{eq:entropy}, \cref{eq:tick}}
\end{proposition}

\begin{proof}
Divide \cref{eq:eboundary} by \cref{eq:entropy}:
\begin{equation}\label{eq:tempcompute}
T_{\partial} = \frac{n\,\varepsilon_{\mathrm P}/2}{k_{\mathrm B}\,\pi n^{2}} = \frac{\varepsilon_{\mathrm P}}{2\pi k_{\mathrm B}\,n},
\end{equation}
the frame index cancelling to first order. For the time form, substitute $\varepsilon_{\mathrm P} = \hbar/t_{\mathrm P}$ (\cref{def:epsilon}) and $n = t/t_{\mathrm P}$ (\cref{eq:tick}):
\begin{equation}\label{eq:temptime}
T_{\partial} = \frac{\hbar/t_{\mathrm P}}{2\pi k_{\mathrm B}\,(t/t_{\mathrm P})} = \frac{\hbar}{2\pi k_{\mathrm B}\,t},
\end{equation}
the tick cancelling against its reciprocal.
\end{proof}

\begin{eresult}[Gibbons--Hawking temperature]\label{er:gh}
A de Sitter horizon of radius $L$ carries
\begin{equation}\label{eq:gh}
T_{\mathrm{GH}} = \frac{\hbar c}{2\pi k_{\mathrm B} L}.
\end{equation}
\reference{G.~W. Gibbons and S.~W. Hawking, Phys.\ Rev.\ D \textbf{15}, 2738 (1977). Quoted as the external benchmark.}
\end{eresult}

\begin{proposition}[Verification against the benchmark]\label{prop:gh-check}
With the boundary radius $L = ct$ (\cref{lem:boundary-sphere}, \cref{eq:stepping}), $T_{\partial}$ of \cref{eq:temperature} equals $T_{\mathrm{GH}}$ of \cref{eq:gh}, term by term.
\uses{\cref{prop:temperature}, \cref{er:gh}, \cref{lem:boundary-sphere}}
\end{proposition}

\begin{proof}
Substitute $L = ct$ into \cref{eq:gh}: $T_{\mathrm{GH}} = \hbar c/(2\pi k_{\mathrm B}\,ct) = \hbar/(2\pi k_{\mathrm B}t)$, the factor $c$ cancelling against itself, which is exactly the second form of \cref{eq:temperature}.
\end{proof}

\begin{remark}[The ratio quantity, flagged]\label{rem:ratio-flag}
$T_{\partial}$ is a ratio of two state quantities, and \cref{rem:integrability} applies to it verbatim: it may be quoted at a frame, but it may not be integrated over frames --- its formal integral would be a continuum object with no frame-level meaning. This chapter respects the flag everywhere (for instance in \cref{sec:identities}, where only state quantities are accumulated), and every later chapter respects it as well.
\end{remark}

\begin{remark}[Asymptotic per-bit allotment]\label{rem:allotment}
The allotment of new content per new bit at large $n$,
\begin{equation}\label{eq:allotment}
\frac{\Delta E_{\partial}}{\Delta \mathcal N} = \frac{\varepsilon_{\mathrm P}/2}{2\pi n} = \frac{\varepsilon_{\mathrm P}}{4\pi n}
\;\xrightarrow{\ \cref{eq:temperature}\ }\;
\frac{k_{\mathrm B} T_{\partial}}{2},
\end{equation}
approaches one half of the thermal unit $k_{\mathrm B}T_{\partial}$ per bit. The factor $1/2$ is the same structural feature as the factor $2$ of \cref{eq:microcanonical}: a linear content against a quadratic count divides every new quantum between new bits and old ones in a fixed proportion, and \cref{sec:identities} records that proportion exactly.
\end{remark}

\section{Two consistency paths}\label{sec:consistency-paths}

The ratio definition \cref{eq:ratio-temp} is the definitional source of $T_{\partial}$. Two further routes reach the same relation from independent inputs, and both are recorded --- as checks, not as derivations, and the difference is stated once here and enforced throughout.

\begin{eresult}[Clausius relation]\label{er:clausius}
For reversible heat exchange,
\begin{equation}\label{eq:clausius}
dS = \frac{\delta Q_{\mathrm{rev}}}{T}.
\end{equation}
\evidence{R.~J.~E.~Clausius (1854/1865);\footnote{R.~J.~E.~Clausius, \emph{\"Uber eine ver\"anderte Form des zweiten Hauptsatzes der mechanischen W\"armetheorie}, Ann.\ Phys.\ \textbf{169}, 481 (1854); consolidated in \emph{Mechanische W\"armetheorie} (1865). Used here as the operational definition of heat, nothing more.} the relation is used as an experimental input, at exactly the level at which it was established.}
\scope{Reversible processes; the framework applies it only to the differential split below.}
\end{eresult}

\begin{proposition}[Clausius consistency]\label{prop:clausius-check}
Assume the identity $E_{\partial} = T_{\partial}S_{\partial}$ (in whatever form it is given --- here, as the product rule)
\begin{equation}\label{eq:productrule}
dE_{\partial} = T_{\partial}\,dS_{\partial} + S_{\partial}\,dT_{\partial},
\end{equation}
in which the first term is the heat exchanged (per \cref{eq:clausius}) and the second is the dilution of the thermal unit as $T_{\partial}$ falls. With the growth rates
\begin{equation}\label{eq:growthrates}
\frac{dE_{\partial}}{dt} = \mathcal P = \frac{c^{5}}{2G},
\qquad
\frac{dS_{\partial}}{dt} = \frac{2S_{\partial}}{t}
\quad\Bigl(\text{from } S_{\partial} = \frac{k_{\mathrm B}\pi t^{2}}{t_{\mathrm P}^{2}},\ \frac{d}{dt}t^{2} = 2t\Bigr),
\end{equation}
the product rule integrates to
\begin{equation}\label{eq:integrated}
T_{\partial}S_{\partial} = \mathcal P\,t + C,
\end{equation}
and regularity at $t \to 0$ forces $C = 0$, returning $T_{\partial} = E_{\partial}/S_{\partial}$.
\uses{\cref{er:clausius}, \cref{eq:entropy}, \cref{def:power}, \cref{prop:temperature} (for the identification of terms)}
\end{proposition}

\begin{proof}
Two steps, both displayed. \emph{Integration.} Along any path with $E_{\partial} = T_{\partial}S_{\partial}$ as a functional relation, \cref{eq:productrule} is the total differential of the product, so
\begin{equation}\label{eq:diff}
\frac{d}{dt}\bigl(T_{\partial}S_{\partial}\bigr) = \frac{dE_{\partial}}{dt} = \mathcal P
\quad\Longrightarrow\quad
T_{\partial}S_{\partial} = \mathcal P t + C
\end{equation}
by one integration, \cref{eq:growthrates} supplying $\mathcal P$; the entropy's own growth rate $\dot S_{\partial} = 2S_{\partial}/t$ is shown in \cref{eq:growthrates} from $S_{\partial} = k_{\mathrm B}\pi n^{2}$ with $n = t/t_{\mathrm P}$: $S_{\partial} = k_{\mathrm B}\pi t^{2}/t_{\mathrm P}^{2}$, and $d(t^{2})/dt = 2t$. \emph{Regularity.} As $t \to 0$: $S_{\partial} \propto t^{2}$ vanishes quadratically and $T_{\partial} \propto t^{-1}$ diverges inversely (\cref{eq:temperature}), so the product $T_{\partial}S_{\partial} \propto t$ vanishes; \cref{eq:integrated} at $t \to 0$ then reads $0 = 0 + C$, forcing $C = 0$. Hence $T_{\partial}S_{\partial} = \mathcal P t = E_{\partial}$ (\cref{eq:eboundary}), the ratio definition.
\end{proof}

\begin{remark}[Status of the Clausius path: a check, stated exactly]\label{rem:clausius-status}
\Cref{prop:clausius-check} assumes what it concludes --- the identity $E_{\partial} = T_{\partial}S_{\partial}$ enters through \cref{eq:productrule}, so the path cannot generate the ratio definition, only verify that the definition is \emph{consistent} with the Clausius relation and the growth rates. Its content, stated exactly: \emph{if} the boundary content is the product of its entropy and temperature, \emph{then} the heat/dilution split of the injected power is $T_{\partial}\dot S_{\partial} = 2\mathcal P$ and $S_{\partial}\dot T_{\partial} = -\mathcal P$ (substituting \cref{eq:growthrates} into \cref{eq:productrule}: $\dot T_{\partial} = -T_{\partial}/t$, so $S_{\partial}\dot T_{\partial} = -S_{\partial}T_{\partial}/t = -E_{\partial}/t = -\mathcal P$), a fixed $2{:}1$ partition of the injection between organisation into new bits and dilution of old ones. The standing status of the path is thereby fixed: the check is kept, and it is not promoted to a derivation.
\end{remark}

\begin{proposition}[Zero-point consistency]\label{prop:zeropoint-check}
The zero-point floor of the reference oscillator equals the per-frame injection:
\begin{equation}\label{eq:zeropoint-check}
E_{0} = \alpha_{0}\,\omega_{0} = \frac{\hbar}{2t_{\mathrm P}} = \frac{\varepsilon_{\mathrm P}}{2} = \Delta E,
\end{equation}
with $\Delta E$ the injection theorem's per-frame quantum.
\uses{\cref{thm:spectrum}, \cref{as:p}, \cref{thm:injection}}
\end{proposition}

\begin{proof}
$E_{0} = \alpha_{0}\omega_{0}$ is the oscillator floor in the form of \cref{rem:half}; $\omega_{0} = 1/t_{\mathrm P}$ is \cref{eq:framerate}; so $E_{0} = (\hbar/2)(1/t_{\mathrm P}) = \hbar/2t_{\mathrm P} = \varepsilon_{\mathrm P}/2$ by \cref{cor:epmp}, which is $\Delta E$ by \cref{eq:inj-theorem}.
\end{proof}

\begin{remark}[The three-way meeting, and its exact structure]\label{rem:three-way}
Three routes now stand at the relation $E_{\partial} = T_{\partial}S_{\partial}$: the ratio definition (definitional), the Clausius differential (a consistency check), and the zero-point path (a cross-chapter consistency check). Their statuses are unequal and are kept unequal: one defines, two verify. What the meeting adds is not a stronger definition but a redundancy check of the framework's joints --- the ratio temperature, the injection rate, and the oscillator floor were built in different chapters from different inputs, and they agree. No route uses another's conclusion as a premise, and none is used below as if it were the definition.
\end{remark}

\section{A negative check: not a blackbody}\label{sec:sb-check}

The constructions so far are positive: they build, and they verify. One candidate description of the boundary can be excluded, and the exclusion is worth recording because it is the chapter's only falsification-shaped statement.

\begin{eresult}[Stefan--Boltzmann law]\label{er:sb}
A blackbody source of area $A$ radiates power
\begin{equation}\label{eq:sb}
P = \sigma A T^{4},
\qquad
\sigma = \frac{\pi^{2} k_{\mathrm B}^{4}}{60\,\hbar^{3} c^{2}}.
\end{equation}
\reference{J.~Stefan (1879, empirical); L.~Boltzmann (1884, theoretical); the constant $\sigma$ in the form \cref{eq:sb} follows from Planck's law. Phenomenology of ordinary radiating bodies --- which is exactly the status the check below probes.}
\end{eresult}

\begin{proposition}[The blackbody candidate fails]\label{prop:sb}
If the boundary were a blackbody source radiating the injected power $\mathcal P$ at temperature $T$, then
\begin{equation}\label{eq:sb-temp}
T_{\mathrm{SB}} = \Bigl(\frac{\mathcal P}{\sigma A}\Bigr)^{1/4} = 2\pi\Bigl(\frac{15}{2\pi^{3}}\Bigr)^{1/4}\sqrt{n}\; T_{\partial} \approx 4.4\,\sqrt{n}\;T_{\partial},
\end{equation}
a law $T_{\mathrm{SB}} \propto t^{-1/2}$ that diverges from $T_{\partial} \propto t^{-1}$ by $\sqrt{n}$ at every frame.
\uses{\cref{er:sb}, \cref{def:power}, \cref{lem:boundary-sphere}, \cref{prop:temperature}, \cref{def:units}}
\end{proposition}

\begin{proof}
\emph{Denominator.} Substitute $A = 4\pi c^{2}t^{2}$ (\cref{eq:area} with $L = ct$) and the constant \cref{eq:sb}:
\begin{equation}\label{eq:sb-denom}
\sigma A = \frac{\pi^{2}k_{\mathrm B}^{4}}{60\hbar^{3}c^{2}}\cdot 4\pi c^{2}t^{2} = \frac{\pi^{3} k_{\mathrm B}^{4}\,t^{2}}{15\,\hbar^{3}},
\end{equation}
the $c^{2}$ cancelling. \emph{Quotient.} Divide $\mathcal P = c^{5}/2G$ (\cref{def:power}) by \cref{eq:sb-denom}:
\begin{equation}\label{eq:sb-quotient}
\frac{\mathcal P}{\sigma A} = \frac{c^{5}}{2G}\cdot\frac{15\hbar^{3}}{\pi^{3}k_{\mathrm B}^{4}t^{2}} = \frac{15\,\hbar^{3}c^{5}}{2\pi^{3} G k_{\mathrm B}^{4} t^{2}},
\end{equation}
the $2$ and $15$ from the two constants combining in the only place they can. \emph{Fourth root, and the ratio.} Taking the fourth root of \cref{eq:sb-quotient} and dividing by $T_{\partial} = \hbar/2\pi k_{\mathrm B}t$ (\cref{eq:temperature}):
\begin{equation}\label{eq:sb-ratio}
\frac{T_{\mathrm{SB}}}{T_{\partial}} = \Bigl(\frac{15}{2\pi^{3}}\Bigr)^{1/4}\frac{\hbar^{3/4}c^{5/4}}{G^{1/4}k_{\mathrm B}t^{1/2}}\cdot\frac{2\pi k_{\mathrm B}t}{\hbar}
= 2\pi\Bigl(\frac{15}{2\pi^{3}}\Bigr)^{1/4}\frac{c^{5/4}t^{1/2}}{\hbar^{1/4}G^{1/4}},
\end{equation}
the $k_{\mathrm B}$ cancelling and the powers of $\hbar$, $t$ reducing to $\hbar^{-1/4}t^{1/2}$. Now $\hbar G = c^{3}\ell_{\mathrm P}^{2}$ (\cref{def:units}) turns the unit combination into
\begin{equation}\label{eq:sb-final}
\frac{c^{5/4}t^{1/2}}{\hbar^{1/4}G^{1/4}} = \Bigl(\frac{c^{5}t^{2}}{\hbar G}\Bigr)^{1/4} = \Bigl(\frac{c^{2}t^{2}}{\ell_{\mathrm P}^{2}}\Bigr)^{1/4} = \Bigl(\frac{ct}{\ell_{\mathrm P}}\Bigr)^{1/2} = \sqrt{n},
\end{equation}
using $ct = n\ell_{\mathrm P}$ (\cref{eq:stepping}). Assembling \cref{eq:sb-ratio,eq:sb-final} gives the coefficient of \cref{eq:sb-temp}; its numerical value, $2\pi(15/2\pi^{3})^{1/4} \approx 4.4$, is an evaluation of the exact constant, not an independent input. \emph{Divergence.} $T_{\mathrm{SB}} \propto t^{-1/2}$ and $T_{\partial} \propto t^{-1}$ are power laws of different exponents: their ratio $\propto t^{1/2} = \sqrt{n}$ grows without bound, and no constant rescaling reconciles them at more than one instant.
\end{proof}

\begin{remark}[Verdict, in the negative]\label{rem:sb-verdict}
The check excludes one candidate: the boundary, as a source, does not radiate its injected power as a blackbody. The exclusion is stated in the framework's own terms --- the injection is an event of the content layer (\cref{as:b3}), not a radiation process --- and it asserts no mechanism about what the boundary does instead. It is the chapter's only falsification-shaped statement, and its shape is worth keeping in view: the framework survives it not by accommodating blackbody behaviour but by never having claimed it.
\end{remark}

\section{The identity \texorpdfstring{$E = TS$}{E = TS}, and the root of one-way conduction}\label{sec:identity-direction}

\subsection{Identity and direction}\label{sec:ets-direction}

\begin{remark}[The direction of \texorpdfstring{$E = TS$}{E = TS}]\label{rem:direction-ets}
The relation $E = TS$ enters this framework as an \emph{identity}, not as an assumption --- and the direction is the point: $E_{\partial}$ is generated by the Newtonian leg (\cref{def:eboundary}), $S_{\partial}$ by the information anchors (\cref{cor:entropy}), and $T_{\partial}$ by the ratio definition (\cref{as:ratio-temp}). The relation $E_{\partial} = T_{\partial}S_{\partial}$ is then not an assumption but a \emph{rewrite} --- it is the ratio definition read backwards, an identity of the form $E = (E/S)\,S$. A consistency loop remains, recorded here and never reversed in use: \emph{if} one assumes $E = TS$ \emph{and} quotes the benchmarks $T_{\mathrm{GH}}$ (\cref{er:gh}) and $S_{\mathrm{BH}}$ (\cref{er:bh}), one recovers the same content, since
\begin{equation}\label{eq:oldloop}
T_{\mathrm{GH}}\,S_{\mathrm{BH}} = \frac{\hbar}{2\pi k_{\mathrm B}t}\cdot\frac{k_{\mathrm B}c^{3}A}{4\hbar G} = \frac{c^{3}A}{8\pi G\,t} = \frac{c^{5}t}{2G},
\end{equation}
using $A = 4\pi c^{2}t^{2}$ in the last step; this loop computes $E$ from benchmarks and may serve as a consistency check, but it is not a premise for anything below --- the direction is single. The free energy $F = E_{\partial} - T_{\partial}S_{\partial}$ vanishes identically, and its content is correspondingly definitional: \emph{all} of the boundary's energy is thermal, in the precise sense that none of it is free to do work at fixed $(S_{\partial}, T_{\partial})$; the vanishing is a definition expressing that, not a coincidence to be explained.
\end{remark}

\subsection{The thermodynamic root of one-way conduction}\label{sec:b1-root}

This subsection discharges the forward reference in the motivation of \cref{as:b1}: the absence of a reverse channel is not merely asserted, it is the shadow of a divergence. The symbol $\beta$ enters here for the first time, in its standard statistical-mechanical meaning $\beta = 1/(k_{\mathrm B}T)$; the present use of the symbol is its only one in this document.

\begin{lemma}[Quadratic entropy]\label{lem:quadratic}
Eliminating $n$ between \cref{eq:entropy,eq:eboundary} gives
\begin{equation}\label{eq:quadratic}
S_{\partial} = a\,E_{\partial}^{2}, \qquad a = \frac{4\pi k_{\mathrm B}}{\varepsilon_{\mathrm P}^{2}},
\end{equation}
i.e.\ the boundary entropy is quadratic in its energy.
\uses{\cref{eq:entropy}, \cref{eq:eboundary}}
\end{lemma}

\begin{proof}
From \cref{eq:eboundary}, $n = 2E_{\partial}/\varepsilon_{\mathrm P}$; substitute into \cref{eq:entropy}: $S_{\partial} = k_{\mathrm B}\pi\bigl(2E_{\partial}/\varepsilon_{\mathrm P}\bigr)^{2} = k_{\mathrm B}\pi\cdot 4E_{\partial}^{2}/\varepsilon_{\mathrm P}^{2} = \frac{4\pi k_{\mathrm B}}{\varepsilon_{\mathrm P}^{2}}E_{\partial}^{2}$.
\end{proof}

\begin{theorem}[No equilibrium, hence no reverse channel]\label{thm:no-equilibrium}
The canonical partition function of a system with the spectrum of \cref{lem:quadratic},
\begin{equation}\label{eq:partition}
Z(\beta) = \int_{0}^{\infty} e^{\,S(E)/k_{\mathrm B} - \beta E}\,dE = \int_{0}^{\infty} e^{\,\alpha E^{2} - \beta E}\,dE,
\qquad \alpha = \frac{4\pi}{\varepsilon_{\mathrm P}^{2}} > 0,
\end{equation}
diverges for every $\beta > 0$. Such a system has no equilibrium state at any temperature, and therefore no channel through which a reverse flow could equilibrate --- the thermodynamic root of \cref{as:b1}.
\uses{\cref{lem:quadratic}; integral bound}
\end{theorem}

\begin{proof}
The exponent factors as $\alpha E^{2} - \beta E = E(\alpha E - \beta)$. For $E \geq \beta/\alpha$, the second factor is non-negative, so the integrand satisfies $e^{\alpha E^{2} - \beta E} \geq 1$ on the whole half-line $[\beta/\alpha, \infty)$, and
\begin{equation}\label{eq:divergence}
Z(\beta) \;\geq\; \int_{\beta/\alpha}^{\infty} 1\;dE \;=\; \infty
\end{equation}
by comparison with a constant integrand. The divergence is not a borderline failure at small $\beta$: it occurs at \emph{every} inverse temperature, and a divergent partition function is the statement that no Gibbs state exists. Two honest boundaries close the argument. (i) What is proved is the \emph{absence of an equilibration channel}: a quadratic-entropy object cannot sit in equilibrium with anything, so no reverse flow can be an approach to equilibrium. (ii) What is not proved is the \emph{direction}: why the flow runs $\mathrm B \to \mathrm r \to \mathrm m$ rather than the reverse is carried by \cref{as:b1} as a hypothesis, and this theorem gives it no direction. Both boundaries are kept.
\end{proof}

\subsection{Status, cleared}\label{sec:clearing}

\begin{remark}[The clearing table]\label{rem:clearing}
The standing items, and what this framework makes of them:
\begin{center}
\begin{tabular}{@{}p{2.2cm}p{1.8cm}p{2.8cm}p{4.3cm}@{}}
\toprule
Standing item & Old status & Present status & Settled by \\
\midrule
$E = TS$ & axiom & identity / rewrite & $E$ from \cref{ch:newton}, $S$ from anchors, $T$ a ratio \\
One-way conduction & assumption & mostly theorem-ized & no equilibrium (\cref{thm:no-equilibrium}); delay (\cref{cor:delay}); $\gamma = 1$ uniqueness (\cref{thm:forced-cascade}) \\
Discrete time & assumption & the energy--time hypothesis & A1 $+$ A2, magnitude supported by \cref{prop:tau} \\
\bottomrule
\end{tabular}
\end{center}
The asymmetric division of load, stated once: if the entropy--temperature layer (anchors or benchmarks) were to fail, $E_{\partial}$ and $\mathcal P$ would stand on \cref{ch:newton} alone, and the framework would survive with them; if the Newtonian leg's standing hypotheses (\cref{as:stepping}, \cref{as:pump}) were to fail, everything above them fails with them. The honesty is symmetric: the weak points of each layer are recorded with the layer. And the loop's direction is irreversible in use: it computes $T$ from $E$, never $E$ from $T$.
\end{remark}

\begin{remark}[Two free energies, one convention apart]\label{rem:dual-free}
The thermal free energy vanishes identically, $F_{\mathrm{th}} = E_{\partial} - T_{\partial}S_{\partial} \equiv 0$ (\cref{rem:direction-ets}), while the statistical free energy of the same state would read $F_{\mathrm{stat}} = -k_{\mathrm B}T_{\partial}\ln W_{\mathrm{st}} = -k_{\mathrm B}T_{\partial}\,\pi n^{2} = -E_{\partial}$. The two differ by the energy-zero convention of the counting formula, and only differences of free energy are observable; the note is recorded to prevent a confusion, and not developed further.
\end{remark}

\section{The identity set, and where the shares live}\label{sec:identities}

\begin{proposition}[Thermodynamic identities]\label{prop:identities}
The following hold at every $t > 0$ (continuum forms, exact):
\begin{enumerate}
\item zero free energy: $F = E_{\partial} - T_{\partial}S_{\partial} = 0$;
\item negative heat capacity: $C = \dfrac{dE_{\partial}/dt}{dT_{\partial}/dt} = -S_{\partial} < 0$;
\item the $2{:}1$ split of the injection: $T_{\partial}\dot S_{\partial} = 2\mathcal P$, $\;S_{\partial}\dot T_{\partial} = -\mathcal P$;
\item entropy production: $\sigma = \dot S_{\partial}/2 = S_{\partial}/t > 0$;
\item second difference of the information count: $\Delta^{2}\mathcal N = 2\pi$ per frame;
\item constant product: $E_{\partial}T_{\partial} = \dfrac{\hbar c^{5}}{4\pi k_{\mathrm B}G}$;
\item boundary tension: $\dfrac{\partial E_{\partial}}{\partial L} = \dfrac{c^{4}}{2G}$;
\item per-bit allotment: $\dfrac{\Delta E_{\partial}}{\Delta \mathcal N} \to \dfrac{k_{\mathrm B}T_{\partial}}{2}$ as $n \to \infty$.
\end{enumerate}
\uses{\cref{eq:entropy,eq:eboundary,eq:temperature,def:power,rem:integrability} (only state quantities are differentiated or differenced)}
\end{proposition}

\begin{proof}
Each item one to three lines; every differentiation is of a state quantity, per \cref{rem:integrability}.
\begin{enumerate}
\item The ratio definition, read backwards (\cref{rem:direction-ets}).
\item $\dot T_{\partial} = -T_{\partial}/t$ from \cref{eq:temperature}, so $C = \mathcal P/(-T_{\partial}/t) = -\mathcal P t/T_{\partial} = -E_{\partial}/T_{\partial} = -S_{\partial}$, using $E_{\partial} = \mathcal P t$ and the ratio definition.
\item $\dot S_{\partial} = 2S_{\partial}/t$ from $S_{\partial} \propto t^{2}$ (\cref{eq:growthrates}); then $T_{\partial}\dot S_{\partial} = 2T_{\partial}S_{\partial}/t = 2E_{\partial}/t = 2\mathcal P$, and $S_{\partial}\dot T_{\partial} = -S_{\partial}T_{\partial}/t = -\mathcal P$.
\item From the same derivative, $\sigma := \dot S_{\partial}/2 = S_{\partial}/t$, positive since $S_{\partial} > 0$.
\item $\Delta \mathcal N(n) = \pi(n+1)^{2} - \pi n^{2} = \pi(2n+1)$; $\Delta^{2}\mathcal N = \pi(2n+3) - \pi(2n+1) = 2\pi$ --- a constant because the second difference of a quadratic is twice its leading coefficient, with the linear term annihilated.
\item $E_{\partial}T_{\partial} = \dfrac{n\varepsilon_{\mathrm P}}{2}\cdot\dfrac{\varepsilon_{\mathrm P}}{2\pi k_{\mathrm B}n} = \dfrac{\varepsilon_{\mathrm P}^{2}}{4\pi k_{\mathrm B}}$, and $\varepsilon_{\mathrm P}^{2} = (\hbar/t_{\mathrm P})^{2} = \hbar c^{5}/G$ by \cref{eq:reducefinal}; the frame index cancels, leaving the $t$-independent constant.
\item With $E_{\partial} = c^{5}t/2G$ and $L = ct$: $E_{\partial} = c^{4}L/2G$; differentiating, $\partial E_{\partial}/\partial L = c^{4}/2G$.
\item \Cref{rem:allotment}.
\end{enumerate}
\end{proof}

\begin{remark}[The discrete residual, and what it is made of]\label{rem:discrete-residual}
The exact per-frame entropy increment in frame units is: $\Delta S_{\partial}(n) = k_{\mathrm B}\pi(2n+1) = 2\pi k_{\mathrm B}\,n + \pi k_{\mathrm B}$ --- a linear part and a constant part. The $2{:}1$ split of item 3 is the leading, $n \gg 1$ form; the constant remainder $\pi k_{\mathrm B}$ per frame is not an anomaly but a second-difference object, of exactly the same origin as item 5: the count $\mathcal N = \pi n^{2}$ has its growth \emph{rate} linear in $n$ and its growth \emph{acceleration} constant, and $\pi k_{\mathrm B}$ is that constant seen in the first difference. If the acceleration of the information count were ever observed to drift with $n$, item 5 --- and with it this residual --- would be the first place the drift showed. Both the exact frame arithmetic and the asymptotic split are kept, and neither is confused with the other.
\end{remark}

\subsection{Equation of state, recalled}\label{sec:eos-recap}

The effective equation of state was derived once, in \cref{sec:density-readout}: $w = -\tfrac13$ from $\rho \propto t^{-2}$ together with the conservation scaling \cref{er:conservation}, with the effective/instantaneous distinction drawn in \cref{rem:instantaneous} --- the injected content is pressureless as deposited material, and $w = -\tfrac13$ is the effective shape of its dilution. This chapter adds nothing to that derivation; it records the boundary-chapter form of the density, $\rho_{\partial} = E_{\partial}/V = \tfrac{3c^{2}}{8\pi G t^{2}}$, the energy-density twin of \cref{prop:density}, and the zero active mass $\rho_{\partial}c^{2} + 3p_{\partial} = 0$ (\cref{prop:eos}). The standing prohibition: $w = -\tfrac13$ is the effective state of the boundary content, not an identification with any dark-energy sector, and this chapter makes no cosmological-constant statement of any kind.

\subsection{Where the bulk share lives}\label{sec:residence}

\begin{proposition}[Geometric residence of the bulk share]\label{prop:residence}
\begin{equation}\label{eq:residence}
M\Bigl(\frac{L}{2}\Bigr) = \frac{c^{2}L}{4G} = \frac{E_{\mathrm B}}{c^{2}},
\end{equation}
the mass inside half the boundary radius equals the bulk content of the chain.
\uses{\cref{lem:deposited}, \cref{cor:twoj}, \cref{eq:eboundary}}
\end{proposition}

\begin{proof}
Two readings of the same arithmetic. \emph{By the profile:} $M(L/2) = c^{2}(L/2)/2G = c^{2}L/4G$ from \cref{eq:profile}. \emph{By the shares:} the deposited shells have equal mass ($dM/dr = c^{2}/2G$ constant, \cref{lem:deposited}), so halving the radius halves the enclosed mass: $M(L/2) = \tfrac12 M(L) = \tfrac12 E_{\partial}/c^{2} = E_{\mathrm B}/c^{2}$, since $E_{\mathrm B} = n\varepsilon_{\mathrm P}/4 = E_{\partial}/2$ (\cref{eq:eboundary}, \cref{cor:twoj}). Both give $c^{2}L/4G$.
\end{proof}

\begin{remark}[Reading]\label{rem:withdrawn-share}
The deposited content is an equal-mass onion: every shell the same mass, doubling the radius doubles the enclosed mass. The chain's bulk share --- one half of the total --- is thereby not a volume fraction but a \emph{radius} fraction: it occupies the inner half by radius, which is the inner one-eighth by volume, and the two statements are the same fact read through \cref{eq:profile}. No hemisphere or $S^{3}$ identification of the shares is made here: such an identification would require a global space and a preferred hemisphere --- assumptions the framework does not make --- together with a geometry in tension with the flat kinematics of the reading layer; whatever residence the deeper chain shares have is deferred to \cref{ch:symmetries}, and nothing here depends on it.
\end{remark}

\section{Summary of the boundary}\label{sec:boundary-summary}

\begin{remark}[What Chapter 7 has established]\label{rem:ch7summary}
The boundary is a sphere of radius $L = n\ell_{\mathrm P}$ (\cref{lem:boundary-sphere}), enclosing volume $V = \tfrac{4}{3}\pi L^{3}$ (\cref{lem:volume}), and carrying the deposited profile $M(r) = c^{2}r/2G$ of equal-mass shells (\cref{lem:deposited}), with the rest default stated (\cref{rem:firstframe-check}). From two empirical inputs --- the holographic bit density and the Boltzmann count --- the boundary's information content $\mathcal N = \pi n^{2}$ and entropy $S_{\partial} = k_{\mathrm B}\pi n^{2}$ follow (\cref{prop:information,cor:entropy}), verified term by term against the Bekenstein--Hawking benchmark with the identity of anchor and benchmark stated openly (\cref{prop:bh-check,rem:anchor-vs-verification}). From the sixth hypothesis --- the ratio definition, forced by the diagnosed failure of the microcanonical derivative (\cref{as:ratio-temp}) --- the temperature $T_{\partial} = \hbar/2\pi k_{\mathrm B}t$ follows (\cref{prop:temperature}), verified against the Gibbons--Hawking benchmark (\cref{prop:gh-check}) and flagged as non-integrable (\cref{rem:ratio-flag}). Two consistency paths reach $E_{\partial} = T_{\partial}S_{\partial}$ from independent inputs, with their statuses fixed as checks rather than derivations (\cref{prop:clausius-check,prop:zeropoint-check,rem:three-way}); the identity itself is a rewrite of the ratio definition, and the old closed loop is kept irreversible in use (\cref{rem:direction-ets}). One candidate description is excluded: the boundary is not a blackbody source (\cref{prop:sb}). One forward reference is discharged: one-way conduction has a thermodynamic root in the divergence of the canonical partition function (\cref{thm:no-equilibrium}), while its direction remains a hypothesis (\cref{as:b1}). The identity set of \cref{prop:identities} --- free energy, heat capacity, the $2{:}1$ split with its discrete residual, entropy production, the second difference $2\pi$, the constant product, the tension, and the per-bit allotment --- is verified line by line, and the bulk share is located geometrically as the inner half by radius (\cref{prop:residence}), the older hemisphere images withdrawn (\cref{rem:withdrawn-share}). The effective equation of state $w = -\tfrac13$ stands as derived in \cref{sec:density-readout}, with the prohibition on dark-energy identifications permanent (\cref{sec:eos-recap}).

The chapter's input account: three empirical inputs (\cref{er:bitdensity,er:boltzmann,er:clausius}), one hypothesis (\cref{as:ratio-temp}), three external benchmarks (\cref{er:bh,er:gh,er:sb}); every other symbol in the chapter is derived. Precision inheritance: the anchors' domain restricts $S_{\partial}$, $T_{\partial}$, and everything differentiated or differenced from them; $E_{\partial}$ and $\mathcal P$ stand on \cref{ch:newton} and are unaffected (\cref{rem:clearing}). Open items, each with its owner: correction of the anchors --- inherited by the boundary quantities alone; the first-frame residuals, owned by the boundary layer of \cref{cor:annihilation}; the residence of the deeper shares, deferred to \cref{ch:symmetries}; and the mechanism of the injection, open in both directions (\cref{rem:mechanismregistry}). No cosmological model entered the chapter, and none was needed.
\end{remark}

\chapter{The Friedmann--Einstein reading}\label{ch:gr}

This chapter assembles the points of contact between the framework's reading and the Friedmann equations of general relativity. There are two such contacts --- the constraint and the acceleration --- and one kinematic contact through the age--Hubble product. Each contact is stated where it occurs, with its direction; none of them derives the other side, and no observational claim is made in this chapter.

\section{The constraint}\label{sec:constraint}

Two ingredients enter, both earned earlier. The first is the content of \cref{ch:newton}: the ball's content grows linearly, $M(t) = c^{3}t/2G$ (\cref{eq:ballmass}), so that the radius advances one step per tick, $R(t) = ct$ (\cref{eq:ct}). The second is the ball geometry of \cref{sec:geometry}: the uniform ball has volume $V = \tfrac{4}{3}\pi R^{3}$ (\cref{lem:volume}). The density of the reading was computed in full in \cref{prop:density}: $\rho(t) = M(t)/V(R(t)) = 3/(8\pi G t^{2})$, the constant $8\pi/3$ being the product of the geometry's $4\pi/3$ and the energetics' $1/2$.

\begin{proposition}[The constraint of the reading]\label{prop:constraint}
With $H = \dot R/R$ (\cref{def:hubble}),
\begin{equation}\label{eq:constraint}
H^{2} = \frac{8\pi G}{3}\,\rho = \frac{1}{t^{2}},
\end{equation}
identically in $t$.
\uses{\cref{prop:hubblelaw}, \cref{prop:density}}
\end{proposition}

\begin{proof}
Two evaluations, each one line. \emph{Hubble side:} $H = \dot R/R = c/(ct) = 1/t$ by \cref{eq:hubblelaw}, so $H^{2} = 1/t^{2}$. \emph{Density side:}
\begin{equation}\label{eq:constraint-calc}
\frac{8\pi G}{3}\,\rho = \frac{8\pi G}{3}\cdot\frac{3}{8\pi G t^{2}} = \frac{1}{t^{2}},
\end{equation}
the factors $8\pi G/3$ and $3/8\pi G$ cancelling term by term --- they are reciprocals by construction, since the $8\pi/3$ of the left is the very constant born in \cref{eq:density}. The two sides are therefore equal to each other and to $1/t^{2}$, at every $t$. The same identity was registered as a self-check in \cref{prop:selfcheck}; it is restated here because the present chapter places it against the external equation.
\end{proof}

Before stating what this coincides with, one clarification about the ingredients' provenance. Every quantity in \cref{eq:constraint} is a reading-layer quantity: $\rho$ is total content over total volume, $H$ is a ratio of readings, and both were fixed by \cref{ch:newton} before this chapter began. No cosmological model entered the derivation, and none could: the arithmetic of \cref{eq:density} involves only $c$, $G$, and $t$.

\begin{eresult}[Friedmann constraint]\label{er:friedmann}
For a spatially flat ($k = 0$), cosmological-constant-free ($\Lambda = 0$) FLRW spacetime with energy density $u$,
\begin{equation}\label{eq:friedmann}
H^{2} = \frac{8\pi G}{3c^{2}}\,u.
\end{equation}
\reference{A.~Friedmann, Z.\ Phys.\ \textbf{10}, 377 (1922); \textbf{21}, 326 (1924). The $G_{00}$ (Hamiltonian) constraint of general relativity, quoted at the level at which it is used: as a relation among measurable kinematic quantities, independent of what drives the expansion.}
\end{eresult}

\begin{remark}[The coincidence, stated exactly and bounded]\label{rem:coincidence}
With $u = \rho c^{2}$ --- the definition of the energy density --- the right side of \cref{eq:friedmann} is $\frac{8\pi G}{3c^{2}}\,\rho c^{2} = \frac{8\pi G}{3}\rho$, identical to the left side of \cref{eq:constraint}. The two relations therefore coincide on the boundary content: the Friedmann constraint evaluated on the reading's density \cref{eq:density} returns exactly $1/t^{2}$, and the reading's own constraint \cref{eq:constraint} is the same arithmetic. The coincidence is constructed from two independent machines --- the Newtonian energetics of \cref{ch:newton} on the framework side, the geometry of \cref{er:friedmann} quoted externally --- and it is one-way in each use: the framework does not derive \cref{eq:friedmann}, and \cref{eq:friedmann} does not generate the framework's injection structure or its content. It is a coincidence of form at the level of the constraint alone: nothing is asserted about the $G_{ij}$ components of relativistic dynamics, and nothing about the framework's deeper structure is inferred from the match. In the framework's own accounting the constraint is an identity; on the relativistic side it is a dynamical equation, and that an identity of one description coincides with an equation of another is the precise content of the contact --- no more, and the chapter keeps it at no more.
\end{remark}

One word, finally, about what exactly is borrowed in this chapter, since the borrowing recurs below. The relations quoted from FLRW cosmology are its \emph{kinematics}: statements about how measurable quantities --- redshift, scale factor, time --- are related, given the symmetry. Kinematics is mathematics once the symmetry is granted, and it keeps that validity whatever the physical picture behind the expansion turns out to be; it is precisely this double status that makes a borrowed kinematics useful as a standard. A definite mathematical prediction turns a departure into a measurable quantity, and it has happened before that the mismatch, not the match, did the work: the Newtonian kinematics of Mercury's orbit predicted a perihelion advance that observation refused to supply,\footnote{U.~J.~Le Verrier, \emph{Th\'eorie du mouvement de Mercure}, Annuaire de l'Observatoire Imp\'erial de Paris (1859); the residual advance of about $43$ arcseconds per century stood against the Newtonian prediction until Einstein's account of 1915.} and the anomaly was the instrument that led to general relativity. The kinematic borrowings of \cref{sec:age-hubble} below serve the same function here: a ruler whose deviations are themselves the observable content.

\section{The acceleration component}\label{sec:acceleration-component}

The second contact concerns the second derivative of the scale factor. On the framework side there is nothing to solve --- only one function is available, and its derivatives are arithmetic.

\subsection{The scale-factor bridge}\label{sec:bridge}

For the comparison with FLRW language, the framework's radius must be identified with the scale factor of the relativistic description. There is exactly one homogeneous length in the framework, and the identification is declared where it is first needed:
\begin{equation}\label{eq:scalebridge}
a(t) := R(t).
\end{equation}
The bridge \cref{eq:scalebridge} is a definitional reading, not a derived relation: the framework defines $R$ from the frame sequence, and the FLRW formalism defines $a$ from the metric; the two are identified for the comparison, and if the identification were ever to fail, the failure would be registered as a failure of the reading, not argued away.

\begin{remark}[Rescaling invariance]\label{rem:rescale}
Only ratios survive the bridge. A constant rescaling $a \to \lambda a$ ($\lambda > 0$) leaves every quantity below unchanged:
\begin{equation}\label{eq:rescale}
\frac{\dot{(\lambda a)}}{\lambda a} = \frac{\dot a}{a},
\qquad
\frac{\ddot{(\lambda a)}}{\lambda a} = \frac{\ddot a}{a},
\end{equation}
the constants $\lambda$ cancelling in numerator and denominator separately. Statements about $\dot a/a$ and $\ddot a/a$ are therefore statements about the framework's shape, not about its normalisation, and the comparison is insensitive to how the two normalisations are chosen.
\end{remark}

\begin{proposition}[Zero acceleration, framework side]\label{prop:zero-accel}
With $a = R = ct$ (\cref{eq:scalebridge}, \cref{eq:ct}),
\begin{equation}\label{eq:zero-accel}
\dot a = c = \text{const}, \qquad \ddot a = 0.
\end{equation}
\uses{\cref{eq:ct}, \cref{eq:scalebridge}}
\end{proposition}

\begin{proof}
Two differentiations, each one step. First: $\dot a = d(ct)/dt = c$, a constant, since $c$ is the same at every $t$. Second: the derivative of a constant is zero, $\ddot a = dc/dt = 0$. No equation of motion is solved in either step; both are the arithmetic of the reading $R = ct$.
\end{proof}

\subsection{The relativistic statement and the match}\label{sec:accel-match}

\begin{eresult}[Acceleration equation]\label{er:accel}
For an FLRW spacetime with energy density $u$ and pressure $p$,
\begin{equation}\label{eq:accel}
\frac{\ddot a}{a} = -\frac{4\pi G}{3c^{2}}\bigl(u + 3p\bigr).
\end{equation}
\reference{Second Friedmann equation; A.~Friedmann (1922/1924), and standard thereafter. Quoted as the external statement of relativistic dynamics at the level used: a relation among the kinematic quantity $\ddot a/a$ and the content $(u, p)$. The framework borrows the form; nothing on the framework side generates it.}
\end{eresult}

The units of \cref{eq:accel} close as they should, and the closure is one line worth writing: $G/c^{2}$ carries $[\,{\rm m/kg}\,]$, the combination $(u + 3p)$ carries $[\,{\rm J/m^{3}} = {\rm kg/(m\,s^{2})}\,]$, and the product is $[\,{\rm s^{-2}}\,]$ --- the units of $\ddot a/a$. The constant $4\pi/3$ is the same geometric constant that appeared in the constraint, and it appears here for the same reason: both equations are moments of the same spherical averaging.

\begin{proposition}[Coincidence of the acceleration component]\label{prop:accel-check}
With the effective equation of state $p = w\,u$ and $w = -\tfrac13$ (\cref{prop:eos}),
\begin{equation}\label{eq:substitute}
u + 3p = u + 3wu = u\,(1 + 3w) = u\Bigl(1 + 3\cdot\bigl(-\tfrac13\bigr)\Bigr) = u\,(1 - 1) = 0,
\end{equation}
so the right side of \cref{eq:accel} vanishes identically, coinciding with the framework side \cref{eq:zero-accel}.
\uses{\cref{prop:eos}, \cref{er:accel}, \cref{eq:ct}}
\end{proposition}

One honesty flag is planted with the substitution itself. The $w$ of \cref{eq:substitute} is the \emph{effective} index of \cref{rem:instantaneous}: the deposited material is pressureless, and $w = -\tfrac13$ belongs to the effective dilution of the averaged content together with the conservation scaling. The acceleration equation \cref{eq:accel}, quoted from relativistic dynamics, is a statement about an effective fluid $(u, p)$; the match of \cref{prop:accel-check} is therefore a match at the level of effective description, which is the only level at which the quoted equation applies. Nothing weaker or stronger is claimed.

\begin{remark}[Two components, one status each]\label{rem:two-components}
The boundary content now satisfies both relativistic components at once: the constraint (\cref{rem:coincidence}) and the acceleration (\cref{prop:accel-check}). The two matches are registered separately, each with its own direction: the framework does not derive either equation, and neither equation is used below as a premise for framework content. The acceleration match is the stronger-looking of the two, since \cref{eq:accel} is a genuinely dynamical statement --- and its strength is exactly bounded by the flag above: it is a match of the effective fluid description, on both sides, nothing more.
\end{remark}

\section{The chain}\label{sec:chain}

The two components of \cref{sec:constraint,sec:acceleration-component} are the ends of a longer object: a chain of six equivalent statements, each of which is derivable from its neighbours by one move, whose only external content is a single link. The chain is the chapter's structural result; its links are proved here one at a time, with the conditions of each stated as they are used.

\begin{theorem}[The uniform-expansion chain]\label{thm:chain}
For the single boundary content of the framework, with $w = p/u$ the effective equation-of-state parameter and $u = \rho c^{2}$ the energy density,
\begin{equation}\label{eq:chain}
\begin{gathered}
w = -\tfrac{1}{3}
\;\Leftrightarrow\;
u + 3p = 0
\;\Leftrightarrow\;
\ddot a = 0
\;\Leftrightarrow\;
a \propto t\\
\;\Leftrightarrow\;
H = \tfrac{1}{t}\ \forall t
\;\Rightarrow\;
L_{0} = 1
\;\Leftrightarrow\;
L_{\partial} = \tfrac{c}{H} = ct .
\end{gathered}
\end{equation}
\uses{link by link below; the only external input is the acceleration equation (link 2)}
\end{theorem}

\begin{proof}
\textbf{Link 1 --- definition.} $w = -\tfrac13 \Leftrightarrow u + 3p = 0$, with no conditions beyond the definition $w := p/u$ (\cref{prop:eos}). Both directions are one substitution: if $w = p/u$, then $u + 3p = u(1 + 3w)$, which vanishes iff $w = -\tfrac13$. The combination $u + 3p$ is the \emph{active mass}: the part of the content that sources the second derivative. This link is pure arithmetic.

\textbf{Link 2 --- external dynamics.} $u + 3p = 0 \Leftrightarrow \ddot a = 0$ \emph{via the relativistic acceleration equation} \cref{eq:accel}, quoted from \cref{er:accel}. Both directions: $u + 3p = 0$ puts the right side of \cref{eq:accel} to zero, giving $\ddot a = 0$; conversely, $\ddot a = 0$ with $a > 0$ forces $u + 3p = 0$ through the same equation. \emph{Condition and relativistic content:} this is the only link whose statement is relativistic dynamics; the framework side of the equivalence is verified separately in \cref{prop:accel-check}, and nothing below the present section uses the implication in the $\Leftarrow$ direction.

\textbf{Link 3 --- integration, with one boundary condition.} $\ddot a = 0 \Rightarrow a \propto t$, in two integrations, each written out:
\begin{equation}\label{eq:twice}
\ddot a = 0 \;\xrightarrow{\ \text{integrate once}\ }\; \dot a = \alpha \ (\text{const})
\;\xrightarrow{\ \text{integrate twice}\ }\; a(t) = \alpha t + b,
\end{equation}
the constants $\alpha$, $b$ fixed by data. The single condition enters here: the zeroth frame is empty, $a(0) = 0$ (\cref{cor:counting} --- the frame count starts at $n = 1$, and $t = 0$ is the limit of an empty frame sequence), so $a(0) = b = 0$, leaving $a(t) = \alpha t$, i.e.\ $a \propto t$. The condition is registered as a condition: without $a(0) = 0$, link 3 would deliver the affine family $a = \alpha t + b$, and the chain would stop at ``constant velocity plus offset''. The framework's content fixes $b = 0$.

\textbf{Link 4 --- definition.} $a \propto t \Rightarrow H = \tfrac1t$ for all $t$: with $a = \alpha t$,
\begin{equation}\label{eq:link4}
H = \frac{\dot a}{a} = \frac{\alpha}{\alpha t} = \frac{1}{t},
\end{equation}
the constant $\alpha$ cancelling against itself --- the rescaling freedom of \cref{rem:rescale} dying exactly here, which is why link 4 holds for \emph{every} $t$ and not as an asymptotic statement: there is no exceptional instant, no early-time or late-time regime, in the identity $H = 1/t$.

\textbf{Link 5 --- value and naming.} $H = \tfrac1t \Rightarrow t_{0}H_{0} = 1$: evaluate at the present age $t = t_{0}$ (\cref{def:hubble}),
\begin{equation}\label{eq:link5}
H_{0} = \frac{1}{t_{0}} \;\Leftrightarrow\; t_{0}H_{0} = 1 .
\end{equation}
The product is named
\begin{equation}\label{eq:L0}
L_{0} := t_{0}H_{0},
\end{equation}
the \emph{age--Hubble product}, a dimensionless derived quantity ($[T]\times[T^{-1}]$). In frame units the same statement reads $H(n) = 1/(n t_{\mathrm P})$ for every $n$ (\cref{eq:tick}): the per-frame content of link 5 is that today's frame count equals the inverse of the current frame rate, a fact about the whole trajectory and not only about today. The development of $L_{0}$ is \cref{sec:age-hubble}; the chain itself needs only \cref{eq:link5}.

\textbf{Link 6 --- stepping.} $L_{\partial} = c/H = ct$ for all $t$: with $H = 1/t$,
\begin{equation}\label{eq:link6}
L_{\partial} = \frac{c}{H} = \frac{c}{1/t} = ct,
\end{equation}
one division, and the right side is the boundary radius itself (\cref{lem:boundary-sphere}: $L_{\partial} = n\ell_{\mathrm P}$ with $n\ell_{\mathrm P} = ct$ by $\ell_{\mathrm P} = ct_{\mathrm P}$ and $t = nt_{\mathrm P}$). The last equivalence of \cref{eq:chain} is therefore a restatement of the stepping structure, read through the Hubble ratio.
\end{proof}

\begin{remark}[Where the external content sits]\label{rem:gr-content}
Exactly one link of the six carries general-relativistic dynamics: link 2. Removing it leaves the internal chain
\begin{equation}\label{eq:internal-chain}
w = -\tfrac13,\ \ a = ct \ \Longrightarrow\ H = \tfrac1t\ \forall t\ \Longrightarrow\ L_{0} = 1,
\end{equation}
in which every implication is framework-internal: $a = ct$ is the reading (\cref{eq:ct}), proved without \cref{er:accel}; $H = 1/t$ is link 4's arithmetic; $L_{0} = 1$ is link 5's evaluation. The relativistic content of the full chain is thus confined to the single equivalence $u + 3p = 0 \Leftrightarrow \ddot a = 0$, and the $\Rightarrow$ direction of the internal chain is free of relativistic input by construction.
\end{remark}

\begin{remark}[What the chain cancels, and what remains]\label{rem:what-remains}
Read end to end, the chain performs one long cancellation. The scale $\alpha$ of the expansion cancels in link 4; the frame count converts to time through the tick in link 5; the boundary radius re-emerges through $\ell_{\mathrm P} = ct_{\mathrm P}$ in link 6; and the active mass vanishes through $1 + 3w = 0$ at the start. What survives the cancellation is not a number but a \emph{trajectory}: $H = 1/t$ at every $t$, or equivalently the per-frame statement $H(n) = 1/(nt_{\mathrm P})$ --- a curve, not a point. A single value of $L_{0}$ is a compression of that curve, and compressions lose information that only the curve carries.
\end{remark}

\section{The age--Hubble product}\label{sec:age-hubble}

\begin{definition}[The referents]\label{def:referents}
\begin{center}
\begin{tabular}{@{}p{2.6cm}p{3.8cm}p{2.8cm}p{2.0cm}@{}}
\toprule
Symbol & Referent & Layer & Character \\
\midrule
$t_{0}$ & age: count of all frames, $t_{0} = n_{0}t_{\mathrm P}$ & frame sequence & accumulated quantity \\
$H_{0}$ & current expansion rate & observation (distance ladder) & instantaneous reading \\
$L_{0} := t_{0}H_{0}$ & age--Hubble product & derived & dimensionless \\
$n_{\mathrm H} := 1/(H_{0}t_{\mathrm P})$ & Hubble frame count & derived & integer in the frame picture \\
\bottomrule
\end{tabular}
\end{center}
The last row carries a flag: the integrality of $n_{\mathrm H}$ is a statement of the frame sequence, not a measurable fact --- the number is too large to count --- and it is registered as such, used nowhere below.
\end{definition}

\begin{proposition}[Frame-counting identity]\label{prop:framecount}
\begin{equation}\label{eq:framecount}
L_{0} = t_{0}H_{0} = n_{0}\,t_{\mathrm P}\,H_{0} = n_{0}\,(H_{0}t_{\mathrm P}) = \frac{n_{0}}{n_{\mathrm H}},
\end{equation}
three substitutions, each shown: $t_{0} = n_{0}t_{\mathrm P}$ (\cref{eq:tick}), then the grouping $t_{\mathrm P}H_{0} = H_{0}t_{\mathrm P} = 1/n_{\mathrm H}$ (\cref{def:referents}). Hence
\begin{equation}\label{eq:framecount-equiv}
L_{0} = 1 \;\Leftrightarrow\; n_{0} = n_{\mathrm H} \;\Leftrightarrow\; t_{0} = \frac{1}{H_{0}}:
\end{equation}
the age--Hubble product is the ratio of the age frame count to the Hubble frame count, and it equals $1$ exactly when today's frame count is the count of frames the current rate would need --- a statement about the whole trajectory, since both $n_{0}$ and $n_{\mathrm H}$ are read off it.
\uses{\cref{eq:tick}, \cref{eq:L0}, \cref{def:referents}}
\end{proposition}

The second identity for $L_{0}$ borrows the kinematics of the FLRW symmetry: the change of variables between cosmic time and redshift is a statement of that symmetry, valid as mathematics whatever drives the expansion, and valuable exactly because a departure from it would be measurable.

\begin{eresult}[Age integral]\label{er:age-integral}
Along the light cone of an FLRW spacetime, $dt = -\frac{dz}{(1+z)H(z)}$, so
\begin{equation}\label{eq:age-integral}
t_{0} = \int_{0}^{\infty}\frac{dz}{(1+z)H(z)},
\qquad
L_{0} = \int_{0}^{\infty}\frac{dz}{(1+z)\,E(z)},
\quad E(z) := \frac{H(z)}{H_{0}}.
\end{equation}
\reference{Standard FLRW kinematics; the change of variables uses only the symmetry (redshift--scale-factor relation), no dynamics.}
\end{eresult}

\begin{remark}[What the integrand refers to]\label{rem:integrand}
The kernel $\bigl((1+z)E(z)\bigr)^{-1}$ is the weight of the redshift shell $dz$ in the total age: $L_{0}$ is a \emph{compression of the whole history} into one number, the sum of weights over all shells. Compressions lose information --- two different histories with the same total weight are indistinguishable to $L_{0}$ --- and the framework's output, by \cref{rem:what-remains}, is the curve $H(n)$, not the number.
\end{remark}

\begin{proposition}[Uniform case and the sufficiency directions]\label{prop:uniform-integral}
For the uniform trajectory $H = 1/t$: $t = t_{0}/(1+z)$, so $E(z) = H(z)/H_{0} = 1+z$, the weight is $(1+z)^{-2}$, and
\begin{equation}\label{eq:uniform-integral}
L_{0} = \int_{0}^{\infty}\frac{dz}{(1+z)^{2}} = \Bigl[-\frac{1}{1+z}\Bigr]_{0}^{\infty} = 1,
\end{equation}
the antiderivative evaluated at both limits. Conversely, \emph{if} $E(z) \geq 1+z$ at every $z$ (strictly somewhere), then the weight is dominated, $\bigl((1+z)E(z)\bigr)^{-1} \leq (1+z)^{-2}$, and $L_{0} < 1$ --- a history whose expansion ran faster than uniform on average is younger than uniform at the same $H_{0}$; symmetrically $E(z) \leq 1+z$ gives $L_{0} > 1$.
\uses{\cref{er:age-integral}, \cref{eq:hubblelaw}}
\end{proposition}

\begin{remark}[Necessary, not sufficient --- and why]\label{rem:necessary-not-sufficient}
The converse of the uniform case fails: $L_{0} = 1$ does not force the uniform history. A history decelerating early and accelerating late can have its weights arranged to sum to exactly $1$ while $E(z)$ differs from $1+z$ on every shell --- the cancellation is exact in the integral and invisible to it. The single number is therefore \emph{necessary} for uniformity --- any non-unit value excludes the trajectory outright --- and \emph{not sufficient}; sufficiency lives in the curve. The framework predicts the curve $H(n) = 1/(n t_{\mathrm P})$ (\cref{thm:chain}); the product is its shadow.
\end{remark}

\begin{remark}[The three cases and the parameter-free references]\label{rem:three-cases}
\begin{center}
\begin{tabular}{@{}lll@{}}
\toprule
$L_{0}$ & Name & Meaning \\
\midrule
$= 1$ & uniform (the framework's value) & $H = 1/t$ at every frame \\
$< 1$ & deceleration-dominated & at equal $H_{0}$, younger than uniform \\
$> 1$ & acceleration-dominated & at equal $H_{0}$, older than uniform \\
\bottomrule
\end{tabular}
\end{center}
For orientation, two parameter-free kinematic references, both derived in one line from their trajectories: dust, $a \propto t^{2/3}$, gives $H = \dot a/a = 2/3t$, hence $H_{0}t_{0} = 2/3$; radiation, $a \propto t^{1/2}$, gives $H = 1/2t$, hence $H_{0}t_{0} = 1/2$. These are values of the \emph{kinematics}, quoted as external standards, with no dynamical model attached on either side; model-dependent fits are deliberately not listed here, and the reason is the discipline of this chapter: it borrows forms, not parameterisations. The framework's value is not ``approximately $1$''; it is exactly $1$, by \cref{thm:chain}.
\end{remark}

\begin{remark}[Direction]\label{rem:l0-direction}
$L_{0} > 1$ is not a framework prediction, and the framework's own value is not a fit: the product is fixed to $1$ at every frame, and any confrontation with observation would have to address the trajectory $H(n)$ and the inequality form --- the universe must be older than its oldest content --- never ``$L_{0} \approx 1$''. Writing the contact as ``approximate agreement'' would already concede the parameterisation this chapter refuses. No such confrontation is undertaken in this paper.
\end{remark}

\section{Degrees of freedom without a reading}\label{sec:no-dof}

The relativistic dynamics quoted in \cref{sec:constraint,sec:acceleration-component} is the $G_{00}$ constraint and its $G_{ij}$ companion used only through the acceleration equation. Observables that live on other components of the field equations --- metric perturbations, gravitational waves, gravitational lensing --- have no matching term anywhere in the framework's algebra, and this section registers that fact, first as it stands and then as what it follows from.

\begin{remark}[The absence, registered as a blank]\label{rem:blank}
The absence of a reading for the $G_{ij}$-sector observables is a \textbf{blank of the reading layer, not an ontological assertion}: the framework has no registered quantity corresponding to them, and it asserts neither their existence nor their non-existence. Gravitational radiation and the growth of structure are accordingly outside this chapter's jurisdiction, at the same specification as the prohibitions of \cref{ch:boundary} --- the framework declines the statements, it does not deny the phenomena.
\end{remark}

\begin{remark}[Where the blank comes from: the hypothesis layer, not an omission]\label{rem:no-dynamics-why}
The blank is not a piece the framework forgot to install; it is a consequence of what the hypothesis layer fixed, and the line of consequence is worth retracing, because the energy--time hypothesis was not an arbitrary definition. Three fixations, each made in \cref{sec:hypotheses}, jointly remove the slot into which dynamics would have to go:
\begin{enumerate}
\item \emph{Time is a sequence marker.} \cref{as:a1} registers processes on a discrete index, and time enters every formula only as $t = n\,t_{\mathrm P}$ (\cref{eq:tick}). A sequence marker has no rate: there is nothing to differentiate, and no initial-value problem can be posed on the frame sequence itself.
\item \emph{Action is quantized per tick.} \cref{as:a2} quantizes every exchange into integer cells of $\alpha_{0}$ (\cref{eq:quantization}), so the contents the framework tracks are counts --- and counts evolve by stepping, not by equations. The injection theorem (\cref{thm:injection}) then fixed the step's size: one minimal action quantum per frame, no more and no less.
\item \emph{The matching degenerates the would-be equation.} When the Newtonian energetics was matched to the frame sequence (\cref{prop:compatibility}), the energy equation collapsed to an identity: $E = \tfrac12 c^{2} - GM/R \equiv 0$. What would have been the dynamics of the radius became the arithmetic of frame counting (\cref{rem:resolved}); $R = ct$ is a reading, and the question of what dynamics drives it has no slot in which to be posed.
\end{enumerate}
The integrability discipline (\cref{rem:integrability}) completes the picture: even where the framework does accumulate --- state quantities only --- the accumulation is the telescope of a frame difference, never the solution of a differential equation. A framework whose time is a marker, whose contents are counts, and whose would-be equation of motion is an identity, registers; it does not evolve. The absence registered in \cref{rem:blank} is the same fact, seen from the relativistic side: the sector of the field equations that describes evolution has no counterpart here, because nothing here evolves. And the consequence runs in the direction the abstract fixed: dynamics, if it is to enter, must enter as additional interactions or hypotheses (\cref{sec:hypotheses}) --- outside the hypothesis layer, never derivable from it.
\end{remark}

\section{The cosmological constant has no slot}\label{sec:lambda}

The two Friedmann components quoted in this chapter belong to the $\Lambda = 0$ member of the relativistic family. This section makes that selection explicit, by the same device as the rest of the chapter: the two sides of the equation are placed next to each other, and the framework's side leaves no room for the additive constant.

\begin{eresult}[Friedmann constraint with cosmological constant]\label{er:friedmann-lambda}
\begin{equation}\label{eq:friedmann-lambda}
H^{2} = \frac{8\pi G}{3c^{2}}\,u + \frac{\Lambda}{3},
\end{equation}
with $\Lambda$ the cosmological constant, the single free additive constant of the family.
\reference{The $\Lambda$-extension of \cref{eq:friedmann}, introduced by Einstein (1917) for a static universe model;\footnote{A.~Einstein, \emph{Kosmologische Betrachtungen zur allgemeinen Relativit\"atstheorie}, Sitzungsber.\ Preuss.\ Akad.\ Wiss.\ (1917). The constant was absent from the field equations of 1915 (\emph{Die Feldgleichungen der Gravitation}, Sitzungsber.\ Preuss.\ Akad.\ Wiss.\ 1915), which contained only the coupling fixed by the Newtonian limit.} quoted as the external family of which \cref{eq:friedmann} is the $\Lambda = 0$ member.}
\end{eresult}

The framework's constraint is an identity fixed point by point (\cref{prop:constraint}): at every $t$,
\begin{equation}\label{eq:framework-side}
H^{2} - \frac{8\pi G}{3c^{2}}\,u = \frac{1}{t^{2}} - \frac{1}{t^{2}} = 0,
\end{equation}
the two terms cancelling as proved in \cref{eq:constraint-calc}. The relativistic family \cref{eq:friedmann-lambda} placed against the same quantities gives, on its side,
\begin{equation}\label{eq:gr-side}
H^{2} - \frac{8\pi G}{3c^{2}}\,u = \frac{\Lambda}{3}.
\end{equation}
Matching the two sides --- the only move this chapter ever makes --- requires $\Lambda/3 = 0$ at every $t$: any non-zero $\Lambda$ would leave a non-zero residual in \cref{eq:framework-side} where the framework has none. The framework's value of the cosmological constant is therefore
\begin{equation}\label{eq:lambda-zero}
\Lambda = 0,
\end{equation}
not by assumption but by comparison: the reading fixes every term of the constraint, and the one free constant of the relativistic family is squeezed to zero between two quantities that coincide identically.

The deeper content of the squeeze is structural, and it is worth one paragraph. In the field equations, $\Lambda$ is an \emph{integration-constant-like} addition: the 1915 equations contained no free constant, the coupling alone being fixed by the Newtonian limit, and the 1917 constant entered as a model-driven modification, chosen to hold a static universe. The framework's reading sides with the parameter-free form: its constraint is the parameter-free arithmetic of \cref{ch:newton}, and the additive constant has no slot in it --- just as the $G_{ij}$ sector had no reading in \cref{sec:no-dof}, the $\Lambda$ term has none here, and for the same reason: there is nothing in the frame sequence for it to be measured against. The reading thereby selects, of the one-parameter relativistic family, precisely the member Einstein first wrote down.

\begin{remark}[Status and jurisdiction]\label{rem:lambda-status}
Three boundaries, per the chapter's discipline. (i) \cref{eq:lambda-zero} is the value forced by matching the reading; it is a registered coincidence-extension of the same kind as \cref{rem:coincidence}, one-way in each use, and nothing in the framework generates the field equations with or without $\Lambda$. (ii) The framework has no cosmological-constant \emph{problem}: the question why a constant is so small does not arise where the constant has no slot; this is the same jurisdictional silence as \cref{ch:boundary}'s prohibition on dark-energy identifications, and it asserts no solution to any constant problem. (iii) If observation were ever to require $\Lambda \neq 0$, the failure would be the reading's, registered as a failure of the reading and of nothing else.
\end{remark}

\section{Summary of the reading}\label{sec:gr-summary}

\begin{remark}[What Chapter 8 has established]\label{rem:ch8summary}
Three contacts between the framework and the Friedmann equations of general relativity, each stated where it occurs and each bounded the same way.

\emph{The constraint.} The boundary content satisfies $H^{2} = \frac{8\pi G}{3}\rho = 1/t^{2}$ identically (\cref{prop:constraint}), with the density computed in full in \cref{prop:density}; the Friedmann constraint, quoted externally at $k = 0$, $\Lambda = 0$, coincides with it term by term (\cref{rem:coincidence}) --- an identity of the framework matched against an equation of relativistic dynamics, recorded and not explained.

\emph{The acceleration.} The scale factor of the reading obeys $\ddot a = 0$ as pure arithmetic (\cref{prop:zero-accel}), and the relativistic acceleration equation vanishes on the boundary content's effective fluid, $u + 3p = 0$ at $w = -\tfrac13$ (\cref{prop:accel-check}) --- a match at the level of effective description, which is the only level at which the quoted equation applies, flagged where it is made.

\emph{The chain.} Six links, from the active mass $u + 3p$ through $H = 1/t$ for all $t$ to $L_{\partial} = c/H = ct$ (\cref{thm:chain}), carry the entire relativistic content of the chapter in a single link (link 2); removed, the internal chain $w = -\tfrac13,\ a = ct \Rightarrow H = 1/t \Rightarrow L_{0} = 1$ stands framework-internal (\cref{rem:gr-content}), and what survives every cancellation is the trajectory $H(n) = 1/(nt_{\mathrm P})$ --- a curve, not a point (\cref{rem:what-remains}).

\emph{The age--Hubble product.} $L_{0} = n_{0}/n_{\mathrm H}$ as frame arithmetic (\cref{prop:framecount}); as a history-integral it borrows FLRW kinematics --- valid mathematics under the symmetry, a standard whose deviations are the observable content, in the precedent of every kinematic anomaly from Mercury onward (\cref{rem:integrand,rem:necessary-not-sufficient}); the uniform trajectory gives $L_{0} = 1$ exactly, the single value is necessary but not sufficient for uniformity, and the parameter-free kinematic references of \cref{rem:three-cases} are quoted without any model fit entering the chapter.

\emph{The two silences.} The $G_{ij}$-sector observables have no reading (\cref{rem:blank}), and the cosmological constant has no slot (\cref{eq:lambda-zero}); both silences trace to the hypothesis layer --- time as a marker, contents as counts, the would-be equation an identity (\cref{rem:no-dynamics-why}) --- and the reading thereby selects, of the one-parameter relativistic family, the parameter-free form of 1915.

The chapter's input account: four empirical inputs (\cref{er:friedmann,er:accel,er:age-integral,er:friedmann-lambda}), all quoted one-way and none generating framework content; one definitional reading (\cref{eq:scalebridge}); no new hypothesis; and numerical content limited to the parameter-free kinematic references --- no observation, no model parameterisation, no fit, entered the chapter. In one sentence: the uniform-expansion member of the Friedmann equations is the arithmetic of this framework's reading --- the rest of the relativistic theory is present here only as silence, and the silence is registered.
\end{remark}

\chapter{Symmetries and geometry}\label{ch:symmetries}

This chapter records the framework's symmetries and their geometric residences. Each object is defined before any theorem about it is stated, and nothing in the chapter is assumed that the census table of \cref{sec:sym-register} does not list.

\section{The census table}\label{sec:sym-register}

Every symmetry is entered with four items: its \emph{source construction} (which established structure it grows from), its \emph{existence status} (exists / does not exist / exists but trivial), its \emph{conserved object}, and its \emph{status label}. Every geometric residence is entered with three: its \emph{construction}, its \emph{status} ($\checkmark$ built, $\times$ proved absent, N/A no such object, or registered open), and its \emph{verdict}, with proof where a verdict is given. Speculative content is confined to remarks labelled as such, under the doctrine of \cref{rem:register}; the census table is the chapter's own device, and the chapter's remaining sections are its entries.

The working list of symmetries is the standard one of relativistic physics --- time translation, time rotation (boost), spatial rotation, spatial translation --- and the working list of residences is read off the chapters already written: the frame sequence (\cref{as:a1}), the reading layer (\cref{sec:units,sec:steps}), the boundary sphere (\cref{ch:boundary}), and the rapidity plane constructed in \cref{sec:rapidity-plane}. The contents of the lists are not assumed in advance; they are what the entries below establish. Where each entry lives: \cref{sec:time-translation} (time translation), \cref{sec:lattice-lorentz,sec:rotation,sec:translation} (the two rotations and the translation), \cref{sec:lattice-theorems} (the plane's lattice theorems), \cref{sec:census} (the residences), \cref{sec:sixteen} (the sixteen-fold grid), \cref{sec:repair} (the repair registry), \cref{sec:correspondence} (the force--geometry correspondence and the hypothesis layer's own symmetry record), and \cref{sec:backfill} (the backfill list).

\section{The geometry discipline}\label{sec:geometry-discipline}

One discipline governs the chapter's use of geometry.

\begin{remark}[What geometry the framework needs, and what it does not]\label{rem:geometry-discipline}
The framework's geometry is exhausted by four objects: the lattice $\mathbb{Z}^{4}$ of counting (\cref{as:a1,as:stepping}), the flat kinematics of the reading layer (\cref{sec:units}), the boundary sphere $S^{2}$ (\cref{lem:boundary-sphere}), and the rapidity plane of \cref{sec:rapidity-plane}. All four are low-dimensional, concrete, and equipped with explicit formulas; none requires a metric tensor field, and none requires the apparatus of Riemannian or pseudo-Riemannian geometry.

The reason is not that the Riemannian machine would output zeros here --- it would not. In a flat FLRW spacetime with $a(t) = ct$, the connection coefficients $\Gamma^{0}_{ij} = a\dot a\,\delta_{ij} = c^{2}t\,\delta_{ij}$ are non-zero, and the scalar curvature $R = 6\bigl(\ddot a/a + (\dot a/a)^{2}\bigr) = 6\bigl(0 + 1/t^{2}\bigr) = 6/t^{2}$ is non-zero as well: the curvature machinery is a box of live objects, not of zeros. The correct statement is sharper and entirely positive: the framework's constraint is one line of algebra (\cref{prop:constraint}), its reading kinematics is arithmetic (\cref{sec:steps}), and its boundary thermodynamics is sphere geometry (\cref{sec:geometry}) --- none of these computations ever calls for a curvature tensor, and none would be shortened by having one. Geometry enters where it is used, and it is not used here.

If a future reading ever involves non-zero spatial curvature, metric perturbations, or a discrete--smooth correspondence, Riemannian geometry may enter as an external comparison --- as it entered, one-way, in \cref{ch:gr}. Today no such reading exists (\cref{rem:blank}), and the discipline is simply the refusal to stock tools no current computation reaches for.
\end{remark}

\section{The order declaration}\label{sec:order-declaration}

\begin{remark}[Construction first, definition after]\label{rem:wigner}
The chapter writes each symmetry in the order \emph{construction} $\to$ \emph{definition} $\to$ \emph{consequence}. This is an epistemic ordering --- it reflects how the objects are most effectively described --- and not an ontological one: whether the mathematical structure or the physical phenomenon ``comes first'' is a question this chapter registers without answering, as the standing problem of the ``unreasonable effectiveness'' of mathematics in the natural sciences (Wigner, 1960).\footnote{E.~P. Wigner, \emph{The unreasonable effectiveness of mathematics in the natural sciences}, Comm.\ Pure Appl.\ Math.\ \textbf{13}, 1 (1960). Registered, not answered.} The same registration covers the correspondence between forces and geometric configurations examined in \cref{sec:correspondence}: the chapter cannot know which came first, and it does not pretend to; it records that the correspondence is effective, and argues from effectiveness alone.
\end{remark}

\section{The rapidity plane}\label{sec:rapidity-plane}

The boost --- the change to a frame in relative motion --- will turn out, in \cref{sec:lattice-lorentz}, to have no exact seat on the frame sequence. Before that theorem can even be stated, the seat it \emph{does} have must be built, because the matching of \cref{def:units} turns a pair of physical quantities into a plane on which the boost is an exact, identity-level symmetry. This section constructs that plane and records its invariant.

\begin{definition}[Plane coordinates]\label{def:plane}
\begin{equation}\label{eq:plane}
X := E\,t_{\mathrm P} \quad (\text{time side}), \qquad Y := p\,\ell_{\mathrm P} \quad (\text{space side}),
\end{equation}
where $E$ is an energy and $p$ its associated momentum, defined by the convention
\begin{equation}\label{eq:momentum}
p := \frac{E v}{c^{2}},
\end{equation}
with the endpoint check $v = c \Rightarrow p = E/c$ (\cref{er:efreq}: the photon relation $E = pc$) and the interpolation at $v < c$ recorded as a convention under the doctrine of \cref{rem:register}.
\end{definition}

The two coordinates share one dimension, and the sharing is the Planck relations: $[E t_{\mathrm P}] = [\hbar] = [p\ell_{\mathrm P}]$ by \cref{def:units}. The time axis of the plane is the quantization axis of \cref{eq:quantization} in units of $\alpha_{0}$: at $E = \tfrac{k}{2}\varepsilon_{\mathrm P}$,
\begin{equation}\label{eq:plane-grid}
\frac{X}{\hbar} = \frac{k}{2},
\end{equation}
so the time axis of the plane carries the half-integer grid of the energy--time hypothesis --- a fact returned to in \cref{sec:lattice-theorems}.

\begin{lemma}[Velocity as slope]\label{lem:slope}
\begin{equation}\label{eq:slope}
Y = \frac{v}{c}\,X,
\end{equation}
\uses{\cref{eq:plane}, \cref{def:units} (first Planck relation)}
\end{lemma}

\begin{proof}
Substitute \cref{eq:momentum} and convert one unit: $Y = p\ell_{\mathrm P} = \dfrac{Ev}{c^{2}}\cdot c\,t_{\mathrm P} = \dfrac{v}{c}\,E t_{\mathrm P} = \dfrac{v}{c}X$, where $\ell_{\mathrm P} = c\,t_{\mathrm P}$ does the only work.
\end{proof}

\Cref{eq:slope} is the statement that velocity is the slope of a ray through the origin of the plane: each state of motion is a ray, and the ray's slope is $v/c$.

\begin{definition}[The invariant]\label{def:invariant}
\begin{equation}\label{eq:invariant}
I := \sqrt{X^{2} - Y^{2}} = X\sqrt{1 - v^{2}/c^{2}},
\end{equation}
the second equality from \cref{eq:slope} by factoring $X > 0$ out of the square root.
\end{definition}

\begin{theorem}[The bilateral identity]\label{thm:bilateral}
With the mass shell $E\sqrt{1 - v^{2}/c^{2}} = mc^{2}$ \textnormal{(quoted externally)},
\begin{equation}\label{eq:bilateral}
\boxed{\;\sqrt{\bigl(E\,t_{\mathrm P}\bigr)^{2} - \bigl(p\,\ell_{\mathrm P}\bigr)^{2}} = \frac{\hbar\,m}{M_{\mathrm P}}\;}
\end{equation}
identically in $v$.
\uses{\cref{def:plane}, \cref{eq:invariant}, \cref{def:units}, mass shell (external)}
\end{theorem}

\begin{proof}
\emph{Step 1 (invariance).} Under a change of frame, $(E, pc)$ transforms as a hyperbolic rotation (the boost, quoted externally as the passive transformation between inertial frames),
\begin{equation}\label{eq:boost-quotes}
E' = E\cosh\theta - pc\sinh\theta, \qquad p'c = pc\cosh\theta - E\sinh\theta,
\end{equation}
with $\tanh\theta = v/c$. Multiply the first line by $t_{\mathrm P}$: $X' = E't_{\mathrm P} = E t_{\mathrm P}\cosh\theta - pc\,t_{\mathrm P}\sinh\theta = X\cosh\theta - Y\sinh\theta$, using $pc\,t_{\mathrm P} = p\ell_{\mathrm P} = Y$ --- one step, and the step where the matching is load-bearing. Similarly $Y' = Y\cosh\theta - X\sinh\theta$. Then
\begin{equation}\label{eq:casimir-inv}
\begin{gathered}
(X')^{2} - (Y')^{2} = (X^{2} - Y^{2})\bigl(\cosh^{2}\theta - \sinh^{2}\theta\bigr)\\
+ \text{cross terms } (-2XY + 2YX)\cosh\theta\sinh\theta,
\end{gathered}
\end{equation}
and the cross terms cancel identically while $\cosh^{2}\theta - \sinh^{2}\theta = 1$ (hyperbolic identity), leaving $X^{2} - Y^{2}$ unchanged: the combination is an invariant of the boost, carried by the group alone and requiring no action principle (\cref{rem:casimir} below).

\emph{Step 2 (evaluation).} An invariant is evaluated where it is easiest, at $v = 0$: there $p = 0$, $E = mc^{2}$, so $I = X = mc^{2}t_{\mathrm P} = m\,\dfrac{\hbar}{M_{\mathrm P}}$, using $M_{\mathrm P}c^{2} = \hbar/t_{\mathrm P}$ (\cref{cor:epmp} with \cref{def:units}). Hence \cref{eq:bilateral} at $v = 0$, and by invariance at every $v$.
\end{proof}

\begin{remark}[Why the minus sign]\label{rem:minus-sign}
Four geometries, one question --- what transformation a change of frame is:
\begin{center}
\begin{tabular}{@{}p{1.8cm}p{1.8cm}p{3.8cm}p{3.8cm}@{}}
\toprule
Geometry & Invariant & Change of frame & Framework \\
\midrule
Minkowski & $X^{2} - Y^{2}$ & hyperbolic rotation, symmetry & the rapidity plane (this section) \\
Euclidean & $X^{2} + Y^{2}$ & ordinary rotation & limit $c \to \infty$ (degenerate) \\
sphere ($K > 0$) & --- & no translation-like boost & no such object \\
saddle ($K < 0$) & --- & --- & --- \\
\bottomrule
\end{tabular}
\end{center}
Only the Minkowski signature makes the change of frame a symmetry of the plane's geometry; in the limit $c \to \infty$, $X^{2} - Y^{2} \to X^{2}$ and the plane degenerates to a single axis. The minus sign of \cref{eq:invariant} is, in this precise sense, the algebraic form of the statement that the speed of light is finite.
\end{remark}

\begin{remark}[Reading, and the Casimir remark]\label{rem:casimir}
Read conservatively: \cref{eq:bilateral} says that the plane's geometry assigns every motion a hyperbolic radius, $\hbar m/M_{\mathrm P}$, fixed by the rest energy alone; the photon lies on the lightlike ray $Y = X$ with $I = 0$, and subluminal motions have $I > 0$, so the invariant measures the imbalance between the time and space axes that a nonzero rest energy sustains. As a structural remark, certified by \cref{eq:casimir-inv}: $I$ depends only on the group of the plane and not on any action --- it is a Casimir, one level deeper than a Noether charge, which would require an action to define. Mass enters this framework as the value of a group invariant; nothing deeper than the group itself is invoked, and nothing deeper is claimed.
\end{remark}

\begin{remark}[Rays and mass shells]\label{rem:rays-shells}
Two families of curves organise the plane, and both are read off what is already proved. The \emph{rays} $Y = \tfrac{v}{c}X$ (\cref{eq:slope}) are the states of constant velocity: moving along a ray changes the energy at fixed mass and fixed speed --- a change of frame, not of content. The \emph{mass shells}
\begin{equation}\label{eq:shells}
X^{2} - Y^{2} = \Bigl(\frac{\hbar m}{M_{\mathrm P}}\Bigr)^{2}
\end{equation}
are the hyperbolae of \cref{thm:bilateral}, one per mass: moving along a shell changes the mass assignment at fixed invariant. The two families are transverse --- every point of the plane lies on exactly one ray and one shell --- and the plane's mixed texture is now visible: the time axis of the plane carries the discrete half-integer grid of \cref{eq:plane-grid}, while the mass shells are continuous, so the masses of free entities do not sit on the grid. The consequences of that mixed texture --- which velocities are quantizable, which masses are --- are the lattice theorems of \cref{sec:lattice-theorems}.
\end{remark}

\section{Time translation}\label{sec:time-translation}

The first symmetry is the one the framework modifies most severely, because the boundary is not a closed system: its content grows. What time translation means here must be assembled from three independent facts before any symmetry statement is made.

\begin{proposition}[Decomposition of the content]\label{prop:stock}
\begin{equation}\label{eq:stock}
E_{\partial}(n) = \underbrace{\mathcal P}_{\text{power}} \times \underbrace{n\,t_{\mathrm P}}_{\text{order}},
\qquad
\begin{cases}
\mathcal P = c^{5}/2G = \text{const} & \text{\cref{def:power}}\\
n\,t_{\mathrm P} = t & \text{\cref{eq:tick}}\\
\mathcal P\,t_{\mathrm P} = \varepsilon_{\mathrm P}/2 = \text{per-frame quantum} & \text{\cref{thm:injection}}
\end{cases}
\end{equation}
The stock $E_{\partial}(n)$ is not conserved; the flow $\mathcal P$ is. The three displayed facts are independent: the first is the fixed injection rate, the second the definition of the tick, the third the injection theorem --- no one of them is a consequence of the other two.
\uses{\cref{def:power}, \cref{eq:tick}, \cref{thm:injection}}
\end{proposition}

\begin{proof}
$E_{\partial}(n) = n\,\varepsilon_{\mathrm P}/2$ (\cref{eq:eboundary}); the injection theorem (\cref{eq:inj-theorem}) gives $\varepsilon_{\mathrm P}/2 = \mathcal P\,t_{\mathrm P}$; hence $E_{\partial}(n) = n\,\mathcal P\,t_{\mathrm P} = \mathcal P\cdot(nt_{\mathrm P})$, the last regrouping purely formal. Independence: \cref{def:power} is fixed from the stepping matching without the chain; \cref{eq:tick} is definitional; \cref{thm:injection} is the chain's theorem --- three different chapters, three different inputs.
\end{proof}

In the vocabulary of closed systems, \cref{prop:stock} would be the statement ``energy is conserved''; here it is decomposed into ``the rate is constant'' and ``the count advances'', and the conservation migrates from the stock to the flow.

\begin{proposition}[Function level: broken]\label{prop:function-broken}
For the frame translation $T_{k}: n \mapsto n + k$ ($k \geq 1$),
\begin{equation}\label{eq:broken}
(T_{k}E_{\partial})(n) = E_{\partial}(n+k) = \frac{(n+k)\,\varepsilon_{\mathrm P}}{2} = E_{\partial}(n) + \frac{k\,\varepsilon_{\mathrm P}}{2},
\end{equation}
step by step: \cref{eq:eboundary} at $n+k$, then distribution over the sum. The difference $k\,\varepsilon_{\mathrm P}/2$ is non-zero at every $n$ and every $k \geq 1$: the content function is changed by every non-trivial translation --- time translation is \emph{broken at the function level}.
\uses{\cref{eq:eboundary}}
\end{proposition}

\begin{proposition}[Difference level: survives]\label{prop:difference-survives}
The difference quotient is translation-invariant:
\begin{equation}\label{eq:difference}
\Delta E_{\partial}(n) := E_{\partial}(n+1) - E_{\partial}(n) = \frac{(n+1)\varepsilon_{\mathrm P}}{2} - \frac{n\varepsilon_{\mathrm P}}{2} = \frac{\varepsilon_{\mathrm P}}{2} = \mathcal P\,t_{\mathrm P},
\end{equation}
independent of $n$ by the cancellation of the two halves, and for any $k$,
\begin{equation}\label{eq:difference-t}
\Delta (T_{k}E_{\partial})(n) = E_{\partial}(n+k+1) - E_{\partial}(n+k) = \frac{\varepsilon_{\mathrm P}}{2},
\end{equation}
the same value, shifted. The broken symmetry therefore leaves an invariant behind: not a charge (no action, no Noether theorem is invoked --- the collision with that vocabulary is exactly what the present account avoids), but a \emph{difference invariant}, the constant per-frame increment.
\uses{\cref{eq:eboundary}, \cref{thm:injection}}
\end{proposition}

\begin{remark}[Root of the breaking]\label{rem:root-breaking}
\emph{Why} the function level must break: a symmetry of the content function would move every frame's content rigidly, $E_{\partial}(n+k) = E_{\partial}(n)$, and that is precisely a postponement of the injection --- no net content may enter while the origin is shifted. The injection, however, is defined as external to the system and non-postponable (\cref{as:pump}): each frame receives its quantum regardless of any relabelling. The symmetry that would require the source to wait cannot exist; what exists is the symmetry of the increment, which does not ask the source to wait. The breaking is therefore not a defect of the framework but a shadow of the injection's externality --- and the difference invariant that survives is the injection theorem wearing its symmetry clothes. The same fact, seen from the reading layer rather than the frame sequence, is \cref{rem:no-dynamics-why}: there, nothing evolves; here, nothing is invariant except what arrives.
\end{remark}

\begin{remark}[Registered note: a borrowed isomorphism]\label{rem:b1-registered}
The two-level structure --- constant norm-like increment against a growing stock --- is isomorphic, as a piece of mathematics, to the spectral situation of a linearly time-dependent Hamiltonian $\hat H \propto n$ in standard kinematics. The isomorphism is recorded here and used nowhere: the framework's injection is not a Hamiltonian, and no spectral theorem is invoked. Registered per \cref{rem:wigner}.
\end{remark}

\section{Time rotation: the lattice--Lorentz theorem}\label{sec:lattice-lorentz}

The boost --- the change to a frame in relative motion, the ``rotation'' that mixes time and space --- must first be denied its apparent home before its true home can be claimed. The theorem below is the denial, and its statement requires four objects to be defined with care, because the question ``does the boost preserve the lattice?'' is not even well-posed until each of them is fixed. The four objects come first.

\subsection{What the objects are}\label{sec:audit}

\emph{1. The frame sequence (the stage).} \cref{as:a1} registers events on a discrete sequence: integer frame $n$ and integer spatial site $j$, with spacings $t_{\mathrm P}$ and $\ell_{\mathrm P}$. The Planck relation $\ell_{\mathrm P} = c\,t_{\mathrm P}$ (\cref{def:units}) makes the two axes commensurable: $c\cdot n t_{\mathrm P} = n\ell_{\mathrm P}$, so in units of $\ell_{\mathrm P}$ the physical frame sequence \emph{is} the lattice $\mathbb{Z}^{2}$. Without this commensurability the preservation question below would have no meaning --- two axes of incomparable units admit no common transformation law. The physical sequence occupies the half-plane $n \geq 1$ (\cref{cor:counting}); the mathematical object of the question is the full $\mathbb{Z}^{2}$, and the theorem's conclusion is unaffected by the extension, as noted after the proof.

\emph{2. What a boost is.} On the reading layer, a boost is the change between inertial frames --- a hyperbolic rotation of the $(ct, x)$ plane through a rapidity $\theta \in \mathbb{R}$, with matrix
\begin{equation}\label{eq:boostmatrix}
B(\theta) = \begin{pmatrix} \cosh\theta & \sinh\theta \\ \sinh\theta & \cosh\theta \end{pmatrix},
\qquad \det B = \cosh^{2}\theta - \sinh^{2}\theta = 1,
\end{equation}
the unit determinant being built into the form (quoted externally as the passive Lorentz transformation, as in \cref{eq:boost-quotes}).

\emph{3. The lattice-preserving condition.} $B(\theta)$ preserves $\mathbb{Z}^{2}$ if and only if it carries lattice points to lattice points, and by linearity it suffices to check the basis: preservation requires
\begin{equation}\label{eq:preserve}
B(\theta)\binom{1}{0} = \binom{\cosh\theta}{\sinh\theta} \in \mathbb{Z}^{2}
\quad\Longrightarrow\quad
\cosh\theta = m \in \mathbb{Z},\qquad \sinh\theta = s \in \mathbb{Z}.
\end{equation}
The second basis vector then follows for free, $B(\theta)\binom{0}{1} = \binom{s}{m} \in \mathbb{Z}^{2}$, and the inverse $B(-\theta)$, having the same integer form, carries $\mathbb{Z}^{2}$ \emph{onto} itself --- surjectivity is automatic once \cref{eq:preserve} holds, and the condition is exactly \cref{eq:preserve}.

\emph{4. The theorem, stated exactly.} There is no non-trivial real rapidity $\theta$ for which $B(\theta)$ preserves $\mathbb{Z}^{2}$: among exact symmetries of the frame sequence, the boost is the identity alone.

\subsection{The proof: a squeeze between consecutive squares}\label{sec:squeeze}

\begin{theorem}[Lattice--Lorentz]\label{thm:lattice-lorentz}
If $B(\theta)$ preserves $\mathbb{Z}^{2}$, then $\theta = 0$.
\uses{\cref{eq:preserve}; the hyperbolic identity; one inequality}
\end{theorem}

\begin{proof}
By \cref{eq:preserve}, preservation forces $m = \cosh\theta$ and $s = \sinh\theta$ to be integers, and the hyperbolic identity becomes a Diophantine equation:
\begin{equation}\label{eq:diophantine}
\cosh^{2}\theta - \sinh^{2}\theta = 1
\quad\Longrightarrow\quad
m^{2} - s^{2} = 1
\quad\Longleftrightarrow\quad
(m - s)(m + s) = 1.
\end{equation}
Two cases exhaust the integers $m \geq 1$ (note $m = \cosh\theta \geq 1$ for real $\theta$).

\emph{Case $m \geq 2$ (the non-identity boost).} Then $s^{2} = m^{2} - 1$, and $m^{2} - 1$ lies strictly between consecutive squares:
\begin{equation}\label{eq:squeeze}
(m-1)^{2} = m^{2} - 2m + 1 \;<\; m^{2} - 1 \;<\; m^{2} \qquad (m \geq 2),
\end{equation}
where the left inequality is $m^{2} - 2m + 1 < m^{2} - 1 \Leftrightarrow -2m + 1 < -1 \Leftrightarrow 2m > 2$, true for every $m \geq 2$; the right inequality is $-1 < 0$. A number strictly between two consecutive squares is not a square, so $m^{2} - 1$ is not a perfect square, and $s = \sqrt{m^{2}-1} \notin \mathbb{Z}$ --- contradicting \cref{eq:preserve}.

\emph{Case $m = 1$.} Then $s^{2} = 1 - 1 = 0$, so $\sinh\theta = 0$, hence $\theta = 0$: the identity.

The version on the physical frame sequence $\mathbb{Z}^{4}$ is verbatim identical: a boost on any $(t, x_{i})$ plane imposes the same integer pair $(m, s)$ on that plane, and the same squeeze applies plane by plane.
\end{proof}

No step used anything but \cref{eq:preserve} and one inequality between integers: the theorem is pure number theory, and its physical content --- analysed next --- is entirely the content of the two assumptions that built the stage.

\subsection{What the theorem is, and is not}\label{sec:witness}

\begin{remark}[A theorem of the frame sequence, not a prohibition of physics]\label{rem:not-prohibition}
\Cref{thm:lattice-lorentz} concerns the frame sequence as a mathematical object: among the \emph{exact} symmetries of $\mathbb{Z}^{2}$, the boost is absent. It forbids nothing physical. On the reading layer, Lorentz kinematics is exact mathematics --- quoted as such throughout \cref{sec:rapidity-plane} --- and on the rapidity plane the boost is an identity-level symmetry (\cref{thm:bilateral}). The three-layer verdict, with the frame sequence's cell named by the theorem itself:
\begin{center}
\begin{tabular}{@{}lll@{}}
\toprule
Layer & Boost & Status \\
\midrule
frame sequence $\mathbb{Z}^{2}$ & $\times$: identity only & \cref{thm:lattice-lorentz} \\
reading layer & emerges & flagged: mechanism not constructed \\
rapidity plane & $\checkmark\checkmark$ identity-level & \cref{thm:bilateral} \\
\bottomrule
\end{tabular}
\end{center}
The boost is a tenant of the plane, a guest of the reading, and a stranger to the frame sequence --- and each residence is certified separately, by its own argument.
\end{remark}

\begin{remark}[The theorem witnesses \cref{as:a1}, without a circle]\label{rem:witness}
\cref{as:a1} posits that all reference frames register against one common sequence. If an exact boost acted on the frame sequence, it would mix the time axis into the space axes frame by frame, and the common sequence would lose its meaning at the level of the sequence itself. \Cref{thm:lattice-lorentz} proves that no such symmetry exists: the common sequence is \emph{compatible} with the geometry of the lattice. The direction of this witness is fixed once and is not reversible: the theorem treats the sequence \emph{defined by} \cref{as:a1} as a number-theoretic object; it does not derive \cref{as:a1} (which remains a hypothesis), and it does not use \cref{as:a2} at any step (the quantization rule never entered \cref{eq:preserve} or the squeeze). What the theorem certifies is consistency, and consistency only: the hypothesis and the geometry fit.
\end{remark}

\begin{remark}[The physical reading: the absolute inside the relative]\label{rem:absolute-ruler}
The reading belongs to the interpretive side (\cref{rem:register}) and is stated once, exactly. Relativity deprived the time axis of absoluteness; \cref{thm:lattice-lorentz} restores a residue of it on the frame sequence: the common sequence is not an additional postulate but a consequence locked by the integers --- no exact transformation of reference frames can stir the time axis into the space axes on $\mathbb{Z}^{2}$, because the stirring transformations do not exist there. The standard ruler of \cref{rem:reading} is absolute for a stated reason: the transformations that would disturb it are absent from the stage it rules.
\end{remark}

\begin{remark}[Not a new hypothesis]\label{rem:not-axiom}
\Cref{thm:lattice-lorentz} introduces no new hypothesis: its entire physical input is \cref{as:a1} (the frame sequence) and \cref{def:units} (the commensurability), and its mathematical content is the single squeeze \cref{eq:squeeze}. Nothing is assumed in its statement that the earlier definitions did not already contain; the theorem computes, it does not legislate.
\end{remark}

\section{Spatial rotation and dimension}\label{sec:rotation}

\subsection{Why rotation is tied to dimension}\label{sec:rigidity}

The angular momentum of a system is presented, in three-dimensional physics, as a vector $\mathbf L = \mathbf r \times \mathbf p$ --- a single arrow, with a magnitude and a direction, from which conservation statements and quantisation rules are built. That presentation is not automatic: it requires the cross product $\mathbf r \times \mathbf p$ to exist as a vector at all, and the existence of a cross product is a condition on the dimension of space.

\begin{lemma}[Dimension of the exterior square]\label{lem:exterior-dim}
For a real vector space $V$ of dimension $n$, the second exterior power $\Lambda^{2}V$ has dimension
\begin{equation}\label{eq:exterior-dim}
\dim \Lambda^{2}V = \frac{n(n-1)}{2},
\end{equation}
with basis $e_{i} \wedge e_{j}$ over pairs $i < j$: choosing $i$ in $n$ ways and then $j > i$ in $n-1$ ways counts each pair once, giving $n(n-1)/2$.
\end{lemma}

\begin{theorem}[Dimensional rigidity]\label{thm:rigidity}
An antisymmetric bilinear map $V \times V \to V$ (a cross product) can be an isomorphism $V \cong \Lambda^{2}V$ --- so that the angular momentum of a rotation acts on $V$ itself --- \textbf{if and only if}
\begin{equation}\label{eq:rigidity}
\frac{n(n-1)}{2} = n \;\Longleftrightarrow\; n(n-3) = 0 \;\Longleftrightarrow\; n = 3,
\end{equation}
the root $n = 0$ being the trivial space and excluded. Spatial rotation, as a symmetry whose charge is a vector of the space itself, exists in three dimensions and nowhere else.
\uses{\cref{lem:exterior-dim}}
\end{theorem}

\begin{proof}
One computation chain, each link written: the angular momentum $\mathbf L = \mathbf r \times \mathbf p$ is an antisymmetric function of two vectors of $V$, hence a linear map on $\Lambda^{2}V$; for $\mathbf L$ to be a vector of $V$ (a ``pure rotation charge''), $\Lambda^{2}V$ must be isomorphic to $V$, which requires equality of dimensions, \cref{lem:exterior-dim}; $\frac{n(n-1)}{2} = n \Leftrightarrow n^{2} - n = 2n \Leftrightarrow n^{2} - 3n = 0 \Leftrightarrow n(n-3) = 0$, whose non-negative roots are $n = 0$ and $n = 3$. For $n \neq 3$ (and $n \geq 2$), the angular momentum is a two-form, not a vector: rotation has no arrow-shaped charge.
\end{proof}

\begin{remark}[What is rigid, and what is topological]\label{rem:topological}
Two registrations, one per direction of the dependence. (i) \emph{Effective description}: the vector form of angular momentum is the unique shape in which a rotation charge can be an arrow --- \cref{thm:rigidity} --- and in every other dimension the same physics, if present at all, would present its rotation charge as a matrix of components. (ii) \emph{Structural}: the dimension itself is a topological invariant --- preserved by homeomorphisms, as registered by homology theory (quoted externally, used nowhere below)\footnote{The ranks of the homology groups $H_{k}$ are homotopy invariants, and $H_{n} \neq 0$, $H_{k>n} = 0$ fixes the dimension topologically; quoted as the standard fact, no proof reproduced.} --- so $n = 3$ cannot be deformed into or out of: no continuous process changes the dimension, and \cref{thm:rigidity} cannot be evaded by reparametrisation. The experimental side is the operational one: an experiment is rotated by an element of the rotation group, its measurements are checked for invariance, and the conserved quantity thereby certified is angular momentum (Noether, quoted externally). None of these external registrations is used as a premise for \cref{thm:rigidity}, which stands on \cref{lem:exterior-dim} alone.
\end{remark}

\subsection{The rotation groups, counted}\label{sec:so-count}

\begin{lemma}[Dimension of $\mathrm{SO}(n)$]\label{lem:so-dim}
\begin{equation}\label{eq:so-dim}
\dim \mathrm{SO}(n) = n^{2} - \frac{n(n+1)}{2} = \frac{n(n-1)}{2},
\end{equation}
\uses{pure counting of constraints}
\end{lemma}

\begin{proof}
An $n \times n$ matrix has $n^{2}$ entries. The orthogonality constraint $R^{\mathrm T}R = \mathbb{1}$ is symmetric --- $(R^{\mathrm T}R)^{\mathrm T} = R^{\mathrm T}R$ --- so it imposes as many independent conditions as a symmetric matrix has entries: $n$ diagonal and $\binom{n}{2} = n(n-1)/2$ off-diagonal, totalling $n(n+1)/2$. The constraint surface in the $n^{2}$-dimensional matrix space therefore has dimension $n^{2} - n(n+1)/2 = n(n-1)/2$. (The unit-determinant condition is discrete and does not change the dimension.)
\end{proof}

The enumeration, dimension by dimension, is the content of the lemma, and each row is worth one line:
\begin{equation}\label{eq:so-table}
\begin{array}{c|l|l}
n & \dim \mathrm{SO}(n) & \text{structure}\\
\hline
1 & 0 & \mathrm{SO}(1) = \{\mathbb{1}\}: \text{no rotation}\\
2 & 1 & \mathrm{SO}(2) \cong U(1),\ \text{abelian: } [J, J] = 0\\
3 & 3 & [J_{i}, J_{j}] = i\epsilon_{ijk}J_{k},\ \text{non-abelian}\\
4 & 6 & \mathrm{SO}(4) \cong \bigl(\mathrm{SU}(2) \times \mathrm{SU}(2)\bigr)/\mathbb{Z}_{2}
\end{array}
\end{equation}
the isomorphism in the last row holding only modulo the central identification of $\pm$ (the unquotiented product double-covers $\mathrm{SO}(4)$).

\begin{proposition}[No pure rotation in $1{+}1$ dimensions]\label{prop:no-pure-rotation}
The propagation seat --- a spacetime with coordinates $(t, x)$ only --- admits no pure spatial rotation: rotating the single spatial coordinate alone is the action of $\mathrm{SO}(1) = \{\mathbb{1}\}$, the trivial group (\cref{eq:so-table}, first row); rotating $(t, x)$ together mixes time with space and is, by the definition of the boost (\cref{eq:boostmatrix}), a boost rather than a rotation. Pure spatial rotation therefore has no seat in the propagation layer.
\uses{\cref{eq:so-table}, \cref{eq:boostmatrix}}
\end{proposition}

\begin{remark}[Reading]\label{rem:rotation-reading}
Conservative reading of the two results. \Cref{thm:rigidity}: the framework's spatial world is three-dimensional not by decree but by the arithmetic of the cross product --- the vector shape of angular momentum is the fingerprint of $n = 3$, and the framework records the fingerprint. \Cref{prop:no-pure-rotation}: what the propagation layer lacks is not rotation in general but the \emph{purely spatial} kind; the boost remains available as the mixed transformation, and the full rotation group survives where \cref{eq:so-table} says it can --- on the boundary sphere (\cref{ch:boundary}), whose symmetry group is $\mathrm{SO}(3)$ and whose angular momentum lives in three dimensions precisely because the boundary does. The grid of \cref{sec:sixteen} will place each of these statements in its cell.
\end{remark}

\section{Spatial translation}\label{sec:translation}

\begin{lemma}[Integer translations are symmetries]\label{lem:integer-translation}
The shifts
\begin{equation}\label{eq:shift}
x \longmapsto x + m\,\ell_{\mathrm P}, \qquad m \in \mathbb{Z},
\end{equation}
are bijections of the spatial frame sequence carrying each site to a site: they are the automorphisms of the one-dimensional lattice, and the sequence is translation-invariant under exactly them --- no smaller or larger step preserves the step structure.
\uses{\cref{as:a1,as:stepping} (the frame sequence)}
\end{lemma}

\begin{proof}
A bijection of $\mathbb{Z}$ (in units of $\ell_{\mathrm P}$) that maps the step structure to itself must send $j$ to $j + m$ for a single fixed $m$: bijectivity forces the difference between consecutive images to be constantly $1$ or constantly $-1$, and preservation of the orientation chosen by the sequence fixes the sign. Both steps of the argument are one line each.
\end{proof}

\begin{eresult}[Fourier duality on the lattice]\label{er:fourier}
The conserved quantity dual to a translation symmetry on a periodic lattice is the quasimomentum $Q$, with plane waves $e^{iQj\ell_{\mathrm P}}$; periodicity in $Q$ with period $2\pi/\ell_{\mathrm P}$ restricts $Q$ to the Brillouin zone $Q\ell_{\mathrm P} \in [-\pi, \pi)$.
\reference{Standard lattice kinematics (Bloch, 1928); quoted as the external dual of \cref{lem:integer-translation}.}
\end{eresult}

The relation between the quasimomentum and the naive momentum lattice is fixed by three clauses, each worth stating exactly:
\begin{equation}\label{eq:correction}
\begin{gathered}
p \in \tfrac{\pi\hbar}{\ell_{\mathrm P}}\mathbb{Z}
\quad\text{(single-cell periodic convention)}
\qquad\text{vs.}\\
Q \in [-\tfrac{\pi}{\ell_{\mathrm P}}, \tfrac{\pi}{\ell_{\mathrm P}})\ \text{continuous on the infinite lattice}
\quad\text{(\cref{er:fourier})}.
\end{gathered}
\end{equation}
(i) The formula $\tfrac{\pi\hbar}{\ell_{\mathrm P}}\mathbb{Z}$ belongs to the \emph{single-cell periodic} convention --- one site per period --- under which the allowed momenta are the reciprocal-lattice points of that convention. (ii) On the infinite lattice, the quasimomentum of \cref{er:fourier} is continuous throughout the zone; no discrete spectrum exists without a finite volume. (iii) Finite volume is itself a further convention --- box normalisation --- and the framework registers it as a convention wherever it is used, never as a derived fact. The three clauses together are the exact status of momentum on the frame sequence: a convention where a lattice appears, continuity on the infinite lattice, and a further convention wherever discreteness is needed.

\begin{lemma}[Centres and broken translation]\label{lem:center-breaking}
Wherever the geometry selects a centre, translation is broken: a centre is a distinguished point, a translation moves it, and the translated configuration differs from the original one. Residence by residence, from the structure built in \cref{ch:boundary,ch:newton}: the boundary sphere has a centre --- $\times$; a mass configuration has its centre of mass --- $\times$; the free entity of the rapidity plane (\cref{sec:rapidity-plane}) carries no intrinsic centre --- $\checkmark$.
\uses{\cref{eq:shift}; geometric enumeration}
\end{lemma}

\begin{remark}[Reading]\label{rem:translation-reading}
Conservative reading: translation is the symmetry the frame sequence most freely admits (\cref{lem:integer-translation}) and the one its centred objects most freely break (\cref{lem:center-breaking}); the conserved quasimomentum lives on the free entity alone, and even there it answers to the zone of \cref{er:fourier} rather than to any momentum lattice. The two lemmas between them fix the symmetry's exact extent, no more and no less: a lattice-full of translations, a zone-full of conserved quantity, and no claim beyond.
\end{remark}

\section{Lattice theorems of the plane}\label{sec:lattice-theorems}

The plane's mixed texture --- discrete grid on the time axis (\cref{eq:plane-grid}), continuous mass shells (\cref{eq:shells}) --- is now made precise. The theorems of this section are classification statements: which motions can be quantised, which masses can, and what their arithmetic looks like. Conventions are registered where they enter, and two of the five results carry the convention label throughout.

\begin{remark}[The two grids and the shell]\label{rem:two-grids}
Two discrete structures meet the continuous shells. The \emph{quantization grid} on the time axis is framework-fixed: $X = \tfrac{k}{2}\hbar$, $k \in \mathbb{Z}^{+}$ (\cref{eq:plane-grid}, from \cref{as:a2} with \cref{def:alpha0}). The \emph{momentum grid} on the space axis is convention-fixed: $Y = j\,\pi\hbar$, $j \in \mathbb{Z}$ (single-cell periodic convention, \cref{eq:correction}); it enters \cref{rem:t5} and the corollary alone, and every conclusion that uses it carries the convention flag. The shell is \cref{eq:shells}: $X^{2} - Y^{2} = (\hbar m/M_{\mathrm P})^{2}$, i.e.\ in grid units
\begin{equation}\label{eq:shell-units}
\Bigl(\frac{k}{2}\Bigr)^{2} - (j\pi)^{2} = \Bigl(\frac{m}{M_{\mathrm P}}\Bigr)^{2}.
\end{equation}
\end{remark}

\begin{theorem}[T1: velocity quantisation]\label{thm:t1}
For a mass $m$, the motions whose energy quantises (\cref{eq:quantization}) are exactly the discrete set
\begin{equation}\label{eq:t1}
\cosh\theta_{k} = \frac{k\,M_{\mathrm P}}{2m}, \qquad k \in \mathbb{Z}^{+},\ k \geq \frac{2m}{M_{\mathrm P}},
\end{equation}
one rapidity per admissible $k$. Check at the boundary of the condition: $m = 2M_{\mathrm P}$, $k = 4$ gives $\cosh\theta = 4M_{\mathrm P}/4M_{\mathrm P} = 1$, the rest state, as the condition $k \geq 2m/M_{\mathrm P} = 4$ saturates.
\uses{\cref{eq:quantization}, \cref{thm:bilateral}}
\end{theorem}

\begin{proof}
On a mass shell, $E = mc^{2}\cosh\theta$ (\cref{eq:boost-quotes} at $v = 0$ composed with the boost, quoted externally as before). Quantisation requires $E = \tfrac{k}{2}\varepsilon_{\mathrm P} = \tfrac{k}{2}M_{\mathrm P}c^{2}$ (\cref{eq:quantization}, \cref{cor:epmp}); equating, $mc^{2}\cosh\theta = \tfrac{k}{2}M_{\mathrm P}c^{2}$, and cancelling $c^{2}$ gives \cref{eq:t1}. The condition $\cosh\theta \geq 1$ is $kM_{\mathrm P}/2m \geq 1 \Leftrightarrow k \geq 2m/M_{\mathrm P}$; each admissible $k$ fixes exactly one $\theta_{k} = \operatorname{arccosh}(kM_{\mathrm P}/2m) \geq 0$, and distinct $k$ give distinct $\theta$ since $\operatorname{arccosh}$ is strictly increasing.
\end{proof}

\begin{theorem}[T2: the quantisable velocities have measure zero]\label{thm:t2}
For fixed $m$, the set $\{\theta_{k}\}$ of \cref{thm:t1} is countable --- a subset of $\mathbb{Q}$-indexed points in the continuum of rapidities. A generic velocity is never quantisable; the quantization rule sees almost no motion.
\uses{\cref{thm:t1}}
\end{theorem}

\begin{proof}
Countability is inherited from $\mathbb{Z}^{+}$: $k \mapsto \theta_{k}$ is an injection of a countable set. Measure zero follows from countability in the usual way (a countable union of points has Lebesgue measure zero), or more concretely: the points are discrete, separated by $\theta_{k+1} - \theta_{k} > 0$ for each $k$.
\end{proof}

\begin{theorem}[T3: the rest section]\label{thm:t3}
At $j = 0$ (zero momentum), quantisation is possible if and only if the mass is a half-integer multiple of the Planck mass:
\begin{equation}\label{eq:t3}
\theta = 0 \text{ quantisable} \quad\Longleftrightarrow\quad \Bigl(\frac{k}{2}\Bigr)^{2} = \Bigl(\frac{m}{M_{\mathrm P}}\Bigr)^{2} \quad\Longleftrightarrow\quad m = \frac{k}{2}M_{\mathrm P},
\end{equation}
from \cref{eq:shell-units} at $j = 0$: the shell crosses the time axis exactly at the half-integer masses. The masses of free entities are, in general, off this lattice --- cf.\ \cref{rem:rays-shells}.
\uses{\cref{eq:shell-units}}
\end{theorem}

\begin{theorem}[T4: the boundary case $m = M_{\mathrm P}/2$]\label{thm:t4}
For $m = M_{\mathrm P}/2$, the condition of \cref{thm:t1} reads $k \geq 1$ with $\cosh\theta_{k} = k$: \emph{every} positive integer $k$ quantises, at the rapidity $\theta_{k} = \operatorname{arccosh} k$. The smallest admissible mass is the most densely quantisable object of the plane --- the case in which the quantization grid and the mass shell are most nearly commensurate.
\uses{\cref{thm:t1}}
\end{theorem}

\begin{remark}[T5: the two grids are incommensurable --- convention-flagged]\label{rem:t5}
Joint quantization of energy and momentum asks for integer solutions of \cref{eq:shell-units}. The grid spacing ratio is
\begin{equation}\label{eq:incommensurate}
\frac{\text{momentum grid spacing}}{\text{quantization grid spacing}} = \frac{\pi\hbar}{\hbar/2} = 2\pi,
\end{equation}
irrational, so the two lattices share no common period: a generic mass has no moving state quantised on both grids simultaneously. Illustration at the boundary mass: for $m = M_{\mathrm P}/2$, \cref{eq:shell-units} reads $\bigl(\tfrac{k}{2}\bigr)^{2} - (j\pi)^{2} = \tfrac14$, whose only integer solution is $(k, j) = (1, 0)$ --- the rest state, verified directly: $\tfrac14 - 0 = \tfrac14$. \emph{Convention status}: the momentum grid is the single-cell convention of \cref{eq:correction}; under a different convention the $j$-lattice moves, and with it every conclusion that mentions $j$. This remark is a convention remark, registered as such, not a standing theorem.
\end{remark}

\begin{corollary}[Degeneracy of a fixed quantum number --- convention-flagged]\label{cor:degeneracy}
For fixed $k$, the condition $(j\pi)^{2} \leq \bigl(\tfrac{k}{2}\bigr)^{2} - \bigl(\tfrac{m}{M_{\mathrm P}}\bigr)^{2}$ allows
\begin{equation}\label{eq:degeneracy}
j = 0, 1, \ldots, \Bigl\lfloor \frac{\sqrt{(k/2)^{2} - (m/M_{\mathrm P})^{2}}}{\pi} \Bigr\rfloor,
\end{equation}
so a single quantum number $k$ carries up to $\bigl\lfloor\sqrt{(k/2)^{2} - (m/M_{\mathrm P})^{2}}\bigr\rfloor/\pi + 1$ momentum assignments --- a multiplicity of realisations per quantum, convention-dependent as \cref{rem:t5}.
\uses{\cref{eq:shell-units}}
\end{corollary}

\subsection{The dual identities}\label{sec:dual-identities}

\begin{proposition}[Rapidity duality]\label{prop:rapidity-dual}
With the quantum number $k = 2E/\varepsilon_{\mathrm P}$ (\cref{eq:quantization}) and its rest value $k_{0} = 2mc^{2}/\varepsilon_{\mathrm P} = 2m/M_{\mathrm P}$ (\cref{cor:epmp}):
\begin{equation}\label{eq:dual1}
\frac{k}{k_{0}} = \frac{E}{mc^{2}} = \cosh\theta.
\end{equation}
Counting oscillations per frame in \emph{radians} (angular-frequency convention throughout, as in \cref{as:p}): the rest phase advance per frame is $\omega_{\mathrm C}t_{\mathrm P} = \tfrac{mc^{2}}{\hbar}t_{\mathrm P} = \tfrac{m}{M_{\mathrm P}}$ rad (\cref{def:units}), and time dilation stretches the Compton interval to $\gamma\,t_{\mathrm C}$ with $\gamma = \cosh\theta$ and $t_{\mathrm C} = \hbar/mc^{2}$, so the moving count is
\begin{equation}\label{eq:dual2}
\frac{t_{\mathrm P}}{\gamma t_{\mathrm C}} = \frac{t_{\mathrm P} mc^{2}}{\hbar\cosh\theta} = \frac{m}{M_{\mathrm P}\cosh\theta} \quad\text{rad/frame},
\end{equation}
and the ratio of the two counts is $1/\cosh\theta$. The product is
\begin{equation}\label{eq:dual3}
k \times \frac{m}{M_{\mathrm P}\cosh\theta} = \frac{2m}{M_{\mathrm P}}\cosh\theta \cdot \frac{m}{M_{\mathrm P}\cosh\theta} = 2\Bigl(\frac{m}{M_{\mathrm P}}\Bigr)^{2},
\end{equation}
convention-independent --- the frame's quantum number and its oscillation count trade through $\cosh\theta$ in opposite directions, leaving the product fixed. Numerical check at $\cosh\theta = 2$: $k/k_{0} = 2$, oscillation ratio $= \tfrac12$, product $= 2(m/M_{\mathrm P})^{2}$, as \cref{eq:dual1,eq:dual2,eq:dual3} give.
\uses{\cref{eq:quantization}, \cref{eq:boost-quotes} (external boost), \cref{def:units}, angular convention. Not used: \cref{as:a2} beyond the definition of $k$}
\end{proposition}

\begin{theorem}[The energy--frequency product, and its uniqueness]\label{thm:ef}
\begin{equation}\label{eq:ef}
E \cdot f = mc^{2}\cosh\theta \cdot \frac{\omega_{\mathrm C}}{\cosh\theta} = mc^{2}\omega_{\mathrm C} = E_{0}f_{0},
\end{equation}
with $f \equiv \omega$ in the angular convention (the product uses the same frequency as \cref{eq:dual2}). The right side is
\begin{equation}\label{eq:ef-value}
E_{0}f_{0} = mc^{2}\cdot\frac{mc^{2}}{\hbar} = \frac{m^{2}c^{4}}{\hbar} = 2\Bigl(\frac{m}{M_{\mathrm P}}\Bigr)^{2}\mathcal P,
\end{equation}
using $\mathcal P = c^{5}/2G$ (\cref{def:power}) and $M_{\mathrm P}^{2} = \hbar c/G$ (\cref{def:units}). \emph{Uniqueness} (the content that forbids over-reading): among all combinations $m^{a}c^{b}\hbar^{d}$, the units of power force
\begin{equation}\label{eq:ef-uniqueness}
a + d = 1\ (\mathrm{kg}), \qquad b + 2d = 2\ (\mathrm{m}), \qquad b + d = 3\ (\mathrm{s}),
\end{equation}
a linear system whose unique solution is $(a, b, d) = (2, 4, -1)$: subtracting the last two equations gives $d = -1$, then $b = 4$, $a = 2$. There is exactly one power built from $\{m, c, \hbar\}$, and \cref{eq:ef} is it; the identity carries no content beyond its own arithmetic.
\uses{\cref{eq:boost-quotes} (external), \cref{def:power,def:units}}
\end{theorem}

\begin{proposition}[The quantization rule holds at every velocity]\label{prop:quantization-loop}
\begin{equation}\label{eq:quantization-loop}
E\,t_{\mathrm P} = mc^{2}\cosh\theta \cdot t_{\mathrm P} = \frac{2m}{M_{\mathrm P}}\cosh\theta\cdot\frac{\hbar}{2} = \frac{k}{2}\hbar,
\end{equation}
three equalities, each one substitution: the energy on shell, $mc^{2}t_{\mathrm P} = 2m\hbar/2M_{\mathrm P}$ (\cref{def:units} rearranged), and \cref{eq:dual1}. The exchange rate $\hbar$ between action and quantum number is the same in every frame --- the quantization rule is covariant as it stands, and no separate ``moving-observer'' rule is needed. \emph{Endpoint excluded}: the derivation presumes $m \neq 0$ (the Compton interval diverges at $m = 0$); the photon endpoint $Y = X$, $I = 0$ (\cref{rem:rays-shells}) lies outside the formulas of this section.
\uses{\cref{eq:boost-quotes} (external), \cref{def:units}, \cref{prop:rapidity-dual}}
\end{proposition}

\begin{remark}[The two-frame comparison, with a demotion on record]\label{rem:two-frames}
Side by side, laboratory and rest frame (all entries from \cref{eq:boost-quotes} and \cref{def:units}):
\begin{center}
\begin{tabular}{@{}llll@{}}
\toprule
Quantity & Laboratory & Rest frame & Per frame \\
\midrule
time & $t = nt_{\mathrm P}$ & $\tau = t/\cosh\theta$ & $t_{\mathrm P}$ \\
energy & $E = mc^{2}\cosh\theta$ & $E_{0} = mc^{2}$ & quantum number $k = 2E/\varepsilon_{\mathrm P}$ \\
momentum & $p = mc\sinh\theta$ & $0$ & --- \\
frequency & $\omega = \omega_{\mathrm C}/\cosh\theta$ & $\omega_{\mathrm C} = mc^{2}/\hbar$ & oscillations $m/(M_{\mathrm P}\cosh\theta)$ \\
invariant & \multicolumn{2}{l}{$I = \hbar m/M_{\mathrm P}$} & --- \\
product & \multicolumn{2}{l}{$E\omega = E_{0}\omega_{\mathrm C}$ (\cref{thm:ef})} & --- \\
\bottomrule
\end{tabular}
\end{center}
Every entry of the first column is driven by $\cosh\theta$ off a constant second column. One demotion, recorded: the relation $t/E = \tau/E_{0}$ is \emph{not} a new invariant --- it is the quotient of the two known statements $t = \tau\cosh\theta$ and $E = E_{0}\cosh\theta$, two invariant relations divided; it carries the content of its factors and nothing beyond, and the table lists it as such.
\end{remark}

\section{The census of residences}\label{sec:census}

\subsection{The classification underneath}\label{sec:three-way}

\begin{theorem}[Three constant-curvature geometries]\label{thm:three-way}
A complete, simply connected, two-dimensional geometry of constant curvature $K$ is, up to scale, exactly one of three: the sphere ($K > 0$), the Euclidean plane ($K = 0$), or the hyperbolic plane ($K < 0$). The sign of $K$ alone decides.
\uses{the Jacobi equation, quoted hypotheses (complete, simply connected, constant $K$)}
\end{theorem}

\begin{proof}
In geodesic polar coordinates a rotationally symmetric metric reads $ds^{2} = dr^{2} + f(r)^{2}d\theta^{2}$, and constant curvature $K$ is equivalent to the Jacobi equation
\begin{equation}\label{eq:jacobi}
f''(r) + K f(r) = 0, \qquad f(0) = 0,\quad f'(0) = 1,
\end{equation}
the initial conditions fixing the normalization of the radial coordinate (unit speed at the origin). The equation is solved in closed form in one line per sign:
\begin{equation}\label{eq:jacobi-sol}
K > 0:\ f = \frac{\sin(\sqrt{K}\,r)}{\sqrt{K}};\qquad
K = 0:\ f = r;\qquad
K < 0:\ f = \frac{\sinh(\sqrt{|K|}\,r)}{\sqrt{|K|}},
\end{equation}
each satisfying \cref{eq:jacobi} by one differentiation. The sign of $K$ therefore determines $f$ uniquely, and with $f$ the metric: three classes, and the hypotheses (complete, simply connected) exclude quotients, so no further cases exist.\footnote{The statement is the classical classification of constant-curvature surfaces, associated with Killing (1891) and Hopf (1925); the proof above is self-contained and the quoted hypotheses are only the completeness and simple connectivity.}
\end{proof}

\subsection{The double-axis census}\label{sec:double-axis}

Two axes classify geometry independently: the \emph{signature} axis (Lorentzian or Euclidean) and the \emph{curvature} axis ($K$'s sign). The axes are independent and must not be conflated: signature classifies the metric's indefinite character, curvature its Gaussian curvature, and a seat's proper cell is fixed by both. The census of the framework's geometric seats:
\begin{center}
\begin{tabular}{@{}p{1.8cm}|p{3.8cm}p{3.5cm}p{1.8cm}@{}}
\toprule
Signature $\backslash$ curvature & flat $K=0$ & sphere $K>0$ & hyperbolic $K<0$ \\
\midrule
Lorentzian & propagation seat (\cref{ch:lattice}); rapidity plane (\cref{sec:rapidity-plane}) & --- no content: curved spacetime is not in the framework & --- \\
Euclidean & static tangent planes of the boundary & boundary $S^{2}$ (\cref{ch:boundary}) & --- \\
\bottomrule
\end{tabular}
\end{center}

Two clarifications accompany the table, because both are the kind of point that is easily misread.

\emph{First, signature is not curvature.} The Minkowski form $X^{2} - Y^{2}$ of the rapidity plane concerns its \emph{signature}, and the plane itself is flat: its curvature is $K = 0$, the same geometric type as the propagation seat's local slices. Placing the plane in the hyperbolic-curvature cell would conflate the plane's \emph{rotation group}, which is hyperbolic ($\mathrm{SO}(1,1)$, the boost), with its \emph{space}, which is not: the plane is flat, of the same geometric type as the propagation seat's local slices, and its seat in the census is a functional one in the Lorentzian--flat cell: same geometric type as the propagation slices, different function (kinematics of $(E, p)$ rather than of $(t, x)$).

\emph{Second, what is hyperbolic here is one-dimensional.} The framework's genuinely hyperbolic objects are the mass shells of \cref{eq:shells} --- the curves $X^{2} - Y^{2} = (\hbar m/M_{\mathrm P})^{2}$, one hyperbola per mass, embedded in the flat plane. There is no two-dimensional hyperbolic geometry anywhere in the framework: the cell is empty. The third row of \cref{eq:jacobi-sol} describes no framework object.

\begin{remark}[The census knows nothing of forces]\label{rem:no-forces-here}
The table classifies geometric \emph{seats}, and nothing in it assigns a seat to a force; three clarifications keep the table honest. First, the wave operator of the propagation seat carries no gauge structure, no charge, and no coupling constant, so the framework names no occupant for that seat, and the vacancy is registered (\cref{sec:correspondence}). Second, the boundary's relation to relativistic gravity is not an identification: it consists solely of the formal coincidences already registered --- one-way and bounded --- in \cref{ch:boundary,ch:gr}. Third, the strong and weak structures have no candidate seat, and the census registers the absence as information rather than as failure: the absence of candidates is data about the framework, not a defect in it. The full statement is \cref{sec:correspondence}; nothing in this section presupposes it, and nothing in this section needs it.
\end{remark}

\section{The sixteen-fold grid}\label{sec:sixteen}

Four residences, four symmetries, sixteen cells. Each cell below carries four items: the existence verdict, the proof or the recorded reason for absence, the condition under which the verdict holds, and a conservative reading. The residences are the four geometric seats of \cref{sec:census}: the propagation seat, the boundary sphere, the repair registry's formal row, and the rapidity plane. Symbols: $T$ time translation, $P$ spatial translation, $R$ spatial rotation, $B$ boost.

\subsection{Residence one: the propagation seat}\label{sec:grid-propagation}

\emph{Cell (1,1), $T$ --- exists at frame level.} \textbf{Proposition.} The integer frame shifts $n \mapsto n+k$ are symmetries of the wave operator's seat; the continuous time translation is a reading-layer statement. \emph{Proof.} The propagation operator is the same at every frame by construction --- its coefficients do not carry $n$ --- so $n \mapsto n+k$ is a legal relabelling that changes no equation; continuous translation requires the reading limit $n \gg 1$, where it holds as the reading of the same relabelling. \emph{Condition:} lattice uniformity (\cref{as:stepping}). \emph{Reading:} the seat is a machine that runs identically at every frame; moving the time origin changes nothing in its operation.

\emph{Cell (1,2), $P$ --- exists (discrete).} \textbf{Proposition.} Integer site shifts are symmetries (\cref{lem:integer-translation}); the conserved quantity is the quasimomentum $Q$ of \cref{er:fourier}, living on the Brillouin zone. \emph{Proof.} \cref{lem:integer-translation} and \cref{er:fourier}. \emph{Condition:} the infinite/finite-volume convention (\cref{eq:correction}). \emph{Reading:} on the lattice, motion is one site at a time, and the conserved momentum is accordingly a zone quantity, not a real line.

\emph{Cell (1,3), $R$ --- exists (discrete version).} \textbf{Proposition.} The lattice preserves a discrete rotation group; continuous $\mathrm{SO}(3)$ is reading-layer. \emph{Proof.} A rotation preserving $\mathbb{Z}^{3}$ must have integer entries and determinant $1$; the integer orthogonal matrices are the signed permutation matrices, $3! \times 2^{3} = 48$ in number, of which the determinant-$1$ half --- $48/2 = 24$ --- forms the rotation group of the cube: $24$ elements, counted by the enumeration itself. The continuous symmetry of the wave equation is $\mathrm{SO}(3)$ (the equation is isotropic, quoted externally); on the lattice only this discrete subgroup survives exactly. \emph{Condition:} cubic lattice structure (the isotropic lattice convention). \emph{Reading:} on graph paper only right-angle turns are exact; continuous rotation is what the reading sees when the mesh is too fine to resolve.

\emph{Cell (1,4), $B$ --- $\times$ on the frame sequence, $\checkmark\checkmark$ on the plane.} \textbf{Proposition.} No exact boost on the counting layer (\cref{thm:lattice-lorentz}); identity-level boost on the rapidity plane (\cref{thm:bilateral}); the reading layer's continuous boost is emergent, registered as unbuilt. \emph{Proof.} The three citations; the middle layer's mechanism is not constructed and is registered as such (\cref{rem:not-prohibition}). \emph{Condition:} countability of the frame sequence. \emph{Reading:} the boost is a tenant of the plane and a stranger to the frame sequence --- the chapter's first and sharpest three-layer verdict.

\subsection{Residence two: the boundary sphere}\label{sec:grid-sphere}

\emph{Cell (2,1), $T$ --- exists, trivially (the frozen cell).} \textbf{Proposition.} Time translation acts on the boundary's intrinsic geometry as the identity; on the family of boundary spheres it acts equivariantly, by re-indexing. \emph{Proof.} The intrinsic geometry of $S^{2}$ --- its metric, its area functional, its thermodynamic readings --- contains no time variable: every $T_{k}$ leaves each such object unchanged, so the group acts through singlets. What \emph{does} vary with the frame is the size reading $L(n) = n\ell_{\mathrm P}$ (\cref{eq:boundaryradius}), and re-indexing $n \mapsto n+k$ carries one member of the family $\{S^{2}_{L(n)}\}$ to another, leaving the construction --- every member is a sphere --- invariant: an equivariant action on the family, a trivial action on each geometry. The cell's verdict: the symmetry exists, and it guards a static object. \emph{Condition:} the statics of the intrinsic geometry (the sizes are readings, not contents of any single frame's geometry). \emph{Reading:} time translation of a photograph leaves the photograph unchanged --- but the family of photographs, one per frame, is ordered by the very shift that does nothing to each.

\emph{Cell (2,2), $P$ --- $\times$.} \textbf{Proposition.} No translation on the sphere. \emph{Proof.} The sphere has a centre (\cref{lem:center-breaking}): a translation moves it, and the translated object differs from the original --- the curvature $1/L^{2}$ is anchored at the centre, so translation is not an isometry. \emph{Condition:} existence of the centre. \emph{Reading:} on a sphere no two points are indistinguishable --- pole and equator are always distinct, and there is nowhere to translate to.

\emph{Cell (2,3), $R$ --- $\checkmark\checkmark$ full symmetry.} \textbf{Proposition.} $\mathrm{SO}(3)$ acts fully on $S^{2}$; the conserved charge is angular momentum. \emph{Proof.} The sphere is the equidistant set from its centre (\cref{lem:boundary-sphere}); the Killing equation $\nabla_{(a}\xi_{b)} = 0$ on $S^{2}$ has a three-dimensional solution space --- one rotational Killing field per axis --- counted by the equation itself, and quoted externally as the standard fact. After quantization the spectrum is $\hbar\sqrt{l(l+1)}$, $l \in \mathbb{Z}^{+}$ (external, quantum mechanics). \emph{Condition:} none --- pure geometry. \emph{Reading:} the sphere is the most symmetric object in the framework: any axis, any angle, and it is itself.

\emph{Cell (2,4), $B$ --- $\times$.} \textbf{Proposition.} No boost on the sphere. \emph{Proof.} The boost mixes a time coordinate into a space coordinate (\cref{eq:boostmatrix}); the sphere is a purely spatial object with no time dimension --- there is nothing to mix (\cref{sec:no-dof} for the framework's standing on time). \emph{Condition:} absence of a time dimension. \emph{Reading:} two observers on a sphere cannot pass each other at constant velocity --- there is no constant-velocity line on a sphere at all.

\subsection{Residence three: the repair registry's formal row}\label{sec:grid-repair}

The third residence is the formal construction of \cref{sec:repair}: a static response $\varphi(r)$ in a flat background. The row is \emph{formal} --- the construction carries no physical application (\cref{sec:repair}) --- and the cells below inherit that status: their verdicts are verdicts about a mathematical object.

\emph{Cell (3,1), $T$ --- exists, trivially (frozen, as in cell (2,1)).} \textbf{Proposition.} As cell (2,1), with the static response in place of the sphere: the response contains no time variable, the group acts through singlets, and any family of such responses is re-indexed equivariantly. \emph{Proof.} The same equivariance argument as cell (2,1), with $\varphi(r)$ for the geometry. \emph{Condition:} statics. \emph{Reading:} another still photograph; the same corrected reason.

\emph{Cell (3,2), $P$ --- $\times$.} \textbf{Proposition.} No translation. \emph{Proof.} The response has a centre --- its source (\cref{lem:center-breaking}) --- and the fixed scale $\mu$ selects that centre's unit; a translation changes the anchor of $\mu|r - r_{0}|$. \emph{Condition:} the source. \emph{Reading:} the response is centred on its source; where the source is, symmetry ends.

\emph{Cell (3,3), $R$ --- $\checkmark$.} \textbf{Proposition.} Rotation symmetry about the centre; the charge is the winding number. \emph{Proof.} $\varphi = \varphi(r)$ carries no direction, so rotations leave it unchanged (the same argument as cell (2,3), one dimension of Killing fields per axis); angular momentum $\mathbf L = \mathbf r \times \mathbf p$ is conserved (Noether, quoted externally). On the lattice the surviving memory of rotation is topological: the winding number $W \in \mathbb{Z}$ of a closed path around the centre --- an integer that no continuous deformation of the path changes, counted by the fundamental group (quoted externally). \emph{Condition:} radiality. \emph{Reading:} the lattice forgets continuous rotation but remembers the count of turns; the winding number is the residue.

\emph{Cell (3,4), $B$ --- $\times$.} \textbf{Proposition.} No boost. \emph{Proof.} The response is instantaneous (the elliptic operator has no time parameter, \cref{sec:repair}) and not Lorentz-covariant: under a boost the form $e^{-\mu r}/r$ does not transform as a scalar, so the boosted configuration is not a symmetry-related copy of the original. \emph{Condition:} statics. \emph{Reading:} this object knows the circle it drew around its source and no frame beyond it.

\subsection{Residence four: the rapidity plane}\label{sec:grid-plane}

\emph{Cell (4,1), $T$ --- $\checkmark$ active.} \textbf{Proposition.} Time translation of the free entity is a symmetry with the energy as conserved quantity. \emph{Proof.} The free entity's generator is time-independent --- the free Hamiltonian of \cref{sec:quantum} carries no explicit frame dependence --- so the evolution over any number of frames is the same relabelled, and the energy reading is frame-independent by \cref{thm:bilateral}. \emph{Condition:} freedom from the injection (the content layer has no such cell). \emph{Reading:} the free entity is the framework's only clock that keeps its own time --- and pays no injection for the privilege.

\emph{Cell (4,2), $P$ --- $\checkmark$ active.} \textbf{Proposition.} Momentum conservation for the free entity. \emph{Proof.} The free generator commutes with $\hat p$ (quoted externally as the standard fact for free kinematics); on the plane, translation is motion along the $Y$-axis, the momentum being the coordinate itself. \emph{Condition:} freedom (no centre, \cref{lem:center-breaking}). \emph{Reading:} on the plane, translation has a geometric face: it is displacement along the space axis, and its conservation is the flatness of that axis.

\emph{Cell (4,3), $R$ --- N/A in $1{+}1$; $\checkmark$ in higher dimensions (little groups).} \textbf{Proposition.} In one space dimension there is no rotation (\cref{prop:no-pure-rotation}); in higher dimensions the rotation symmetry of a state is its stabiliser --- $\mathrm{SO}(3)$ (spin) for massive states, $\mathrm{ISO}(2)$ (helicity) for massless --- quoted externally (Wigner classification). \emph{Condition:} dimension. \emph{Reading:} on the plane only rotation about one's own axis remains --- and that is spin, not orbital motion.

\emph{Cell (4,4), $B$ --- $\checkmark\checkmark$ home.} \textbf{Proposition.} The boost is the plane's defining symmetry, with the Casimir $I = \hbar m/M_{\mathrm P}$ (\cref{thm:bilateral}, \cref{rem:casimir}). \emph{Proof.} \Cref{thm:bilateral}. \emph{Condition:} none --- identity level. \emph{Reading:} the boost is not a tenant here and not a guest: it is the rotation that defines the plane, and its invariant is the mass.

\subsection{The verdict of completeness}\label{sec:verdict}

\begin{remark}[Completeness of the grid]\label{rem:grid-complete}
Sixteen cells, sixteen verdicts: every $\checkmark$ carries its proof above, every $\times$ carries the recorded reason the symmetry fails there, every N/A carries the missing object, and the two frozen cells carry the equivariance argument. No cell is left unargued, and no argument is left unregistered. The completeness claimed is of this form: not that every cell has a symmetry --- the grid is mostly negative --- but that \emph{every negative cell has its reason on record}. A withdrawn row (the two-dimensional hyperbolic seat, \cref{sec:census}) stands outside the grid by the census's own decision, and the grid's fourfold--fourfold shape is exactly the shape of what the framework contains.
\end{remark}

\section{The repair registry, and odd symmetry}\label{sec:repair}

\subsection{The repair registry}\label{sec:repair-registry}

\begin{remark}[What the repair registry is]\label{rem:repair-status}
The repair registry is the single location where the framework's undetermined items are recorded with their exact status (the parameter discipline of \cref{rem:layer-scope} admits no others). One item stands. The mass $m$ is an \emph{anchor-family debt}: every role mass plays in this framework is derived --- the bilateral invariant (\cref{thm:bilateral}), the energy--frequency product (\cref{thm:ef}), the ladder spectrum (\cref{thm:spectrum}) are all roles, derived and proved --- but no role determines the \emph{value} of any mass, exactly as no role in this framework determines the value of $\hbar$, $c$, or $G$. The requirement ``the framework must generate the masses'' is accordingly a category-misplaced pseudo-requirement: anchor families never generate the values of their constants, and registering $m$ beside $\hbar$, $c$, $G$ is the honest treatment, not an omission. The registry's content: one entry, so stated.
\end{remark}

\begin{remark}[The formal construction, registered]\label{rem:formal-construction}
Two registrations stand as pure formalities. (i) \emph{Existence}: the causal-consistency test of the propagation layer applies to time-dependent propagation operators --- a dispersion law must hold at every wavenumber, or some mode outruns the step structure (\cref{sec:steps}) --- and the shielded operator $(-\nabla^{2} + \mu^{2})$ is elliptic: it contains no time derivative, propagates nothing, and offers no dispersion for the test to act on; the test is idle against it, and the operator's legitimacy is the absence of competition with the propagation layer, nothing more. (ii) \emph{Lattice form}: on the frame sequence the operator $(-\delta^{2} + \mu^{2}\ell_{\mathrm P}^{2})$ is well-defined ($\delta^{2}$ the nearest-neighbour second difference), with decay at infinity automatic, no short-distance singularity, and unique solutions for compact sources --- each clause the standard fact, quoted externally. Both registrations carry the same disclaimer, stated once and inherited by each: formal constructions, no physical application in this framework, kept on record so that the registry is complete.
\end{remark}

\begin{remark}[The residence verdict]\label{rem:residence-verdict}
The shielded operator does not reside in a space of constant curvature: the kernel $e^{-\mu r}/r$ is the Green's function of the \emph{flat}-space shielded Poisson operator $(-\nabla^{2}+\mu^{2})\varphi = \rho$ (quoted externally as the standard potential theory), not a solution of any constant-curvature metric's field equations. The correct residence is the one the operator itself states: a flat background carrying a shielded elliptic term. The decomposition of the entry: it belongs to the repair registry ($\checkmark$, this registry), its curvature residence fails ($\times$, this remark), and the label ``strong force'' sometimes attached to it is an anachronism, carrying no content.
\end{remark}

\subsection{Odd symmetry}\label{sec:odd-symmetry}

\begin{proposition}[Odd extension of the content]\label{prop:odd}
The linear growth of the boundary content extends oddly:
\begin{equation}\label{eq:odd}
E_{\partial}(-n) = -E_{\partial}(n),
\end{equation}
a direct substitution into \cref{eq:eboundary}: $E_{\partial}(-n) = (-n)\varepsilon_{\mathrm P}/2 = -n\varepsilon_{\mathrm P}/2$. The share coefficients and every dimensionless ratio are invariant under $n \mapsto -n$ --- ratios are blind to the overall sign (\cref{eq:shares}, \cref{rem:gamma-one}) --- while every quantity carrying absolute dimensions changes sign with the content. The normalization factor $\gamma_{n}$ of \cref{def:normalization} has domain $n \geq 2$ and is not an object of this proposition.
\uses{\cref{eq:eboundary}, \cref{eq:shares}}
\end{proposition}

\begin{remark}[The negative hemisphere: a registered conjecture, nothing more]\label{rem:negative-conjecture}
The extension \cref{eq:odd} is algebra --- it computes --- and nothing beyond algebra is claimed for it. If the complete structure were a closed global object, the oddness \cref{eq:odd} would have a geometric reading: a positive hemisphere, an equatorial boundary, and a negative hemisphere as the mirror image --- and the reading is recorded here \emph{as a conjecture and only as a conjecture}, for the reason the doctrine exists: the framework assigns no physical content to negative frames. There is no mechanism for them, no observable of them, and no interpretation of the negative content; the extension is consistent, and consistency without content is what a conjecture is. The conjecture is therefore registered in both directions at once --- existence and non-existence of the negative side are \emph{both} unasserted --- under the standing doctrine (\cref{rem:truncation-restated} for the template: open status, three possibilities, no claim). The grid of \cref{sec:sixteen} is untouched by the conjecture: no cell of the framework's four residences depends on it, and none is asked to.
\end{remark}

\section{The correspondence record, and the hypothesis layer's account}\label{sec:correspondence}

\subsection{The record}\label{sec:correspondence-register}

The correspondence between this framework's objects and the interactions of established physics is, after the discipline of \cref{rem:register} has done its work, a short statement of what the framework does \emph{not} claim.

\begin{remark}[No identifications are made]\label{rem:autopsy}
Three statements, all negative, complete the record. (i) \emph{The propagation seat is not identified with any named interaction.} The propagation layer ends at a wave operator --- unique among local symmetric operators, proved in \cref{ch:lattice} --- while the named interactions of established physics are characterised by gauge structure, charge, and coupling constants, none of which the framework contains; any identification would rest on nothing the framework has. The seat's occupant is therefore un-named, and the vacancy stands on record. (ii) \emph{The boundary is not identified with gravity.} What \cref{ch:boundary,ch:gr} registered is one-way and bounded --- thermodynamic identities coinciding term by term with the relativistic horizon's, a constraint coinciding with the Friedmann equation's flat member --- and none of it is an identification; nothing in the framework generates the relativistic field equations. (iii) \emph{The strong and weak structures have no candidate seat at all,} and the census records the absence as information: the framework's seat count is three residences and a plane, and what is absent is as much a registered fact as what is present.
\end{remark}

\begin{remark}[The Wigner clause, closed]\label{rem:wigner-closed}
One clause of \cref{rem:wigner} is discharged here: the correspondence between forces and geometric configurations is, in the established theories, strikingly effective, and the framework registers that effectiveness without either sharing or explaining it. Its own objects stand in no identified relation to any force; the effectiveness it borrows from the record of physics remains a fact about that record, not about this framework.
\end{remark}

\subsection{The hypothesis layer's symmetry record}\label{sec:hypothesis-record}

The four symmetries' final status, as the chapter has established it:
\begin{center}
\begin{tabular}{@{}p{2.2cm}p{3.5cm}p{2.7cm}p{2.7cm}@{}}
\toprule
Symmetry & Frame sequence & Reading layer & Plane \\
\midrule
Time translation & difference invariant survives (\cref{sec:time-translation}) & function level broken & --- \\
Time rotation (boost) & identity only (\cref{thm:lattice-lorentz}) & emergence, unbuilt & identity-level (\cref{thm:bilateral}) \\
Spatial rotation & discrete subgroup, 24 elements & $\mathrm{SO}(3)$ reading & N/A in $1{+}1$ (\cref{prop:no-pure-rotation}) \\
Spatial translation & lattice automorphisms & quasimomentum, zone & free entity active (\cref{sec:translation}) \\
\bottomrule
\end{tabular}
\end{center}
Each entry cites the section where it is proved; the table itself asserts nothing new.

\begin{remark}[Frequency status, recorded once]\label{rem:frequency-status}
For closed entities --- the isolated oscillator and free entity of \cref{sec:quantum} --- the phase advances per frame by
\begin{equation}\label{eq:phase-advance}
\Delta\varphi = \frac{E\,t_{\mathrm P}}{\hbar} = \frac{k}{2}\ \mathrm{rad},
\end{equation}
using \cref{eq:quantization}, and the frequency reading is $\omega \equiv \Delta\varphi/t_{\mathrm P} = \tfrac{k}{2}\times(\text{frame rate})$: energy binds frequency for closed entities, with the half-integer lattice of \cref{eq:plane-grid} as the unit. The content layer carries no such binding --- power is not a phase-carrying object, and no frequency is assigned to the injection --- and the boundary, whose thermodynamics is built from contents, is correspondingly frequency-free. The asymmetry is structural, not an omission: phase belongs to amplitudes (\cref{def:quantum-layer}), and the content layer has none.
\end{remark}

\begin{remark}[Geometric association, by reference]\label{rem:geo-association}
The question of which symmetry sits in which geometric configuration is discharged by reference rather than repetition: the census of \cref{sec:census} gives the residences, the grid of \cref{sec:sixteen} gives the sixteen verdicts, and the two are the complete answer. Nothing is added here, because nothing needs to be.
\end{remark}

\section{Backfill, and the chapter's account}\label{sec:backfill}

\subsection{The backfill list}\label{sec:backfill-list}

The earlier chapters made forward references into this one; each is now closed, and the closure is recorded where it happened:
\begin{center}
\begin{tabular}{@{}p{6.2cm}p{6.6cm}@{}}
\toprule
Forward reference, made at & Closed at \\
\midrule
``common sequence'' of \cref{as:a1} (motivation) & witness of compatibility, \cref{rem:witness} \\
frequency convention of \cref{as:p} & angular convention of \cref{prop:rapidity-dual} \\
correspondence clause of \cref{rem:wigner} & \cref{rem:wigner-closed} \\
residence of the deeper shares, deferred in \cref{rem:withdrawn-share} & the census of \cref{sec:census} \\
\bottomrule
\end{tabular}
\end{center}
Two further closures belong to the lattice chapter and are registered here only so the list is complete: the dispersion relation's independent rederivation (\cref{er:dispersion}) awaits \cref{ch:lattice}, and the zero-point route's registration (\cref{rem:meeting}) was closed in \cref{ch:unitarity} --- neither is this chapter's debt, and neither is silently left open.

\subsection{The chapter's account}\label{sec:ch9-account}

\begin{remark}[What Chapter 9 has established]\label{rem:ch9summary}
The census table and its discipline (\cref{sec:sym-register,rem:geometry-discipline}). The rapidity plane (\cref{sec:rapidity-plane}) and its bilateral invariant, identity-level, matching its coordinates into the boost group (\cref{thm:bilateral}). Time translation (\cref{sec:time-translation}): broken at the function level by the non-postponable injection, surviving as the difference invariant of \cref{prop:difference-survives}. Time rotation: the lattice--Lorentz theorem (\cref{thm:lattice-lorentz}) with its four fixed distinctions --- a frame-sequence fact, a consistency witness to \cref{as:a1} without circularity, the absolute locked inside the relative, and no new hypothesis --- and the three-layer verdict of \cref{rem:not-prohibition}. Spatial rotation: dimensional rigidity (\cref{thm:rigidity}), the rotation groups counted (\cref{lem:so-dim}), and the absence of pure rotation in $1{+}1$ (\cref{prop:no-pure-rotation}). Spatial translation: the lattice automorphisms (\cref{lem:integer-translation}), the quasimomentum convention record (\cref{eq:correction}), and centre-breaking (\cref{lem:center-breaking}). The plane's lattice theorems (\cref{sec:lattice-theorems}): velocity quantisation, measure zero, the rest section, the boundary case, the convention-flagged incommensurability, and the dual identities --- rapidity duality with its convention-independent product, the energy--frequency product with its dimensional uniqueness, the quantisation rule at every velocity, and the two-frame comparison with its demotion on record. The census (\cref{sec:census}): three constant-curvature geometries, the double-axis table with signature and curvature kept distinct and the hyperbolic content one-dimensional. The sixteen-fold grid (\cref{sec:sixteen}): sixteen verdicts, each with proof or recorded reason, the two frozen cells in the equivariance form, and the completeness verdict of \cref{rem:grid-complete}. The repair registry (\cref{sec:repair}): one entry, the mass debt, beside the anchor family where it belongs; two formal registrations; one residence verdict; and the odd extension with its two-way conjecture (\cref{prop:odd,rem:negative-conjecture}). The correspondence record (\cref{sec:correspondence}): three negative statements, the Wigner clause closed, and the hypothesis layer's symmetry record tabled.

The input account: no new hypothesis in the chapter; external results quoted in proofs --- the Lorentz boost, the mass shell, Fourier duality, the Wigner classification, the constant-curvature classification --- each cited where used, none generating framework content; registrations without use --- homology, Noether's theorem, quantum angular spectra --- kept in remarks, never in premises. The chapter's open items, each registered with its successor: the reading layer's continuous boost, whose emergence mechanism is unbuilt (\cref{rem:not-prohibition}); the negative hemisphere, conjecture in both directions (\cref{rem:negative-conjecture}); and the mass debt (\cref{rem:repair-status}). Nothing in this chapter is used by an earlier chapter except through the closures of \cref{sec:backfill-list}; the dependency order of the whole text survives the chapter intact.
\end{remark}

\part{The lattice}

\chapter{Quantum mechanics on the lattice}\label{ch:lattice}

The frame sequence carries one field: the amplitude $\varphi$ of the propagation layer. This chapter establishes its two exactly provable properties --- what propagates, and with what dispersion --- and closes, with the first section, the forward reference made at \cref{er:dispersion}.

\section{Propagation}\label{sec:propagation}

\subsection{Setting}\label{sec:exit-line}

The propagation layer lives on the frame sequence: frame index $n \in \mathbb{Z}^{+}$ (\cref{as:a1}), spatial sites $j \in \mathbb{Z}$ in one dimension (the three-dimensional version is verbatim, dimension by dimension, since nothing below couples distinct axes). The amplitude $\varphi_{j}^{(n)} \in \mathbb{R}$ (or $\mathbb{C}$) is a field value, not a quantum state (\cref{def:quantum-layer}) and not a content (\cref{def:contentvector}): it carries no norm and no sum. The frame differences are
\begin{equation}\label{eq:frame-diff}
\begin{aligned}
\Delta_{t}^{2}\varphi_{j}^{(n)} &= \varphi_{j}^{(n+1)} - 2\varphi_{j}^{(n)} + \varphi_{j}^{(n-1)},\\
\Delta_{x}^{2}\varphi_{j}^{(n)} &= \varphi_{j+1}^{(n)} - 2\varphi_{j}^{(n)} + \varphi_{j-1}^{(n)},
\end{aligned}
\end{equation}
dimensionless in site and step units. Time is measured in $t_{\mathrm P}$, space in $\ell_{\mathrm P}$, and the first Planck relation (\cref{def:units}) makes the two interchangeable --- the commensurability on which everything below rests. The symbols and their exit lines:
\begin{center}
\begin{tabular}{@{}p{2.2cm}p{2.2cm}p{3.4cm}p{3.0cm}@{}}
\toprule
Symbol & Identity & Domain & Beyond the domain \\
\midrule
$n, j$ & pure counting & integers & no meaning between frames \\
$\varphi$ & field amplitude & propagation layer & $\neq \Psi$, $\neq \mathcal E_{i}$ \\
$Q$ & wavenumber & Brillouin zone $|Q\ell_{\mathrm P}| \leq \pi$ & folds back (\cref{rem:folding}) \\
$c$ & the Planck relations & matching conversions & the two uses not interchanged \\
$v_{\mathrm{ph}}, v_{\mathrm{g}}, v_{\mathrm{f}}$ & comparison vocabulary & \cref{prop:causal-auto} & outside the proved chain \\
\bottomrule
\end{tabular}
\end{center}
An exit line is enforced throughout: a symbol asked to leave its row is stopped.

\subsection{The wave equation}\label{sec:wave-eq}

In lattice units the propagation layer's equation is
\begin{equation}\label{eq:wave-eq}
\bigl(\Delta_{t}^{2} - \Delta_{x}^{2}\bigr)\varphi_{j}^{(n)} = 0
\quad\Longleftrightarrow\quad
\varphi_{j}^{(n+1)} = 2\varphi_{j}^{(n)} - \varphi_{j}^{(n-1)} + \bigl(\varphi_{j+1}^{(n)} - 2\varphi_{j}^{(n)} + \varphi_{j-1}^{(n)}\bigr),
\end{equation}
a recurrence determining each frame from the two preceding ones. The coefficients $1$ and $-1$ are the matching in arithmetic form: with time in $t_{\mathrm P}$ and space in $\ell_{\mathrm P}$, the ratio $\ell_{\mathrm P}^{2}/t_{\mathrm P}^{2} = c^{2}$ is absorbed into the unit convention, and the two differences enter with equal weight.

\subsection{The exact dispersion relation}\label{sec:dispersion}

\begin{theorem}[Exact dispersion]\label{thm:dispersion}
The plane wave $\varphi_{j}^{(n)} = e^{i(Qj\ell_{\mathrm P} - \omega n t_{\mathrm P})}$ solves \cref{eq:wave-eq} if and only if
\begin{equation}\label{eq:exact-dispersion}
\cos(\omega t_{\mathrm P}) = \cos(Q\ell_{\mathrm P})
\quad\Longrightarrow\quad
\omega = cQ, \qquad Q\ell_{\mathrm P} \in [-\pi, \pi],
\end{equation}
\uses{\cref{eq:wave-eq}, \cref{def:units}}
\end{theorem}

\begin{proof}
\emph{Step 1 (differences on a plane wave).} Each difference returns the plane wave times one factor, by the geometric sum of two phases, i.e.\ by the definition of cosine:
\begin{equation}\label{eq:fourier-action}
\Delta_{x}^{2}e^{iQj\ell_{\mathrm P}} = 2\bigl(\cos Q\ell_{\mathrm P} - 1\bigr)e^{iQj\ell_{\mathrm P}},
\end{equation}
and identically in time, $\Delta_{t}^{2}e^{-i\omega n t_{\mathrm P}} = 2(\cos\omega t_{\mathrm P} - 1)e^{-i\omega n t_{\mathrm P}}$. One line each.

\emph{Step 2 (substitution).} \Cref{eq:wave-eq} becomes $2(\cos\omega t_{\mathrm P} - 1) - 2(\cos Q\ell_{\mathrm P} - 1) = 0$, the two $-2$'s cancelling: $\cos\omega t_{\mathrm P} = \cos Q\ell_{\mathrm P}$. The matching entered in Step 1's silent agreement of the two factors --- without $\ell_{\mathrm P} = ct_{\mathrm P}$ the factors would carry different constants and the cancellation would not exist.

\emph{Step 3 (root and branch).} On $[-\pi, \pi]$ cosine is even and injective on each half, so $\cos(\omega t_{\mathrm P}) = \cos(Q\ell_{\mathrm P})$ forces $\omega t_{\mathrm P} = \pm Q\ell_{\mathrm P}$; the physical branch (positive frequency, the sign of $Q$ carrying the direction) is $\omega t_{\mathrm P} = Q\ell_{\mathrm P}$, and $\ell_{\mathrm P} = ct_{\mathrm P}$ gives $\omega = cQ$. Both directions of the equivalence are covered by reading the same three equalities backwards.
\end{proof}

This is the independent rederivation foreseen at \cref{er:dispersion}: $\omega = ck$, quoted there from standard kinematics, is here exact at every wavenumber of the zone --- not an approximation, not a low-energy limit, an identity. The closure is one-way as registered there: the lattice result is not used before this point, and the earlier uses remain valid as external quotations.

\subsection{The light cone}\label{sec:lightcone}

\begin{lemma}[Light-cone support]\label{lem:lightcone}
The value $\varphi_{j}^{(n)}$ depends only on initial data within the cone $|j' - j| \leq n - n'$.
\uses{\cref{eq:wave-eq}}
\end{lemma}

\begin{proof}
Induction on the frame. The recursion \cref{eq:wave-eq} expresses $\varphi_{j}^{(n+1)}$ through $\varphi_{j\pm1}^{(n)}$, $\varphi_{j}^{(n)}$, and $\varphi_{j}^{(n-1)}$ alone: each contributing site lies within one step of $j$ and one frame of $n$. If every value at frames $\leq n$ respects the cone, the values at frame $n+1$ respect the cone widened by one step --- the cone of slope one step per frame. The induction starts at the initial surface, where the claim is the tautology of the data.
\end{proof}

\begin{proposition}[Causality is automatic; undistortion is not]\label{prop:causal-auto}
Any evolution law coupling only nearest neighbours --- any recursion of the form \cref{eq:wave-eq}, whatever its coefficients --- transmits influence at most one step per frame, by exactly the induction of \cref{lem:lightcone}. Causality in this sense is a property of \emph{locality}, shared by every local operator on the frame sequence. The wave operator's distinctive property is what \cref{thm:dispersion} adds beyond it: every mode propagates \emph{without distortion}, $\omega = cQ$ exactly, waveforms unchanged within the cone. Two comparisons keep the two properties separate: a superluminal \emph{phase} velocity alone would signal nothing (a phase pattern is not a signal), and a local operator with distorted dispersion --- exhibited in \cref{thm:match} --- respects the cone while deforming everything inside it. As a measure of the undistortion: the three standard velocities coincide, $v_{\mathrm{ph}} = \omega/Q = c$, $v_{\mathrm{g}} = d\omega/dQ = c$, and the cone slope $v_{\mathrm{f}} = c$ each by one substitution into $\omega = cQ$, whereas a generic dispersion splits them (quoted externally as standard wave theory). Causality is the floor; undistorted propagation is the wave operator's alone.
\uses{\cref{lem:lightcone}, \cref{thm:dispersion}}
\end{proposition}

\begin{remark}[Brillouin folding]\label{rem:folding}
The general solution of $\cos\theta = \cos\phi$ is $\theta = \pm\phi + 2\pi m$; every $m \neq 0$ branch describes the same physical mode as $m = 0$, folded back by the $2\pi$-periodicity of the lattice's Fourier image. The negative branch is the backward-propagating wave, its direction carried by the sign of $Q$. The modes of the frame sequence are therefore exactly the interval $Q\ell_{\mathrm P} \in [-\pi, \pi)$, counted once each: no hidden excitations, no spectrum beyond the zone.
\end{remark}

\subsection{The matching theorem: undistorted propagation is unique}\label{sec:match}

Locality, translation invariance, and time-reversal symmetry define the class precisely: any operator in it, given the plane wave, produces a dispersion equation
\begin{equation}\label{eq:multiplier}
F(\omega t_{\mathrm P},\; Q\ell_{\mathrm P}) = 0,
\end{equation}
with $F$ a real trigonometric polynomial --- a finite sum of sines and cosines of integer multiples of the two arguments, the direct image of the finite difference stencil. The wave operator's $F$ is $\cos\theta - \cos\phi$ in the shorthand $\theta = \omega t_{\mathrm P}$, $\phi = Q\ell_{\mathrm P}$ used for the duration of this subsection only.

\begin{theorem}[The matching theorem]\label{thm:match}
Let $F = F_{0} + \varepsilon F_{1}$ with $F_{0} = \cos\theta - \cos\phi$ and $F_{1}$ a real trigonometric polynomial. If, for every $\phi \in (-\pi, \pi)$, the real zeros of $F$ in $\theta$ are exactly $\theta = \pm\phi$ --- every mode undistorted --- then $F$ and $F_{0}$ have the same zero set near $\varepsilon = 0$, and the homogeneous equation of $F$ has the same solution space as the wave operator's.
\uses{\cref{eq:multiplier}; implicit differentiation; one change of variables}
\end{theorem}

\begin{proof}
\emph{Step 1 (zeros stay put).} Write a zero branch as $\theta^{\pm}(\phi) = \pm\phi + \varepsilon\theta_{1}^{\pm}(\phi) + O(\varepsilon^{2})$ and impose $F = 0$ order by order:
\begin{equation}\label{eq:implicit}
0 = F_{0}(\pm\phi, \phi) + \varepsilon\Bigl[\theta_{1}^{\pm}(\phi)\,\partial_{\theta}F_{0}(\pm\phi, \phi) + F_{1}(\pm\phi, \phi)\Bigr] + O(\varepsilon^{2}),
\end{equation}
where $F_{0}(\pm\phi,\phi) = 0$ identically and $\partial_{\theta}F_{0} = -\sin\theta = \mp\sin\phi$. The zeros stay put, $\theta_{1}^{\pm} \equiv 0$, if and only if $F_{1}$ vanishes on both diagonals:
\begin{equation}\label{eq:divisibility}
F_{1}(\phi, \phi) \equiv 0, \qquad F_{1}(-\phi, \phi) \equiv 0.
\end{equation}

\emph{Step 2 (factorization, by one change of variables).} Set $u = \tfrac{\theta + \phi}{2}$, $v = \tfrac{\theta - \phi}{2}$: every trigonometric polynomial in $(\theta, \phi)$ becomes one in $(u, v)$. The two diagonals are $v \equiv 0 \pmod{\pi}$ and $u \equiv 0 \pmod{\pi}$; a polynomial in $v$ vanishing at $v \equiv 0 \pmod{\pi}$ is divisible by $\sin v$ (factor out $e^{\pm iv}$ and use that a one-variable polynomial with those roots has them), and symmetrically for $u$. Hence
\begin{equation}\label{eq:factor}
F_{1} = \sin u \,\sin v \,\tau = \tfrac{1}{2}\bigl(\cos\theta - \cos\phi\bigr)\,\tau
\end{equation}
for some trigonometric polynomial $\tau$, where $u + v = \theta$ and $u - v = \phi$ give $2\sin u \sin v = \cos(u - v) - \cos(u + v) = \cos\phi - \cos\theta$, the sign absorbed into $\tau$. Then $F = (\cos\theta - \cos\phi)\bigl(1 + \tfrac{\varepsilon}{2}\tau\bigr)$: the second factor is nonvanishing on the zero set at small $\varepsilon$, and the solution space --- determined by the zero set alone --- is unchanged.

\emph{Step 3 (what non-divisibility does).} An $F_{1}$ violating \cref{eq:divisibility} moves the zeros by
\begin{equation}\label{eq:distortion}
\theta_{1}^{\pm}(\phi) = \frac{F_{1}(\pm\phi, \phi)}{\mp\sin\phi},
\end{equation}
from \cref{eq:implicit}. Example: $F_{1} = (\cos\phi - 1)^{2}$, a fourth-order spatial correction. Near $\phi = 0$: $F_{1} = \phi^{4}/4 + O(\phi^{6})$ and $\mp\sin\phi = \mp(\phi + O(\phi^{3}))$, so
\begin{equation}\label{eq:example}
\theta_{1}^{+}(\phi) = \frac{\phi^{4}/4 + O(\phi^{6})}{\phi + O(\phi^{3})} = \frac{\phi^{3}}{4} + O(\phi^{5}) \neq 0,
\end{equation}
a cubic distortion of the dispersion at small wavenumber; the same computation at $\phi \to \pi$ gives a power of different order, so no cancellation between the two ends of the zone can restore \cref{eq:divisibility}. Undistorted propagation therefore forces \cref{eq:factor}, and the wave operator's solution space is the only one with the property.
\end{proof}

\begin{remark}[Reading of the matching theorem]\label{rem:match-reading}
The theorem is a uniqueness statement, and its two clauses deserve separation. The wave operator is the unique \emph{local, translation-invariant, time-reversal-symmetric} operator whose modes all propagate undistorted; operators outside the class are beyond the theorem's reach, not disproved. And the corrections the theorem excludes from the propagation seat do not vanish from the framework --- they are the formal constructions of the repair registry (\cref{sec:repair}), legitimate there because they claim no propagation. The seat admits one operator; the registry admits the rest; \cref{eq:factor} is the hinge.
\end{remark}

\section{The zero point}\label{sec:zero-point}

The frame sequence's second exact property concerns the oscillator that defines its frame rate. Nothing in this section is re-derived: every ingredient was established where it belongs, and the section's content is the one identity they compose, with two notes that keep its factors apart.

\begin{proposition}[The zero-point cell]\label{prop:zero-point}
The ground-state energy of the reference oscillator (\cref{as:p}) is
\begin{equation}\label{eq:zero-point}
E_{0} = \frac{\hbar\omega_{0}}{2} = \frac{\hbar}{2t_{\mathrm P}} = \frac{\varepsilon_{\mathrm P}}{2} = \frac{\alpha_{0}}{t_{\mathrm P}},
\end{equation}
\uses{\cref{as:p}, \cref{thm:spectrum}, \cref{def:units}, \cref{def:alpha0}, \cref{def:epsilon}}
\end{proposition}

\begin{proof}
Three equalities, each one substitution. First: \cref{thm:spectrum} gives $E_{0} = \hbar\omega_{0}/2$, the oscillator floor of \cref{sec:quantum}. Second: $\omega_{0} = 1/t_{\mathrm P}$ (\cref{eq:framerate}) gives $\hbar\omega_{0}/2 = \hbar/2t_{\mathrm P}$. Third: $\hbar/t_{\mathrm P} = \varepsilon_{\mathrm P}$ (\cref{def:epsilon}) gives $\hbar/2t_{\mathrm P} = \varepsilon_{\mathrm P}/2$. Last: $\alpha_{0} = \hbar/2$ (\cref{def:alpha0}) gives $\alpha_{0}/t_{\mathrm P} = \hbar/2t_{\mathrm P}$, closing the chain.
\end{proof}

The chain \cref{eq:zero-point} is the zero-point content of the framework: the oscillator's floor, expressed in the matched units, is one minimal action quantum per tick --- the same value the injection theorem assigns to the per-frame growth of the content layer (\cref{thm:injection}). That the two constructions meet was registered in \cref{rem:meeting}, with the shared premise declared (the matched unit system) and independence claimed only for the derivation chains; nothing is added to the registration here, and nothing here is used to strengthen it. The identity is recorded as what it is: one value, read in two languages --- the spectral language of the oscillator and the incremental language of the open system --- whose agreement is a consistency of the framework, not a mechanism within it.

\begin{remark}[The two quarters]\label{rem:two-quarters}
The floor's split $\hbar\omega/4 + \hbar\omega/4$ of \cref{eq:twoquarters} is the equal partition of kinetic and potential energy: $\langle \hat T\rangle = \langle \hat V\rangle$ at the minimum, here read directly off the two evaluations rather than through the virial theorem (quoted externally as the general statement for homogeneous potentials of degree $2$).
\end{remark}

\begin{remark}[The two half-factors, kept apart]\label{rem:two-halves}
Two unrelated factors of $\tfrac12$ meet in the neighbourhood of \cref{eq:zero-point}, and the framework keeps them apart as a matter of record. The half in $E_{0} = \hbar\omega_{0}/2$ is the saturation floor of the uncertainty relation (\cref{rem:half}): it originates in \cref{lem:kennard} and belongs to \cref{sec:quantum}. The half in the Cayley parameter $\vartheta = \hat H t_{\mathrm P}/2\hbar$ is a coefficient-matching artifact of the $[1,1]$ Pad\'e form (\cref{rem:approx}): it originates in \cref{def:cayley} and belongs to \cref{ch:unitarity}. The two share a numeral and nothing else; no narrative merges them, and the $2\pi$ convention that fixes $\omega_{0} = 1/t_{\mathrm P}$ against $2\pi/t_{\mathrm P}$ was settled by the consistency check of \cref{thm:injection}, recorded at \cref{as:p}, not by either half.
\end{remark}

\section{Dependencies and closures}\label{sec:lattice-backfill}

\begin{remark}[What this chapter stands on]\label{rem:ch10-premises}
The chapter's premises are three: the frame sequence (\cref{as:a1}), the stepping structure (\cref{as:stepping}), and the Planck relations (\cref{def:units}). Everything else in \cref{sec:propagation,sec:zero-point} is proved inside the chapter: the dispersion relation (\cref{thm:dispersion}), the light cone (\cref{lem:lightcone}), the automaticity and the uniqueness (\cref{prop:causal-auto}, \cref{thm:match}), and the zero-point identity (\cref{prop:zero-point}) use no external theorem within their proofs. The chapter's only external content is comparative --- the splitting of the three velocities in generic dispersions, quoted as standard wave theory in one sentence of \cref{prop:causal-auto} --- and is marked as comparison wherever it appears.
\end{remark}

\begin{remark}[The closures made here]\label{rem:ch10-closures}
Three forward references end in this chapter, and the endings are recorded. (i) The dispersion relation $\omega = ck$, quoted from standard kinematics at \cref{er:dispersion} of \cref{sec:light} with a promise of independent rederivation, is rederived as \cref{thm:dispersion} --- exact at every wavenumber of the zone; the closure is one-way, as registered there, and the earlier uses stand as external quotations. (ii) The propagation layer's causal test, cited by the repair registry (\cref{rem:formal-construction}) as the criterion idle against elliptic operators, is \cref{prop:causal-auto}: locality's automatic cone, with undistorted propagation as the wave operator's additional, non-automatic property. (iii) The $2\pi$ convention that fixes the frame rate against its cycle-reading alternative, settled by the consistency check of \cref{thm:injection} and recorded at \cref{as:p}, is mutually referenced by \cref{rem:two-halves}; the two references state the same fact from their two sides, and no further closure is owed.
\end{remark}

\begin{remark}[What Chapter 10 has established]\label{rem:ch10summary}
On the frame sequence, with the matching as the only arithmetic: one field (\cref{sec:propagation}), one equation (\cref{eq:wave-eq}), one dispersion (\cref{thm:dispersion} --- $\omega = cQ$ exactly, throughout the Brillouin zone), one cone (\cref{lem:lightcone}), and one uniqueness (\cref{thm:match} --- the wave operator alone among local symmetric operators propagates every mode without distortion). And one zero point (\cref{prop:zero-point}): the reference oscillator's floor is one minimal action quantum per tick, consistent with the injection theorem and adding nothing to it. The chapter's numerics are none; its assumptions are none beyond the frame sequence and the matching; its open items are none --- the one registration touching it, the repair registry's formal constructions, claims no propagation and is unaffected by everything proved here.
\end{remark}

\chapter{Survival without memory}\label{ch:decay}

A point-like entity --- a component of the frame sequence with no internal structure to remember by --- either survives a frame or terminates. This chapter classifies what survival laws such an entity can have, and reads the classification in continuous time.

\section{Uniqueness of the memoryless survival law}\label{sec:survival}

\begin{definition}[Survival function and memorylessness]\label{def:survival}
A \emph{survival function} is $S: \mathbb{Z}^{+} \to [0,1]$ with $S(0) = 1$, $S$ non-increasing; $S(n)$ is the probability of surviving $n$ frames. It is \emph{memoryless} if
\begin{equation}\label{eq:memoryless}
S(m + k) = S(m)\,S(k) \qquad \text{for all } m, k \geq 0.
\end{equation}
\end{definition}

\begin{theorem}[The geometric law is the only memoryless one]\label{thm:geometric}
If $S$ is memoryless and $S(1) = 1 - p$ with $0 < p < 1$, then
\begin{equation}\label{eq:geometric}
S(n) = (1-p)^{n} \qquad \text{for all } n \geq 0,
\end{equation}
and conversely, every geometric law \cref{eq:geometric} is memoryless with $S(1) = 1-p$. The parameter $p$ --- the per-frame termination probability --- labels the entire family: one number, one law.
\uses{\cref{eq:memoryless}; induction}
\end{theorem}

\begin{proof}
\emph{Forward.} Memorylessness with $m = 1$, $k = n$ gives $S(n+1) = S(1)\,S(n)$, a first-order recursion in $n$; iterating from $S(0) = 1$,
\begin{equation}\label{eq:iterate}
S(n) = S(1)^{n} = (1-p)^{n},
\end{equation}
the iteration written out rather than presumed: $S(2) = S(1)S(1)$, $S(3) = S(1)S(2) = S(1)^{3}$, and induction closes the chain for all $n$.

\emph{Converse.} $(1-p)^{m+k} = (1-p)^{m}(1-p)^{k}$ holds identically, one line; the geometric law satisfies \cref{eq:memoryless} with $S(1) = 1 - p$ as claimed.
\end{proof}

\begin{remark}[What is classified, and what is not]\label{rem:classification}
\Cref{thm:geometric} classifies the \emph{shape} of survival under the hypothesis of memorylessness; it does not derive that a given entity is memoryless --- that remains a physical input, stated here as the hypothesis it is. What the classification removes is every other shape: under memorylessness there is no alternative survival law, no free functional form to fit, nothing between the geometric law and its negation. The hypothesis is one sentence long; the theorem is its complete analysis.
\end{remark}

\section{The continuous reading, and the two layers}\label{sec:continuous-decay}

\begin{proposition}[Exponential reading]\label{prop:exponential}
At per-frame probabilities with $p \ll 1$ --- many frames per termination --- the geometric law's continuous reading is exponential:
\begin{equation}\label{eq:exponential}
S(t) = (1-p)^{t/t_{\mathrm P}} = e^{\,\frac{t}{t_{\mathrm P}}\ln(1-p)} \simeq e^{-pt/t_{\mathrm P}},
\end{equation}
using $\ln(1-p) = -p + O(p^{2})$; the rate is
\begin{equation}\label{eq:rate}
\Gamma := -\frac{\ln(1-p)}{t_{\mathrm P}} = \frac{p}{t_{\mathrm P}} + O\Bigl(\frac{p^{2}}{t_{\mathrm P}}\Bigr),
\end{equation}
whose leading term is the per-frame probability over the tick.
\uses{\cref{thm:geometric}; one Taylor step}
\end{proposition}

\begin{proof}
Two steps, both displayed. Regroup the exponent: $t/t_{\mathrm P}$ frames elapse in time $t$, so $S(t) = S(1)^{t/t_{\mathrm P}}$ by \cref{eq:geometric} with $n = t/t_{\mathrm P}$ --- the reading of the frame count, not a new law. Then $\ln(1-p) = -p - p^{2}/2 - \cdots$, one term kept at $p \ll 1$ with the remainder recorded; substituting gives \cref{eq:exponential}, and \cref{eq:rate} is the same expansion named.
\end{proof}

\begin{remark}[Constant hazard is structural]\label{rem:hazard}
The geometric law's hazard --- the conditional termination probability per frame, $S(n) - S(n+1)$ over $S(n)$ --- equals $p$ at every frame by \cref{eq:iterate}: constant hazard is not an additional assumption but a restatement of memorylessness, and the exponential reading inherits it. The short-time behaviour of the frame sequence's law is therefore linear-in-time from the first frame, by \cref{eq:memoryless} and not by approximation; the framework states this as a structural feature and compares it with no other theory's short-time regime.
\end{remark}

\begin{remark}[Two layers, two laws, no conflict]\label{rem:two-layers-decay}
The point layer's geometric law and the content layer's algebraic kernel \cref{eq:kernel} coexist without tension, because they are answers to different questions. The kernel $\mathcal{W}(n;m) = (m-1)/(n-1)$ tracks a \emph{perturbation of a content} --- a deviation from the exact solution, attenuated by the conductance --- while $S(n)$ tracks a \emph{termination event} of a point-like entity. One is linear relaxation of deviations, the other constant-hazard survival; different objects, different function classes, and no statement of either chapter is asked to reproduce the other.
\end{remark}

\begin{remark}[Counting noise, quoted]\label{rem:counting-noise}
For $N$ independent termination events the relative fluctuation is $1/\sqrt{N}$ (quoted externally as the standard counting statistics of independent events). The remark is registered as the experimental counterpart of the framework's point-layer law and used nowhere below; it carries no numerical content.
\end{remark}

\begin{remark}[What Chapter 11 has established]\label{rem:ch11summary}
Under one hypothesis --- memorylessness --- exactly one survival law exists (\cref{thm:geometric}), read continuously as the exponential (\cref{prop:exponential}) with a structurally constant hazard (\cref{rem:hazard}); it is layered against the content chain's kernel without conflict (\cref{rem:two-layers-decay}), and its experimental face is one line of standard counting statistics (\cref{rem:counting-noise}). The chapter's premises are \cref{eq:memoryless} alone; its externals are two one-line quotations, both comparative; its open items are none; and the chapter adds no new object, parameter, or scale to the framework.
\end{remark}

\chapter{Two roads to the continuum}\label{ch:limits}

The frame sequence is discrete, and the two fundamental equations of quantum physics are differential. Two external results bridge the gap, each connecting a discrete model to one equation, and this chapter records both --- with their status stated first: both theorems are established outside this framework, quoted with their conditions as stated in the sources, and neither is derived here. What the framework contributes is only the carrier: the lattice with its matched units (\cref{def:units}), and the question of which discrete dynamics lives on that carrier is the chapter's one open item.

\section{The quantum walk and the Dirac equation}\label{sec:walk-dirac}

\begin{eresult}[Quantum walks in the continuous limit]\label{er:walk-dirac}
A $1{+}1$-dimensional quantum walk with the Hadamard-type coin, under the parameter and measurement-setup approach of the source, converges in the continuous limit to the Weyl equation, and with the appropriate two-component structure to the Dirac equation; the convergence is made rigorous in the strictification cited below.
\reference{D.~A. Meyer, J.\ Stat.\ Phys.\ \textbf{85}, 551 (1996); F.~W. Strauch, Phys.\ Rev.\ A \textbf{74}, 030301(R) (2006). Conditions as stated there: the walk parameters and the measurement setup must follow the source's approach; the result is a theorem about the limit, not an identification of the frame sequence's dynamics with any particular walk.}
\end{eresult}

The interface with the framework is two numbers: the walk's step is the lattice spacing $\ell_{\mathrm P}$, and its clock is the frame rate $1/t_{\mathrm P}$ (\cref{eq:tick}) --- the matched units are exactly the scale at which a walk's continuum limit is taken, so \cref{er:walk-dirac} is applicable to the frame sequence as a carrier the moment a walk on it is specified. That specification is the chapter's open item, registered here and used nowhere: the coin and shift operators of the sequence's walk are \textbf{not determined by the framework's conservation structure} --- nothing in \cref{ch:eth,ch:newton,ch:unitarity,ch:cascade} selects a coin --- and their choice is an additional hypothesis, to be made where the walk is defined, not before.

\section{Diffusion and the Schr\"odinger equation}\label{sec:nelson}

\begin{eresult}[Nelson's stochastic mechanics]\label{er:nelson}
A diffusion with coefficient $D = \hbar/2m$, together with Newton-type dynamics for the osmotic velocity, implies the Schr\"odinger equation with mass $m$; the derivation is made rigorous in the modern strictification cited below.
\reference{E.~Nelson, Phys.\ Rev.\ \textbf{150}, 1079 (1966); M.~Hall and M.~Reginatto, quant-ph/0002050 and subsequent rigorisation. Conditions as stated there: the osmotic condition is part of the hypothesis, and the result is a theorem about the continuum limit of the stochastic model.}
\end{eresult}

The interface is definitional: the relation $m = \hbar/2D$ is used as a definition of the mass in terms of the diffusion coefficient wherever the two descriptions meet, the same anchor-family status as every other use of $m$ in this framework (\cref{rem:repair-status}). The two roads then arrange themselves by layer: the \emph{coherent} walk --- ballistic, variance growing as $n^{2}$ --- has the Dirac equation as its continuum limit (\cref{er:walk-dirac}); the \emph{decohered} count --- variance growing as $n$ --- diffuses, and diffusion has the Schr\"odinger equation as its limit (\cref{er:nelson}). The amplitude description and the counting description are thus the two continuum faces of the frame sequence's two transport regimes, and which face a given phenomenon presents is a matter of measurement, not of the sequence.

\begin{remark}[What Chapter 12 has established]\label{rem:ch12summary}
Two interfaces, both external and both one-way: the frame sequence's matched units are the natural carrier for a quantum walk whose limit is the Dirac equation (\cref{er:walk-dirac}), and the sequence's decohered counting is the natural carrier for a diffusion whose limit is the Schr\"odinger equation (\cref{er:nelson}). Nothing in the chapter is derived; nothing in the chapter is used as a premise elsewhere; and the chapter's single open item --- the coin and shift operators, undetermined by the conservation structure --- stands registered until a walk on the sequence is defined. The chapter's premises are the two citations; its contributions are two interfaces and one registration.
\end{remark}

\part{Summary}

\chapter{The account, closed}\label{ch:conclusion}

\section{The chain, end to end}\label{sec:chain-closure}

The paper's content compresses into one chain, each link earned before the next is forged. Measurement supports the magnitude (\cref{prop:tau}): no apparatus of the empirical inputs resolves below $t_{\mathrm P}$. The hypothesis layer stands on that support (\cref{sec:hypotheses}): discrete registration, action quantised per tick, a reference oscillator at the frame rate. The matching fixes the units (\cref{sec:units}), the oscillator floor gives the half-unit $\alpha_{0}$ and the odd quantum numbers (\cref{sec:quantum,sec:oddquantum}). The open three-component chain runs on three hypotheses, and its law is forced --- the $1{:}1$ split and the exponent $\gamma = 1$ are the unique non-degenerate steady state (\cref{thm:forced-cascade}). The Newtonian leg, matched to the frame sequence, fixes the injection rate, the power, and the per-frame quantum (\cref{ch:newton}); unitarity holds where it belongs and the injection theorem makes the growth one minimal action quantum per frame (\cref{thm:injection}); the open system's environment is explicit and minimal (\cref{thm:minimal}). The boundary acquires entropy and temperature from two anchors, and $E = TS$ is a rewrite, not an assumption (\cref{ch:boundary}); the relativistic reading is a set of one-way formal coincidences, with no slot for dynamics and none for $\Lambda$ (\cref{ch:gr}). Symmetry is catalogued seat by seat: the boost is a stranger to the frame sequence and a symmetry of the plane (\cref{thm:lattice-lorentz,thm:bilateral}), rotation is three-dimensional rigidity (\cref{thm:rigidity}), and sixteen cells of verdicts are complete (\cref{sec:sixteen}). On the frame sequence itself, one field propagates without distortion at exactly $\omega = cQ$ (\cref{thm:dispersion}), the unique such local operator (\cref{thm:match}), and the oscillator's floor equals the per-frame quantum (\cref{prop:zero-point}). Survival without memory is geometric and nothing else (\cref{thm:geometric}), and two external bridges connect the sequence's two transport regimes to the Dirac equation and the Schr\"odinger equation (\cref{er:walk-dirac,er:nelson}), with the walk's own operators registered open.

\section{The master input table}\label{sec:master-table}

The paper's foundational input layer, so that it can be checked in one place:
\begin{center}
\begin{tabular}{@{}p{2.4cm}p{3.3cm}p{2.9cm}p{2.6cm}@{}}
\toprule
Empirical input & Content & Role & Source \\
\midrule
\cref{er:light} & invariance of $c$ & time--space conversion & \cref{sec:units} \\
\cref{er:action} & quantum of action $\hbar$ & the quantum unit & \cref{sec:units} \\
\cref{er:grav} & Newtonian constant $G$ & the gravitational unit & \cref{sec:units} \\
\cref{er:ccr} & $[\hat x, \hat p] = i\hbar$ & quantum kinematics & \cref{sec:quantum} \\
\cref{er:efreq} & $E = \hbar\omega$ & energy--frequency & \cref{sec:quantum} \\
\cref{er:etuncertainty} & $\Delta E \geq \hbar/2\tau$ & measurement resolution & \cref{sec:measurement} \\
\bottomrule
\end{tabular}
\end{center}
\begin{center}
\begin{tabular}{@{}p{2.4cm}p{3.3cm}p{2.7cm}p{2.6cm}@{}}
\toprule
Hypothesis & Content & Role & Source \\
\midrule
\cref{as:a1} & discrete time & the frame sequence & \cref{sec:hypotheses} \\
\cref{as:a2} & action quantised per tick & the quantisation grid & \cref{sec:hypotheses} \\
\cref{as:p} & $\omega_{0} = 1/t_{\mathrm P}$ & the reference oscillator & \cref{sec:hypotheses} \\
\cref{as:b1,as:b2,as:b3} & conduction, conductance, splitting & the open chain & \cref{sec:cascade-axioms} \\
\cref{as:stepping} & $R(n) = n\ell_{\mathrm P}$ & the spatial sequence & \cref{sec:steps} \\
\cref{as:pump} & constant injection & the open system & \cref{sec:pumped-family} \\
\cref{as:ratio-temp} & $T := E/S$ & boundary temperature & \cref{sec:temperature-path} \\
\bottomrule
\end{tabular}
\end{center}
Chapter-level empirical inputs are tabulated in each chapter's closing account: the measurement inputs of \cref{sec:measurement}; the dispersion quotation of \cref{ch:light} (\cref{er:dispersion}); the Newtonian ball and the conservation scaling of \cref{ch:newton} (\cref{er:ball,er:conservation}); the measurement bridge, the dilation theorem, and the collision model of \cref{ch:unitarity,ch:dilation} (\cref{er:bridge,er:stinespring,er:collision}); the three anchors and three benchmarks of \cref{ch:boundary} (\cref{er:bitdensity,er:boltzmann,er:clausius,er:bh,er:gh,er:sb}); the four relativistic quotations of \cref{ch:gr} (\cref{er:friedmann,er:accel,er:age-integral,er:friedmann-lambda}); Fourier duality in \cref{ch:symmetries} (\cref{er:fourier}); and the two bridges of \cref{ch:limits} (\cref{er:walk-dirac,er:nelson}). Below the hypotheses sit conventions, registered wherever they are used: the first-frame clause (\cref{eq:firstframe}), the scale-factor bridge $a := R$ (\cref{eq:scalebridge}), the momentum-grid convention (\cref{eq:correction}), box normalisation, and the angular-frequency convention (\cref{rem:twopi-light}). Everything else in the paper --- every theorem, identity, and verdict of Chapters \ref{sec:units}--\ref{ch:limits} --- is derived from the tables above plus cited external results, or it is an explicitly registered open item.

\section{The open items}\label{sec:open-registry}

What the framework does not contain, collected once, each entry in the place it has always occupied:
\begin{enumerate}
\item the mechanism of the injection (\cref{rem:mechanismregistry}) --- open in both directions, by doctrine;
\item the even action sites (\cref{rem:evensites}) --- registered with the non-bound processes;
\item the flight of a free quantum on the quantisation grid (the scope clause of \cref{rem:layer-scope}) --- deferred to where the propagation layer meets the exchange layer;
\item the emergence of the continuous boost on the reading layer (\cref{rem:not-prohibition}) --- mechanism unbuilt;
\item the negative hemisphere (\cref{rem:negative-conjecture}) --- conjecture in both directions at once;
\item the mass debt $m$ (\cref{rem:repair-status}) --- beside $\hbar$, $c$, $G$, where it belongs;
\item the value of $G$ --- an anchor-family debt of the whole field (\cref{er:grav});
\item the coin and shift operators of the frame sequence's walk (\cref{ch:limits}) --- undetermined by the conservation structure;
\item the coupling rules of any interaction on the propagation seat --- no gauge structure, charge, or coupling constant is fixed anywhere in the framework (\cref{rem:autopsy});
\item the counting layer's ontological status (\cref{rem:wigner}) --- registered, not answered.
\end{enumerate}
No entry is concealed; no entry is promised a closure this paper does not deliver.

\section{The disciplines, finally}\label{sec:disciplines}

Three disciplines governed the writing, and they are the paper's claim on the reader's trust. The \emph{parameter discipline}: no parameter is introduced whose value the framework cannot in principle determine --- one number per structure, and a structure that would require an undetermined constant is redesigned rather than parametrised. The \emph{attribution discipline}: every statement wears its status --- empirical input, hypothesis, definition, derivation, external result, registration --- and no statement is used as a premise for its own ancestor; hypotheses wear no experimental clothing, and conclusions are never quoted as evidence for their premises. The \emph{integrability discipline}: state quantities may be accumulated over frames, ratio quantities may not, and every integral in the paper is a telescope of a frame difference or is flagged as reading-layer shorthand (\cref{rem:integrability}).

What the framework is, in one closing sentence: a frame sequence --- time as a sequence, contents as counts, and a single matched calibration --- from which the Planck scale, the injection quantum, the boundary's thermodynamics, the Friedmann coincidences, the symmetry verdicts, and the exact dispersion follow as arithmetic, with every physical reading labelled as a reading, every mechanism question registered as open, and every contact with the continuous theories made one-way, formal, and explicit.

\end{document}
